\section{One mean volume and the degree of short sections}
\label{rh-geometry}

The geometric task is to give lower bounds for the degrees of the
actual section lattices used by the arithmetic maps. Three reductions
are useful: average the normalized intersection numbers before
estimating them, preserve the component-degree weights, and use the
existing short sections to control negative degree. These are
comparison statements; their geometric hypotheses are retained.

Arithmetic degrees in this section are divided by the degree of the
ground number field. For a nonzero Hermitian lattice $E$, its ordinary
slope is $\mu(E)=\widehat{\deg}(E)/\rank E$. A convention using logarithms
of squared norms gives slope $2\mu(E)$. We use ordinary slopes and
state every factor accordingly.

\subsection{The exact mean-volume input}

Fix a finite nonempty actual family of components $Z$, each with the prescribed
period-normalized metric and integral conditions. Let $V_Z(t)$ be
its normalized asymptotic count volume of sections whose norm in
power $k$ is at most $e^{kt}$. A real number $s_Z$ is an admissible
lower volume slope if
\begin{equation}\label{rh-volume-slope}
 V_Z(t)\ge\max\{\delta_Z(s_Z+t),0\}\quad(t\in\mathbb R),
 \qquad \delta_Z>0.
\end{equation}
The required mean estimate has the form
$|\mathcal Z|^{-1}\sum_Zs_Z\ge-X$, with an explicit $X\ge0$.
An estimate for each separate component is stronger than necessary.

\begin{lemma}[A finite-prime allowance]\label{rh-prime-cost}
Let $E$ be a number field of degree $n_E$, absolute discriminant
$\Delta_E$, and root discriminant $r_E=\Delta_E^{1/n_E}$.
Let $Q$ be the prime divisors of $2\Delta_EM'$ for a positive
integer $M'$. Suppose $B\ge128$ and $|Q|\le\lfloor B\rfloor$.
With $S=\{p\text{ prime}:p\le B\}\cup Q$, put
\[
 c_2=\sum_{p\in S}\left(\frac1{p-1}
       +\frac{40v_p(\Delta_E)}{n_E}+\mathbf1_{p=2}\right)\log p.
\]
Then $0\le c_2\le(59/7)\log B+40\log r_E
\le40(\log B+\log r_E)$.
\end{lemma}

\begin{proof}
The discriminant and the prime two contribute exactly
$40\log r_E+\log2$. Partial summation and \eqref{rh-chebyshev} give
\[
 \sum_{p\le B}\frac{\log p}{p-1}
 \le2\sum_{m\le B}\frac{\Lambda(m)}m
 \le8\log2(1+\log B).
\]
The decreasing function $\log t/(t-1)$ bounds the remaining primes
of $Q$ by $B\log B/(B-1)\le(8/7)\log B$.
Since $\log B\ge7\log2$ and $\log2<3/4$, the total
nondiscriminant contribution is at most
$(6+9/7+8/7)\log B=(59/7)\log B$.
Finally $r_E\ge1$ and $59/7<40$.
\end{proof}

\begin{proposition}[From total intersection to mean volume]
\label{rh-intersection-mean}
Suppose the components have equal positive volumes $v_Z$, and put
$v_T=\sum_Zv_Z$. Let $D>0$, $d\ge0$, and let $I_Z$ be their actual
arithmetic intersection numbers. For nonnegative $L,W,q$ and real
$\lambda$, set
\[
 \mathcal E=(\log B+\log r_E+|\lambda|)L+qW.
\]
Assume, for $C_T\ge0$,
\[
 \sum_ZI_Z\ge-C_T(d+1)v_TD\mathcal E,
\]
and assume the prescribed section-counting comparison makes
$s_Z=I_Z/((d+1)v_ZD)-2c_2L$ admissible in
\eqref{rh-volume-slope}. Under Lemma~\ref{rh-prime-cost},
\begin{equation}\label{rh-intersection-result}
                  \frac1{|\mathcal Z|}\sum_Zs_Z
                         \ge-(C_T+80)\mathcal E.
\end{equation}
\end{proposition}

\begin{proof}
Equal volumes give the exact identity
\[
 \frac1{|\mathcal Z|}\sum_Z\frac{I_Z}{(d+1)v_ZD}
                  =\frac{\sum_ZI_Z}{(d+1)v_TD}.
\]
Its denominator is positive. Divide the assumed inequality and use
$2c_2L\le80(\log B+\log r_E)L\le80\mathcal E$.
\end{proof}

In the intended application, $L$ is the sum of the absolute values
of a weight, $W$ is its total spread, and $\lambda$ is the specified
logarithmic derivative at one. The actual period factor in the norm
of weight $u$ is $(2\pi\xi)^{-\frac12\sum u}$ for a prescribed
$\xi>0$. At the finite primes in $S$ the lattice uses its specified
Hodge--Tate valuations, and elsewhere its original integral valuations.
The proposition does not supply the comparison for an arbitrary
replacement of these norms or lattices.

This normalization is essential. Multiplying the norm in power $k$
by $e^{kc}$ changes $V_Z(t)$ to $V_Z(t-c)$ and every admissible
lower volume slope by $-c$, while leaving the field and its
$L$-function unchanged. A bound for $\lambda$ alone cannot prevent
that change. The uniform arithmetic intersection bound is a separate
geometric input.

The asymptotic comparison fixes each datum before taking the section
power to infinity. A fixed shift of that power and fixed finite
constants disappear in the normalized limit. When the volume uses
powers divisible by eight, the shift must also be divisible by eight.
No uniform convergence threshold is needed for this fixed-data limit.
An estimate at a prescribed finite section power remains a different
quantitative requirement.

The next comparison sharpens the intersection input itself. Under its
displayed polynomial and metric hypotheses, a single normalized test
intersection controls the weight-dependent lower bound. The test and
the chosen weight must use the same coefficients and prescribed metric.

\subsection{One block-sensitive test intersection}

Let $\Psi$ be a finite nonempty set of size $f$, let $N=p+q$ with
$1\le q\le p$, and put $d=fN(N-1)/2$. For a dominant weight $u$, set
\[
 S=\sum_{\tau,i<p}u_{\tau i},\qquad
 T=\sum_{\tau,i<N}u_{\tau i},\qquad C=S-\frac pN T,\qquad
 W=\sum_\tau(u_{\tau0}-u_{\tau,N-1}),\quad L=\sum_{\tau,i}|u_{\tau i}|.
\]
The exact identity
\[
 C=\frac1N\sum_\tau\sum_{a<p\le b}(u_{\tau a}-u_{\tau b})
 \quad\text{gives}\quad
 0\le C\le\min\left\{\frac{pq}{N}W,\frac pN L\right\}.
\]
In particular, $C$ is unchanged by adding a constant to each
embedding's weight. All intersections below use the prescribed metric.

Write $\ell_u=-\sum_{\tau,i}u_{\tau i}X_{\tau i}$ and use the flag
degree convention
\[
 \deg_{\rm Fl}F=(-1)^d\sum_{(\sigma_\tau)\in(\mathfrak S_N)^\Psi}
 \left(\prod_\tau\operatorname{sgn}\sigma_\tau\right)
 \left[\prod_{\tau,i}X_{\tau i}^{N-1-\sigma_\tau(i)}\right]F.
\]
The Vandermonde identity gives
\[
 D(u)=\deg_{\rm Fl}(\ell_u^d)
 =\frac{d!}{(\prod_{i=0}^{N-1}i!)^f}
             \prod_{\tau,a<b}(u_{\tau a}-u_{\tau b}).
\]
Thus $D(u)>0$ for a strictly dominant weight, as assumed below.
For a polynomial of degree $d+1$, put
\[
 \mathcal S_\tau(F)=\deg_{\rm Fl}\sum_{i<p}\partial_{\tau i}F,
 \qquad
 \mathcal T_\tau(F)=\deg_{\rm Fl}\sum_{i<N}\partial_{\tau i}F.
\]

Fix a cutoff $B$, total component volume $v>0$, and an arithmetic
intersection functional $I_B$. Suppose its comparison has the form
\begin{equation}\label{geo-two-polynomial-comparison}
 I_B(F)/v
 =H(F)-A_B\sum_\tau\mathcal S_\tau(F)
       -b_B\sum_\tau\mathcal T_\tau(F)+R_B(F).
\end{equation}
Only two polynomials are needed, with the same coefficients and metric:
$\ell_u^{d+1}$ and the following test. Put $d_0=fpq$,
\[
 \mu=\sum_{\tau,i\ge p}X_{\tau i},\qquad
 V=\prod_\tau\prod_{\substack{a<b\\a,b<p\ {\rm or}\ a,b\ge p}}
                        (X_{\tau b}-X_{\tau a}),\qquad
 F_0=\mu^{d_0+1}V.
\]
Assume $J=\deg_{\rm Fl}(\mu^{d_0}V)>0$ and put $K_0=(d_0+1)J$.
Define
\[
 \beta_B=\frac{I_B(F_0)}{vfqK_0},\quad
 h_0=\frac{H(F_0)}{fqK_0},\quad r_0=\frac{R_B(F_0)}{fqK_0},
 \qquad c_B=\frac{pA_B+Nb_B}{N},
\]
\[
 \sigma_B(u)=\frac{I_B(\ell_u^{d+1})}{v(d+1)D(u)},\qquad
 h_u=\frac{H(\ell_u^{d+1})}{(d+1)D(u)},\quad
 r_u=\frac{R_B(\ell_u^{d+1})}{(d+1)D(u)}.
\]

\begin{proposition}[Eliminating the two scalar coefficients]
\label{geo-tower-two-polynomial}
Under these hypotheses,
\[
 \boxed{\quad
 \sigma_B(u)=h_u+\frac Np
       \{C(\beta_B-h_0-r_0)+c_BS\}+r_u.\quad}
\]
\end{proposition}

\begin{proof}
Differentiation in \eqref{geo-two-polynomial-comparison} gives
$\sigma_B(u)=h_u+A_BC+c_BT+r_u$.
The polynomial $V$ is invariant under translating either block at
a fixed $\tau$. Hence $\mathcal S_\tau(F_0)=0$ and
$\mathcal T_\tau(F_0)=qK_0$, so $\beta_B=h_0-b_B+r_0$.
Thus $A_B=(N/p)(c_B+\beta_B-h_0-r_0)$. Substitution and
$T+(N/p)C=(N/p)S$ prove the formula.
\end{proof}

\begin{corollary}[A single total test suffices]\label{geo-tower-block-bound}
Let $\Lambda\ge0$, $C_0,C_F\ge0$, and suppose
\[
 \beta_B\ge-C_0\Lambda,\quad
 h_u\ge0,\quad h_0\le\tfrac12\log(N+1),\quad
 |r_0|\le\frac8{5fq},\quad r_u\ge-2qW,\quad |c_B|\le C_F\Lambda.
\]
With $Z=C_0\Lambda+\tfrac12\log(N+1)+8/(5fq)$, one has
\[
 \sigma_B(u)\ge-\frac Np\{CZ+C_F\Lambda|S|\}-2qW
 \ge-\max\left\{Z,\frac qpZ+\frac Np C_F\Lambda\right\}L-2qW.
\]
The first inequality may retain the signed term $(N/p)c_BS$
in place of its absolute-value bound.
\end{corollary}

\begin{proof}
Use $C\ge0$ in Proposition~\ref{geo-tower-two-polynomial}.
For the second inequality write $R=T-S$ and
$(N/p)C=(q/p)S-R$, then use $|S|+|R|\le L$ and $Z\ge0$.
\end{proof}

The remaining geometric input is the single lower bound
$I_B(F_0)/(vfqK_0)\ge-C_0\Lambda$ on the selected towers, with the
required uniformity. The comparison does not establish that bound.
The exact factor $N/p$ is less than $3/2$ when $3q+1\le N$;
the remainder retains its denominator $fq$.

Finally suppose the same functional satisfies
\[
 I_B\bigl((\sum_iX_{\tau i})Q\bigr)
   =v\bar h_{B,\tau}\deg_{\rm Fl}Q,\qquad
 \bar h_{B,\tau}=-Nc_B
 \quad(\deg Q=d).
\]
For any constants $z_\tau$, differentiation along $u+tz$ gives the
exact translation law
\[
             \sigma_B(u+z)=\sigma_B(u)+Nc_B\sum_\tau z_\tau.
\]
Thus centering the weight incurs a signed determinant term. An
anchoring condition $u_{\tau a}=0$ does not make its total sum zero.
Removing this term requires an identity for the actual flat-section
metric, including its finite-place corrections; its cancellation is
not among the hypotheses proved here.


\paragraph{The classical height theorem and the cusp limit.}
For a projective integral variety $X$ of dimension $n$ over a number
field $K$, normalize arithmetic intersections by $[K:\mathbb Q]$ and
point heights by the absolute degree of their field. Write
\[
 D(\bar L)=\deg_LX,\qquad
 I(\bar L)=\frac{\widehat{\deg}(\bar L^{n+1})}{[K:\mathbb Q]}.
\]
Zhang's successive-minima theorem, in its relative-semipositive form
\cite[Theorems A.4.2 and A.5.2]{yuan-zhang-adelic}, implies
\begin{equation}\label{geo-classical-height-lower}
 h_{\bar L}(P)\ge a\quad(P\in X(\overline K))
 \quad\Longrightarrow\quad
 I(\bar L)\ge(n+1)aD(\bar L).
\end{equation}
Here $L$ is ample. For a smooth Hermitian model, relative semipositivity
means semipositive curvature and nonnegative degrees on curves in
finite fibres; the metric is conjugation invariant. Nonnegative
horizontal point heights are not assumed. The point-height lower
bound makes every successive minimum at least $a$. Rational tensor
powers are included by homogeneity.

\begin{proposition}[Reciprocal perturbations suffice]
\label{geo-reciprocal-height-limit}
Fix a component and its prescribed cutoff and metrics, and finite
real numbers $a,\gamma,u_B$.
Let $\bar L_{m,j}$, $m,j\ge1$, be prescribed smooth Hermitian
$\mathbb Q$-line bundles on projective models of $X$, with perturbation
$1/m$ and smoothing index $j$. Suppose that, for each $m$ and all
sufficiently large $j$, they satisfy the positivity and ampleness
conditions above and
\[
             h_{\bar L_{m,j}}(P)\ge a-\gamma/m
                         \quad(P\in X(\overline K)).
\]
The threshold may depend on $m$ but must work for every point.
Suppose finite real $d_m,A_m,V_m$ satisfy
\[
 D(\bar L_{m,j})\longrightarrow d_m,\qquad
 \limsup_{j\to\infty}I(\bar L_{m,j})\le A_m+(n+1)V_m,
\]
and, for the original prescribed intersection $I_B$ and degree $d_B>0$,
\[
 A_m\longrightarrow I_B,\qquad d_m\longrightarrow d_B,\qquad
 V_m\longrightarrow u_Bd_B/2.
\]
Then
\[
                   \frac{I_B}{(n+1)d_B}\ge a-u_B/2.
\]
\end{proposition}

\begin{proof}
Apply \eqref{geo-classical-height-lower} and take the lower limit in
$j$. It is at most the assumed upper limit, so
\[
 (n+1)(a-\gamma/m)d_m\le A_m+(n+1)V_m.
\]
Let $m\to\infty$ and divide by $(n+1)d_B>0$. The argument permits
negative $a$, $u_B$ or $V_m$.
\end{proof}

For a finite nonempty set of actual components, suppose the same-tower test
identification is
\[
 I^{\rm test}_{B,Z}
 =s_{B,Z}(d_0+1)v_ZJ,\qquad
 s_{B,Z}=\frac{I_{B,Z}}{(n_Z+1)d_{B,Z}},\qquad
 v_Z,J>0,\quad d_0=fpq.
\]
With $v=\sum_Zv_Z$, the normalized total test is
\[
 \beta_B=\frac{\sum_ZI^{\rm test}_{B,Z}}{vfq(d_0+1)J}
         =\frac1{fq}\sum_Z\frac{v_Z}{v}s_{B,Z}.
\]
If Proposition~\ref{geo-reciprocal-height-limit} applies with
$a_Z=-C_hfq\Lambda_B$ and $u_{B,Z}\le C_ufq\Lambda_B$, then
\[
                   \beta_B\ge-(C_h+C_u/2)\Lambda_B.
\]
The signed component means may be retained before estimating them.
This is the input of Corollary~\ref{geo-tower-block-bound}. Its
remaining geometric premises are the prescribed positive
regularizations, the all-point height floor including the boundary,
and the one-sided intersection comparison and signed cusp limit.
A lower bound for the essential minimum alone does not imply
\eqref{geo-classical-height-lower}. Each datum is fixed before the
limits, so no convergence rate uniform in the growing family is
required for this deduction.


\subsection{Integrability of the isotropic lattice count}
\label{sec:arithmetic-isotropic-integrability}

The finite-volume and integrability assumptions in
Proposition~\ref{prop:cusp-counting-haar} follow from the arithmetic
unitary quotient itself.

\begin{proposition}\label{prop:arithmetic-isotropic-integrability}
Let $E/F$ be a CM extension with $[F:\mathbf Q]=f$, let $(V,h)$
be a nondegenerate Hermitian $E$-space of dimension $N\ge3$, and
let $L$ be a full $\mathcal O_E$-lattice. Let $\Gamma$ have finite
index in its unitary lattice stabilizer, and put
$G=\prod_{\sigma:F\hookrightarrow\mathbf R}U(h_\sigma)$.
The quotient $\Gamma\backslash G$ has finite Haar volume. For its
normalized Haar measure $\nu$ and any fixed positive norm, the count
\[
 N_t(g)=\#\{v\in L\setminus\{0\}:h(v,v)=0,\ \|g^{-1}v\|<t\}
\]
is integrable for every $t>0$, and
\[
            \int N_t\,d\nu=t^{2f(N-1)}\int N_1\,d\nu.
\]
The same holds for the special unitary group.
\end{proposition}

\begin{proof}
The rational group $\operatorname{Res}_{F/\mathbf Q}U(h)$ has
semisimple derived group and a norm-one center with compact real
points. It therefore has no rational characters. The
Borel--Harish-Chandra theorem gives finite arithmetic covolume
\cite{arithmetic-borel-harish-chandra}. If no rational isotropic
vector exists, all counts vanish. This includes any case with a
definite archimedean factor.

Otherwise, reduction theory gives finitely many $\Gamma$-orbits
of rational isotropic lines \cite{arithmetic-platonov-rapinchuk}.
Choose representatives $Ee_j$ with $e_j\in L$, and put
\[
 J_j=\{a\in E:ae_j\in L\},\qquad
 U_j=\{u:\gamma e_j=ue_j\text{ for some }\gamma\in\Gamma_{Ee_j}\}.
\]
Then $J_j$ is a fractional ideal and $U_j\subset\mathcal O_E^\times$
has finite index. To see the latter directly, complete $e_j$ to a
hyperbolic pair and clear the denominators of its two projections by
an integer $m$. Every unit $u\equiv1\pmod{m\mathcal O_E}$ then
gives a lattice isometry
$\operatorname{diag}(u,1,\bar u^{-1})$. Its inverse also preserves
the lattice. The preimage of the prescribed finite-index group
$\Gamma$ still has finite index in these congruence units.

Let $H_j$ fix the vector $e_j$. Fixing the vector removes the
multiplicative factor present in its line stabilizer: $H_j$ is an
extension of $U(e_j^\perp/Ee_j)$ by a unipotent group. It is
connected and has no rational characters. Thus
$\Gamma_{e_j}\backslash H_j$ has finite volume by the same theorem.
Tonelli and unfolding give
\[
 A(e_j;t):=\int\sum_{v\in\Gamma e_j}
           \mathbf1_{\{\|g^{-1}v\|<t\}}\,d\nu(g)
 =\frac{\operatorname{vol}(\Gamma_{e_j}\backslash H_j)}
        {\operatorname{vol}(\Gamma\backslash G)}
   \int_{H_j\backslash G}\mathbf1_{\{\|g^{-1}e_j\|<t\}}\,d\bar g.
\]
The last quotient is the product of nonzero complex null cones.
Its invariant measure is proportional to
$\prod_\sigma dS_\sigma/\|\nabla h_\sigma(v,v)\|$.
Its radial density in each factor is $r^{2N-3}\,dr$, so bounded
norm balls have finite measure. Hence $A(e_j;1)<\infty$.

For $a\in E^\times$, the vectors $e_j$ and $ae_j$ have
\emph{the same} real and arithmetic stabilizers. Keep these same
Haar measures. Multiplication by $a$ scales the cone measure by
$|N_{E/\mathbf Q}(a)|^{N-1}$. Therefore, with $s=N-1$,
\[
 A(ae_j;t)=t^{2fs}|N_{E/\mathbf Q}(a)|^{-s}A(e_j;1).
\]
The vector orbits above $Ee_j$ are $(J_j\setminus\{0\})/U_j$.
The ideal $\mathfrak a=a\mathcal O_EJ_j^{-1}$ is integral, has
class $[J_j]^{-1}$ and norm $|N(a)|/NJ_j$. Every such ideal has
$[\mathcal O_E^\times:U_j]$ preimages. Thus
\[
 \int N_t\,d\nu
 =t^{2fs}\sum_j A(e_j;1)[\mathcal O_E^\times:U_j]
           (NJ_j)^{-s}\zeta_{[J_j]^{-1}}(s).
\]
All terms are finite: $s>1$, and each partial ideal zeta series is
bounded by $\zeta_E(s)\le\zeta(s)^{2f}$. This proves both assertions.

For $SU$, determinant correction on a nonisotropic line in the
hyperbolic complement gives rational and real transitivity; here
$N-2\ge1$. Use real congruence units in the scalar lift. They
have finite index in the CM units, since $u/\bar u$ is a root of
unity for every CM unit $u$. The vector stabilizer has special
unitary reductive quotient, and the remaining proof is unchanged.
\end{proof}

On $\Gamma\backslash G/K$, with $K$ compact and preserving the
norm, use the normalized Haar lift. Apply the proposition on each
of finitely many arithmetic components with its original weight.
A bounded, invariant finite-adelic weight of compact support is
dominated by a constant times the count on $m^{-1}L$ for some
integer $m$, so it is integrable too; Haar invariance supplies its
scaling law. The original packing estimate uses the original count.

In the cusp bound, $N\ge3j+1$ and $j\ge1$, so the dimension
condition holds. No numerical covolume or cusp count is required.
The remaining geometric inputs are the actual arithmetic model,
its lattice and norm, equality of its geometric mean with this
normalized Haar mean, and the signed boundary intersection limit.

\subsection{The geometric mean is a Haar mean}
\label{sec:curvature-haar-mean}

We now identify the mean defined by the original Hodge metric, once
its complex uniformization and local metric identifications are given.
For the matrix domain
\[
 \mathcal D_{p,q}=\{Z\in M_{p,q}(\mathbf C):I_q-Z^*Z>0\},
 \qquad \Phi(Z)=-\log\det(I_q-Z^*Z),
\]
the determinant quotient Hodge norm has local logarithm $\Phi/2$,
up to a pluriharmonic term. Indeed, on the positive plane
$\operatorname{graph}(W)$ the quotient metric has Gram matrix
$(I_q-WW^*)^{-1}$; its determinant norm is the square root of
its determinant. The transpose coordinate $Z=W^t$ gives the
displayed potential.

Put $\Phi=\sum_\sigma\Phi_\sigma$ on the product domain and
$n=\sum_\sigma p_\sigma q_\sigma$. Direct differentiation gives
\[
 (\partial\bar\partial\Phi_\sigma)_Z(T,\bar T)
 =\operatorname{tr}\bigl((I-ZZ^*)^{-1}
                       T(I-Z^*Z)^{-1}T^*\bigr).
\]
The determinant of this Hessian is
$\det(I-Z^*Z)^{-(p_\sigma+q_\sigma)}$. Therefore, with
$dd^c=(i/\pi)\partial\bar\partial$,
\[
 (dd^c(\Phi/2))^n=c_n b\,dV,\qquad
 c_n=\frac{n!}{\pi^n},\qquad
 b=\prod_\sigma\det(I-Z_\sigma^*Z_\sigma)^{-(p_\sigma+q_\sigma)}.
\]
This includes zero-dimensional factors by the empty-determinant
convention. In our equal-signature case, $b=\exp(N\Phi)$.

The fractional linear unitary action changes $\Phi$ by
$2\log|\det(CZ+D)|$. Its Levi form is unchanged, so $b\,dV$
is invariant. Since the stabilizer $K$ is compact, the pushforward
of Haar measure under $G\to G/K$ is $c\,b\,dV$ for some
$0<c<\infty$.

\begin{proposition}\label{prop:geometric-haar-mean}
Retain the original finite collection of arithmetic components of
complex dimension $n$, their measurable fundamental sets $\mathcal F_k$,
and restrictions $\gamma_m:B_m\to X$ of local holomorphic sheets, with
$B_m\subset\mathcal F_{k(m)}$. Suppose these sheets are disjoint
and cover up to chart-null sets, their multiplicity over almost
every point of $\mathcal F_k$ is $d_k\in\mathbf Z_{\ge0}$ with
$\sum_k d_k>0$, and their original Hodge
potential is $\Phi_k/2+\Re H_m$ locally. Then the original
normalized geometric mean is
\[
 \operatorname{mean}_X(g)=
 \sum_{k:d_k>0}\frac{d_kV_k}{\sum_\ell d_\ell V_\ell}
                 \int\bar g_k\,d\nu_k,\qquad
 V_k=\int_{\mathcal F_k}b_k\,dV,
\]
where $\nu_k$ is normalized arithmetic Haar measure and
$\bar g_k$ averages the $d_k$ sheet values. The assertion holds
for every absolutely integrable function $g$ on the projective variety.
\end{proposition}

\begin{proof}
Holomorphic change of variables gives the original analytic
integral as
\[
 I(g)=c_n\sum_m\int_{B_m}g(\gamma_m(t))b_{k(m)}(t)\,dV(t).
\]
Use nonnegative measures first; absolute integrability then
justifies signed sums. The arithmetic finite-volume theorem gives
$0<V_k<\infty$. Consequently
$I(1)=c_n\sum_k d_kV_k>0$, and division gives the formula.
The original denominator is defined by this same Hodge-potential
integral against $1$, so it is exactly $I(1)$.

For precision, let $s_k$ be the order of the kernel of
$\Gamma_k$ on its domain, and write
$(\pi_k)_*dg_k=c_kb_k\,dV$. Almost every point has stabilizer
of order $s_k$: every additional fixed locus is a proper analytic
subset. The raw quotient volume is
$(c_k/s_k)V_k$. Thus normalized component Haar measure is
$b_k\,dV/V_k$, while the weights retain $d_kV_k$.
No stabilizer or relative volume is suppressed.
\end{proof}

For lattice invariants the prescribed fibre isometries make all
sheet values equal. Proposition~\ref{prop:arithmetic-isotropic-integrability}
therefore gives finite mean and the same radius scaling for the
original isotropic count. The cusp bound applies with its original
normalization. The required geometric input is the actual
uniformization with these sheets and this metric; their existence
is not inferred from a probability-measure contract. The arithmetic
model and signed boundary intersection limit remain separate.

\subsection{Bounding the average cusp correction by counting}
\label{sec:cusp-counting-mean}

The averaging estimate has a short proof once the original quotient
measure and counting identity are identified. Write $\log^+x=\max(0,\log x)$.

\begin{lemma}\label{lem:cusp-counting-tail}
On a probability space let $\lambda>0$ and $N_t\ge0$ be measurable,
and fix $r>0$ and $0<a<k$. Suppose $N_r$ is integrable and, for
$0<t\le r$,
\[
 \mathbf1_{\{\lambda<t\}}\le N_t,\qquad
 N_r\le(3r/\lambda)^a\mathbf1_{\{\lambda<r\}},\qquad
 \int N_t\le(t/r)^k\int N_r.
\]
Then $\log^+(r/\lambda)$ is integrable and
\[
 \int\log^+(r/\lambda)
 \le\frac{a\log3+\log(k/(k-a))}{k-a}+\frac1k.
\]
For $a\ge2$ and $k\ge3a$, this is at most $23/24<1$.
\end{lemma}

\begin{proof}
Put $M=\int N_r$. The hypotheses give
$\Pr(\lambda<t)\le\min(1,M(t/r)^k)$. If $M\ge1$, integrating this
tail, with the split at $u=M^{a/k}$, gives
\[
 M\le3^a\int(r/\lambda)^a\mathbf1_{\{\lambda<r\}}
 \le3^a\left(1+\int_1^\infty\min(1,Mu^{-k/a})\,du\right)
 =\frac{3^ak}{k-a}M^{a/k}.
\]
Thus $M\le M_*=(3^ak/(k-a))^{k/(k-a)}$, also when $M<1$.
Tonelli's theorem now gives
\[
 \int\log^+(r/\lambda)
 \le\int_0^\infty\min(1,M_*e^{-ks})\,ds
 =\frac{\log M_*+1}{k}.
\]
This proves integrability as well as the claimed bound. For $a\ge2$,
$k\ge3a$, it is at most
$\frac12\log3+\frac14\log(3/2)+\frac16\le23/24$;
use $\log3\le4/3$ and $\log(3/2)\le1/2$.
\end{proof}

Here is how the pointwise hypotheses arise. Let $E$ be CM of degree
$d=2f$, let $L$ be an $\mathcal O_E$-stable full lattice in a Hermitian
$E$-space of Witt index at most $j$, and let $\lambda$ be its shortest
nonzero vector length in the original Hodge norm. Suppose
$b(v,w)=\operatorname{Tr}_{E/\mathbb Q}(\xi h(v,w))$ is integral on
$L$, with $\xi\ne0$, and
\[
 |b(v,w)|\le\|v\|\|w\|,\qquad
 \|\alpha v\|\le\max_\sigma|\sigma\alpha|\,\|v\|.
\]
Choose integral elements spanning $E$ over $\mathbb Q$, all of whose
conjugates have absolute value at most $v_r\ge1$, and put $r=v_r^{-1/2}$.
For $\|v\|,\|w\|<r$, the integers $b(\alpha v,w)$ have absolute
value less than one. They vanish; the spanning property and the trace
pairing imply $h(v,w)=0$. All short vectors therefore span a totally
isotropic $E$-space, of real dimension at most $2fj$. Balls of radius
$\lambda/2$ around distinct vectors are disjoint. Comparing their
volumes inside the ball of radius $t+\lambda/2$ gives, for the number
$N_t$ of nonzero isotropic vectors of norm less than $t$,
\[
 1\le N_t\le(3t/\lambda)^{2fj}\quad(\lambda<t\le r),
 \qquad N_t=0\quad(t\le\lambda).
\]
Thus the lemma applies with $a=2fj$ and $k=2f(N-1)$ when $N\ge3j+1$,
provided the original quotient probability measure makes $N_r$
integrable and satisfies $\int N_t=(t/r)^k\int N_r$.

For abelian $E/\mathbb Q$, the spanning elements admit a simple bound.
Let $c$ be its conductor. The conductor--discriminant formula gives
$c\le\operatorname{rd}(E)^2$: at each prime dividing $c$, the
characters with smaller conductor exponent form a proper subgroup,
so at least half the $d$ characters have the maximal exponent.
The integral elements
\[
 \operatorname{Tr}_{\mathbb Q(\zeta_c)/E}(\zeta_c^i),
 \qquad0\le i<\varphi(c),
\]
span $E$ over $\mathbb Q$ and have all conjugates bounded by
$v_r=\varphi(c)/d$. This is a rational spanning assertion; their
integral span may have nontrivial index in $\mathcal O_E$.
The classical conductor identities and cyclotomic inclusion are in
Milne, \emph{Class Field Theory}, v4.03,
\href{https://www.jmilne.org/math/CourseNotes/CFT.pdf}{Theorems I.4.16
and V.3.27, and Example V.3.10}.

Assume the same measurable, nonnegative cusp correction $U$ satisfies
the original pointwise comparison
\[
 U\le fj\bigl(2\log^+(1/\lambda)+\log(2fj\operatorname{rd}(E))\bigr).
\]
Since $\log^+(1/\lambda)\le\log^+(r/\lambda)+\tfrac12\log v_r$,
the lemma proves its integrability and
\[
 \int U\le fj\bigl(2+\log v_r+\log(2fj\operatorname{rd}(E))\bigr)
 \le fj\bigl(2+\log j+3\log\operatorname{rd}(E)\bigr).
\]
For $B\ge4fN^2$ and $j\ge1$, one has $B\ge64j$; hence, whenever
$\Lambda_B\ge\log B+\log\operatorname{rd}(E)$,
\[
                         \int U\le3fj\Lambda_B.
\]
This supplies the numerical allowance $C_u=3$. Application to the
arithmetic inequality still requires the original lattice and
covolume, the normalized quotient measure and its counting identity,
and the identification of this same integral with the signed cusp
term. The arithmetic model and the one-sided boundary limit are
retained hypotheses.

\subsection{Metric variation and its normalization}
\label{sec:exact-arithmetic-metric-variation}

Write \(I(L_0,\ldots,L_n)\) for an arithmetic intersection satisfying
symmetry, tensor additivity and the divisor recursion.  Use metric
denominators \(h_L\), so that \(\|s\|=|s|/h_L\); the archimedean
term in the recursion is minus the integral of \(\log\|s\|\).
Denote its original mixed integral by
\(\langle u;L_1,\ldots,L_n\rangle_\infty\), retaining every complex
component and its weight.

For two smooth metrics \(A,B\) of the same multidegree, put
\(\delta=\log h_B-\log h_A\).  Then
\begin{equation}\label{eq:metric-variation-single-slot}
 I(L_1,\ldots,L_n,B)-I(L_1,\ldots,L_n,A)
       =\langle\delta;L_1,\ldots,L_n\rangle_\infty .
\end{equation}
Indeed, the trivial line with denominator \(e^\delta\) has section \(1\),
empty divisor and logarithmic section norm \(-\delta\).
The two minus signs in the recursion give the positive sign in
\eqref{eq:metric-variation-single-slot}.

There is a short way to justify locality even when homogeneous metric
data have arbitrary values at degenerate coordinate blocks.
Let \(\mathcal C\) be the set where every block is nonzero and use
\[
                    h_C(z)=\prod_k
                       \left(\sum_i|z_{ki}|^2\right)^{1/2}.
\]
This is a smooth positive metric on \(\mathcal C\), and vanishes
outside it.  If two same-degree metrics agree on \(\mathcal C\),
their tensor products with \(C\) agree everywhere; tensor additivity
and cancellation give equal intersections.  Extend \(\delta\) by zero
outside \(\mathcal C\) and apply this observation to
\(A\otimes(\mathcal O,e^\delta)\) and \(B\), proving
\eqref{eq:metric-variation-single-slot}.
If \(A,B\) agree only on the variety, \(\delta\) vanishes at every
point of its mixed integral, so \eqref{eq:metric-variation-single-slot}
also gives locality on the variety.
Thus changing a smooth representative does not change the intersection.

Replacing the factors one at a time now gives
\begin{equation}\label{eq:metric-variation-top-intersection}
 I(B^{n+1})-I(A^{n+1})
   =\sum_{j=0}^{n}
          \langle\delta;B^{[j]},A^{[n-j]}\rangle_\infty .
\end{equation}
Here brackets indicate repeated arguments.  If the original mixed
integrals are identified on a smooth compact complex model with the
corresponding Chern integrals, write
\(h_A=h_0e^\varphi\), \(h_B=h_0e^\psi\) and
\(\theta=dd^c\log h_0\), with \(dd^c=i\partial\bar\partial/\pi\).
Equation \eqref{eq:metric-variation-top-intersection} is precisely
\begin{equation}\label{eq:metric-variation-energy-factor}
 I(B^{n+1})-I(A^{n+1})
                =(n+1)\bigl(E_\theta(\psi)-E_\theta(\varphi)\bigr),
\end{equation}
with the original weighted sum over components understood.
The energy is unnormalized; the mixed difference formula is
\cite[Proposition~2.1]{boundary-bbgz}.  There is no further factor \(1/2\):
the half in the cusp comparison is the given difference
\(\psi-\varphi=U/2\).

For integers \(p\geq0\) and \(q\geq1\), multilinearity also gives
\(q^{-(n+1)}I((L^{\otimes q}\otimes A^{\otimes p})^{n+1})
 =I((L+(p/q)A)^{n+1})\), interpreted as the usual polynomial,
whose constant term is \(I(L^{n+1})\).
This establishes the algebraic normalization.  Identifying the original
mixed integrals and common vertical curvature with the specified compact
model remains a separate geometric requirement.

\subsection{Logarithmic boundary growth and energy}
\label{sec:boundary-energy}

Let $X$ be a compact Kähler manifold of dimension $n$, let $\omega$
be Kähler, and let $D$ be a normal-crossing divisor. We use unnormalized
Monge--Ampère energy. Its monotonicity and upper semicontinuity, and the
full-mass property of finite-energy potentials, are
\cite[Propositions~2.10--2.11]{boundary-begz} and
\cite[Propositions~2.4 and~2.8]{boundary-bbgz}.

\begin{lemma}\label{lem:boundary-energy}
Let $\theta$ be smooth and closed, and suppose
$T=\theta+dd^c\varphi\geq0$, with $\varphi$ smooth off $D$,
bounded above, and
\[
 \varphi\geq-C\log r-C,\qquad
 r=O(1)+\sum_i\log|z_i|^{-2}\geq1
\]
in normal-crossing coordinates. Suppose smooth $\varphi_t$ satisfy
$\varphi\leq\varphi_t\leq C_0$, converge locally in $C^2$ off $D$,
and $\theta+dd^c\varphi_t+\varepsilon\omega\geq0$ eventually for
every $\varepsilon>0$. Then
\[
 \int_{X\setminus D}T^k\wedge\omega^{n-k}
      =\int_X\theta^k\wedge\omega^{n-k}\quad(0\leq k\leq n).
\]
All the corresponding integrals weighted by $|\varphi|$ are finite.
The original signed analytic energies also converge.
\end{lemma}

\begin{proof}
For $n=0$ the conclusion is immediate; assume $n\geq1$.
First take a Kähler reference form $\Omega$ with
$\Omega+dd^c\varphi\geq0$. Choose
$\rho=\log\|s_D\|^2\ll-1$, so that
$\log(-\rho)=\log r+O(1)$. A barrier
$v=-b\log(-\rho)-c\leq\varphi$ is $\Omega$-plurisubharmonic
after decreasing the norm of $s_D$. Indeed, if $dd^c\rho=-\eta$
off $D$, then
$dd^cv=-b\eta/(-\rho)+b\,d\rho\wedge d^c\rho/\rho^2$.
Smooth convex truncations $v_j=h_j(\rho)$ have
$0\leq h_j'\leq b/(-\rho)$ and $h_j''\geq0$.
In
$dv_j\wedge d^cv_j\wedge(\Omega+dd^cv_j)^{k-1}$ the rank-one
term containing $h_j''$ vanishes. Integration by parts therefore
bounds their energies using
\[
 \int|\log(-\rho)|\,\Omega^n<\infty,\qquad
 \int\frac{d\rho\wedge d^c\rho\wedge\Omega^{n-1}}{\rho^2}<\infty.
\]
At a crossing the second assertion follows by integrating
$ds_i/(A+\sum s_j)^2$ first, with exponential volume factors in
the other variables. Thus $v$, and by domination $\varphi$, have
finite energy.

Apply this with $\Omega=\theta+B\omega$ for large $B$.
Fatou's lemma for the positively corrected smooth approximations
bounds every mixed mass by its cohomological value. Full mass for
$T+B\omega$ makes the positive polynomial combination of these
deficits zero, so every deficit vanishes. Finite energy also gives
all the asserted weighted integrability.

For any fixed positive correction, choose a smooth positive reference
in its class. Then $\varphi_t\geq\varphi$ and $L^1$ convergence,
together with monotonicity and upper semicontinuity, give energy
convergence. Finally
\[
 Q_t(B)=\frac1{n+1}\sum_{k=0}^n
  \int_X\varphi_t
       (\theta+dd^c\varphi_t+B\omega)^k
                       \wedge(\theta+B\omega)^{n-k}
\]
is a polynomial of degree at most $n$ in $B$. Its limit is the same
expression off $D$ with $\varphi$, whose coefficients are absolutely
integrable. Convergence at $n+1$ distinct large values of $B$
implies coefficientwise convergence, including the original signed
quantity $Q_t(0)$.
\end{proof}

For the bounded correction, let $\psi$ be bounded, and take finitely
many charts $V_j$ with
$\overline I_j\subset M_j$, $\overline M_j\subset V_j$ and
$\bigcup_jI_j=X$. Let $q_j$ be smooth local potentials of $\theta$,
with $q_j+\psi$ plurisubharmonic. Use a nonnegative smooth kernel
of integral one, supported in the unit ball and invariant under
simultaneous complex rotations. Put $f_{j,t}=M_t(q_j+\psi)-q_j$.
The cutoffs satisfy $-1\leq\eta_j\leq0$ and $\eta_j=0$ on $I_j$.
Let $\epsilon(t)\geq t$ tend to zero, with
$|f_{j,t}-f_{k,t}|\leq\epsilon(t)$ on overlaps. The original gluing
satisfies $\max a_j\leq\mathcal G_\epsilon(a)\leq\max a_j+\epsilon$,
with maxima over the active regions $M_j$.
Its smooth output and curvature bound are
\[
 \psi_t=\mathcal G_{\epsilon(t)}(f_{j,t}+4\epsilon(t)\eta_j),
 \qquad \theta+dd^c\psi_t\geq-C\epsilon(t)\omega.
\]

\begin{lemma}\label{lem:bounded-cusp-energy}
Under these finite-chart regularization conditions, $\psi_t\to\psi$
pointwise and in $L^1$, and their energies converge in every
positively corrected class.
\end{lemma}

\begin{proof}
Averaging over complex rotations and using the complex-line
mean-value inequality
gives $u_j\leq M_tu_j\leq\sup_{\overline B(\cdot,t)}u_j$.
Here $u_j=q_j+\psi$, so $f_{j,t}$ tends pointwise to $\psi$.

At each point choose an inner chart, where its cutoff is zero.
The original overlap error $\epsilon(t)\to0$, the nonpositive
cutoffs, and the regularized maximum give
\[
             \psi\leq f_{j,t}\leq\psi_t
                            \leq f_{j,t}+2\epsilon(t).
\]
Compactness gives a common bound, hence $L^1$ convergence.
The curvature bound gives a positive smooth reference in every
positively corrected class. Energy is then squeezed between its value at $\psi$
and its upper semicontinuous limit. This proves convergence without
requiring $\psi$ to be continuous or $\psi_t$ to be monotone.
\end{proof}

For the original cusp correction, $\psi-\varphi=U/2$. Concavity gives
\[
 E_z(\psi)-E_z(\varphi)
       \leq\frac12\int_X U\,(T+z\omega)^n,\qquad z>0.
\]
These arguments supply the analytic boundary limits. Their application
to the original arithmetic height still requires its metric-variation
identity, the compact-model integration dictionary, and the specified
vertical and rational twists.

\subsection{The original integral and the signed boundary bound}
\label{sec:compact-model-signed-stokes}

\begin{lemma}\label{lem:compact-model-integral}
Let $Y$ be the original smooth complex projective variety, and let
$X$ be a supplied compact complex manifold of the same dimension,
whose projection from homogeneous coordinates has precisely the
coordinate-scaling orbits as fibres and admits local holomorphic
sections. Then $X\simeq Y$ biholomorphically. For an absolutely
integrable mixed curvature density, the original sum over inverse
projection branches equals its integral on $X$, with all complex
components included.

Moreover, an original vertical metric factor defined by a rational
section that is a regular unit away from one arithmetic prime has
zero curvature on $X$, including its complex boundary.
\end{lemma}

\begin{proof}
Local sections define a holomorphic bijection $X\to Y$. It is
biholomorphic by the injective holomorphic mapping theorem
\cite[Theorem~1.6.6]{compact-stokes-lebl}. Smoothness of $Y$ here is
the retained smooth-generic-fibre hypothesis.

A linear projection whose centre misses the Segre embedding of $Y$
is finite: its fibres are proper subvarieties of affine spaces,
hence finite \cite{compact-stokes-projection}. Outside its critical
values, infinity and the prescribed null boundary image, apply
change of variables to every inverse holomorphic branch. Different
homogeneous lifts change logarithmic metric potentials by
pluriharmonic functions. With $dd^c=(i/\pi)\partial\bar\partial$,
their mixed Hessian density is exactly the wedge of the curvature
forms. Applying this argument to the absolute value first justifies
the signed identity. Each projective point occurs once; any later
finite-cover multiplicities and stabilizer factors retain their
original values. In dimension zero this is a finite sum.

At every complex point the arithmetic prime is invertible. The
original unit relation supplies a local homogeneous quotient $G/H$
with $G,H$ nonzero and vertical factor $|G/H|$. Its logarithm is
locally the real part of a holomorphic logarithm, so its curvature
vanishes there. This argument applies at boundary points as well
\cite{compact-stokes-source}.
\end{proof}

Fix the original finite vertical cutoff and put $V$ for its product
of vertical factors. Let $h_0$ be the common reference denominator,
$\theta$ its curvature, and let the prescribed positive metric
$A'$ have Kähler curvature $\omega$. Suppose the tied potentials
$\varphi_t,\psi_{t'}$ satisfy the preceding boundary-energy lemmas,
with limits $\varphi,\psi$ and
$\psi-\varphi=U/2$. Write $T=\theta+dd^c\varphi\geq0$.
The actual smooth lines have restricted denominators
\[
 h_{\mu_t}=h_0e^{\varphi_t}V,\qquad
 h_{N_{p,q,t}}=h_{\mu_t}^{q}h_{A'}^p,\qquad
 h_{M_{p,q,t'}}=(h_0e^{\psi_{t'}}V)^qh_{A'}^p.
\]
Retain eventual smoothness of $\mu_t$ and, for each
$1\leq p\leq q$, of $M_{p,q,t'}$. The full intersection pairing
and all the displayed geometric data are supplied hypotheses.

\begin{proposition}[Signed arithmetic boundary comparison]
\label{prop:compact-model-signed-stokes}
Put $d(z)=\int_X(T+z\omega)^n$ and assume $d(0)>0$.
Suppose the original normalized height inequalities give
\[
 q^{-(n+1)}I(M_{p,q,t'}^{n+1})
       \geq(n+1)(a-(p/q)K_A)d(p/q)
\]
eventually for each $1\leq p\leq q$.
Then the original arithmetic limit exists and satisfies
\[
 \mathcal A:=\lim_{t\to0}I(\mu_t^{n+1})
 \geq(n+1)\left(a-\frac{\overline U}{2}\right)d(0),
 \qquad \overline U=\frac{\int_XU\,T^n}{d(0)}.
\]
\end{proposition}

\begin{proof}
Metric locality and multilinearity give one polynomial
\[
 F_t(z)=\sum_{j=0}^{n+1}\binom{n+1}{j}z^j
          I(\mu_t^{[n+1-j]},(A')^{[j]}),
 \qquad F_t(p/q)=q^{-(n+1)}I(N_{p,q,t}^{n+1}).
\]
The model, cutoff, vertical factors and metrics in this polynomial
are fixed independently of $p,q$. At each fixed rational
$z\in(0,1]$, exact metric variation and energy convergence give
a limit for $F_t(z)$. Values at $n+2$ distinct such points recover
all coefficients by interpolation. Thus $F_t\to F$ coefficientwise,
and $F(0)=\mathcal A$.

For fixed $z=p/q>0$, the bounded energy convergence gives a limit
$G(z)$ for the corresponding normalized bounded intersection.
No bounded line with $p=0$ is needed. The integral dictionary and
flatness of $V$ give the exact difference
\[
 G(z)-F(z)=(n+1)\bigl(E_z(\psi)-E_z(\varphi)\bigr).
\]
Here $E_z$ is unnormalized energy in the class of $\theta+z\omega$.
Near semipositivity supplies a smooth positive reference in that
class; $\theta$ itself need not be positive. The logarithmic metric
ratio contributes $q$ and the $n$ curvature factors contribute $q^n$,
exactly cancelling the original $q^{n+1}$.

Finite-energy concavity and $\psi-\varphi=U/2$ yield
\[
 E_z(\psi)-E_z(\varphi)
       \leq\frac12\int_XU\,(T+z\omega)^n.
\]
Together with the retained height inequality this gives
\[
 F(z)\geq(n+1)\left((a-zK_A)d(z)
                    -\frac12\int_XU\,(T+z\omega)^n\right).
\]
Full mixed mass identifies the geometric degree with $d(z)$.
Weighted mixed integrability makes the last integral a finite
polynomial in $z$. Letting positive rational $z$ tend to zero proves
the assertion with the original signed $U$. This also covers $n=0$.
\end{proof}

This proves the signed comparison for the given arithmetic and
regularization data. The construction of those data and the
arithmetic height and section estimates remain separate inputs.


\subsection{Selected tangent degrees and integral theta sections}

The all-point height estimate can be reduced further to the selected
tangent eigenspaces. The following comparison retains the exact
integral indices and separates the toric and abelian ranks. Its
theta argument supplies the tangent estimate. The determinant calculation
then constructs the integral line with its exact correction at ramified
coefficient primes. The prescribed pure-period point is algebraic
because the moduli space separates its full labelled fibres; this
preserves the exact component and original rigidified line. Properness
then identifies the integral CM model. The actual arithmetic model,
its original boundary-fibre data and pointwise period presentations
now supply the flat-degree theorem below. The counting argument above bounds the
cusp mean once its original quotient and lattice identities are supplied;
the signed boundary comparison is supplied by
Proposition~\ref{prop:compact-model-signed-stokes} under its stated
model and height hypotheses. The actual arithmetic section remainder
is supplied below by Theorem~\ref{thm:actual-component-section-degree};
its mean height and rank losses remain to be bounded.

\paragraph{The selected tangent degrees.}
All degrees in this paragraph are ordinary Arakelov degrees divided by
the degree of the number field. Let $\bar V$ have rank $n$ and degree
$d$, with pairwise orthogonal, generically complementary integral
sublattices $W_\tau,W'_\tau$. Let $V_\tau\subset W_\tau$ have the same
generic fibre and the inherited metric. Put
\[
 X=\sum_\tau\widehat{\deg}_{\rm n}W_\tau,\quad
 Y=\sum_\tau\widehat{\deg}_{\rm n}W'_\tau,\quad
 T=\sum_\tau\widehat{\deg}_{\rm n}V_\tau,\quad
 R=\sum_\tau\operatorname{rk}V_\tau .
\]
The two integral indices give
\begin{equation}\label{geo-selected-index-identity}
 \delta=d-X-Y\ge0,\qquad \nu=X-T\ge0,\qquad
 b=X-Y,\qquad d=2T-b+\delta+2\nu .
\end{equation}
There is no archimedean correction, since all summands are orthogonal
and retain the given metrics.

\begin{proposition}\label{geo-selected-tangent-sum}
If $C_s\ge0$ and
\[
 T\le RC_s\log(2+n+(-d)_+),
\]
then
\[
 T\le RC_s\log\bigl(2+n+(b-\delta-2\nu)_+\bigr).
\]
If also $\Lambda\ge1$, $C_f\ge0$, $n\le G\le e^\Lambda$ and
$b\le C_fG\Lambda$, then
\begin{equation}\label{geo-selected-log-budget}
 T\le C_A R\Lambda,\qquad C_A=C_s\{1+\log(3+C_f)\}.
\end{equation}
\end{proposition}
\begin{proof}
For $T\le0$ the first conclusion is immediate. For $T>0$,
\eqref{geo-selected-index-identity} gives
$(-d)_+=(b-\delta-2\nu-2T)_+\le(b-\delta-2\nu)_+$.
For the second conclusion, bound the logarithm's argument by
$e^\Lambda(3+C_f\Lambda)$ and use concavity:
\[
 \log(3+C_f\Lambda)
 \le\log(3+C_f)+\frac{C_f}{3+C_f}(\Lambda-1)
 \le\Lambda\log(3+C_f).
\]
The proof includes $R=0$ and $C_s=0$.
\end{proof}

Only the signed total $b$ occurs; both integral defects are favourable.
If the selected eigenspace has dimension $m$, an adapted family has
$r$ toric vectors and $m-r$ abelian vectors, and there are $f$ conjugate
pairs, then the exact selected rank is $R=f(m-r)$. Thus a comparison
\[
 \Phi\le fmX_0+T-fr\log\operatorname{rd}(E)
                         +\frac{fr}{2}\log(4\pi)
\]
together with $X_0\le C_X\Lambda$, $\log\operatorname{rd}(E)\ge0$ and
$\log(4\pi)\le\Lambda$ gives
\[
 \Phi\le\{C_X+\max(C_A,\tfrac12)\}fm\Lambda .
\]
In the purely toric case the selected abelian sum is zero.

\paragraph{A short integral jet argument.}
Write $\operatorname{av}_K$ for the average over all complex embeddings
of $K$, counting conjugate embeddings separately. The next lemma
isolates the elementary determinant argument.

\begin{lemma}\label{geo-leading-jet-determinant}
Let $J_0,\ldots,J_k$ be independent $K$-linear forms on $K^r$ and
$W_i=\bigcap_{j<i}\ker J_j$. Suppose
$J_i(W_i\cap\mathcal O_K^r)\subset\mathcal O_K$ and, for every embedding,
\[
 \|J_i|_{W_i\otimes_\sigma\mathbb C}\|_{\rm op}\le A_{i,\sigma},
 \qquad A_{i,\sigma}>0 .
\]
Then $0\le\sum_{i=0}^k\operatorname{av}_K\log A_{i,\sigma}$.
\end{lemma}
\begin{proof}
At each finite prime, choose an integral basis adapted to the saturated
kernel flag. The row matrix is triangular in its first $k+1$ columns,
zero in the remaining columns, and has integral diagonal entries.
Every coordinate of $J_0\wedge\cdots\wedge J_k$ is therefore integral.
For a nonzero such coordinate $c$, row Gram--Schmidt gives
\[
 |\sigma(c)|\le\|J_0\wedge\cdots\wedge J_k\|_\sigma
       =\prod_{i=0}^k\|J_i|_{W_i\otimes_\sigma\mathbb C}\|_{\rm op}
       \le\prod_{i=0}^kA_{i,\sigma}.
\]
Now use $|N_{K/\mathbb Q}(c)|\ge1$.
\end{proof}

The required integrality has a direct geometric interpretation.
Let $\mathcal G/\mathcal O_K$ be a group scheme, let $\mathcal L$
be a line bundle rigidified at its identity, and let $x$ be an integral
invariant tangent vector. For integral sections of $\mathcal L^{\otimes n}$,
take successive directional derivatives in a generic local frame
with that prescribed value at the identity. At a finite prime choose
an integral frame with the same value. Derivatives of its regular
representing functions are integral. When the earlier derivatives
vanish, the Leibniz rule shows that the leading derivative is unchanged
by the transition function, whose value is one. It is consequently
integral at every finite prime. These are undivided derivatives:
no assertion about division by $i!$ is needed. The same argument
identifies the leading derivative in a holomorphic frame with the
same rigidification.

\begin{proposition}\label{geo-theta-jet-bound}
Suppose the generic fibre is an abelian variety $A/K$ of dimension
$g\ge1$, $L$ is ample, and $x\ne0$ is an integral tangent vector.
Let $n=3k$, $k\ge1$. Suppose $r$ integral sections of $L^{\otimes n}$
span its generic section space. Write $\ell_\sigma=\|x\|_\sigma>0$.
Assume their functions in the rigidified theta frame satisfy
\[
 |f_c(zx)|\le\sqrt r\,M_\sigma\|c\|_2
                    \exp(n\ell_\sigma^2|z|^2/2),
 \qquad M_\sigma>0,\quad c\in\mathbb C^r .
\]
For $U=\operatorname{av}_K\log M_\sigma$ one has
\begin{equation}\label{geo-exact-jet-cost}
 -\operatorname{av}_K\log\ell_\sigma
       \le\frac12\log(enk)+\frac{\log r+2U}{k}.
\end{equation}
\end{proposition}
\begin{proof}
The third power of an ample line bundle on an abelian variety is
very ample \cite[Theorem~2.2.3]{genestier-ngo-shimura}.
Choose $u_0,u_1\in H^0(A,L^{\otimes3})$ with $u_0(0)\ne0$,
$u_1(0)=0$, and $D_x(u_1/u_0)(0)\ne0$. Then
$u_0^{k-i}u_1^i$ has exactly order $i$ along $\exp(zx)$,
for every $0\le i\le k$. The supplied spanning sections therefore
have independent jet rows through order $k$. Only this generic
rank argument uses $u_0,u_1$; their integrality is unnecessary.

For $1\le i\le k$, apply Cauchy's formula at radius
$\sqrt{i/n}/\ell_\sigma$ and use $i!\le i^i$. The operator bound
on the $i$-th leading jet is
\[
 A_{i,\sigma}=\sqrt r\,M_\sigma
                    (\sqrt{enk}\,\ell_\sigma)^i .
\]
The same formula holds for $i=0$. Leading-jet integrality and
Lemma~\ref{geo-leading-jet-determinant} give
\[
 0\le(k+1)(\tfrac12\log r+U)
       +\frac{k(k+1)}2
          \{\tfrac12\log(enk)+\operatorname{av}_K\log\ell_\sigma\}.
\]
Division proves \eqref{geo-exact-jet-cost}.
\end{proof}

The norm convention is essential: for the prescribed Hodge metric,
$\ell_\sigma^2=\pi H_\sigma(x,x)$ and the theta weight is
$\exp(-\pi nH_\sigma(w,w)/2)$. These give exactly the Gaussian
bound above. Also, $h^0(A,L^{\otimes n})=n^g h^0(A,L)$.
A spanning family of that cardinality contributes
$g\log n+\log h^0(A,L)$ to $\log r$; the principal case has
$h^0(A,L)=1$.

\paragraph{An economical theta power and the arithmetic input.}
Suppose $L$ gives a principal polarization. Assume that for every
integer $m\ge1$, after a finite number field extension, the prescribed
construction gives $m^{2g}$ integral sections spanning $L^{\otimes m^2}$
with the same bounds $M_\sigma$, pulled back under the embeddings.
Assume the integral tangent lattice and its metric pull back
isometrically, and that for fixed $A_0,B_0\ge0$,
\[
 U\le h/2+A_0g+B_0g\log Y,\qquad Y=2+g+h,\qquad h\ge0.
\]
This includes the uniform finite-prime integrality of those sections;
it is not implied by their complex translation identities.

Put
\[
 Z=1+g+h+2A_0g+2B_0g\log Y,\qquad
 a=\lceil\sqrt Z\rceil,\quad m=3a,\quad k=3a^2,\quad n=9a^2 .
\]
Thus $n=m^2=3k$, $3Z\le k\le12Z$ and $n\le36Z$.
The tensor exponent is $n$; the torsion order is $m$.
With $r=n^g$, Proposition~\ref{geo-theta-jet-bound} yields
\[
 -\operatorname{av}_K\log\ell_\sigma
 \le\log k+\tfrac12\log(3e)+\frac{g\log(3k)+2U}{k}
 \le\frac43\log(12Z)+\frac56\log3+\frac12 .
\]
For the last inequality, use $\log(3k)>1$ and
$g+2U_+\le Z$. If $C_0=1+2A_0+2B_0$, then
$Z\le C_0Y\log Y$ and $\log Y\ge1$. Consequently
\begin{equation}\label{geo-logarithmic-tangent-line}
 -\operatorname{av}_K\log\|x\|
 \le C_{\rm line}\log(2+g+h),\qquad
 C_{\rm line}=\frac83+\frac43\log(12C_0)+\frac56\log3+\frac12 .
\end{equation}

The passage to degrees is supplied by the classical absolute Siegel
lemma. For a nonzero Hermitian sublattice $F$ of rank $r$, let $h_F$
be its logarithmic adelic height, using the given archimedean norms
and the finite-place unit balls $F_v$. Gaudron--R\'emond give
\[
 \inf_{0\ne y\in F_K\otimes_K\overline K}h_F(y)
       \le-\widehat\mu(F)+c_r,\qquad
 c_r=\frac12\sum_{j=2}^r\frac1j\le\frac12\log r .
\]
These are the same ordinary normalized degree and norm conventions
\cite[Sections~2.1--2.3 and Lemma~3.2]{gaudron-remond-minima}.

For completeness, the algebraic vector can be made integral without
changing its height. Choose $y$ within $\epsilon>0$ of the displayed
infimum, over a finite extension $K'$. Its saturated integral line
is $P=Iy$ for a fractional ideal $I$ of $\mathcal O_{K'}$.
Choose $a\ge1$ with $I^a=(b)$ and adjoin $\alpha$ with $\alpha^a=b$.
Fractional ideals have no torsion, so $I\mathcal O_L=(\alpha)$ in
$L=K'(\alpha)$. The vector $x=\alpha y$ is an integral generator
of $P\otimes\mathcal O_L$. The quotient by $P$ is projective,
so saturation survives base change and every finite gauge norm
of $x$ is one. The product formula now gives exactly
\[
 \operatorname{av}_{L}\log\|x\|=h_F(x)=h_F(y),\qquad
 \widehat\mu(F)-c_r-\epsilon
                  \le-\operatorname{av}_{L}\log\|x\| .
\]
No metric has changed.

Suppose the theta construction and \eqref{geo-logarithmic-tangent-line}
apply after every finite extension, with unchanged constants, rank,
and $h=(-d)_+$. Apply that line bound to $x$ and let
$\epsilon\downarrow0$. Since $r\le g$,
\[
 \widehat\mu(F)
 \le C_{\rm line}\log(2+g+(-d)_+)+c_r
 \le(C_{\rm line}+\tfrac12)\log(2+g+(-d)_+).
\]
The extension may depend on $\epsilon$; no bound on its degree is
needed. Multiplying by the actual selected ranks and summing supplies
Proposition~\ref{geo-selected-tangent-sum} with
$C_s=C_{\rm line}+\tfrac12$.
If the ambient abelian tangent rank is zero, the selected sum is zero
and none of these theta or Siegel inputs is needed.

The arithmetic inputs isolated by this argument are the primitive theta
section with its mean norm bound and integral normalized torsion
translates with the same supremum bounds at the chosen levels,
compatible with the prescribed field extensions and metric. The following
classical construction supplies them at positive even levels, which
suffice for the parameter choice.
At a degeneration the final valuation inequality is the elementary
identity
\[
 Q(\chi)+\frac{b_Q(y,\chi)}N+\frac{Q(y)}{N^2}
             =\frac{Q(N\chi+y)}{N^2}\ge0
\]
for a nonnegative quadratic form $Q$ and $N\ge1$. Applying it to
theta values still requires the coefficient valuation bound and the
identification of the prescribed rigidified torsion frame.

\paragraph{Classical theta sections supply the tangent estimate.}
The arithmetic inputs to the preceding jet argument can be supplied on
the prescribed lattice. The distinction between an even theta level and
an arbitrary positive level is retained in the parameter choice below.

\begin{theorem}\label{geo-classical-tangent-slope}
Let $G/\mathcal O_K$ be a smooth separated group scheme of finite type
with geometrically connected semiabelian fibres and abelian generic
fibre $A$ of dimension $g\ge1$. Let $\mathcal L_0$ be a symmetric
cubical line bundle, rigidified at zero, whose generic fibre $L$
defines a principal polarization. Give it the canonical metric and
give $V=\operatorname{Lie}(G)$ the metric
$\|v\|_\sigma^2=\pi H_\sigma(v,v)$, where $H_\sigma$ is the
Appell--Humbert form of $L$. With ordinary normalized Arakelov degree
$d=\widehat{\deg}_{\rm n}V$, one has
\begin{equation}\label{geo-classical-tangent-twelve}
 \widehat\mu_{\max}(V)
       \le12\log\bigl(2+g+(-d)_+\bigr).
\end{equation}
Here maximal slope is the supremum of the mean degrees of nonzero
integral submodules with inherited metrics. The bound has the same
constant after every finite extension of $K$.
\end{theorem}
\begin{proof}
\emph{The same integral lattice and line.}
The semistable comparison theorem identifies $G$ with the identity
component of the N\'eron model of $A$, and preserves that identity
component under finite extension
\cite[Theorem~4.4 and Corollary~4.5]{conrad-semistable}.
Thus the identity component of a Moret--Bailly model over an extension
$K'$ is $G_{\mathcal O_{K'}}$; its tangent module is precisely
$V\otimes\mathcal O_{K'}$.

The Moret--Bailly extension of the specified $L$ restricts to
$\mathcal L_0$ on this component. Indeed the divisor of the generic
rigidified isomorphism is vertical. Locally over a discrete valuation
ring, the total space is regular and its special fibre is integral,
so that divisor is a multiple of the special fibre. Pullback at zero
shows that its multiplicity is zero. Normality extends the isomorphism,
and generic density makes it unique. Its canonical metrics agree:
their logarithmic ratio is harmonic on the compact complex torus and
vanishes at the zero rigidification. No integral index or metric
change is introduced.

\emph{One integral theta section.}
In symplectic coordinates with period matrix $\tau$, the squared norm
on the dual tangent determinant is
\[
 \det(\operatorname{Im}\tau)/\pi^g
  =\frac{i^{g^2}}{(2\pi)^g}\int_A
        dz_1\wedge\cdots\wedge dz_g\wedge
        d\bar z_1\wedge\cdots\wedge d\bar z_g.
\]
At finite places this is the invariant Hodge line of the same model.
Consequently the Faltings height in this convention is $h_F(A)=-d$.
Write $h=(-d)_+$ and $Y=2+g+h$.

For a Moret--Bailly model $\pi:\mathcal A\to\operatorname{Spec}
\mathcal O_{K'}$, the rank-one module $\pi_*\mathcal L$, with the
$L^2$ norm for Haar probability measure, has degree
$-h_F(A)/2-g\log(2\pi)/4$. It is preserved isometrically under
change of model and finite extension
\cite[Definition~2.1 and Theorem~3.4(i),(iii)--(v)]{pazuki-theta-faltings}.
Principalize this rank-one module by the ideal argument used above
for the absolute Siegel lemma. An integral generator $\theta$ then
satisfies
\begin{equation}\label{geo-theta-ltwo-supplier}
 \operatorname{av}\log\|\theta\|_2
          \le h/2+g\log(2\pi)/4.
\end{equation}
Restriction gives an integral section of the prescribed line bundle.

For every positive even integer $m$, extend the field to contain
$A[m^2]$ and use a model containing these torsion sections. The normalized
theta-group maps $[m]^*L\simeq L^{m^2}$ and
$t_a^*L^{m^2}\simeq L^{m^2}$ extend integrally and isometrically
\cite[Propositions~2.5 and~5.1]{pazuki-theta-faltings}.
Representatives $a\in A[m^2]/A[m]$ produce $m^{2g}$ integral
sections $s_a$ spanning $H^0(A,L^{m^2})$, with
\[
 \|s_a(z)\|=\|\theta([m](z+a))\|,
 \qquad \|s_a\|_{\infty,\sigma}=\|\theta\|_{\infty,\sigma}.
\]
The first section remains integral on each new model by model
independence. The integral spanning assertion is only generic spanning;
an integral basis of the whole section module is unnecessary.

\emph{A uniform supremum bound.}
Let $\rho_\sigma$ be the shortest nonzero period in the metric
$H_\sigma$, and put $r_\sigma=\min(1,\rho_\sigma/2)$.
The principal period lattice has real $H_\sigma$-covolume one.
In a Gaussian holomorphic frame centred at any point, the squared
norm is $|f(w)|^2e^{-\pi H_\sigma(w,w)}$. The ball of radius
$r_\sigma$ embeds in the torus. Its volume in Haar probability units
is $\pi^g r_\sigma^{2g}/g!$, so the holomorphic mean-value inequality gives
\[
 \frac{\|\theta\|_{\infty,\sigma}^2}{\|\theta\|_{2,\sigma}^2}
       \le\frac{g!e^{\pi r_\sigma^2}}{\pi^g r_\sigma^{2g}}.
\]
Autissier's matrix lemma uses the height
$h_{\rm Fa}=h_F-g\log\pi/2$. Its corollary at $\varepsilon=1/2$
gives exactly
\cite[Corollary~1.4 and the height convention in Section~2]{autissier-matrix}
\[
 \operatorname{av}\rho_\sigma^{-2}
 \le\frac{12}{\pi}h_F+\frac{6g}{\pi}\log(4\pi)
 \le4h+6g,\qquad
 \operatorname{av}r_\sigma^{-2}\le24Y.
\]
Take logarithms, average, and use Jensen's inequality together with
$\log(g!)\le g\log g$. Combining with
\eqref{geo-theta-ltwo-supplier} yields, for
$M_\sigma=\|\theta\|_{\infty,\sigma}$,
\begin{equation}\label{geo-theta-sup-supplier}
 U:=\operatorname{av}\log M_\sigma
                 \le h/2+5g+g\log Y.
\end{equation}
For the constant, it suffices that
$2+\tfrac12\log24+\tfrac14\log(2\pi)<5$.
All these averages are unchanged under finite extension.

\emph{Even level and the slope.}
Set
\[
 Z=1+11g+h+2g\log Y,\quad a=\lceil\sqrt Z\rceil,
 \quad m=6a,\quad k=12a^2,\quad n=m^2=3k.
\]
Then $12Z\le k\le48Z$ and $Z\le13Y\log Y$. The level $m$
is even, so the preceding integral sections apply. In
\eqref{geo-exact-jet-cost} their number is $n^g$, giving every
nonzero integral tangent vector the bound
\[
 -\operatorname{av}\log\|x\|
 \le\log k+\tfrac12\log(3e)+\frac{g\log(3k)+2U}{k}
 \le\frac{13}{12}\log(48Z)+\frac7{12}\log3+\frac12.
\]
Here $g+2U_+\le Z$, $Z/k\le1/12$, and $\log(3k)>1$.
Thus the line bound has coefficient
\[
 C_{\rm line}=\frac{13}{6}+\frac{13}{12}\log624
                              +\frac7{12}\log3+\frac12
\]
in front of $\log Y$. Apply the already established absolute Siegel
comparison, after an isometric finite extension, to each nonzero
integral submodule. Its additional coefficient is at most $1/2$.
Since $\log624<7$ and $\log3<2$,
$C_{\rm line}+1/2<143/12<12$, proving
\eqref{geo-classical-tangent-twelve}.
\end{proof}

For selected submodules of total rank $R$, the theorem supplies
$T\le12R\log(2+g+(-d)_+)$ on these same lattices. In particular
Proposition~\ref{geo-selected-tangent-sum} applies with $C_s=12$.
If the abelian rank is zero, the selected degree sum is zero without
using theta sections. The flat determinant, point-height, cusp and
finite section-degree comparisons are separate geometric requirements.

\paragraph{The exact conjugate determinant comparison.}
The passage from local flat norms to the degrees of the prescribed
eigenlattices has a short direct proof. Let $K$ be a number field and
$\bar V$ an integral Hermitian lattice with an action of $\mathcal O_E$.
For two conjugate embeddings into $K$, let $W,W'$ be the integral
eigenlattices, of ranks $a,b$, with their inherited metrics. Put
$\mathcal D=\det W\otimes(\det W')^{-1}$.

\begin{lemma}\label{geo-flat-determinant-bridge}
After one common finite extension if necessary, suppose a scalar
$\varphi(x,y)\in K^\times$ attached to rational eigenframes commutes
with field embeddings and field extension, and satisfies
\[
 \varphi(Ax,By)=\varphi(x,y)\frac{\det B}{\det A}.
\]
At complex embeddings suppose its specified logarithmic norm is
\[
 \log|\varphi(x,y)|+
 \tfrac12\log\det\operatorname{Gram}(x)-
 \tfrac12\log\det\operatorname{Gram}(y).
\]
At every finite place suppose it is $\log|\varphi(x_0,y_0)|$
on integral local frames of these same eigenlattices. Then its normalized
adelic degree is
$\widehat{\deg}_{\rm n}W-\widehat{\deg}_{\rm n}W'$.
The statement remains valid after a common finite field extension,
including at ramified places.
\end{lemma}
\begin{proof}
The section
$s=\varphi(x,y)(\bigwedge x)\otimes(\bigwedge y)^{-1}$ of
$\mathcal D_K$ is independent of the frames. Its complex norm is the
displayed Gram expression. At a finite place, writing $x=Ax_0$ and
$y=By_0$ gives
\[
 \log\|s\|=\log|\varphi(x,y)|+\log|\det A|-\log|\det B|
                         =\log|\varphi(x_0,y_0)|.
\]
Thus the supplied norms are exactly those of $\mathcal D$. Sum their
negatives over all embeddings and finite places, and divide by
$[K:\mathbb Q]$. The product formula removes any change of rational
section. For integral base change, choose an integral basis $e_i$ of
$\mathcal O_E$: an eigenlattice is the kernel of
$v\mapsto((e_i-\tau(e_i))v)_i$. Flat extension preserves this kernel,
so no denominator or new index occurs. Normalized degree is unchanged
because the embeddings occur with their extension multiplicities.
For rank zero the determinant is the trivial line.
\end{proof}

The geometric input is a flat family with these local norms and their
compatibility with the prescribed complex multiplication points.
Proposition~\ref{geo-cm-flat-degree} below computes its value from one
elliptic curve under precisely that compatibility. The finite-place
normalization must be retained. Write
\[
 P_E=\gamma+\log(2\pi)-\log\operatorname{rd}(E)
                      -\frac{L'(1,\chi_{-d})}{L(1,\chi_{-d})}.
\]
Suppose the corrected flat degree is
$(a-b)P_E/2+b(\log d-\log\operatorname{rd}(E))$ and the actual
integral-frame logarithmic norms add
$b(\log|\mathfrak D_E|-\log|d|)$ at finite places. Their normalized
sum is $b(\log d-\log\operatorname{rd}(E))$, by the ideal norm
of the different and the product formula. Degree is the negative
sum of logarithmic norms, so this term cancels exactly. The lemma gives
\begin{equation}\label{geo-exact-conjugate-degree}
 \widehat{\deg}_{\rm n}W-\widehat{\deg}_{\rm n}W'
                         =\frac{a-b}{2}P_E.
\end{equation}
These are hypotheses about the same semistable model and its prescribed
metric. They do not follow from a scalar estimate for $P_E$.

\begin{corollary}\label{geo-signed-signature-budget}
In the notation of \eqref{geo-selected-index-identity}, assume the
flat-family and period inputs above for each conjugate pair with the
same $P_E$. Assume also the tangent-slope estimate
$T\le12R\log(2+n+(-d)_+)$ for these same lattices, with $d$ their
total degree. Then
\[
 b=\frac{2R-n}{2}P_E,\qquad
 T\le12R\log\left(2+n+
       \left(\frac{2R-n}{2}P_E-\delta-2\nu\right)_+\right).
\]
For constant signature $(p,q)$, toric rank $r$, and a selected half
of size $m\in\{p,q\}$, one has
$2R-n=f(2m-p-q)$, independently of $r$. Equal signatures remove
the period cost. Selecting the smaller half also removes it when
$P_E\ge0$.
\end{corollary}
\begin{proof}
Sum \eqref{geo-exact-conjugate-degree}; the complementary ranks sum
to $n-R$. Apply the selected-degree argument with the tangent-slope
constant $12$. Finally use $R=f(m-r)$ and $n=f(p+q-2r)$.
\end{proof}

\subsection{Determinant lines and their integral norms}
\label{shared-integral-determinants}

The prescribed level gives determinant monodromy one, and a direct period
argument extends its generator in the actual boundary tangent frames.
The finite norm must use the same integral lattice as the complex norm.
First-order sections of the biextension
identify the coefficient module and construct both of its filtrations.
A universal vector extension and a short toric filtration control the
ramified coefficient action. Frobenius then forces determinant torsion, which makes
the finite order constant on a connected fibre. At a complex
multiplication point, relative trace duality computes the exact
eigenframe scaling. One elliptic curve and the Kronecker limit formula
then evaluate the signed degree. The final assembly retains the exact
integral determinant line and the CM maps to the selected component.
For the finite-place degree argument, compatible formal charts on test
curves suffice; no additional global family on those curves is required.

\paragraph{The level already kills determinant monodromy.}
There is a short direct argument at the prescribed level.
Let $E$ be a CM field, let $\Lambda$ be a finite projective
$\mathcal O_E$-module, and let $B$ be a rationally nondegenerate
integral alternating form with
$B(ax,y)=B(x,\bar a y)$.

\begin{proposition}\label{geo-half-level-determinant}
Let $I\subset\mathcal O_E$ and $M\ge3$ satisfy
$I\bar I=M\mathcal O_E$ and $I+\bar I=\mathcal O_E$.
If an $\mathcal O_E$-linear automorphism $g$ of $\Lambda$
preserves $B$ and
\[
                 (g-1)(\bar I\Lambda)\subset M\Lambda,
\]
then $\det_E(g)=1$.
\end{proposition}
\begin{proof}
Trace duality defines a nondegenerate skew-Hermitian form by
$\operatorname{Tr}_{E/\mathbb Q}(a h(x,y))=B(ax,y)$.
It is linear in its first argument and is preserved by $g$.
Hence $\delta=\det_E(g)$ satisfies
$\delta\bar\delta=1$. Local bases of the projective lattice give
$\delta\in\mathcal O_E^\times$.

Localizing at each finite prime and cancelling the invertible
ideal $\bar I$ gives $(g-1)\Lambda\subset I\Lambda$, whence
$\delta-1\in I$. Also
$\delta-1=-\delta(\bar\delta-1)\in\bar I$. Therefore
\[
                  \delta-1\in I\cap\bar I
                            =M\mathcal O_E.
\]
If $\delta\ne1$, the norm of $\delta-1$ has absolute value at least
$M^{[E:\mathbb Q]}$. At every complex embedding
$|\sigma(\delta)|=1$, so this absolute value is at most
$2^{[E:\mathbb Q]}$, a contradiction. Rank zero is immediate.
\end{proof}

This proves that the corresponding determinant local system on
the abelian locus has trivial monodromy, without a further level
cover. The lower threshold is sharp: for
$E=\mathbb Q(\sqrt{-7})$, $\omega=(1+\sqrt{-7})/2$,
$I=(\omega)$ and $M=2$, multiplication by $-1$ on
$\mathcal O_E$ satisfies the stated conditions and has determinant
$-1$. One may take
$B(x,y)=\operatorname{Tr}_{E/\mathbb Q}(x\bar y/\sqrt{-7})$.

The next argument verifies extension in the actual tangent bundle.
Trivial monodromy on the open set alone permits a boundary twist;
controlling both the section and its reciprocal excludes it.

\paragraph{The determinant extends in the prescribed tangent frames.}
The analytic boundary estimate can be proved directly from periods
of the actual semiabelian family. The following local statement uses
the original polarization and its prescribed tangent metric.

\begin{lemma}\label{geo-boundary-gram}
Let a smooth family of semiabelian varieties over a complex
polydisc be the pullback of an algebraic family with a line bundle
whose restriction off $z_1\cdots z_a=0$ is a principal polarization
on an abelian scheme. Trivialize its actual Lie bundle, of rank $g$.
For every holomorphic row frame $F$ of rank $k$ at zero, its Hodge
Gram determinant satisfies, after shrinking,
\[
 c_F\left(1+\sum_{j=1}^a|\log|z_j||\right)^{-k}
              \le \det\Gamma_F(z)\le C_F
\]
on the punctured polydisc, for positive constants $c_F,C_F$.
The metric is the prescribed $\pi$ times the Riemann Hermitian form.
\end{lemma}
\begin{proof}
Write the central fibre as an extension of an abelian variety
$A$ of dimension $b$ by a torus of dimension $t$, where $g=b+t$.
Its period lattice $\Gamma_0$ has an exact sequence
\[
 0\longrightarrow2\pi i\mathbb Z^t\longrightarrow\Gamma_0
                \longrightarrow H_1(A,\mathbb Z)\longrightarrow0.
\]
It spans the central Lie algebra over $\mathbb C$.
Every algebraic line on this semiabelian variety pulls back from
$A$: its restriction to the generic split torus has trivial
Picard group, so a rational section can be chosen with vertical
divisor; smoothness and integral fibres make each such prime
divisor the multiplicity-one pullback of a divisor of $A$.
Thus its Chern pairing kills the toric periods and induces an
integral alternating form of type $(1,1)$ on $H_1(A,\mathbb Z)$.
For the latter assertion, average the Chern curvature over
translations of the compact complex torus.

Such a form, even if degenerate or indefinite, admits a rational
isotropic half-dimensional real subspace transverse to its
imaginary translate. Indeed its rational kernel is stable under
the complex structure. On the nondegenerate quotient, a Hermitian
diagonal basis gives a real Lagrangian with this transversality;
rational Lagrangians are dense, since their Darboux charts are
symmetric matrices. Transversality is open. Choose a rational
transverse half-subspace in the kernel as well, and lift and
adjoin the two choices. Intersecting with the integral lattice
makes the choice primitive. Lift it to $\Gamma_0$ and adjoin
the toric periods. This gives a primitive isotropic rank-$g$
subgroup whose periods form a complex basis.

The family exponential is holomorphic: integrate the invariant
vector fields near zero and extend using multiplication by an
integer. Its derivative at a period is invertible, so the implicit
function theorem continues the finitely many chosen periods
holomorphically. Their alternating pairings remain constant,
being continuous integer-valued signed Chern numbers. Their matrix $A(z)$
is invertible at zero.

These periods remain primitive in every nearby abelian period
lattice. For a continued central period $S_\gamma$ and $n\ge1$,
the section $\exp(S_\gamma/n)$ is $n$-torsion on the entire
polydisc. Multiplication by $n$ is étale in characteristic zero.
Its separated étale kernel makes equality of this section with
zero an open and closed condition. Thus divisibility is detected
at the central fibre, for every $n$ on the same connected
polydisc.

Complete the continued periods to an integral symplectic basis
on the universal cover of the punctured polydisc, with
$B(A_i,B_j)=-\delta_{ij}$ in the linear-first convention
$B=\operatorname{Im}H$. Write the complementary period rows as
$\Omega A$. With
$z_j=e^{2\pi i w_j}$, the normalized matrix $\Omega$ is symmetric
with positive imaginary part and
\[
 \Omega(w+e_j)=\Omega(w)+N_j,\qquad
                    N_j\in\operatorname{Sym}_g(\mathbb Z).
\]
For $v\in\mathbb Z^g$, put $f=v^t\Omega v$ and
$c=v^tN_jv$. The bounded single-valued function $e^{2\pi i f}$
extends across $z_j=0$. Its zero order equals its winding
number $c$, hence $c\ge0$. Dividing by $z_j^c$ and taking a
local logarithm extends $f-cw_j$. Integer vectors and their
pairwise sums recover every entry. First extend with parameters
across the smooth parts of the coordinate divisors, then use
Hartogs extension across their intersections. We obtain
\[
       N_j\ge0,\qquad
       \Omega(w)=\sum_j w_jN_j+\Omega_0(z),
\]
with $\Omega_0$ holomorphic on the full polydisc.

Put $N=\sum_jN_j$ and $K_0=\ker N$. For every nonzero
$v\in K_0$, the function $v^t\operatorname{Im}\Omega_0v$ is positive
off the boundary and harmonic everywhere. The minimum principle
makes it strictly positive at zero. Thus
$\operatorname{Im}\Omega_0(0)|_{K_0}$ is positive definite.
On $K_0^\perp$, $N$ is positive definite, so a Schur-complement
estimate makes $\operatorname{Im}\Omega_0(0)+TN$ positive
definite for sufficiently large $T$. Shrink until
$-\log|z_j|/(2\pi)\ge T$ and retain positivity by continuity.
For $Y=\operatorname{Im}\Omega$ this gives
\[
       c\,1_g\le Y\le C D(z)\,1_g,\qquad
                  D(z)=1+\sum_j|\log|z_j||.
\]
In row-vector notation the Riemann relations give
$AHA^*=Y^{-1}$ for the actual tangent Hermitian form $H$.
Consequently
\[
             \Gamma_F=\pi(FA^{-1})Y^{-1}(FA^{-1})^*.
\]
The nonzero singular values of $FA^{-1}$ are uniformly bounded
above and away from zero. The asserted determinant bounds
follow. Empty frames have determinant one.
\end{proof}

\begin{theorem}\label{geo-flat-family-generator}
Let $K$ be a number field, $E$ a CM field with
$\tau:E\to K$, and $\bar\tau=\tau\circ\mathrm{conj}$.
Let $Y/K$ be a normal geometrically integral proper variety and
$U\subset Y$ a dense smooth open. Suppose a smooth separated
commutative group scheme $G/Y$ has semiabelian geometric fibres
of constant dimension, is abelian over $U$, and has an original
line bundle inducing a principal polarization there.
Suppose its $\mathcal O_E$-action has the polarization
adjunction, and its integral period lattice and level satisfy
Proposition~\ref{geo-half-level-determinant}.
For the actual bundle $\mathcal V=\operatorname{Lie}(G/Y)$, put
\[
              \mathcal D_\tau=
              \det\mathcal V_\tau\otimes
                    (\det\mathcal V_{\bar\tau})^{-1}.
\]
Then $\mathcal D_\tau$ has a nowhere-vanishing algebraic
section, unique up to $K^\times$. Its prescribed Hodge norm is
constant on each complex abelian locus.
\end{theorem}
\begin{proof}
The level proposition makes the determinant of every period
monodromy operator one. On the $\tau$-summand of complexified
homology, the Hodge exact sequence has quotient $\mathcal V_\tau$
and kernel identified by polarization with
$\mathcal V_{\bar\tau}^*$. Its determinant is therefore
$\mathcal D_\tau$, with a single-valued parallel generator.

The metric assertion also has a direct check. Use the
linear-first convention $B=\operatorname{Im}H$ on real homology,
and put $e_\pm(v)=(v\mp iJv)/2$. The period map sends $e_+(v)$
to $v$ and kills $e_-(v)$. Normalize the analytic kernel
identification by
\[
 e_-(v)\longmapsto
       \bigl(w\longmapsto 2iB(e_+(w),e_-(v))\bigr)=H(-,v).
\]
It is independent of the lift of $w$ and is holomorphic.
The flat Hermitian form $2iB(z,\bar w)$ is $H$ on the
positive summand and minus its conjugate dual on the kernel.
Its absolute determinant norm is the ratio of the two tangent
determinant norms. Replacing $H$ by the prescribed $\pi H$
multiplies the squared norm by the fixed factor $\pi^{p-q}$,
where $p,q$ are the two eigenranks. Thus the parallel generator
has constant norm; normalize it analytically to norm one.

In actual local eigenframes write
$s=\varphi\,\det(F_+)\det(F_-)^{-1}$. Its unit norm gives
\[
       |\varphi|^2\det\Gamma_{F_+}=\det\Gamma_{F_-}.
\]
Lemma~\ref{geo-boundary-gram} yields
$|\varphi|\le C D^{p/2}$ and
$|\varphi^{-1}|\le C'D^{q/2}$.
A holomorphic function of logarithmic growth has no negative
Laurent coefficient: Cauchy's estimate bounds the coefficient
of $z^{-m}$ by $r^m(1+|\log r|)^N$, tending to zero.
Apply this at each smooth boundary divisor and use Hartogs
extension at the intersections. Both $\varphi$ and its inverse
extend holomorphically, so the section extends as a unit in the
actual frames.

For general proper $Y_{\mathbb C}$, Chow's lemma followed by
resolution and principalization gives a smooth projective
modification with a simple normal crossing boundary.
Pull back the actual family and eigenline. The preceding
argument gives a nowhere-vanishing analytic section there;
projective GAGA makes it algebraic
\cite[Theorems 1.0.1 and 1.0.3]{WlodarczykResolution}
\cite[\S12, Theorem 1]{SerreGAGA}.
For the proper birational map $f:Z\to Y_{\mathbb C}$,
normality gives $f_*\mathcal O_Z=\mathcal O_{Y_{\mathbb C}}$:
properness tests regularity of a rational function upstairs
over every codimension-one discrete valuation ring downstairs,
and normality then makes it regular. The projection formula
descends the section, and surjectivity preserves its
nonvanishing.

Finally, flat base change identifies
$H^0(Y,\mathcal D_\tau)\otimes_K\mathbb C$ with the
one-dimensional space of sections of the now trivial line
on $Y_{\mathbb C}$. A nonzero $K$-section is therefore a
generator, unique up to $K^\times$. At every complex embedding
it is a constant multiple of the parallel section just
constructed, so its norm is constant.
\end{proof}

\paragraph{Freeness at ramified coefficient primes.}
The missing integral condition is stronger than a characteristic-polynomial
identity. It can be proved on geometric fibres by using the universal
extension of the abelian quotient, followed by a three-step filtration.
The argument uses the full ring of integers, including its nilpotents
after reduction.

\begin{lemma}\label{geo-witt-coefficient-freeness}
Let $k$ be an algebraically closed field of characteristic $p>0$,
$W=W(k)$, and $\mathcal O$ the ring of integers of a number field $E$.
If a finite free $W$-module $D$ has an $\mathcal O$-action, then $D/pD$
is projective over $C=k\otimes_{\mathbf Z}\mathcal O$.
If moreover, for every $a\in\mathcal O$,
\begin{equation}\label{geo-regular-character-freeness}
 \det(X-a\mid D/pD)
 =\det(X-a\mid C)^r,
\end{equation}
for an integer $r\geq0$, then $D/pD$ is free of rank $r$ over $C$.
\end{lemma}
\begin{proof}
The ring $W\otimes_{\mathbf Z}\mathcal O$ is a product of discrete
valuation rings. Indeed, decompose $\mathcal O\otimes\mathbf Z_p$ by
its primes over $p$. The maximal unramified subextension of each local
field splits over $W$, and a uniformizer of the remaining totally
ramified extension has an Eisenstein polynomial. Each resulting factor
is a discrete valuation ring with uniformizer $\varpi$ and
$p=u\varpi^e$, where $u$ is a unit.
The corresponding summand of $D$ has no $\varpi$-torsion, since it has
no $p$-torsion. It is therefore free over this factor. Reduction modulo
$p$ proves projectivity over $C$.

Write $C=\prod_i C_i$, where $C_i\simeq k[t]/(t^{e_i})$ and the
residue character is $\tau_i:\mathcal O\to k$. If the rank on $C_i$
is $r_i$, the characteristic polynomial of $a$ is
$\prod_i(X-\tau_i(a))^{e_i r_i}$.
For a selected prime of $\mathcal O$ above $p$, the Chinese remainder
theorem gives an $a$ which is zero at the other primes and is a
nonzero generator of its residue field at that prime. Its distinct
residue embeddings give distinct nonzero roots. Comparing their
multiplicities in \eqref{geo-regular-character-freeness} gives
$e_i r_i=e_i r$, hence $r_i=r$ for every factor.
\end{proof}

For the next two propositions, retain an algebraically closed field
$k$ of characteristic $p>0$, the full ring of integers $\mathcal O$
of the number field $E$, and $C=k\otimes_{\mathbf Z}\mathcal O$.

\begin{proposition}\label{geo-proper-formal-hodge-comparison}
Let $B/k$ be an abelian variety with an $\mathcal O$-action and a
rigidified line bundle $M$ whose associated map $B\to B^\vee$ is an
isomorphism. Suppose its adjunction identifies the endomorphism
$a$ with the dual of $\bar a$, for an involution
$a\mapsto\bar a$ of $\mathcal O$.
Put $V=\operatorname{Lie}B$.
Choose formal coordinates at zero and a frame $e(x,y)$, rigidified
on both axes, of the associated Poincar\'e biextension.
Write its adjunction in this frame as
\[
             e(ax,y)\longmapsto u_a(x,y)e(x,\bar a y).
\]
If $\rho(a)$ is the tangent action and
$b_a$ is the mixed bilinear coefficient of $u_a$, then the module
$V\oplus V^*$ with action
\begin{equation}\label{geo-formal-hodge-action}
       a(v,\lambda)=
       \bigl(\rho(a)v,\rho(\bar a)^t\lambda-b_a^t v\bigr)
\end{equation}
is projective over $C=k\otimes_{\mathbf Z}\mathcal O$.
\end{proposition}
\begin{proof}
The universal vector extension of $B^\vee$ parametrizes extensions
of $B$ by $\mathbf G_m$ together with a pointed section over the
first infinitesimal neighbourhood of zero. This description is
functorial and commutes with arbitrary base change
\cite[I, \S2.1 and Proposition 2.6.7]{mazur-messing-universal}.
Use $M$ to identify the parameter space $B^\vee$ with $B$.
Over $k[\epsilon]/(\epsilon^2)$, a tangent point above $\epsilon v$
has a unique rigidification of the form
\[
                  e(\epsilon v,y)(1+\epsilon\lambda(y)),
                 \qquad \lambda\in V^*,
\]
on that first neighbourhood. This identifies the tangent space
of the universal extension with $V\oplus V^*$.
Pulling the rigidification back by $\bar a$ and then through the
inverse adjunction gives
\[
 e(a\epsilon v,y)
 \bigl(1+\epsilon\{\rho(\bar a)^t\lambda-b_a^t v\}(y)\bigr).
\]
This proves precisely \eqref{geo-formal-hodge-action}, with its
minus sign. It uses only first infinitesimal neighbourhoods, so
no logarithm or division by $p$ occurs.

The Lie algebra of the universal extension is the reduction of
the finite free $W(k)$-module furnished by its Dieudonn\'e crystal;
the comparison is functorial in the abelian variety and its
endomorphisms
\cite[II, \S9, \S13 and \S15]{mazur-messing-universal}.
Consequently the $\mathcal O$-action lifts to this module.
Lemma~\ref{geo-witt-coefficient-freeness} applies. In particular,
neither the abelian variety nor its endomorphisms need to lift
together to characteristic zero.
\end{proof}

\begin{proposition}\label{geo-toric-hodge-filtration}
Let $G/k$ be a semiabelian variety with an $\mathcal O$-action,
and let $q:G\to B$ be an equivariant quotient with torus kernel.
Put $V=\operatorname{Lie}G$, $U=\operatorname{Lie}B$, and
$T=k\otimes_{\mathbf Z}\Lambda$, where $\Lambda$ is the
cocharacter lattice of the torus. Write its exact tangent sequence
as
\[
                 0\longrightarrow T\xrightarrow{j}V
                    \xrightarrow{P}U\longrightarrow0.
\]
Suppose $B$ has the self-dual line bundle of
Proposition~\ref{geo-proper-formal-hodge-comparison}, and that the
rigidified line bundle defining the formal module on $G$ is the
pullback of this line bundle. Give both biextensions their
rigidified adjunctions.
Then the module $V\oplus V^*$ with action
\eqref{geo-formal-hodge-action} is projective over $C$.
If its characteristic polynomials satisfy
the $r$-fold regular identity for an integer $r\geq0$,
it is free of rank $r$.
\end{proposition}
\begin{proof}
Let $e_G,e_B$ be the frames, and write their comparison as
$e_G(x,y)=\ell(x,y)^{-1}e_B(qx,qy)$, with $\ell$ equal to one
on both axes. The two induced adjunctions agree: their ratio is
a global unit on $G\times G$ equal to one on both axes, and hence
is one by the theorem on units on products
\cite[Theorem 1.1]{conrad-units-product}.
If $b_a,c_a$ are the mixed coefficients of the two adjunctions
and $Z$ is that of $\ell$, their relation is therefore
\[
 u_a^G(x,y)=u_a^B(qx,qy)\ell(x,\bar a y)/\ell(ax,y).
\]
Taking the mixed bilinear coefficient gives, for every $a$,
\begin{equation}\label{geo-toric-gauge-comparison}
 b_a+\rho(a)^tZ-Z\rho(\bar a)=P^t c_aP.
\end{equation}
The change $(v,\lambda)\mapsto(v,\lambda+Z^tv)$ identifies the
action with bilinear term $P^tc_aP$ with the original action.
In these new coordinates, the map
\[
       (v,\lambda)\longmapsto j^*\lambda
          \quad\hbox{from }V\oplus V^*
          \hbox{ to }T^\dagger
\]
is equivariant and surjective; the dagger denotes the
conjugate dual action. Its kernel is $V\oplus P^*U^*$.
The equivariant map
\[
           (v,P^*\mu)\longmapsto(Pv,\mu)
\]
from this kernel onto the proper abelian module $U\oplus U^*$
has kernel $T\oplus0$.
Thus the three successive quotients are
$T$, $U\oplus U^*$, and $T^\dagger$.

The $\mathcal O$-module $\Lambda$ is torsion-free: if $a\ne0$
kills an element, multiplying by an integral $b$ with
$ab=N_{E/\mathbf Q}(a)\ne0$ would give integer torsion.
It is therefore projective over the Dedekind domain $\mathcal O$.
The same argument applies to its integer dual, with conjugate
action. Their scalar extensions give the two projective toric
quotients. The middle quotient is projective by
Proposition~\ref{geo-proper-formal-hodge-comparison}.
Both exact sequences therefore split as $C$-modules, proving
projectivity of the full module. The residue-character argument
in Lemma~\ref{geo-witt-coefficient-freeness} gives the common rank.
No trace pairing is used at a ramified prime.
\end{proof}

These are statements about the prescribed formal module, rather
than its semisimplification. They pass to the actual local
lattice as follows. Let $S$ be a local coefficient ring and let
$\mathcal H$ be the underlying finite free $S$-module with its
$\mathcal O$-action. If its geometric closed fibre is free of rank
$r$ over $\mathcal O\otimes k$, then its residue-field fibre is
free of that rank: on each Artin local factor, choose $r$ generators
modulo the maximal ideal, apply Nakayama, and compare dimensions
after algebraic closure. Lift these generators to $\mathcal H$.
The resulting map
$(S\otimes_{\mathbf Z}\mathcal O)^r\to\mathcal H$ is a map between
free $S$-modules of the same rank and is an isomorphism modulo
the maximal ideal. Its determinant is a unit, so it is an
isomorphism over $S$.

For the boundary application, identify the prescribed integral
formal charts and their residue-fibre semiabelian quotients with
the original self-dual pullback line bundle. The propositions
then supply the frame comparison and coefficient freeness,
including at ramified primes. The next construction supplies
both filtrations on these same charts. The actual boundary
model and its complex multiplication evaluation remain separate.

\paragraph{Constructing both filtrations from the biextension.}
The two filtrations can be built on the original family itself.
The construction uses first infinitesimal neighbourhoods and
multiplication by $p$.

\begin{theorem}\label{geo-biextension-conjugate-filtration}
Let $p$ be a prime, let $S$ be a scheme over $\mathbf F_p$,
and let $G/S$ be a
smooth separated commutative group scheme of relative dimension $d$.
Let $\mathcal P$ be a symmetric biextension of $G\times_S G$ by
$\mathbf G_m$, with its rigidifications on the axes and its symmetry
equal to the identity on the diagonal.
Suppose that, at every geometric point of $S$, the group is
semiabelian and its biextension is the pullback from the Poincar\'e
biextension of a line bundle inducing an isomorphism between its
abelian quotient and the dual quotient.

Put $\mathcal V=\operatorname{Lie}(G/S)$ and
$\mathcal U=\mathcal V^*$.
There is a vector bundle $\mathcal H$ on this same base with exact
sequences
\begin{equation}\label{geo-actual-two-filtrations}
 0\longrightarrow\mathcal U\longrightarrow\mathcal H
       \longrightarrow\mathcal V\longrightarrow0,\qquad
 0\longrightarrow F_S^*\mathcal V\longrightarrow\mathcal H
       \longrightarrow F_S^*\mathcal U\longrightarrow0.
\end{equation}
It has the formal-coordinate transition functions and integral
lattices defined by the first derivatives and mixed bilinear
coefficients of $\mathcal P$.
An action of a commutative coefficient ring with involution
$a\mapsto\bar a$ on $G$, with compatible adjunction on $\mathcal P$,
makes both sequences equivariant; the action on $\mathcal U$ is
the conjugate dual action.
\end{theorem}
\begin{proof}
Let $G_1$ be the first infinitesimal neighbourhood of the zero
section. It is finite locally free over $S$ of rank $d+1$, with
square-zero ideal $\mathcal U$.
Over a point $x$ of $G$, consider pointed sections of
$\mathcal P_x|_{G_1}$. They form a torsor under the additive vector
group $\mathcal U$.
These torsors are representable: push the line bundle on
$G\times_S G_1$ forward to $G$, and take the inverse image of the
unit section under restriction to $G\times_S S$.
The first biextension law gives this torsor a commutative group
law. Denote the resulting vector extension by $\mathcal E$:
\[
        0\longrightarrow\mathcal U_{\mathrm{add}}
          \longrightarrow\mathcal E\xrightarrow{\pi}G
          \longrightarrow0.
\]
Taking Lie algebras gives the first sequence, with
$\mathcal H=\operatorname{Lie}(\mathcal E/S)$.
All these constructions commute with base change.

Multiplication by $p$ on $\mathcal E$ kills its additive kernel.
It therefore descends to a group homomorphism $j:G\to\mathcal E$
with $j\pi=[p]_{\mathcal E}$.
Differentiation at the identity gives $dj=0$, since $d\pi$ is
surjective and $d[p]_{\mathcal E}=0$. Translation gives zero
relative differential everywhere.
For a smooth scheme in characteristic $p$, functions with zero
relative differential are precisely the image under relative
Frobenius; this follows in polynomial coordinates and passes to
\'etale coordinates. Hence there is a unique homomorphism
\[
       \Phi:G^{(p)}\longrightarrow\mathcal E,\qquad
       \Phi F_{G/S}=j.
\]
Its differential is a map
$\alpha:F_S^*\mathcal V\to\mathcal H$.

We first verify its exact coordinate formula and isotropy.
Choose coordinates $x$ on the formal group and a frame $e(x,y)$
of $\mathcal P$, rigidified on both axes. A pointed section on
$G_1$ is written $e(x,y)(1+\lambda(y))$.
Write $C(x,y)$ and $D(x,y)$ for the scalars of $p$-fold
biextension multiplication in the first and second variable,
respectively, and write the symmetry as
\[
               e(x,y)\longmapsto s(x,y)^{-1}e(y,x).
\]
Put
\[
 h_{ij}=[x_j^p]\,[p]_G(x)_i,\qquad
 A_{ij}=[x_i^py_j]\,C(x,y),\qquad
 \Sigma_{ij}=[x_i y_j]\,s(x,y).
\]
Multiplication in $\mathcal E$ adds the derivative in $y$ of the
first biextension cocycle to the vector coordinate. Thus
\[
 j(x)=\bigl([p]_G(x),d_yC(x,0)\bigr),\qquad
                         \alpha(v)=(hv,A^tv).
\]
In particular $dj=0$ gives $[x_i y_j]C=0$.
The biextension symmetry and interchange laws give
\begin{align*}
 D(x,y)&=C(y,x)s(x,[p]y)s(x,y)^{-p},\\
 C(x,[p]y)D(x,y)^p&=D([p]x,y)C(x,y)^p.
\end{align*}
The first identity also gives $[x_i y_j]D=0$.
Since $[p]_G$ factors through Frobenius, its terms have degree at
least $p$. Taking the coefficient of $x_i^py_j^p$ in the second
identity, and then using the first, yields
\begin{equation}\label{geo-biextension-isotropy-identity}
                  Ah=h^tA^t+h^t\Sigma h.
\end{equation}
Indeed the $p$-th powers of $C,D$ contribute zero to that
coefficient, while $s([p]x,y)^{-p}$ has first-variable degree at
least $p^2$. The three remaining contributions are exactly those
displayed in \eqref{geo-biextension-isotropy-identity}.

The symmetry of $\mathcal P$ defines a perfect alternating
pairing on $\mathcal H$. Intrinsically, compare the two pointed
sections at a pair of first-order tangent points after swapping
the two factors. In the chosen frame its value is
\begin{equation}\label{geo-biextension-pairing}
 \langle(v,\lambda),(w,\mu)\rangle
              =\mu(v)-\lambda(w)+v^t\Sigma w.
\end{equation}
The diagonal condition gives $s(x,x)=1$ and makes $\Sigma$
alternating, also when $p=2$. This description proves that the pairing is independent
of the frame and is perfect. Equation
\eqref{geo-biextension-isotropy-identity} says precisely that
$\operatorname{im}\alpha$ is isotropic.

It remains to prove injectivity of $\alpha$ on geometric fibres.
For a proper abelian fibre $B$, the extension just constructed is
its universal vector extension, after the prescribed self-duality.
The universal extension of $V_B:B^{(p)}\to B$ kills its additive
kernel: on that kernel the map is pullback by the dual
homomorphism $F_{B^\vee}$, whose differential is zero.
It consequently factors through a map $B^{(p)}\to\mathcal E$.
Functoriality and $V_BF_B=[p]_B$ identify this map with $\Phi$.

To fix the Frobenius convention, take the contravariant
Dieudonn\'e module of $B^\vee[p^\infty]$.
Its reduction is the Lie algebra of the universal extension of
$B$, and the reduction of crystalline Frobenius is the differential
of the universal extension of
$(F_{B^\vee})^\vee=V_B$.
Its kernel is the Frobenius pullback of the additive Hodge
subspace. Therefore its induced map on $F_S^*\operatorname{Lie}B$
is injective, as required
\cite[II, \S\S9 and 13]{mazur-messing-universal}
\cite[\S1.1.4, equation (1.1.4.1)]{madapusi-pera-realizations}.

For a semiabelian geometric fibre let
$0\to T\to G\xrightarrow{q}B\to0$ be its torus quotient.
The pullback biextension is canonically trivial on $G\times T$.
Restricting pointed sections to $T_1$ gives a homomorphism
$\mathcal E\to\omega_T$. Its kernel is
$\mathcal E'=G\times_B E_B$, where $E_B$ is the extension for $B$.
Since $p$ kills $\omega_T$, the map $\Phi$ lands in $\mathcal E'$.
The composite to $E_B$ is $\Phi_Bq^{(p)}$; this follows by
composing with relative Frobenius and using multiplication by $p$.
On $T^{(p)}$, the map is the toric Verschiebung into
$T\subset\mathcal E'$, which is an isomorphism over the
algebraically closed field.
If $\alpha(v)=0$, its image in the proper quotient gives
$dq^{(p)}(v)=0$. Thus $v$ lies in $\operatorname{Lie}T^{(p)}$,
where toric Verschiebung is injective. Hence $v=0$.

The fibrewise rank criterion now makes $\alpha$ a locally split
injection of vector bundles over $S$.
Define $\beta:\mathcal H\to(F_S^*\mathcal V)^*$ by
$\beta(z)(v)=\langle\alpha(v),z\rangle$.
Perfection of the pairing makes $\beta$ surjective; isotropy and
the ranks give $\ker\beta=\operatorname{im}\alpha$.
This is the second sequence in
\eqref{geo-actual-two-filtrations}.

Finally, an adjunction of the coefficient action induces
endomorphisms of $\mathcal E$. They commute with multiplication
by $p$, so the uniqueness of $\Phi$ proves equivariance of
$\alpha$ with the pulled-back action. In coordinates, compatibility
of adjunction with symmetry says
\[
 u_a(x,y)u_{\bar a}(y,x)
                    =s(x,\bar a y)/s(ax,y).
\]
Its mixed coefficient is
$b_a+b_{\bar a}^t=\Sigma\rho(\bar a)-\rho(a)^t\Sigma$.
Together with \eqref{geo-formal-hodge-action}, this proves
$\langle az,z'\rangle=\langle z,\bar a z'\rangle$.
It gives the conjugate dual action on the quotient.
The construction uses the original infinitesimal
neighbourhood and line bundle at every stalk, so its lattices are
the original integral lattices.
\end{proof}

For clarity, the coordinate formulas also specify the comparison
with the initial formal construction. Write the two frames as
$e_2(x,y)=\ell(x,y)e_1(\phi x,\phi y)$, with
$L=d\phi$ and $Z$ the mixed matrix of $\ell$. Then the transition
$T:\mathcal H_2\to\mathcal H_1$ satisfies
\begin{align*}
        T(v,\lambda)&=\bigl(Lv,L^{-t}(\lambda+Z^tv)\bigr),\\
        \Sigma_2&=L^t\Sigma_1L+Z^t-Z.
\end{align*}
The coefficient action is \eqref{geo-formal-hodge-action},
$\alpha=(h,A^t)$, and
\[
                  \beta(v,\lambda)
                        =h^t\lambda-(A-h^t\Sigma)v.
\]
The matrix $\Sigma$ must be retained, including when $Z$ is
nonsymmetric.
These formulas preserve the projection to the actual Lie lattice
and the actual dual sublattice.

\begin{corollary}[Compatible formal charts on a curve]
\label{geo-biextension-formal-curve}
The same conclusion holds on a smooth curve in characteristic $p$
when one is given
compatible formal biextension charts at its generic point and
over its local rings, with the displayed integral transition
matrices, and each geometric residue chart comes from a
semiabelian variety satisfying the theorem's fibre hypothesis.
Here no semiabelian scheme over the whole curve is required.
\end{corollary}
\begin{proof}
The first-neighbourhood construction and its multiplication law
make sense formally. Over a ring of characteristic $p$, a formal
power series with zero partial derivatives involves only powers
divisible by $p$, so the same relative Frobenius factorization
defines $\alpha$. Its formula, the pairing, and $\beta$ commute
with all the displayed transitions. They therefore give maps of
the prescribed integral vector bundles. The fibrewise argument
proves injectivity of $\alpha$ at each geometric residue point.
The locally split rank criterion, followed by isotropy, gives both
exact sequences. This also proves their coefficient equivariance.
\end{proof}

In applying this result, the line bundle itself must be tracked.
If the original line $L$ on a geometric fibre is isomorphic to
$q^*M$ and the translation homomorphism
$\lambda_M:B\to B^\vee$ has trivial scheme-theoretic kernel, then
$\lambda_M$ is an isomorphism: its kernel is finite, so its image
has the full dimension of $B^\vee$, and its degree is one.
The associated biextension is the pullback of that of $M$.
Indeed, after normalizing the line isomorphism at zero, any
discrepancy in the multiplication laws is a unit on a product
of connected reduced varieties, trivial on the coordinate axes,
and is therefore one by \cite[Theorem~1.1]{conrad-units-product}.
The canonical exchange of the two factors is identity on the
diagonal. This supplies the required symmetry also in
characteristic two. Tensoring $L$ by an untracked positive power
would not justify this argument.

Together with coefficient freeness, the result supplies the two
filtered bundles from the prescribed compatible charts and
geometric fibres. Constructing those data on the chosen integral
model, and evaluating its integral complex multiplication
determinant, remain separate inputs.

\paragraph{Determinants before degrees.}
The two filtrations give more than a degree identity. They force a
finite tensor power of the determinant to be trivial, even on a
singular base.

\begin{proposition}\label{geo-frobenius-determinant-torsion}
Let $X$ be a scheme over $k=\mathbf F_{p^f}$, let $F$ be absolute
Frobenius, and let $A$ be a finite commutative $\mathbf F_p$-algebra.
Put $B=k\otimes_{\mathbf F_p}A$. Suppose finite locally free
$\mathcal O_X$-modules with $A$-actions have $A$-linear exact sequences
\begin{equation}\label{geo-hodge-conjugate-sequences}
\begin{gathered}
0\longrightarrow\mathcal U\longrightarrow\mathcal H
 \longrightarrow\mathcal V\longrightarrow0,\\
0\longrightarrow F^*\mathcal V\longrightarrow\mathcal H
 \longrightarrow F^*\mathcal U\longrightarrow0.
\end{gathered}
\end{equation}
The pullbacks carry the pulled-back $A$-actions. For any idempotent
$e\in B$, put $\mathcal D_e=\det(e\mathcal H)$. Then
\[
                    \mathcal D_e^{\otimes(p^f-1)}
                         \simeq\mathcal O_X.
\]
Let $\chi:B\to k$ be a $k$-algebra character and $B_e=eB$ its
Artin local factor, of length $n=\dim_k B_e$. If $e\mathcal H$ is
locally free over $\mathcal O_X\otimes_k B_e$, then
$\mathcal Q_\chi=\mathcal H\otimes_{B,\chi}k$ is a vector bundle and
\begin{equation}\label{geo-character-determinant-torsion}
       (\det\mathcal Q_\chi)^{\otimes n(p^f-1)}
                             \simeq\mathcal O_X.
\end{equation}
In particular every pullback of $\mathcal Q_\chi$ to a proper curve
has degree zero.
\end{proposition}
\begin{proof}
The automorphism $\sigma(c\otimes a)=c^p\otimes a$ of $B$ fixes $A$;
it need not equal Frobenius on the possibly nonreduced algebra $B$.
For any bundle $\mathcal W$ with this action,
$F^*(e\mathcal W)=\sigma(e)F^*\mathcal W$. Project the second
sequence by $\sigma(e)$ and the first by $e$, and take determinants:
\[
 \mathcal D_{\sigma(e)}
 \simeq F^*\{\det(e\mathcal V)\otimes\det(e\mathcal U)\}
 \simeq F^*\mathcal D_e
 \simeq\mathcal D_e^{\otimes p}.
\]
The last identity follows directly by raising transition functions
to their $p$-th powers. Since $\sigma^f=1$, iteration proves the
first assertion. Idempotent images are direct summands, and the
bundle sequences are locally split; no flatness of $F$ is required.

The residue field of $B_e$ is $k$. Tensor a composition series of
$B_e$ with the locally free coefficient module $e\mathcal H$.
Its $n$ quotient bundles are all isomorphic to $\mathcal Q_\chi$,
so $\mathcal D_e\simeq(\det\mathcal Q_\chi)^{\otimes n}$.
This proves \eqref{geo-character-determinant-torsion}.
Degree on a proper curve is additive under tensor products and
takes integer values, giving the final assertion. Rank zero causes
no exception.
\end{proof}

The factor $n$ retains the nilpotent multiplicity at a ramified
coefficient prime. The proposition applies to the entire base only
if its filtrations and coefficient module are defined there. If
they are constructed only on test curves, its conclusion is
curvewise numerical triviality, which already suffices below.

\paragraph{One finite valuation on the whole fibre.}
We use the elementary fact that a numerically trivial line bundle on
a proper integral variety has no nonzero section with a zero.
Indeed Chow's lemma reduces to a projective integral modification.
A nonempty effective zero divisor then meets a projective curve
not contained in it, after a finite field extension if necessary.
The restricted section has positive degree. Pullback preserves
numerical triviality, by the degree formula for maps of proper
curves, giving a contradiction
\cite[Tags 0200 and 0B40]{stacks-finite-place}.

\begin{theorem}\label{geo-vertical-line-rigidity}
Let $R$ be a discrete valuation ring with uniformizer $\pi$,
fraction field $K$ and residue field $k$. Let $X$ be a normal
integral proper flat $R$-scheme with nonempty connected special fibre.
Let $\mathcal L$ be a line bundle whose restriction to every proper
integral curve in $X_k$ has degree zero. Suppose $s$ generates
$\mathcal L$ everywhere on $X_K$. Then there are integers $M>0$
and $a$ for which
\begin{equation}\label{geo-vertical-generator}
                    s^{\otimes M}/\pi^a
          \quad\hbox{generates }\mathcal L^{\otimes M}
          \quad\hbox{on all of }X.
\end{equation}
Write the fibre divisor as $\operatorname{div}(\pi)=\sum_i m_iX_i$.
The normalized order $q=a/M$ is unique, and its denominator in
lowest terms divides $\gcd_i m_i$. In particular $M=1$ is possible
when the fibre multiplicities are relatively prime.
\end{theorem}
\begin{proof}
The divisor $D=\operatorname{div}(s)$ is vertical, so write
$D=\sum_i d_iX_i$ and set $q=\min_i(d_i/m_i)=a/M$ in lowest terms,
with $M>0$. The Cartier divisor
\[
                  E=MD-a\operatorname{div}(\pi)
\]
has nonnegative coefficients and omits at least one component.
Normality implies that it is effective: locally a rational
function with nonnegative order at every prime divisor is regular.

Suppose $E$ is nonempty. The intersection graph of the irreducible
components of the connected special fibre is connected. Hence
some component $X_i$ omitted from $E$ meets its support. Restricting
$E$ to the reduced integral scheme $X_i$ gives a nonempty effective
Cartier divisor. Its line bundle is
$\mathcal L^{\otimes M}|_{X_i}$, because the whole fibre is
principal. This line bundle is numerically trivial, contradicting
the preceding fact. Thus $E=0$. Normality again implies that a
rational section with neither zeros nor poles is a generator,
proving \eqref{geo-vertical-generator}.

All ratios $d_i/m_i$ are therefore $q$. In lowest terms,
$M d_i=a m_i$ and $\gcd(a,M)=1$, so $M$ divides every $m_i$.
These identities also prove uniqueness.
\end{proof}

The theorem does not assert $q=0$. For example, the trivial line
with generic generator $\pi$ has normalized order one. A totally
ramified extension of degree $m$, with a uniformizer as generator,
can give order $1/m$, so the multiplicities cannot be discarded.
Connectedness is also essential: the generators $1$ and $\pi$
on two disjoint copies of $\operatorname{Spec}R$ have different
orders.

For the prescribed conjugate determinant, the characteristic-$p$
proposition supplies the curve degrees once the actual Hodge
filtrations and coefficient freeness are identified. The theorem
then makes the finite norm constant across the connected fibre.
Its value still requires the same integral complex multiplication
evaluation and different correction as in
\eqref{geo-exact-conjugate-degree}. Neither that evaluation nor
the boundary lattice identification follows from the abstract
determinant argument.

\paragraph{Degrees of character quotients in characteristic $p$.}
The degree calculation on a special-fibre curve has a short proof that
also explains the extra hypothesis needed at ramified coefficient primes.

\begin{proposition}\label{geo-frobenius-character-degree}
Let $k$ be a finite field of characteristic $p$, let $C/k$ be a smooth
projective geometrically connected curve, and let $F:C\to C$ be absolute
Frobenius. Let $A$ be a finite commutative $\mathbf F_p$-algebra and put
$B=k\otimes_{\mathbf F_p}A$. Suppose vector bundles
$\mathcal H,\mathcal U,\mathcal V$ on $C$ have $A$-actions and
$A$-linear exact sequences
\begin{equation}\label{geo-hodge-conjugate-sequences}
 0\longrightarrow\mathcal U\longrightarrow\mathcal H
       \longrightarrow\mathcal V\longrightarrow0,
 \qquad
 0\longrightarrow F^*\mathcal V\longrightarrow\mathcal H
       \longrightarrow F^*\mathcal U\longrightarrow0.
\end{equation}
The actions on the Frobenius pullbacks are the pulled-back $A$-actions.
Then $\deg(e\mathcal H)=0$ for every idempotent $e\in B$.

Let $\chi:B\to k$ be a $k$-algebra character, and let $e$ be the
idempotent of its Artin local factor $B_e=eB$. If $e\mathcal H$ is
locally free as a module over $\mathcal O_C\otimes_k B_e$, then
\[
 \mathcal Q_\chi=\mathcal H\otimes_{B,\chi}k
 \quad\hbox{is a vector bundle and}\quad
                         \deg\mathcal Q_\chi=0.
\]
Here local freeness over the coefficient algebra is stronger than
local freeness merely over $\mathcal O_C$. If $B_e=k$, the stronger
condition is automatic.
\end{proposition}
\begin{proof}
Write $\sigma(c\otimes a)=c^p\otimes a$ on $B$; this acts on the
finite field coefficients and leaves $A$ fixed. It is not the
absolute Frobenius of the possibly nonreduced algebra $B$.
For a bundle $\mathcal W$ with $A$-action, pulling back the action
of $e=\sum_i c_i\otimes a_i$ raises the scalars $c_i$ to their
$p$-th powers. Hence
\[
 F^*(e\mathcal W)=\sigma(e)F^*\mathcal W.
\]
Idempotent images are direct summand bundles, so applying an
idempotent preserves both exact sequences. Frobenius is finite flat
on the smooth curve, and $\deg F^*\mathcal W=p\deg\mathcal W$.
Taking the $\sigma(e)$ part of the second sequence and the $e$ part
of the first gives
\[
 \deg\bigl(\sigma(e)\mathcal H\bigr)
       =p\{\deg(e\mathcal V)+\deg(e\mathcal U)\}
       =p\deg(e\mathcal H).
\]
If $k=\mathbf F_{p^f}$, then $\sigma^f=1$. Iteration yields
$\deg(e\mathcal H)=p^f\deg(e\mathcal H)$, proving the first claim.

The residue field of $B_e$ is $k$. A composition series of its
ideals has
$n=\operatorname{length}_{B_e}B_e=\dim_kB_e>0$ successive quotients
isomorphic to $k$. Local freeness over
$\mathcal O_C\otimes_k B_e$ makes tensoring this series with
$e\mathcal H$ exact. It gives a filtration by subbundles whose
$n$ successive quotients are each isomorphic to $\mathcal Q_\chi$.
In particular $\mathcal Q_\chi$ is locally free, and additivity of
degree gives
\[
 0=\deg(e\mathcal H)=n\deg\mathcal Q_\chi.
\]
This proves the assertion, including rank zero.
\end{proof}

For $A=\mathcal O_E/p\mathcal O_E$ and $p$ unramified in $E$,
the algebra becomes a product of copies of $k$ after a finite extension
of the constant field. The character conclusion then follows from
the two exact sequences alone. At a ramified prime the nilpotents
must be retained, and local freeness over the Artin factor supplies
their multiplicity in the proof.

The following example shows why a zero degree for the whole bundle
cannot replace that hypothesis. On $C=\mathbf P^1_k$ take
$B=k[\epsilon]/(\epsilon^2)$ and
$\mathcal H=\mathcal O(1)\oplus\mathcal O(-1)$, with
$\epsilon(u,v)=(sv,0)$ for a nonzero
$s\in H^0(C,\mathcal O(2))$. Although $\deg\mathcal H=0$, the
torsion-free quotient by the image of $\epsilon$ is
$\mathcal O(-1)$, of degree $-1$. The unsaturated quotient also
has a torsion summand of length $2$. The $B$-module is not locally
free at the zero divisor of $s$. This example concerns the inference
from total degree alone; it is not claimed to satisfy
\eqref{geo-hodge-conjugate-sequences} or to arise from the prescribed
abelian family.

To apply the proposition to the geometric construction, the two
filtrations and the required coefficient freeness must hold for its
actual integral Hodge bundle along each special-fibre curve. Its
generic character determinant must also be identified with the ratio
of the prescribed conjugate tangent determinants. These model
comparisons, including their behaviour at the boundary, remain to be
established; the abstract degree calculation does not supply them.

\paragraph{The different gives the exact integral eigenframe.}
The scalar needed at a complex multiplication point is an integral
lattice calculation. Its local form does not require the CM extension
to be unramified over its maximal real subfield.

\begin{proposition}\label{geo-relative-different-eigenframe}
Let $F\subset E$ be number fields, let $\ell$ be a rational prime,
and put $A=\mathcal O_F\otimes\mathbf Z_{(\ell)}$ and
$B=\mathcal O_E\otimes\mathbf Z_{(\ell)}$.
Let $R$ be a discrete valuation ring with a map $A\to R$ and
fraction field $K$ of characteristic zero splitting $E/F$ relative
to that map. For an embedding $\sigma:E\to K$ extending the map
on $F$, let $e_\sigma$ be its primitive idempotent in $K\otimes_F E$.
Then the integral character module is exactly
\begin{equation}\label{geo-relative-different-lattice}
 (R\otimes_A B)[\sigma]
          =\bigl(\sigma(\mathfrak D_{E/F})R\bigr)e_\sigma.
\end{equation}
The image of the relative different denotes its extended ideal in $R$.
\end{proposition}
\begin{proof}
The finite projective module $B$ is free over the semilocal
Dedekind domain $A$. Choose a basis $b_i$ and its trace-dual basis
$b_i^*$, so $\operatorname{Tr}_{E/F}(b_i b_j^*)=\delta_{ij}$.
The $A$-span of the dual basis is the inverse different
\cite[Tag 0BW0]{stacks-finite-place}.
In the splitting algebra,
\[
                  e_\sigma=\sum_i\sigma(b_i^*)b_i.
\]
Indeed the embedding matrix and the transpose of the dual-basis
embedding matrix are inverses. Evaluating the displayed sum at
another embedding gives the corresponding Kronecker delta.
The integral basis $b_i$ now gives
\[
 c e_\sigma\in R\otimes_A B
 \quad\Longleftrightarrow\quad
 c\sigma(\mathfrak D_{E/F}^{-1})R\subset R
 \quad\Longleftrightarrow\quad
 c\in\sigma(\mathfrak D_{E/F})R.
\]
Here images of integral elements lie in $R$ by normality, and
inverse fractional ideals remain inverse after extension to $R$.
Every generic character vector is $c e_\sigma$, proving the claim.
\end{proof}

Suppose now that $E$ is CM and contains
$F=\mathbf Q(\sqrt{-d})$, with
$\mathfrak D_{F/\mathbf Q}=(\sqrt{-d})$.
The relative different is conjugation-invariant: its inverse is
preserved because
\[
 \operatorname{Tr}_{E/F}(b\bar x)
       =\overline{\operatorname{Tr}_{E/F}(\bar b x)}.
\]
Choose a generator $\delta$ of the relative different over the
semilocal ring at all primes above $\ell$. Then
$\bar\delta/\delta$ is a unit of that ring. Every $\ell$-adic
embedding sends this unit and its inverse to integral elements,
so they have absolute value one. Multiplicativity of the different
in the tower $E/F/\mathbf Q$ therefore gives
\begin{equation}\label{geo-cm-local-different-norm}
 \log|j\delta|
  =\log|j\mathfrak D_{E/\mathbf Q}|-\log|j\sqrt{-d}|,
 \qquad |j\bar\delta|=|j\delta|.
\end{equation}
For completeness, tower multiplicativity follows by composing field
traces: locally, a generator of the lower inverse different multiplies
the upper inverse different. Inverting the fractional ideals gives
$\mathfrak D_{E/\mathbf Q}
=\mathfrak D_{E/F}\mathfrak D_{F/\mathbf Q}\mathcal O_E$.
No real generator is needed for
\eqref{geo-cm-local-different-norm}. Under the stronger prescribed
host condition that $E/E^+$ is unramified, one may instead take
$\delta=\delta_0/\sqrt{-d}$ with $\delta_0\in E^+$ generating
the absolute different locally. That choice gives $\bar\delta=-\delta$.
The separate global host conditions remain in force.

To apply the proposition, suppose the actual integral CM tangent
lattice has basis
\[
 \{\rho(b_i)u_a\}_{i,\,1\le a\le r}
       \ \cup\ \{\rho(b_i)u'_b\}_{i,\,1\le b\le s},
\]
where $u_a$ and $u'_b$ have the two conjugate $F$-actions.
These are regular $B$-blocks for the two base embeddings.
Generic character separation kills the opposite block, even if
the residue characters coincide at $\ell=2$.
Writing $P_\sigma$ for the generic eigenprojection, the integral
$\sigma$-eigenframe is therefore exactly
\[
                  \{\sigma(\delta)P_\sigma u_a\}_{a=1}^r.
\]
The conjugate formula uses $\bar\sigma=\sigma\circ\overline{(\cdot)}$
and the $u'_b$. Thus frame covariance and
\eqref{geo-cm-local-different-norm} give the finite scaling in the
determinant comparison. The next argument constructs this same
integral product basis and its CM point, then evaluates the period mean.
A generic basis or an isogeny of degree divisible by $\ell$ does not
supply the integral-basis hypothesis.

\paragraph{The complex multiplication value comes from one elliptic curve.}
The period calculation needs no height formula for higher dimensional
CM abelian varieties. An elliptic curve, a local orthogonal decomposition,
and the product formula suffice. We retain the prescribed host condition:
$E$ is CM, $E/E^+$ is unramified at finite places, and
$F=\mathbf Q(\sqrt{-d})\subset E$, where $-d$ is the prescribed
odd fundamental discriminant. Put $m=[E:F]$.

\begin{lemma}\label{geo-pure-cm-local-basis}
Let $L$ be a projective $\mathcal O_E$-lattice of rank $p+q$,
with a skew-Hermitian form $h$, linear in its first argument,
taking values in
$\mathfrak D_{E/\mathbf Q}^{-1}$. Suppose
$\operatorname{Tr}_{E/\mathbf Q}h$ is a perfect alternating
form on $L$. At the half of the complex embeddings extending
one fixed embedding $F\hookrightarrow\mathbf C$, suppose
$i\tau h$ has constant signature $(p,q)$.
For every rational prime $\ell$ there are orthogonal $E$-subspaces
$W_+,W_-$, positive and negative at all these embeddings, and
vectors $x_1,\ldots,x_p\in L\cap W_+$ and
$x_{p+1},\ldots,x_{p+q}\in L\cap W_-$ which form a basis of
$L\otimes\mathbf Z_{(\ell)}$ over
$\mathcal O_E\otimes\mathbf Z_{(\ell)}$.
The intersections $L\cap W_\pm$ are saturated projective lattices;
their direct sum equals $L$ after this localization.
Writing $G_\pm$ for their Gram blocks, every $\ell$-adic embedding
$j$ satisfies
\[
 \log|j\det G_+|=-pD_\ell(j),\qquad
 \log|j\det G_-|=-qD_\ell(j),\qquad
 D_\ell(j)=\log|j\mathfrak D_{E/\mathbf Q}|.
\]
The assertion includes $\ell=2$ and empty blocks.
\end{lemma}
\begin{proof}
Write $A=\mathcal O_E\otimes\mathbf Z_{(\ell)}$.
Unramifiedness over $E^+$ gives a real generator $\delta_0$ of
the absolute different over $A$. Trace duality and perfection
show that $\delta_0h$ is a unimodular skew-Hermitian form on $L_A$.
There is an anti-invariant unit $u\in A$.
At an inert dyadic place take $a$ with unit residue trace and
$u=a-\bar a=2a-\operatorname{Tr}(a)$; at a split place choose
$a$ with residues $(1,0)$, so $a-\bar a$ has residues $(1,-1)$.
At the other inert places choose
$a-\bar a$ nonzero in the quadratic residue field.
Choose $a$ simultaneously by the Chinese remainder theorem and
set $u=a-\bar a$, preserving exact anti-invariance.
Thus $u\delta_0h$ is a unimodular Hermitian form.

It has an orthogonal basis over the semilocal ring $A$.
Indeed a nondegenerate Hermitian space over a quadratic finite
field extension has a vector of nonzero norm, also in characteristic
two: otherwise the values on $v+aw$ would force the nondegenerate
trace pairing to vanish. At a split factor choose mutually dual
vectors in the two factors. Lift these choices simultaneously,
split off the resulting unit norm line, and induct.
Clearing denominators prime to $\ell$ makes this basis $e_i$ integral.

Choose a positive $p$-frame at every complex embedding in the
specified half. In each coset $e_i+\ell L$, the Minkowski lattice
has bounded covering radius. Approximate a sufficiently large
positive scalar multiple of the prescribed frames. After rescaling,
the errors tend to zero, so the resulting $x_1,\ldots,x_p$ are
positive at every embedding. Their Gram matrix is still invertible
over $A$, since it is unchanged modulo every prime above $\ell$.
Consequently its inverse gives an integral orthogonal projector
$L_A\to\sum A x_i$.
Take $W_+$ to be their $E$-span and $W_-=W_+^\perp$.
The complementary lattice is free over $A$; clear denominators
prime to $\ell$ in a basis to obtain the remaining $x_i$.
The signature gives negativity on $W_-$.
The two Gram blocks of $u\delta_0h$ are unimodular, proving the
valuation formulas. Intersections with $E$-subspaces are saturated
and hence projective over the Dedekind domain $\mathcal O_E$.
For rank zero the determinant is one.
\end{proof}

\begin{proposition}\label{geo-cm-integral-product}
On $W_+$ use the fixed complex embedding of $F$, and on $W_-$
its conjugate. Give the resulting complex torus the Riemann
form $\operatorname{Tr}h$, and keep its specified symmetric
principal line bundle rigidified at zero.
It has a model over a number field, and after a finite extension
it has good reduction everywhere, with that same rigidified line.
For each chosen $\ell$, its tangent lattice is the exact image of
a product of integral elliptic tangent lattices under an isogeny
whose kernel is killed by an integer prime to $\ell$.
\end{proposition}
\begin{proof}
Choose a CM elliptic curve with period lattice homothetic to
$\mathcal O_F$. Use this one curve with the two conjugate
$\mathcal O_F$-actions. The pure vectors in the preceding lemma
give holomorphic homomorphisms from these curves to the torus.
Choose an $\mathcal O_F\otimes\mathbf Z_{(\ell)}$-basis of
$\mathcal O_E\otimes\mathbf Z_{(\ell)}$ and clear denominators
prime to $\ell$. Its products with the $x_i$ span a sublattice
of finite index prime to $\ell$. The product map $\psi$ is
therefore an isogeny, and some integer $N$ prime to $\ell$
kills its homology cokernel.

Polarized complex tori and their homomorphisms are algebraic
\cite[Theorems~2.8--2.9]{cm-milne-av}.
The finite kernel of $\psi$ is contained in the algebraic
$N$-torsion of the elliptic product, so the quotient descends to
$\overline{\mathbf Q}$
\cite[Remark~8.12]{cm-milne-av}.
Homomorphisms between abelian varieties over an algebraically
closed field do not increase after extending that field
\cite[Corollary~16.4]{cm-milne-av}; hence the coefficient action
and the principal polarization descend too. For
$\overline{\mathbf Q}\subset\mathbf C$, define the homomorphism
over a finitely generated regular subextension first.
To retain the specified line, first choose an algebraic line
inducing this polarization
\cite[Proposition~13.6, in characteristic zero]{cm-milne-av}.
Its Tate pairing is alternating because it becomes the prescribed
polarization after complex base change; one may choose an
auxiliary prime different from $2$.
Adjust it by a square root in $\operatorname{Pic}^0$ to make it
symmetric. Its difference from the specified symmetric complex
line is in $\operatorname{Pic}^0[2]$, hence is algebraic.
Choose the comparison to respect the zero rigidification.
All this data, and any prescribed finite level structure, descend
to one number field.

CM elliptic curves acquire good reduction everywhere after a
finite extension \cite[Theorem~12.14]{cm-milne-mf}.
Good reduction is invariant under isogeny
\cite[Theorem~1 and Corollary~2, p.~493]{cm-serre-tate},
so the quotient also has good reduction everywhere.
The Moret--Bailly theorem extends the specified symmetric
rigidified line after another finite extension
\cite[Theorem~3.4(i)]{pazuki-theta-faltings}.
Its identity component is the same abelian scheme, by the
N\'eron comparison already used in
Theorem~\ref{geo-classical-tangent-slope}.
The vertical divisor and zero rigidification argument there
preserves the line and its canonical metric exactly.

Finally $N\psi^{-1}$ sends period lattices into period lattices
and is an algebraic homomorphism $\chi$.
Both maps extend to the good models by the N\'eron mapping
property \cite[Definition~5.1 and Example~5.3]{cm-conrad-minimal}.
Their identities $\chi\psi=[N]$ and $\psi\chi=[N]$ extend.
The coefficient action and its polarization adjunction identities
extend by the same uniqueness. On products of the good model,
the zero rigidification again removes the vertical divisor in
the corresponding cubical line identifications.
At a place above $\ell$, differentiation gives inverse integral
maps $d\psi$ and $N^{-1}d\chi$. Thus the product tangent basis
is a basis of the prescribed integral lattice, without an index.
\end{proof}

\begin{lemma}\label{geo-cm-elliptic-period}
For a CM elliptic curve by $\mathcal O_F$, choose after finite
extension an everywhere good model and an integral tangent basis
$v$. In the metric $\|v\|_\sigma^2=\pi H_\sigma(v,v)$, put
$n_\sigma=\|v\|_\sigma^2$. Then
\[
 -\frac1{[M:\mathbf Q]}\sum_{\sigma:M\hookrightarrow\mathbf C}
       \log n_\sigma
 =x_d:=\gamma+\log\frac{2\pi}{\sqrt d}
                  -\frac{L'(1,\chi_{-d})}{L(1,\chi_{-d})}.
\]
Finite extension and reversal of the CM action preserve this mean.
\end{lemma}
\begin{proof}
The invariant differential line can be principalized after finite
extension: if an ideal $I$ satisfies $I^a=(b)$, adjoining an
$a$th root of $b$ makes its extension principal, since the group
of fractional ideals is torsion free. Let $\omega(v)=1$ and write
its period lattice as $\omega_1(\mathbf Z+\tau\mathbf Z)$,
with $y=\operatorname{Im}\tau>0$. The principal Riemann form gives
\[
 n_\sigma=\frac{\pi}{|\omega_1|^2y},\qquad
 \Delta_\omega(\sigma)
   =(2\pi)^{12}\omega_1^{-12}\eta(\tau)^{24}.
\]
The second formula is the classical discriminant product
\cite[Theorem~4.21]{cm-milne-mf}.
Good reduction and the integral differential make
$\Delta_\omega$ a unit at every finite prime, including $2$ and $3$
\cite[Example~5.3]{cm-conrad-minimal}.
Consequently
\[
 \log n_\sigma
   =-\log(4\pi y|\eta(\tau)|^4)
                         +\tfrac16\log|\sigma\Delta_\omega|,
\]
and the last term has mean zero.

The CM class theorem makes the ideal classes occur with equal
multiplicity among the embeddings
\cite[Theorem~12.18]{cm-milne-mf}.
For
$Z_\tau(s)=\sum' y^s/|m+n\tau|^{2s}$, the Kronecker limit formula is
\[
 Z_\tau(s)=\frac{\pi}{s-1}
       +2\pi\left(\gamma-\log2
                    -\log(\sqrt y|\eta(\tau)|^2)\right)+O(s-1)
 \quad\cite[\S2]{cm-zagier}.
\]
If $h_F$ is the class number and $w_F$ the number of roots of unity,
the ideal norm decomposition gives
\[
 \zeta_F(s)=\frac1{w_F}\left(\frac2{\sqrt d}\right)^s
                       \sum_{[\mathfrak a]}Z_{\tau_{\mathfrak a}}(s).
\]
Compare the constant term divided by the residue with
$\zeta_F(s)=\zeta(s)L(s,\chi_{-d})$. This gives
\[
 \frac1{h_F}\sum_{[\mathfrak a]}
       \log(4\pi y_{\mathfrak a}|\eta(\tau_{\mathfrak a})|^4)=x_d.
\]
Conjugation permutes the ideal classes and preserves this real
modular expression. The product formula and uniform class
multiplicities prove the asserted mean.
\end{proof}

\begin{proposition}\label{geo-cm-flat-degree}
Suppose the determinant family of
Lemma~\ref{geo-flat-determinant-bridge} is algebraically defined,
is constant in the period trivialization, and has the following
local norm compatibility: at each place its prescribed flat norm
at the point under comparison equals the norm at every allowed
CM point above, computed on the same integral eigenlattices and
canonical metric. This equality is independent of which CM
point is chosen at that place.
For constant signature $(p,q)$ its actual normalized degree is
\[
 \frac{p-q}{2}P_E,\qquad
 P_E=\gamma+\log(2\pi)-\log\operatorname{rd}(E)
                     -\frac{L'(1,\chi_{-d})}{L(1,\chi_{-d})}.
\]
If the finite logarithmic norms are corrected by subtracting
$q(\log|\mathfrak D_E|-\log|d|)$, the degree is instead
\[
             \frac{p-q}{2}P_E
                      +q(\log d-\log\operatorname{rd}(E)).
\]
These are the precise inputs of the calculation.
Theorem~\ref{geo-flat-family-generator} constructs the generic family
from the prescribed model, and Corollary~\ref{geo-flat-norm-assembly}
supplies the norm comparison from its exact integral identifications.
\end{proposition}
\begin{proof}
Fix an $E$-basis $e$ of the homology lattice, with Gram matrix $G_e$.
At a pure CM point write $x=Te$ and $t=\det T$.
For the algebraic eigenframes $v,w$ supplied by the elliptic product,
let $\varphi(v,w)$ denote the scalar of the determinant section.
The CM elliptic tangent vectors have period scalars $z_+,z_-$.
In the pure basis, the period map on the positive space and the
dual negative space has blocks
$z_+I_p$ and $z_-^{-1}G_-^{-1}$.
If $\rho$ is the constant in the flat period trivialization, then
\[
 \varphi(v,w)\tau(t)z_+^p=\rho z_-^q\tau(\det G_-).
\]
Thus
$r=\varphi(v,w)\tau(t)/\tau(\det G_-)$ is a nonzero algebraic
number and satisfies $rz_+^p=\rho z_-^q$.
This same $r$ applies at every CM point. One curve with opposite
actions can be used for both signs, so even $z_+=z_-$ is possible.

There is no extra metric constant in the product Gram matrix.
Indeed, with $\omega=(1+\sqrt{-d})/2$, the elliptic Hermitian
pairing is
\[
 \begin{aligned}
 h_F^\epsilon(x,y)
 &=\epsilon\sqrt{-d}\operatorname{Tr}_{E/F}h(x,y)\\
 &=\epsilon\bigl(B(\omega x,y)-B(x,y)+\omega B(x,y)\bigr),
 \qquad B=\operatorname{Tr}_{E/\mathbf Q}h.
 \end{aligned}
\]
Polarization adjunction therefore gives
$\langle d\psi_xv,d\psi_yv\rangle
=n_\sigma\,\sigma(h_F^\epsilon(x,y))$.
For an integral basis $b_i$ of $\mathcal O_E$ and its absolute
trace-dual $b_i^*$, the eigenprojection is
$P_\tau=\sum_i\tau(b_i^*)\rho(b_i)$.
The trace-dual identities collapse the two Gram sums, giving
\[
 \Gamma_+=n_+\epsilon_+\tau(\sqrt{-d}\,G_+),\qquad
 \Gamma_-=n_-\epsilon_-\bar\tau(\sqrt{-d}\,G_-).
\]
Here $\epsilon_\pm=\pm1$ and $\bar\tau=\tau\circ\overline{(\cdot)}$.
As
$\det G_+\det G_-=t\bar t\det G_e$, the complex logarithmic norm is
\[
 a_\sigma=\log|\sigma r|
 +\frac p2(\log n_{+,\sigma}+\tfrac12\log d)
 -\frac q2(\log n_{-,\sigma}+\tfrac12\log d)
 +\tfrac12\log|\sigma\tau(\det G_e)|.
\]

Fix a number field containing the given coefficient field
(in particular $\tau(E)$), the elliptic data and $r$.
At each $\ell$ choose the local pure basis in
Lemma~\ref{geo-pure-cm-local-basis}. Proposition
~\ref{geo-cm-integral-product} supplies its exact integral product
tangent basis. Proposition~\ref{geo-relative-different-eigenframe}
then supplies the integral eigenframes: multiply the positive
and negative projected frames by $\delta$ and $\bar\delta$,
where $\delta$ generates $\mathfrak D_{E/F}$ locally.
Using
$\log|j\delta|=\log|j\bar\delta|
=D_\ell(j)-\log|j\sqrt{-d}|$, their logarithmic norm is exactly
\[
 a_{\ell,j}
 =\log|jr|-\log|jt_\ell|-qD_\ell(j)
             +(q-p)(D_\ell(j)-\log|j\sqrt{-d}|).
\]
The transition determinant also satisfies
\[
 \log|jt_\ell|+\log|j\bar t_\ell|
 +\log|j\det G_e|=-(p+q)D_\ell(j).
\]
It has norm one at all embeddings for all but finitely many
$\ell$, because $e$ is then itself an integral local basis.

Average over all embeddings, with normalized absolute values.
Conjugation permutes the embeddings, so the two transition terms
have equal means. The product formula cancels $r$ and $\det G_e$.
The remaining three sums are
\[
 \overline{\log n_+}=\overline{\log n_-}=-x_d,\qquad
 \sum_\ell\overline{D_\ell}=-\log\operatorname{rd}(E),\qquad
 \sum_\ell\overline{\log|\sqrt{-d}|_\ell}=-\tfrac12\log d.
\]
The negative of the total logarithmic norm is consequently
$(p-q)(x_d+\tfrac12\log d-\log\operatorname{rd}(E))/2$.
This is the asserted degree. Subtracting the specified finite
correction increases the degree by
$q(\log d-\log\operatorname{rd}(E))$.
Each local CM calculation may use its own finite extension of
the fixed number field. Normalized means descend to that field;
there is no need for one extension containing every local CM point.
\end{proof}

\paragraph{Assembling the prescribed local norms.}
\begin{corollary}\label{geo-flat-norm-assembly}
Suppose the data of Theorem~\ref{geo-flat-family-generator} extend
to a normal integral proper flat
model $\mathcal Y/\mathcal O_K$ with nonempty connected special
fibres. Suppose its prescribed integral conjugate determinant
is an invertible sheaf $\mathcal M_\tau$ restricting to
$\mathcal D_\tau$, and that integral point pullbacks give
the same eigenlattices as the prescribed norms.
Assume the actual Hodge and conjugate filtrations, coefficient
freeness and determinant comparison of
Proposition~\ref{geo-frobenius-determinant-torsion} identify
this line on every special-fibre test curve, including at
ramified coefficient primes.
Then a generic generator has point-independent finite norms,
zero logarithmic norm at all but finitely many places.

Retain the other arithmetic and metric hypotheses of
Proposition~\ref{geo-cm-flat-degree}. If the CM constructions give
integral points on this same component, whose original line and tangent
lattice identify with the product-isogeny model of
Proposition~\ref{geo-cm-integral-product}, then the local-norm comparison
required in Proposition~\ref{geo-cm-flat-degree} follows. Its signed determinant degree is therefore
$(p-q)P_E/2$.
\end{corollary}
\begin{proof}
The stated curvewise identification makes the degree of
$\mathcal M_\tau$ zero on every proper integral curve in each
special fibre. Apply Theorem~\ref{geo-vertical-line-rigidity}
at a finite place $v$: for some integers $M_v>0,a_v$,
$s^{\otimes M_v}/\pi_v^{a_v}$ generates the tensor power
on the whole localized model. Every integral point consequently
has logarithmic norm
$(a_v/M_v)\log|\pi_v|_v$, also after a ramified extension.
The divisor of $s$ has finitely many components, so this is zero
at almost all places. Theorem~\ref{geo-flat-family-generator}
gives constant complex norms. Evaluating these same norms
at the supplied integral CM points proves the required
comparison, and the CM degree formula applies.
\end{proof}

The remaining construction concerns the actual arithmetic model,
its curvewise determinant identifications and the integral CM
maps to its selected component. Good reduction of an abstract
CM variety and a generic level point alone do not supply these
identifications.



\subsection{Positivity and birationality for the actual cusp systems}
\label{actual-geometric-positivity}

Use an original constructed Host and its smooth projective toroidal
component $T$ over the prescribed coefficient field $K$. Let
$\pi:Y\to T$ be the original full flag space of the eigenspaces of
the semiabelian Lie bundle, $B$ its reduced toroidal boundary, and
$U=Y\setminus B$. Retain the original automorphic line bundle
$\mathcal L(a)$ for a weight $a=(a_{\tau,i})$, $0\le i<r=p+q$.
In the arithmetic application $p,q\ge1$. Put
\[
 f=|\Psi|,\qquad D=f\binom r2,\qquad
 \operatorname{gap}(a)=\min_{\tau,\ i<r-1}(a_{\tau,i}-a_{\tau,i+1}).
\]
The original canonical weight is independent of $\tau$ and equals
\begin{equation}\label{actual-positive-kappa}
 \kappa_i=\begin{cases}
 q-p+1+2i,&i<p,\\
 2i-3p-q+1,&i\ge p.
 \end{cases}
\end{equation}
This is the source's weight, with its original order and signs.
No tensor power replaces $\mathcal L(a)$ in the conclusions.

\begin{theorem}\label{actual-positive-cusp-theorem}
For these actual models the following assertions hold.
\begin{enumerate}
\item If $\operatorname{gap}(a)\ge0$, then $\mathcal L(a)$ is nef.
If $\operatorname{gap}(a)\ge1$, it is also big, and
$(c_1(\mathcal L(a))^{\dim V}\cdot V)\ge1$ for every
positive-dimensional integral $V\subset Y_{\mathbb C}$ meeting $U$.
\item Put $m_1=\binom{D+1}{2}+1$ and $m_2=\binom{D+2}{2}$.
If $\operatorname{gap}(a-\kappa)\ge m_1$, the original cusp
linear system has no base point on $U$. If
$\operatorname{gap}(a-\kappa)\ge m_2$, that system separates
every pair of distinct complex points of $U$ and is birational.
\item Consequently the original bound
$\operatorname{gap}(a)\ge(D+3)^2+2r$ supplies the source's
birationality condition. It applies to both prescribed weight
systems whenever the original weight-size condition holds.
\end{enumerate}
These are the actual nef, big, open generation and birational
suppliers. They assert no bound for integral generator norms,
coefficient-field degrees or arithmetic heights.
\end{theorem}

\begin{proof}
We may check the geometric assertions after an embedding $K\hookrightarrow
\mathbb C$. The constructed component is geometrically integral.
All line bundles and section spaces below descend to the original $K$.

\paragraph{Nefness with the determinant characters retained.}
The canonical Hodge extension on $T$ is the dual of the original
semiabelian Lie bundle. Lan's algebraic and analytic comparisons
identify these very extensions\cite{positive-lan-kuga,positive-lan-analytic}.
The boundary monodromy is unipotent, and $T$ is smooth projective
with normal-crossing boundary. Hodge semipositivity therefore makes
the Hodge bundle and each of its exterior powers nef: for the latter,
apply the same theorem to the exterior power of the polarized
weight-one variation\cite{positive-fujino-fujisawa}.
This proves nefness on the whole compactification, including curves
contained in the boundary.

For a head step weight, $\mathcal L$ is the dual determinant of
the corresponding tautological subbundle of $\operatorname{Lie}_\tau$.
The eigenspaces are quotient direct summands of the Hodge bundle.
The line in question is therefore a quotient of the corresponding
exterior power of the pulled-back Hodge bundle, hence nef. The conjugate tail steps have
the identical description. A weight constant on all $r$ entries
at $\tau$ gives the power of
\[
 (\det\operatorname{Lie}_\tau)^{-1}
       \otimes\det\operatorname{Lie}_{\bar\tau}
       =\det\mathcal H^1_{\mathrm{dR},\tau}.
\]
The equality follows from the Hodge filtration and polarization.
Its monodromy determinant preserves an $\mathcal O_E$-lattice
and has absolute value one at every embedding. It is thus a root
of unity, of order bounded for the fixed field $E$. Boundary
monodromy has determinant one. Its canonical extension is a
torsion line bundle, so either sign of this step is numerically
trivial. No uniform bound on that order is used.

Here is the exact weight decomposition. Put
$d_i=a_i-a_{i+1}$. For $i<p$ let $h_i$ be one on the first
$i+1$ entries and zero elsewhere; for $i\ge p$ let $t_i$ be
minus one on the entries after $i$, and zero elsewhere. Then
\[
 a=a_p{\bf1}+\sum_{i=0}^{p-1}d_i h_i
+\sum_{i=p}^{r-2}d_i t_i.
\]
Checking each coordinate proves the identity. Every $d_i\ge0$,
so the preceding nef and numerically trivial factors prove the
first assertion. The argument is applied independently at each $\tau$.

\paragraph{Strict gaps are positive at every interior subvariety.}
Let $s_i=r-1-i$, let $\chi_i=1$ for $i<p$ and $\chi_i=-1$
otherwise, and put
\[
 \rho_i=\begin{cases}i-p,&i<p,\\ i-p+1,&i\ge p.\end{cases}
\]
The bundle $\mathcal L(-\rho)$ is relatively ample for $\pi$,
by the Pl\"ucker embeddings of the two full flags. The bundle
$\mathcal L(\chi)$ is the pullback of the determinant Hodge
bundle $\omega$ on $T$. This is an identity of the prescribed
characters, not a change of metric or level.

Choose an ample line bundle $A$ on $T$ and $t>0$ so that
$H=\pi^*A^t\otimes\mathcal L(-\rho)$ is ample on $Y$.
The original toroidal map $v:T\to T^{\min}$ is an isomorphism
over the interior, and $\omega$ is the pullback of an ample
line bundle on $T^{\min}$, with a positive power if the level
is not neat\cite[Theorem 7.2.4.1(1)--(3)]{positive-lan-pel}.
For any $y\in U$, Serre generation applied to
$v_*A^{-t}$ gives some $l>0$, a multiple of that Hodge power, and a section of
$\omega^l\otimes A^{-t}$ nonzero at $\pi(y)$. Its pullback
has effective divisor $E_y$ avoiding $y$.

Choose $b$ so large that $\mu=bs-l\chi+\rho$ is dominant;
for example $b\ge2l+2$ suffices. The exact line identity is
\[
 \mathcal L(s)^b
  =(H\otimes\mathcal L(\mu))\otimes\mathcal O(E_y).
\]
The first factor is ample, since $\mathcal L(\mu)$ is nef.
If $V$ passes through $y$, it is not contained in $E_y$.
Restrict this identity to $V$. Expanding the difference of
the top self-intersections as a telescoping sum of nef classes
intersected with the effective divisor gives
\[
 b^{\dim V}(c_1(\mathcal L(s))^{\dim V}\cdot V)
   \ge(c_1(H\otimes\mathcal L(\mu))^{\dim V}\cdot V)>0.
\]
The intersection on the left before multiplication by $b^{\dim V}$
is an integer, hence at least one. The same identity also proves
bigness. If $\operatorname{gap}(a)\ge1$, then $a-s$ is
dominant. Tensoring with its nef line bundle preserves the
intersection lower bounds and bigness. This proves assertion 1.

\paragraph{The exact canonical bundle and cusp reading.}
The extended Kodaira--Spencer isomorphism on the constructed
toroidal model identifies its logarithmic cotangent bundle with
the prescribed polarized Hodge tensor\cite[Theorem 6.4.1.1(4)]{positive-lan-pel}.
Taking its determinant and the relative flag canonical bundle gives
\begin{equation}\label{actual-positive-canonical}
                 \omega_Y(B)=\mathcal L(\kappa).
\end{equation}
For a direct sign check, on the first Lie flag the relative
canonical exponents are $p-1-2i$, and the base canonical exponent
is $-q$; on the conjugate flag they are $q-1-2j$ and $-p$.
The original reversed conjugate weight convention gives exactly
\eqref{actual-positive-kappa}. Thus
\[
 H^0(Y,\omega_Y\otimes\mathcal L(a-\kappa))
                =H^0(Y,\mathcal L(a)(-B)).
\]
In the original normalized analytic reading these are cusp
sections: the subcanonical extension is the canonical extension
with the negative reduced boundary, and its algebraic sections
read as cuspidal automorphic forms
\cite{positive-lan-analytic,positive-lan-stroh}.
Their fixed holomorphic highest-weight equations give an
infinitesimal character, hence finiteness under the center of
the enveloping algebra; canonical-extension
coordinates give moderate growth. The group has compact real
center, as in the already verified raw-theta cusp argument.
The rapid-decay theorem therefore places these very sections
in the original bounded cusp space\cite{positive-cusp-decay}.
Only this inclusion is needed for generation and the birational
subsystem; no integral section-lattice identification is inferred.

\paragraph{Effective generation and separation.}
Set $N=\mathcal L(a-\kappa)$ and suppose
$\operatorname{gap}(a-\kappa)\ge m$. Then
$N=\mathcal L(s)^m\otimes\mathcal L(a-\kappa-ms)$,
where the second factor is nef. Assertion 1 gives
$(c_1(N)^k\cdot V)\ge m^k$ for every $k$-dimensional
integral $V$ through any prescribed point of $U$.

Apply Koll\'ar's effective adjoint theorems for a smooth variety
and a nef and big $N$\cite[Theorems 5.8--5.9]{positive-kollar}.
For generation take $m=m_1$ and $c(k)=m_1-\tfrac12$.
The strict intersection hypotheses hold, and
\[
 \sum_{k=1}^D\frac{k}{c(k)}
       =\frac{m_1-1}{m_1-\tfrac12}<1.
\]
For separation of two distinct points take $m=m_2$ and
$c(k)=m_2-\tfrac12$. Since
$2^{1/k}\le1+1/k$ for integral $k\ge1$,
\[
 \sum_{k=1}^D\frac{k2^{1/k}}{c(k)}
 \le\frac{D(D+3)/2}{m_2-\tfrac12}
 =\frac{m_2-1}{m_2-\tfrac12}<1.
\]
The conclusions apply to $\omega_Y\otimes N$, which is
exactly the cusp bundle in \eqref{actual-positive-canonical}.
The resulting map on $U$ separates all complex points, so in
characteristic zero it is birational onto its image. Proper
cohomology base change descends generation and the equality
of ratio fields to $K$; no coefficient-field enlargement is
introduced in the conclusion.

Finally, the consecutive gaps of $\kappa$ are $-2$ within
each block and $2r-2$ at the middle break. Hence the original
bound in assertion 3 gives
$\operatorname{gap}(a-\kappa)\ge(D+3)^2+2\ge m_2$.
This proves assertion 3 with room to spare, including $D=1$.
The original weight-size inequalities supply that displayed gap
for both $u_1$ and $u_2$; their exact consuming statement and
size quantifiers are recorded in the companion source certificate.
\end{proof}

For the actual Hilbert--Samuel consumer, a pullback of the
strict-gap line is nef and big under its prescribed birational
modification. Principalizing the other original cusp base ideal
gives a generated line whose map is birational; it too is nef
and big. Thus these geometric hypotheses and both original
birationality conditions are supplied. The common integral
section lattices, their short spanning bounds and the numerical
arithmetic remainder are not inferred from geometric positivity.

\subsection{The actual flat determinant degree}
\label{sec:actual-flat-degree}

The arithmetic model and period presentations constructed above supply
the hypotheses of the determinant calculation.  We record this
application explicitly, because a flat line of the correct generic
weight alone does not supply its prescribed finite norms.

\begin{theorem}\label{thm:actual-flat-degree}
Retain the original host: $E$ is CM, contains
$\mathbf Q(\sqrt{-d})$, $-d$ is its prescribed odd fundamental
quadratic discriminant, and $E/E^+$ is unramified at finite primes.
Let $p+q=r\ge1$ and use an actual sized good cover\cite{current-construction-record}, with its original
coefficient field $K$, level, primitive line, and component marking.
Write
\[
 P_E=\gamma+\log(2\pi)-\log\operatorname{rd}(E)
          -\frac{L'(1,\chi_{-d})}{L(1,\chi_{-d})},\qquad
 a=-\frac{p-q}{2}P_E-q\bigl(\log d-\log\operatorname{rd}(E)\bigr).
\]
For every original component $Z$ and every prescribed half-embedding
$\tau$, the original $\mathrm{FlatSec}(Z,\tau)$ is nonempty, and
every such section has degree $-a$.  Thus the actual source
assertion $\mathrm{FlatExistsAt}(a)$ holds, and its averaged
coefficient satisfies $c_1(\tau)=a$ for every $\tau$.

Under the original numerical bounds
\[
 |P_E|\le3L_b,\qquad 0\le\log d\le L_b,\qquad
 0\le\log\operatorname{rd}(E)\le2L_b,
\]
one has $|c_1(\tau)|\le\tfrac72rL_b<9rL_b$.
There is no extension of $K$ in this conclusion.
\end{theorem}

\begin{proof}
\emph{The same global component.}
The established global domination theorem\cite{current-construction-record} gives the normal projective
flat $\mathcal O_K$-model carrying the original semiabelian family,
coefficient action and primitive cubical line.  Normalize the closure
of the selected generic component.  The original
$\mathrm{FnAlgClosed}$ clause makes its generic fibre geometrically
integral.  Its proper Stein factor is $\operatorname{Spec}\mathcal O_K$:
the finite algebra of global functions embeds in the generic
constant field $K$ and is integral over the integrally closed
$\mathcal O_K$.  Its nonempty special fibres are consequently
connected.  Restriction and normalization retain the same family and
rigidified line.

\emph{The required boundary fibres.}
The original principal polarization gives the exact Raynaud formula
\[
       1=[X:\phi(Y)]^2\deg\lambda_B.
\]
Hence the abelian quotient is principally polarized and the period
map $\phi$ has index one.  The transported primitive identity gives
the actual line $q_B^*(1,w_0)^*\mathcal P_{\lambda_B}$, with its
specified rigidification; this is precisely the calculation in
the actual DVR uniformization theorem\cite{current-construction-record}.  These are the
self-dual-quotient and line-pullback hypotheses in
Theorem~\ref{geo-biextension-conjugate-filtration}.
They hold at every geometric fibre.  Equivalently, they may be
checked on a trait through a closed point with characteristic-zero
generic point: such a trait is obtained by a horizontal curve in the
flat finite-type model and normalization.  The Raynaud quotient of
that trait specializes to the same torus quotient, and the primitive
identity fixes the same quotient line.  Closed points suffice for
these bundle identities and freeness, since the model is Jacobson.

The global normalized adjunctions therefore construct the actual
first-order Hodge extension, with conjugate-dual kernel, as in
Proposition~\ref{geo-ramified-hodge-determinant}.  The proper and
toric comparisons prove coefficient freeness; no freeness of the
individual tangent eigensheaves at ramified primes is assumed.
The proposition supplies the integral determinant line whose
pullback is the exact original tangent eigenlattice at every
integral point, after arbitrary ramified extensions too.

\emph{A section with the original local norms.}
The half-level has $M\ge3$ and satisfies
Proposition~\ref{geo-half-level-determinant}.  Apply
Theorem~\ref{geo-flat-family-generator} on the selected generic
component to obtain an algebraic generator of
\[
      \det\operatorname{Lie}(G)_\tau\otimes
       \det\operatorname{Lie}(G)_{\bar\tau}^{-1}
\]
with constant original complex norm.  The Hodge and conjugate
filtrations make its character determinant have degree zero on
every special-fibre curve.  Theorem~\ref{geo-vertical-line-rigidity}
then makes its normalized finite order independent of the integral
point.  These orders vanish at all but finitely many places.
The original local models have the same orders: their generic
points lift to the proper dominating model by the valuative
criterion, and its family comparisons identify the actual tangent
lattices.  Thus this is the original constant section, not a
substitute norm on its generic line.

\emph{Its degree.}
The fine labelled moduli family and every required pure period
presentation are supplied by
the actual normalized-member theorem\cite{current-construction-record}.  Consequently
Propositions~\ref{geo-cm-period-point-descent}
and~\ref{geo-cm-integral-map-extension} give the exact CM points
on this component and their integral comparisons.  Apply
Corollary~\ref{geo-flat-norm-assembly} and
Proposition~\ref{geo-cm-flat-degree}.  The uncorrected signed
determinant degree is $(p-q)P_E/2$.  The original
$\mathrm{flDefect}$ changes it by
$q(\log d-\log\operatorname{rd}(E))$; hence it is $-a$.
This is exactly the source's passage from a constant frame to
$\mathrm{FlatSec}$, including its sign
\cite{theta-flat-source}.

Any two generators differ by $K^\times$, so the product formula
makes their degrees equal.  More generally, any original constant
section is a nonzero scalar multiple: its quotient has constant
complex norm and is holomorphic on the connected component, so it
is constant, and the original constant-field clause places the
scalar in $K$.  This proves the universal quantifier over sections
in $\mathrm{FlatExistsAt}$.  The definition of $c_1$ is the
negative average of these degrees, so $c_1(\tau)=a$ independently
of the component and half-embedding.

Finally $|p-q|\le r$, $q\le r$, and
$|\log d-\log\operatorname{rd}(E)|\le2L_b$ give
$|a|\le(3/2+2)rL_b$.  The numerical hypotheses are precisely
the retained host and weight bounds in the original consumer.
\end{proof}

The theorem closes the actual flat-degree supplier.  The parallel
integral twists are supplied in Theorem~\ref{thm:actual-short-parallel-twist};
the original theta key section and the full shifted short families
are supplied in Section~\ref{sec:actual-key-short-span}.

\subsection{The actual abelian quotients at boundary points}
\label{sec:actual-point-tangent-degrees}

The flat determinant construction applies to the abelian quotient at
an actual boundary point. This supplies the signed degree comparison
used in the point-height argument. The integral image of the tangent
lattice need not be saturated; its index has a favourable sign and is
retained below.

Keep the original host $E$, with $f=[E^+:\mathbb Q]$, the prescribed
quadratic subfield $\mathbb Q(\sqrt{-d})$, and the hypotheses that
$-d$ is the original odd fundamental discriminant and
$E/E^+$ is unramified at every finite prime. Retain the original
primitive element $w=(1+\sqrt{-d})/2$, with $w+\bar w=1$. Put
\[
 P_E=\gamma+\log(2\pi)-\log\operatorname{rd}(E)
             -\frac{L'(1,\chi_{-d})}{L(1,\chi_{-d})}.
\]
Let an algebraic point of the original compactification have toric
$\mathcal O_E$-rank $r$ and signature $(p,q)$, with
$0\le r\le\min(p,q)$. Its abelian quotient $A$ has signature
$(p-r,q-r)$. After a finite extension $L$ containing the coefficient
embeddings and admitting semistable reduction, let $\bar V$ be the
Lie lattice of the N\'eron identity component of $A$, with the
prescribed metric $\|v\|^2=\pi H(v,v)$. At every conjugate pair of
embeddings, write $W_\tau,W'_\tau$ for its saturated integral
eigenlattices, with inherited metrics. The chosen half has rank
$m-r$, where $m$ is either $p$ or $q$; the other half has rank
$p+q-m-r$. All degrees below are divided by $[L:\mathbb Q]$.

\begin{theorem}[Conjugate degrees at the actual points]
\label{thm:actual-point-conjugate-degrees}
For these lattices,
\begin{equation}\label{eq:actual-point-conjugate-degrees}
 \widehat{\deg}_{\rm n}W_\tau-
 \widehat{\deg}_{\rm n}W'_\tau
                    =\frac{2m-p-q}{2}P_E.
\end{equation}
This is an identity on the actual integral lattices, including the
places ramified in $E$ and the places above two. It is unchanged by
further finite extension. The extension used to realize the point
does not introduce a field-size term in this normalized identity.
\end{theorem}

\begin{proof}
\emph{The quotient and its polarization.}
The original primitive line and polarization restrict to the
semiabelian fibre. Its toric character module is a projective
$\mathcal O_E$-module of rank $r$, so every coefficient embedding
occurs $r$ times in its Lie algebra. Exactness of the Lie sequence
in characteristic zero gives the stated quotient signature. The
principal degeneration formula
\[
                 1=[X:\phi(Y)]^2\deg\lambda_A
\]
makes $\lambda_A$ principal. The original primitive identity gives
the rigidified line
$(1,\lambda_A\iota(w))^*\mathcal P_A$, with
$w+\bar w=1$. These are precisely the quotient and line
identifications in the arithmetic-point uniformization theorem
\cite{current-construction-record}; no independent polarization or
translate of its line is chosen.

\emph{Realization in the appropriate moduli space.}
If $p+q-2r=0$, both eigenlattices are zero and the assertion follows.
Otherwise fix one complex embedding and take the integral Betti
lattice of this principally polarized quotient, with its coefficient
action and perfect alternating form. Betti--\'etale comparison gives
the required complete Tate genus. The moduli construction is rerun
for this quotient's own integral Betti and Tate datum. Choose an
allowed odd half-level,
and pass to a finite field extension over which its finite torsion
labels and the required roots are defined. The generic PEL moduli
theorem then represents this very abelian variety, primitive line and
both entire finite labels. The component-field and global-model
constructions apply to its signature, including either definite
signature \cite{current-construction-record}. Their bounds are not
being used to bound the field of definition of the point. Passing
to the auxiliary cover may enlarge $L$ again; this is permitted
because the conclusion is invariant under normalized semistable
base change, with no bound on that extension required.

The global model is proper. Thus the resulting point extends over
every valuation ring of its number field. Its semiabelian family
has the same generic fibre as the prescribed semistable model of
$A$. The semistable N\'eron comparison, with the zero section fixed,
identifies their identity components and their integral Lie modules
\cite[Theorem 4.4 and Corollary 4.5]{conrad-semistable}.
The rigidified restriction argument in the theta construction
identifies the original lines and their canonical metrics. Hence
the integral eigenlattices and their Hermitian norms are exactly
the ones in the statement. This argument also applies after
ramified field extension; it does not compare only generic fibres.

\emph{Evaluation of the same determinant.}
Apply the actual flat-degree theorem to this component and to the
signature $(a,b)=(m-r,p+q-m-r)$. Its corrected flat degree is
\[
 \frac{a-b}{2}P_E+b(\log d-\log\operatorname{rd}(E)).
\]
On integral eigenframes, the prescribed finite logarithmic norms
add $b(\log|\mathfrak D_E|-\log|d|)$. Their normalized sum is
$b(\log d-\log\operatorname{rd}(E))$. Degree is the negative
sum of logarithmic norms, so this term cancels the displayed finite
correction. The exact determinant comparison
\eqref{geo-exact-conjugate-degree} therefore gives
$(a-b)P_E/2$ for the actual eigenlattice degree difference.
Since $a-b=2m-p-q$, this is
\eqref{eq:actual-point-conjugate-degrees}.

Eigenlattices are kernels of the finitely many integral coefficient
action equations. Flat base change preserves these kernels.
Semistable base change preserves the Lie lattice, and normalized
Hermitian degree is unchanged. This proves the field-extension
assertion, with every original finite correction retained.
\end{proof}

Let $V_\tau\subset W_\tau$ be the image of the original integral
tangent lattice in the selected quotient eigenspace. The quotient
map extends to the N\'eron model by its mapping property, so the
inclusion is integral. It has full generic rank. Set
\[
 g=f(p+q-2r),\qquad R=f(m-r),\qquad
 T=\sum_\tau\widehat{\deg}_{\rm n}V_\tau,\qquad
 b=\frac{f(2m-p-q)}2P_E.
\]
The coefficient summands are orthogonal for the prescribed Hodge
metric. Let $\delta\ge0$ be the normalized logarithmic index of
$\bigoplus_\tau(W_\tau\oplus W'_\tau)$ in $V$, and let
$\nu\ge0$ be the sum of the normalized logarithmic indices
$[W_\tau:V_\tau]$.

\begin{corollary}[The selected degrees for the actual family]
\label{cor:actual-selected-point-degrees}
For the original lattices,
\begin{equation}\label{eq:actual-selected-point-degrees}
 T\le12R\log\bigl(2+g+(b-\delta-2\nu)_+\bigr).
\end{equation}
If $B\ge\max\{128,4f(p+q)^2\}$ and
\[
 \Lambda_B=\log B+\log\operatorname{rd}(E)
                  +\left|\frac{L'(1,\chi_{-d})}{L(1,\chi_{-d})}\right|,
\]
then
\begin{equation}\label{eq:actual-selected-point-budget}
                       T\le36f(m-r)\Lambda_B.
\end{equation}
Neither statement assumes a bound on either integral index.
\end{corollary}

\begin{proof}
Write $d_V=\widehat{\deg}_{\rm n}V$. Orthogonality and the two
integral inclusions give
\[
                        d_V=2T-b+\delta+2\nu.
\]
The classical theta construction supplies the tangent-slope estimate
with constant $12$ on this exact principally polarized semistable
lattice and its metric. Applied to the selected integral images,
it gives $T\le12R\log(2+g+(-d_V)_+)$. If $T>0$, the displayed
identity gives $(-d_V)_+\le(b-\delta-2\nu)_+$; if $T\le0$,
\eqref{eq:actual-selected-point-degrees} is immediate. The case
$g=0$ has $T=R=0$ and needs no slope theorem.

The previously established period estimate is
$|P_E|\le C_f\Lambda_B$, with
$C_f=1+\gamma+|\log(2\pi)|<4$. Moreover
$g\le G=f(p+q)\le B\le e^{\Lambda_B}$ and
$|f(2m-p-q)|\le g$. Thus $b\le C_fG\Lambda_B$.
The selected-degree calculation in
Proposition~\ref{geo-selected-tangent-sum} gives
$T\le12\{1+\log(3+C_f)\}R\Lambda_B<36R\Lambda_B$.
The weak inequality in \eqref{eq:actual-selected-point-budget}
includes all zero-rank cases.
\end{proof}

This removes the conjugate-degree and selected-abelian-degree
hypotheses from the actual point-height application. The corrected fibre comparison is supplied in
Section~\ref{sec:geometry-hodge-specialization}. The regularized
point-height identity and final mean-volume estimate are separate
obligations; they do not follow from a tangent-degree bound alone.

\subsection{Specialization of the corrected tangent lattices}
\label{sec:geometry-hodge-specialization}

The correction used in the point-height argument is made on the original
normal compactification, before specializing to a boundary point. This
order matters. A semiabelian variety of dimension $g$ and torus rank $r$
has $p$-adic Tate module of rank $2g-r$. Its intrinsic torsion cannot be
substituted for the rank-$2g$ torsion of the generic abelian family.
We prove the required comparison by applying the finite-group
Hodge--Tate theorem on the original charts and then specializing the
original locally free modules.

Let $X$ be one of the original normal integral models, with function
field $K$ of characteristic zero, and let $G/X$ be its smooth
semiabelian family of positive relative dimension. Its generic fibre
is abelian. Retain its coefficient
action by $\mathcal O_E$, its rigidified cubical line bundle, and its
generic principal polarization. Write
\[
 \mathcal T=\operatorname{Lie}(G/X),\qquad
 \omega=\mathcal T^\vee.
\]
All coefficient embeddings used below are the original labels. Their
values on $\mathcal O_E$ are regular: they are integral over $\mathbb Z$
in the function field of a normal model.

For $p\in S$, put $n_p=p^{R_p}$, where the prescribed $R_p\ge4$,
and let $T_p=G_K[n_p](K)$. The original level makes all geometric
$n_p$-torsion rational over $K$. For $t\in T_p$, differentiate at the
identity the original Weil character $e_{n_p}(t,-)$ on $G[n_p]$.
As proved below, this gives a canonical local section
\[
 \delta_t\in\omega/n_p\omega.
\]
Define $\mathcal N_p\subset\omega$ to be the inverse image of the
span of these sections. Thus it contains $n_p\omega$, and on a chart
it is generated by arbitrary lifts of the $\delta_t$, together with
$n_p\omega$. This is precisely the cotangent module used in the
original construction; changing a lift changes it by $n_p\omega$.

Let $\mathcal T_{K,\sigma}$ be the $\sigma$-eigenspace of the
coefficient action on $\mathcal T_K$. Set
\[
 \mathcal Q_{p,\sigma}
   =\operatorname{im}\bigl(\mathcal N_p\longrightarrow
                          \mathcal T_{K,\sigma}^{\vee}\bigr),
 \qquad
 \mathcal D_{p,\sigma}=\mathcal Q_{p,\sigma}^{\vee}
       \subset\mathcal T_{K,\sigma}.
 \tag{HS1}
\]
We assume the original local-freeness condition: $\mathcal N_p$ and
each $\mathcal Q_{p,\sigma}$ are locally free of their prescribed
ranks near the fibre over $p$. In the selected half of the labels the
corrected module is
$\sigma(\mathfrak D_E)^{-1}\mathcal D_{p,\sigma}$; in the other
half it is $\mathcal D_{p,\sigma}$. Here $\mathfrak D_E$ is the
different of $E/\mathbb Q$, extended to the chart. Away from $S$ the
original corrected bundle is the ordinary integral eigenbundle.
Denote this globally prescribed bundle by $\mathcal F_\sigma$.

\begin{theorem}[The actual corrected fibre]
\label{thm:actual-hodge-specialization}
Suppose $S$ contains every rational prime dividing
$\operatorname{disc}(E)$. Let $z:\operatorname{Spec}\mathcal O_L\to X$
be an original arithmetic point, including a boundary point, with
its original coefficient labels $\sigma:E\hookrightarrow L$.
Let
\[
 \Lambda_\sigma
   =\operatorname{Lie}(G_z/\mathcal O_L)
                   \cap\operatorname{Lie}(G_{z,L})_\sigma,
 \qquad
 F_\sigma=(z^*\mathcal F_\sigma)
                   \subset\operatorname{Lie}(G_{z,L})_\sigma.
\]
The inclusion defining $F_\sigma$ is the original evaluation of
sections of the corrected bundle. There is a nonzero integral ideal
\[
 \mathfrak a_\sigma
   =\sigma(\mathfrak D_E)\mathcal O_L
                   \prod_{p\in S}(\kappa_{p,z}),
 \qquad
 \kappa_{p,z}=
 \begin{cases}
  4,&p=2,\\
  \zeta_{p,z}-1,&p\ne2,
 \end{cases}
 \tag{HS2}
\]
where $\zeta_{p,z}$ is the specialization of an original primitive
$p$th root of unity, such that
\begin{equation}\label{eq:actual-hodge-inclusions}
 \mathfrak a_\sigma F_\sigma\subset\Lambda_\sigma\subset F_\sigma.
\end{equation}
Consequently $\mathfrak a_\sigma\Lambda_\sigma\subset F_\sigma$ as
well. Its exact normalized norm is
\begin{equation}\label{eq:actual-hodge-ideal-norm}
 \frac{\log N_{L/\mathbb Q}\mathfrak a_\sigma}{[L:\mathbb Q]}
 =\log\operatorname{rd}(E)
  +\sum_{\substack{p\in S\\p\ne2}}\frac{\log p}{p-1}
  +2\,\mathbf 1_{2\in S}\log2
 \le 2\log\operatorname{rd}(E)
       +2\sum_{p\in S}\frac{\log p}{p-1}.
\end{equation}
Every lattice and label here is the original one. In particular,
the assertion applies without replacing an image lattice in an
abelian quotient by its saturation.
\end{theorem}

\begin{proof}
We separate the finite-group estimate from the specialization argument.

\emph{The original characters and their differentials.}
Work first over a normal affine chart $A\subset K$ of $X$ and put
$n=n_p$. The kernel $\Gamma=G[n]$ is separated, flat and quasi-finite
of finite presentation over $A$. At a point of the chart pass to
the henselization of its normal local ring. The finite-part theorem
for a separated quasi-finite scheme over a henselian local base
gives an open and closed finite subscheme $\mathcal H\subset\Gamma$
containing the whole closed fibre\cite{hodge-finite-part}. It is a
subgroup: multiplication and inversion preserve it on the closed
fibre, and the inverse image of the complementary open and closed
part in a finite scheme is empty if its closed fibre is empty.
It is finite flat because $\Gamma$ is flat.

The restriction of $e_n(t,-)$ to $\mathcal H_K$ is a rational point
of the finite Cartier dual $\mathcal H^D$. It extends uniquely over
the normal henselian ring: the images of its finite coordinate
algebra are integral over the base and lie in its fraction field.
The resulting character is regular near the identity of $\Gamma$.
Its generic rational function was defined before henselization;
faithfully flat descent of membership in the local ring therefore
gives a regular germ on the original $\Gamma$ at the identity.
Uniqueness follows from flatness and generic-fibre density.

Lift this germ through the closed immersion $\Gamma\hookrightarrow G$.
The cotangent sequence at the identity identifies
\[
 \omega_\Gamma=\operatorname{coker}
       \bigl([n]^*:\omega\longrightarrow\omega\bigr)
       =\omega/n\omega,
 \tag{HS3}
\]
since the differential of multiplication by $n$ is multiplication
by $n$. The differential of the lift is consequently independent
of its choice modulo $n\omega$. Its value is $\delta_t$.
The character has value $1$ at the identity, so its ordinary
differential there is its logarithmic differential.

This also checks the definition of $\mathcal N_p$ against the
original one: a regular function realizing the original generic
Weil character has differential congruent to $\delta_t$ modulo
$n\omega$, and every lift of $\delta_t$ occurs in this way at the
level of differentials. The zero character supplies $n\omega$.
Thus no new characters or tangent module have been introduced.
For an injective valued extension of $A$, the characters agree
after base change by their generic values and flatness. Their
differentials and the generated module agree as well. We use this
compatibility only for injective test extensions; specialization
with a kernel will be handled separately below.

\emph{A finite flat model of the full generic torsion.}
Let $C$ be a complete algebraically closed rank-one valued extension
of $K$, with valuation ring $\mathcal O_C$ and $v(p)=1$. The
generic fibre of $G_{\mathcal O_C}$ is still abelian. Put
$T=G_C[n](C)$ and let $\mathcal H\subset G_{\mathcal O_C}[n]$ be
its finite part containing the closed fibre. Write $H=\mathcal H(C)$.
Properness identifies $H$ with $\mathcal H(\mathcal O_C)$.

There is a finite flat commutative group scheme $\mathcal B$ whose
generic group is all of $T$: take one copy of $\mathcal H$ for each
coset of $H$ in $T$. Choose representatives $t_c$ with $t_0=0$ and define the sum
on these copies by
\[
 (c,h)+(d,h')=(c+d,h+h'+t_c+t_d-t_{c+d}).
 \tag{HS4}
\]
The last term belongs to $H$ and hence is an integral section of
$\mathcal H$. The cocycle identity proves associativity. Inversion
and the identity are defined by the same formula. The resulting
scheme is independent of representatives by translation, is killed
by $n$, and is finite flat of rank $|T|=n^{2\dim G_K}$.
The identity copy is $\mathcal H$, so
\[
 \omega_{\mathcal B}=\omega_{\mathcal H}
                      =\omega_{G_{\mathcal O_C}}/n\omega_{G_{\mathcal O_C}}.
 \tag{HS5}
\]
This is also the original induced finite Hopf order: a function on
$T$ belongs to it exactly when its restriction to every translate
of $H$ is regular on $\mathcal H$.

Here the original cubical line supplies a necessary integral
compatibility. If $\mathcal L$ is that rigidified line, set
\[
 \mathcal P(x,y)=\mathcal L_{x+y}\otimes\mathcal L_x^{-1}
                     \otimes\mathcal L_y^{-1}\otimes\mathcal L_0.
\]
The cubical structure gives its biextension laws: addition in
either variable tensors the corresponding lines, associatively
and compatibly with addition in the other variable. The
rigidification trivializes $\mathcal P(0,y)$ and
$\mathcal P(x,0)$. On $G[n]\times G[n]$, the two $n$-fold
addition laws consequently give two nowhere-vanishing
trivializations of the same line $\mathcal P(x,y)^{\otimes n}$,
one through $\mathcal P(nx,y)$ and one through
$\mathcal P(x,ny)$. Their ratio, in the order of the original
Weil convention, is a regular unit. This construction is made
with the universal torsion points and commutes with every base
change; it is a regular function on the product scheme, not only
a collection of functions at its rational points.

To check the normalization, choose lifts of $x$ and $y$ to the
theta group of $\mathcal L^{\otimes n}$ on a geometric generic
fibre. The cubical addition law writes the two translation
isomorphisms as the corresponding $n$-fold trivializations,
each multiplied by a scalar. In their ordered commutator the two
scalars cancel, leaving exactly their ratio. Hence the resulting
unit is the original $e_n(x,y)$, rather than an unspecified power
or scalar multiple of it. Its bilinear identities hold on the
generic fibre and extend by flatness. Restricting this construction
gives the required regular pairing on
$\mathcal H\times\mathcal H$\cite{hodge-specialization-source}.
For arbitrary $t,t'\in T$,
the mixed characters $e_n(t,-)|_H$ and $e_n(-,t')|_H$ extend over
$\mathcal H$, since their points in $\mathcal H^D(C)$ extend by
properness. On two coset copies the formula
\[
 e_n(t+h,t'+h')
   =e_n(t,t')e_n(t,h')e_n(h,t')e_n(h,h')
 \tag{HS6}
\]
therefore defines a regular pairing. The constant factor is a
root of unity and hence an integral unit. The identities needed
to glue and to make the pairing bilinear hold generically and
thus hold on the finite flat schemes. We have extended exactly
the original pairing to a homomorphism
$\lambda:\mathcal B\to\mathcal B^D$.

The original generic principal polarization makes $\lambda_C$
an isomorphism. We do not need to claim that $\lambda$ is an
integral isomorphism. The maps from integral points to $C$-points
of both finite proper schemes are bijections. It follows that
\[
 \lambda:\mathcal B(\mathcal O_C)
                       \longrightarrow\mathcal B^D(\mathcal O_C)
 \quad\hbox{is bijective}.
 \tag{HS7}
\]

\emph{The finite-group Hodge--Tate bound.}
For any commutative finite flat group $J/\mathcal O_C$ of $p$-power
order, Fargues's finite-group theorem says that the cokernel of
\[
 J(\mathcal O_C)\otimes_{\mathbb Z}\mathcal O_C
       \longrightarrow\omega_{J^D}
\]
is annihilated by every element of valuation at least $1/(p-1)$
\cite[\S5.4, Theorem~3]{hodge-fargues-finite}. His base-field
convention allows complete rank-one valuations that are not discrete
and imposes no condition on the residue field. The theorem includes
$p=2$; we retain the original weaker annihilator $4$ in that case.
Apply it to $J=\mathcal B^D$. By \textup{(HS5)--(HS7)}, its
Hodge--Tate differentials are exactly the original $\delta_t$ in
$\omega/n\omega$. With
$\kappa_p=\zeta_p-1$ for odd $p$ and $\kappa_2=4$, it gives
\[
 \kappa_p\omega_{\mathcal O_C}
             \subset\mathcal N_{p,\mathcal O_C}
             \subset\omega_{\mathcal O_C}.
 \tag{HS8}
\]
Here $v(\zeta_p-1)=1/(p-1)$ and $v(4)=2$ at $p=2$.

Return to the normal noetherian chart $A$. The generic full level
and the perfect generic pairing supply the chosen primitive roots,
which are integral on $A$ by normality. At a height-one prime over
$p$, embed the fraction field of its completed discrete valuation
ring into a complete algebraically closed valued field. This is an
injective test extension of $A$, so the preceding identification
of $\mathcal N_p$ applies. The extension of valuation rings is
faithfully flat. Thus \textup{(HS8)} descends to $A_{\mathfrak q}$.
At height-one primes not over $p$ we already have
$\mathcal N_p=\omega$, since $n$ is a unit. Near the fibre over
$p$, local freeness makes $\mathcal N_p$ reflexive. Intersecting
at height-one primes in a normal ring proves the original chart
inclusions
\begin{equation}\label{eq:hodge-chart-inclusions}
 \kappa_p\omega\subset\mathcal N_p\subset\omega.
\end{equation}
This argument has used a single finite level. It requires neither
a compatible tower of finite flat groups nor a replacement of the
given boundary point by an abelian variety.

\emph{Specializing the restricted dual.}
We give the algebra needed when $z$ has a kernel on a chart.
Let $A\to R$ be the resulting map to a characteristic-zero
valuation ring, let $C_R=\operatorname{Frac}(R)$, and write
$M=\mathcal T(A)$ and $N=\mathcal N_p(A)$ after shrinking so the
relevant modules are free. For the fixed label put
\[
 Q=\operatorname{im}(N\longrightarrow M_{K,\sigma}^{\vee}),
 \qquad D=Q^\vee.
\]
The original local-freeness condition makes $Q$ projective. Hence
$N\twoheadrightarrow Q$ is a split surjection and
$D\otimes_A R=(Q\otimes_A R)^\vee$.
The chart inclusions imply, in the original generic eigenspace,
\[
 M\cap M_{K,\sigma}\subset D,\qquad \kappa_p D\subset M.
 \tag{HS9}
\]

To identify the specialized subspace, choose a nonzero integer
$b_E$ clearing the denominators in the trace-dual integral basis
of $E$. The $\sigma$-projector, expressed using that dual basis,
is a polynomial in the coefficient action with coefficients in
$A[1/b_E]$; the same formula specializes to the $\sigma$-projector
over $C_R$. After inverting $b_E\kappa_p$, \textup{(HS9)}
identifies $D$ with that eigensummand of $M$. Both constants
specialize to nonzero elements of $C_R$. Thus the map obtained by
specializing $\kappa_p D\to M$ and dividing by $\kappa_{p,z}$
identifies $D\otimes_A C_R$ with the actual specialized
$\sigma$-eigenspace. Since $D\otimes_A R$ is projective, it
embeds in that space. This is the evaluation of the original
sections, because it agrees after multiplication by the same
nonzero scalar with evaluation of integral tangent sections.

Put $\Lambda_{\sigma,R}=M_R\cap M_{C_R,\sigma}$. For
$v\in\Lambda_{\sigma,R}$, evaluation on $N_R$ is integral. It
annihilates $\ker(N\to Q)\otimes_A R$: after passing to $C_R$
this kernel is in the annihilator of the same eigenspace, by the
projector formula, and an element of $R$ vanishing in $C_R$ is
zero. Evaluation therefore descends to an element of
$(Q\otimes_A R)^\vee=D\otimes_A R$. Under the preceding rational
identification this element is $v$ itself. Conversely
$\kappa_{p,z}(D\otimes_A R)$ lies in $M_R$ by \textup{(HS9)}
and remains in the same eigenspace. We obtain
\begin{equation}\label{eq:hodge-dual-specialization}
 \kappa_{p,z}D_R\subset\Lambda_{\sigma,R}\subset D_R.
\end{equation}
The integer $b_E$ was used only to verify identities over $C_R$.
It contributes no integral multiplier to this inclusion.

The embedded different is an invertible integral ideal, so its
inverse commutes with this specialization. For either original
half of the labels, the definition of $\mathcal F_\sigma$ and
\eqref{eq:hodge-dual-specialization} give
\[
 \kappa_{p,z}\sigma(\mathfrak D_E)F_{\sigma,R}
                 \subset\Lambda_{\sigma,R}\subset F_{\sigma,R}.
 \tag{HS10}
\]
At a finite place outside $S$, the discriminant is a unit. The
coefficient algebra is then split etale on the coefficient cover,
and its projectors are integral. The ordinary eigenbundle is a
direct summand and commutes with specialization. Hence
$F_{\sigma,R}=\Lambda_{\sigma,R}$ there.
These are fibres of the original locally free bundles, so their
local section evaluations are exactly the localizations of the
global lattices in the theorem. Intersecting at all finite places
of $L$ proves \eqref{eq:actual-hodge-inclusions} with the ideal
\textup{(HS2)}.

\emph{The ideal norm.}
The specialized odd root is still primitive: its value satisfies
the cyclotomic polynomial $\Phi_p$, and the target has
characteristic zero. Thus every $\kappa_{p,z}$ is nonzero. The
cyclotomic identity gives
\[
 \frac{\log N_{L/\mathbb Q}(\zeta_{p,z}-1)}{[L:\mathbb Q]}
       =\frac{\log p}{p-1},\qquad
 \frac{\log N_{L/\mathbb Q}(4)}{[L:\mathbb Q]}=2\log2.
\]
Extension of the different gives
\[
 \frac{\log N_{L/\mathbb Q}
                  (\sigma(\mathfrak D_E)\mathcal O_L)}{[L:\mathbb Q]}
       =\frac{\log|\operatorname{disc}(E)|}{[E:\mathbb Q]}
       =\log\operatorname{rd}(E).
\]
Multiplicativity proves the equality in
\eqref{eq:actual-hodge-ideal-norm}; its upper bound follows from
$\operatorname{rd}(E)\ge1$ and nonnegativity of the prime terms.
\end{proof}

\paragraph{Passage to the abelian quotient.}
Let $q:G_{z,L}\to A_z$ be the original abelian quotient, and retain
its differential on the original lattices. Distinguish
\[
 W_\sigma=\text{the full integral eigenlattice of }A_z,\qquad
 V_\sigma=q_*\Lambda_\sigma,\qquad
 \overline F_\sigma=q_*F_\sigma.
\]
The theorem gives
\[
 \mathfrak a_\sigma\overline F_\sigma
                       \subset V_\sigma\subset\overline F_\sigma.
 \tag{HS11}
\]
It gives no identification of $V_\sigma$ with $W_\sigma$.
The finite index $[W_\sigma:V_\sigma]$ in the actual point-height
bound remains unchanged. If $r_\sigma=\dim_L\operatorname{Lie}(A_z)_\sigma$
and both image lattices carry the prescribed common quotient
metric, then
\[
 0\le\widehat{\deg}_{\rm n}\overline F_\sigma
          -\widehat{\deg}_{\rm n}V_\sigma
 \le r_\sigma\,
      \frac{\log N_{L/\mathbb Q}\mathfrak a_\sigma}{[L:\mathbb Q]}.
 \tag{HS12}
\]
Indeed the quotient $\overline F_\sigma/V_\sigma$ is annihilated
by $\mathfrak a_\sigma$ and is generated locally by at most
$r_\sigma$ elements. This is the finite correction needed by the
point-height comparison. The regularized intersection identity
is a separate obligation.

\begin{proposition}[Simultaneous local freeness on the original model]
\label{prop:actual-hodge-free-model}
Start with the prescribed normal projective model of finite type
over $\mathbb Z$, the original full levels $p^{R_p}$ for the finite
set $S$, the same generic coefficient field, and the same
semiabelian family and cubical line. There is a normal projective
birational modification $f:X'\to X$, isomorphic over the whole
characteristic-zero fibre, on which all the actual modules
$\mathcal N_p$ and $\mathcal Q_{p,\sigma}$ are locally free.
The original labels, levels, generic field, and generic family
are unchanged. The construction gives precisely the required
corrected bundles on $X'$ and hence supplies the local-freeness
hypothesis of Theorem~\ref{thm:actual-hodge-specialization}.

If some of these modules are already locally free, they are
pulled back as the same locally free modules throughout the
construction. Every original arithmetic point lifts uniquely
to $X'$, because $f$ is an isomorphism on its characteristic-zero
generic point. The comparison in the theorem concerns the actual
modules at this lift and the pullback of the same semiabelian
family.
\end{proposition}

\begin{proof}
We first verify the compatibility that an abstract flattening
argument would not supply. For a dominant morphism $f:X'\to X$
between normal integral models with the same generic field,
the pullback of a regular germ realizing $e_{n_p}(t,-)$ is a
regular germ realizing that same character on the pulled-back
family. Any other such germ on its torsion kernel has the same
generic restriction, and equality follows from flatness over
the integral base. Their differentials modulo $n_p\omega$ are
therefore identical. Since the set of original generic torsion
labels is unchanged, and the original smooth cotangent bundle
commutes with base change, we obtain the equality of actual
modules
\[
 \mathcal N'_p
   =\operatorname{im}(f^*\mathcal N_p\longrightarrow f^*\omega).
 \tag{HS13}
\]
Restriction to the same rational eigenspace gives, with the
original coefficient projector,
\[
 \mathcal Q'_{p,\sigma}
   =\operatorname{im}(f^*\mathcal Q_{p,\sigma}
                      \longrightarrow\mathcal T_{K,\sigma}^{\vee}).
 \tag{HS14}
\]
Thus these modules are their torsion-free pullbacks, not arbitrary
coherent images on the modified model. The displayed equalities
are the original Weil-differential pullback identities.

The local description by finitely many $\delta_t$ and
$n_p\omega$ proves coherence of $\mathcal N_p$; each
$\mathcal Q_{p,\sigma}$ is coherent as the image of a coherent
module in a fixed rational vector space. After inverting $p$,
$\mathcal N_p=\omega$. Over the characteristic-zero fibre, the
coefficient action splits into its eigensummands, so each
$\mathcal Q_{p,\sigma}$ is the corresponding locally free
cotangent summand. Their ranks are the prescribed signature
ranks. Hence all these finitely many coherent modules have
constant generic rank and are locally free on the entire
characteristic-zero fibre.

For a coherent module of generic rank $a$, blow up its $a$th
Fitting ideal. Its strict transform is locally free of rank $a$
\cite[Lemma~31.36.3]{hodge-flattening}. On an integral dominating
model the strict transform is the image in the original rational
vector space: the kernel of the map from the pulled-back module
is exactly its torsion, supported over the exceptional locus.
Normalize this blowup. Normalization is finite because the models
are excellent, and it preserves the locally free strict transform.
Equations \textup{(HS13)--(HS14)} identify this transform with
the actual module calculated from the pulled-back semiabelian
family and its original Weil characters.

Repeat for the finite list of modules. A module already made
locally free stays so: its pullback is locally free and the map
to the original rational vector space is injective, since it is
injective generically on an integral model. Every centre is
vertical, because all the modules are locally free in
characteristic zero. The composite is therefore projective,
birational, normal after the indicated normalizations, and an
isomorphism over the whole characteristic-zero fibre. It remains
flat over $\mathbb Z$: its integral affine coordinate rings have
characteristic zero and are torsion-free over $\mathbb Z$.
Rank-zero
restricted modules are zero and require no modification.

The same generic field and torsion labels remain throughout;
pullback retains smoothness, the semiabelian structure, the
coefficient action, and the rigidified cubical line and its
principal generic polarization. The different is an invertible
coefficient ideal, so dualizing the locally free
$\mathcal Q'_{p,\sigma}$ and applying exactly the original
different correction gives locally free corrected modules.
On the open subset where $p$ is invertible and the other primes
of $S$ have been removed, these agree with the original ordinary
eigenbundle: its coefficient projectors are integral there.
The open subsets obtained by removing all primes of $S$ except
one, together with the complement of $S$, cover $X'$. They glue
these actual modules to the required $\mathcal F'_{\sigma}$.

The construction retains the old projective coordinates by
pullback; adjoining homogeneous coordinates for the projective
modification realizes it as the permitted refinement of the
original model. Every previously prescribed line bundle and
metric is pulled back, with all its original finite factors.
The complex fibre and its metrics are unchanged. We make no
additional numerical claim about the added projective
coordinates. Finally, the valuative criterion for properness
lifts an original arithmetic point at each finite place, from
its unique lift in characteristic zero. Separatedness makes the
lifts agree and glue over $\mathcal O_L$. No extension of $L$
is required.
\end{proof}

\subsection{Exact corrected Hodge flags at the correction primes}
\label{sec:geometry-exact-corrected-flags}

We retain the original full-level model, cubical line, torsion
characters and locally free modules of Section~\ref{sec:geometry-hodge-specialization}.
At a correction prime \(p\), write \(R=R_p\ge4\) and \(m=R-2\).
The following argument uses all torsion levels only over the injective
valued test fields. It does not enlarge the prescribed level or the
function field of the global model.

\begin{lemma}[A compatible polarized test tower]
\label{lem:corrected-ht-test-tower}
At an injective height-one test point above \(p\), let \(C\) be the
complete algebraically closed valued field, \(O_C\) its valuation ring,
and \(G\) the pulled-back semiabelian family of dimension \(g\) with
abelian generic fibre. There is a \(p\)-divisible group \(B/O_C\) of
height \(2g\) such that:
\begin{enumerate}
\item its generic torsion at every level is the original generic
torsion \(G_C[p^n]\), with all transition maps and coefficient actions;
\item its invariant cotangent module is \(\omega_G\), with the
identification at level \(n\) given by (HS5);
\item the original alternating Weil pairings extend compatibly to
a morphism \(\lambda:B\to B^D\) whose Tate-module map is an
isomorphism preserving the original pairing.
\end{enumerate}
The kernel of the characteristic-zero Hodge--Tate map
\[
 H=\alpha_{B^D}\circ T_p(\lambda):
 T_p(B)\otimes_{\mathbb Z_p}C\longrightarrow\omega_G\otimes_{O_C} C
\]
is Lagrangian for the original pairing (using a compatible system of
roots of unity to trivialize the Tate twist).
\end{lemma}

\begin{proof}
Work first over the henselization of the height-one discrete valuation
ring of the original model. For every \(n\ge1\), the finite-part theorem
used in (HS3)--(HS4) gives the finite flat open and closed subgroup
\(\mathcal H_n\) of \(G[p^n]\) containing the entire closed fibre.
Pull these subgroups back to \(O_C\). Their complements still have
empty closed fibre, and the finite-part construction is functorial:
the inclusion and multiplication-by-\(p\) maps preserve the finite
parts. Indeed their restrictions from a finite scheme cannot meet
the complementary open and closed component if they do not meet
it on the closed fibre. The inclusion
\(\mathcal H_n\hookrightarrow\mathcal H_{n+1}\) is closed.

Multiplication by \(p\) maps \(\mathcal H_{n+1}\) faithfully flatly
onto \(\mathcal H_n\). To justify flatness, on the special fibre
these are the full semiabelian torsion schemes. The map there is
finite faithfully flat of rank \(p^{2g-r}\), where \(r\) is the
torus rank of that fibre. Over the henselian DVR both coordinate
algebras are finite free, of ranks \(p^{(n+1)(2g-r)}\) and
\(p^{n(2g-r)}\). The finite special-fibre algebra of the target is
Artinian, and the source is free of the displayed relative rank
over it. Lift a basis over that algebra. Nakayama's lemma gives
a surjection from the free module of the same relative rank over
the integral target coordinate algebra. The two underlying free
DVR modules have equal rank, so the surjection is an isomorphism.
This proves faithful flatness before base change, and hence after it.
In particular the generic finite parts map onto each other.

For each \(n\) apply the explicit induced-coset construction (HS4)
to \(T_n=G_C[p^n](C)\) and \(H_n=\mathcal H_n(C)\), obtaining \(B_n\).
An integral section of \(G[p^n]\) factors through \(\mathcal H_n\):
its closed point belongs to the latter, and its inverse image of
the complementary open and closed part is empty. Consequently
\(H_{n+1}\cap T_n=H_n\). The inclusions of generic labels therefore
induce closed immersions \(B_n\hookrightarrow B_{n+1}\), using the
finite-part inclusions on the individual copies. The maps
\(p:T_{n+1}\twoheadrightarrow T_n\) induce maps
\(B_{n+1}\twoheadrightarrow B_n\), surjective on the coset sets
and on each copy finite faithfully flat, by the preceding paragraph.
The choices of representatives cancel by the translation formula
(HS4). In particular these maps satisfy all their compatibility
relations; alternatively check the relations on the generic fibres
and use flatness and separatedness.

The kernel of \(p^n:B_{n+1}\to B_1\) is finite locally free of
rank \(p^{2gn}\). The inclusion of \(B_n\) factors through it and
is a closed immersion of finite locally free schemes of equal rank,
hence an isomorphism. This proves the Barsotti--Tate axioms for
the system \(B_n\); the same argument gives the usual exact sequences
at arbitrary levels. It is important here that the ranks of
\(B_n\) are those of the full generic abelian torsion, not those
of the intrinsic special-fibre semiabelian torsion.
The identity copy has cotangent \(\omega_G/p^n\omega_G\).
These identifications commute with the cotangent restrictions, so
their limit identifies \(\omega_B\) with \(\omega_G\).
In particular the dimension of \(B\) is \(g\). Its dual has dimension
\(2g-g=g\).

At every level the original cubical pairing extends by (HS6).
The resulting maps \(B_n\to B_n^D\) are compatible: compatibility
of the Weil pairings with inclusion and multiplication by \(p\)
holds on the generic fibres, and each equation extends between the
finite flat schemes. This yields \(\lambda\). The map of Tate modules
is an isomorphism by the original generic principal polarization;
we have not asserted that each finite-level \(\lambda_n\) is a
scheme-theoretic isomorphism. The coefficient action extends by
the same finite-part and coset construction.

Apply the integral Hodge--Tate complex theorem to \(B\) and \(B^D\)
(Fargues, \emph{Ann. Sci. ENS} 44 (2011), \S5.3.2, Theorem 2,
pp.~934--935; see also the finite-group theorem in \S5.4
already used in (HS8)). Over \(C\) it gives the exact sequence
\[
 0\longrightarrow(\omega_B\otimes C)^\vee(1)
 \longrightarrow T_p(B)\otimes C
 \xrightarrow{\alpha_B} \omega_{B^D}\otimes C
 \longrightarrow0 ,
\]
where the first arrow is the transpose of \(\alpha_{B^D}\) via
Cartier duality. In particular the kernel of \(\alpha_B\) is the
annihilator of the kernel of \(\alpha_{B^D}\) under the Tate pairing.
Naturality gives
\(\alpha_{B^D}\circ T_p(\lambda)=(\lambda^D)^*\circ\alpha_B\).
Thus \(\ker\alpha_B\subseteq\ker H\). Both kernels have dimension
\(g\), since \(T_p(\lambda)\) is an isomorphism and both Hodge--Tate
maps are surjective over \(C\) with rank \(g\). They coincide.
Using the alternating pairing given by \(\lambda\), the annihilator
identity becomes \((\ker H)^\perp=\ker H\).
This is the required orthogonality, including at \(p=2\).
Only rational exactness of the displayed complex is used here;
for the integral cokernel bound we retain (HS8) and its factor \(4\).
\end{proof}

\begin{theorem}[An exact flag sequence modulo the original loss]
\label{thm:exact-corrected-flag-sequence}
Keep the original correction prime \(p\), \(\mathrm{HTFree}\) model
and the original flag with \(\mathrm{RootLines}\). On its
\(p^{R-2}\)-thickened fibre, for each CM-type embedding \(\sigma\)
there is an exact sequence of vector bundles
\[
 0\longrightarrow P_\sigma
 \longrightarrow \mathcal O^{N}
 \longrightarrow Q_\sigma
 \longrightarrow0,
 \qquad
 P_\sigma=\sigma(\mathfrak D_E)^{-1}
                  Q_{\bar\sigma}^{\vee},\quad N=p_{\rm sig}+q_{\rm sig}.
 \tag{HF1}
\]
Here \(Q_\sigma\) denotes the restricted image of the original
corrected cotangent module, with the \emph{conjugate} coefficient
action on differentials; all bundles in (HF1) are reduced modulo
\(p^{R-2}\). Thus \(P_\sigma\) is the original selected corrected
tangent bundle and \(Q_\sigma\) is the dual of the other corrected
tangent bundle. The sequence uses the original full-level basis
and normalized Weil pairing. It supplies \(\mathrm{TrivFlagOnHodge}\)
and the correction-prime \(\mathrm{ModelNefStatement}\) on the
actual model.
\end{theorem}

\begin{proof}
Use the convention on the action on differentials just stated, so
the names \(Q_\sigma\) here correspond to the source
\(\mathrm{omHT}\), not to an un-conjugated dual of the
\(\sigma\)-eigenspace. Before reduction work on the original model
over \(\mathbb Z_{(p)}\). By (HS8), and the inclusions proved on
the locally free charts,
\[
 p^R\omega_G\ \subset\ p^{R-2}\mathcal N_p,\qquad
 \mathcal N_p=\langle h(t)\rangle+p^R\omega_G ,
 \tag{HF2}
\]
where \(h(t)\) are local lifts of the actual character differentials.
For (HF2), \(\kappa_p\mid p^2\) and
\(\kappa_p\omega_G\subseteq\mathcal N_p\).
The containment is also valid with the original projected image in
place of \(\mathcal N_p\). In particular all the character
differentials are canonically defined in
\(\mathcal N_p/p^m\mathcal N_p\), irrespective of the lifts.
Their additivity and coefficient linearity hold there: both were
identities modulo \(p^R\omega_G\) before using (HF2).

Choose the original \(N\) generators of the free
\(\mathcal O_E/p^R\mathcal O_E\)-module of generic torsion.
Their projected differentials now give a global surjective map
\[
 q_\sigma:(\mathcal O/p^m)^N\longrightarrow Q_\sigma/p^m Q_\sigma .
 \tag{HF3}
\]
Indeed after projection to the \(\sigma\)-eigenspace the coefficients
in \(\mathcal O_E\) act through \(\sigma\), and the additional
\(p^R\omega_G\) generators disappear by (HF2).
All embeddings of integers are integral on the normal charts.
This argument does not assume that the idempotent projections are
integral on the uncorrected cotangent module. It uses precisely
their locally free images.

Here is the integral normalization of the adjoint. Put
\(D_p=\mathcal O_E\otimes\mathbb Z_{(p)}\).
Its inverse different is the trace dual of \(D_p\), and it is
principal because \(D_p\) is semilocal Dedekind. Choose once a
generator \(d\) of its different. The original perfect alternating
balanced pairing of the full level, after choosing the original
primitive \(p^R\)-th root of unity, is the trace of a unique
sesquilinear pairing with values in
\(\mathfrak D_E^{-1}D_p/p^R\mathfrak D_E^{-1}D_p\).
Indeed the trace-dual identification applied in one variable
constructs that pairing; the balanced identity gives its conjugate
linearity in the other variable. The perfectness of the trace
pairing over \(\mathbb Z/p^R\) implies perfectness of this
sesquilinear pairing over \(D_p/p^R\).
Multiplication by \(d\) makes its matrix invertible over that
ring. Evaluation at \(\sigma\) therefore gives an invertible
bilinear matrix pairing the two conjugate labelled copies of
\((\mathcal O/p^m)^N\).
This is the different-normalized Gram matrix of
\(\mathrm{HTChart.gram}\); it is defined by the constant full-level
pairing, not by choosing a fresh Gram form on each chart.

The dual of \(q_{\bar\sigma}\) composed with the inverse of this
matrix defines an injection \(i_\sigma:P_\sigma/p^m\to
(\mathcal O/p^m)^N\). In this formula a member of
\(P_\sigma=\sigma(d)^{-1}Q_{\bar\sigma}^{\vee}\) is first multiplied
by \(\sigma(d)\) to give a member of the ordinary dual, as required
by the different-normalized Gram. Equivalently the defining
adjoint equality is the one for the unscaled pairing and evaluation
of the member of \(P_\sigma\). This construction is independent of
the choice of different generator. Since \(Q_{\bar\sigma}\) is
locally free and \(q_{\bar\sigma}\) is surjective, \(i_\sigma\) is
locally a split injection, including over a nonreduced thickening.

We claim \(q_\sigma i_\sigma=0\). It is enough to check this at
the height-one primes above \(p\) of the original normal chart.
Indeed locally lift the composite to a map of projective modules
before reduction; its image in the fraction field is divisible
by \(p^m\) in the target if and only if its quotient by \(p^m\)
belongs to the target at every height-one prime. This is the
intersection property of reflexive modules over a normal
noetherian domain. At height-one primes not over \(p\) the
divisibility condition is automatic.

At a height-one prime above \(p\), extend faithfully flatly from
the DVR to the ring of the complete algebraically closed valued
field. Construct \(B\) by Lemma~\ref{lem:corrected-ht-test-tower}.
Lift the original level generators to an
\(\mathcal O_E\otimes\mathbb Z_p\)-basis of \(T_p(B)\);
Nakayama's lemma, using their basis property modulo \(p^R\),
justifies the lift. Lift also the chosen root of unity compatibly
to all levels. The Tate pairing reduces to the original one
modulo \(p^R\); its trace-dual pairing and its normalized Gram
do so as well.

The image of the actual Tate Hodge--Tate map
\(\alpha_{B^D}\circ T_p(\lambda)\) over \(O_C\) is exactly
\(\mathcal N_p\otimes O_C\). Its reduction in \(\omega_G/p^R\)
is the original character differential map, since at each level
the identity copy of \(B_n\) is \(\mathcal H_n\).
Its cokernel is annihilated by \(\kappa_p\): apply the integral
Hodge--Tate cokernel bound (or the compatible finite bounds and
compact finite generation of the image) to \(B^D\).
In particular its image contains \(p^R\omega_G\).
The equality of the two images follows.
The finite free projected modules commute with this flat
base change, as images of a fixed rational linear map.

Over \(O_C\) therefore one can use the actual Tate Hodge--Tate
map and actual Tate Gram matrix before reduction. Their maps
on the projected image lattices are integral and give exactly
(HF3) and its adjoint after reduction modulo \(p^m\).
The kernel of the rational Hodge--Tate map is Lagrangian by
Lemma~\ref{lem:corrected-ht-test-tower}, so the adjoint of the
conjugate labelled quotient has image in the kernel of the
labelled quotient. The different normalization merely scales
that adjoint by a nonzero constant over \(C\).
The composite of the actual integral maps is consequently zero
over \(C\), hence over \(O_C\). The congruence between actual
Tate differentials and the chosen level differentials modulo
\(p^R\omega_G\), together with (HF2), proves zero composition
for the original reduced maps. Faithful flatness and the
reflexive-module argument prove the claim on the original chart.

The ranks of \(P_\sigma\) and \(Q_\sigma\) are respectively
\(p_{\rm sig}\) and \(q_{\rm sig}\), so they sum to \(N\).
The split injection and surjection just proved are exact:
locally quotient the free middle bundle by the split first
bundle. The induced surjection to \(Q_\sigma\) is between free
modules of the same rank, hence is an isomorphism over the
possibly nonreduced local ring. This proves (HF1).
The maps glue because the original torsion labels, Weil pairing,
different ideal and classes in (HF2) are global.

Now use the actual flag in \(P_\sigma\), and in \(Q_\sigma\)
use the reversed dual flag of the other corrected tangent bundle.
The successive inverse images of this second flag under (HF1),
following the first flag, give a full flag of subbundles of the
trivial middle bundle modulo \(p^m\). Its graded line bundles are
exactly the original root lines, including the inverses and reversal
in \(\mathrm{rootIdx},\mathrm{rootSgn}\). The flag restricts to every
special-fibre component without any new base-field extension.
In particular it supplies the source's \(\mathrm{TrivFlagAt}\)
on those components. In the coordinate model one can choose a flag
basis on each original affine chart after localizing as needed;
the prescribed root-line generators fix the diagonal transition
factors. The change matrices are upper triangular modulo \(p^m\),
and the original root monomial ratios are their diagonal entries.
Equivalently the exact sequence of subbundles is the geometric
form of the original trivial-flag assertion; no change to the
coordinate model or to its degree is necessary.

For a complete integral curve on a special-fibre component,
write \(L_i\) for the successive quotients of this full flag and
\(F_i\) for its \(i\)-th subbundle (\(F_0=0,F_N=\mathcal O^N\)).
The determinant of \(F_i\) injects into a trivial exterior-power
bundle, so \(\deg F_i\le0\): a nonzero component of the induced
map to a trivial line gives an effective divisor after normalization
of the curve. Also \(\deg F_N=0\). Thus for decreasing real
weights \(\mu_1\ge\cdots\ge\mu_N\), summation by parts gives
\[
 -\sum_{i=1}^N\mu_i\deg L_i
 =-\sum_{i=1}^{N-1}(\mu_i-\mu_{i+1})\deg F_i\ \ge0 .
 \tag{HF4}
\]
This is the sign convention of
\(\mathrm{RootModel.coef}(\mu,0)=-\mu\).
Sum over the original CM-type labels. It proves precisely
\(\mathrm{ModelNefStatement}\) at the correction primes.
The Whitney formula also gives
\(\prod_i(1+c_1(L_i))=1\) on each component, and the positive-degree
power sums of these roots vanish there by the Newton identities.
\end{proof}

The proof establishes the exact Lagrangian comparison and the
correction-prime general flag positivity, rather than only an
annihilator inclusion. It does not supply the distinct \(\mathrm{M2}\)
condition away from the correction set, or the Frobenius-linked
vertical correction identities comparing two levels. Those remain
necessary for the complete weighted mean. No numerical change to
the base-test bound of Section~\ref{sec:geometry-height-assembly}
is asserted here.

% W10: relative Chern comparison on the original labelled lattice family.
\subsection{The two Frobenius corrections by relative Chern characters}
\label{sec:relative-chern-correction}

The following argument replaces the root-by-root computation of the
two-level correction by a calculation with two short exact sequences.
It is a numerical descent theorem: its input is the original
\emph{pointwise} torsion-lattice identities, not the desired global
Chern comparison. In particular those pointwise identities still need
a proof for the actual torsion orders.

Let \(X\) be an integral, normal, projective flat model of relative
dimension \(d\) over an excellent mixed-characteristic DVR. Let
\(Z=X_{p=0}\), where the rational prime \(p\) is nonzero in the
base, and suppose that a base element \(w\) satisfies
\(w^{p-1}=up\) with \(u\) a unit. The generic fibre is denoted \(X_\eta\).
Write \(T=\mathcal O_X^n\). All lattices below carry their
specified generic identifications, and all maps between them are
the prescribed coordinate maps. We use the same
labels \(\tau\) and Frobenius permutation \(\phi\) as in the source.
The four bundles are written \(L_\tau,W_\tau,O_\tau,O_\tau^{\rm HT}\),
where \(L\) denotes the old Lie bundle and \(O\) the old
cotangent bundle, and
\[
 L_\tau\subset W_\tau\subset w^{-1}L_\tau,\qquad
 O_\tau^{\rm HT}\subset O_\tau\subset w^{-1}O_\tau^{\rm HT}.
\]
No identification of the two different corrected bundles is made.

Here is the lattice input, in coordinate notation. It is exactly the
part of \texttt{LieLatticeData.Base} and \texttt{Facts} used in this
calculation. Let \(\mathbb F\) be the finite field of the torsion
coordinates, \(M\subset\mathcal O_X^{\mathbb F^n}\) the original
function lattice, containing the constants and a power of \(p\) times
the full function lattice. Write \(\chi_\tau\) for the original
character functions with \(\chi_\tau(0)=0,\chi_\tau(1)=1\). Put
\[
 A_\tau=\{a:\ \sum_i a_i\chi_\tau(x_i)
                      \in M+p\mathcal O_X^{\mathbb F^n}\},
 \qquad H_\tau=A_\tau^\vee .
\]
Here and below all constructions on lattices are sheaf constructions;
duals use the given standard pairing. The lattice \(A_\tau\)
lies in \(T\), by evaluation at the coordinate vectors. We require
\(wT\subset A_\tau\). The original two graphs
\(\Gamma_\tau\subset L_{\tau,\eta}\oplus T_\eta\) and
\(\Delta_\tau\subset T_\eta\oplus O_{\tau,\eta}\)
have the following projections and coordinate kernels:
\[
\begin{array}{c|cc|cc}
 &\operatorname{pr}_1&\operatorname{pr}_2&
 \ker(\operatorname{pr}_1)&\ker(\operatorname{pr}_2)\\ \hline
 \Gamma_\tau &w^{-1}W_\tau&C_\tau^{\rm all}&T&W_\tau\\
 \Delta_\tau &w^{-1}T&w^{-1}O_\tau^{\rm HT}
      &O_\tau^{\rm HT}&C_\tau^{\rm all}
\end{array}
\]
where \(C_\tau^{\rm all}=\operatorname{pr}_2\Gamma_\tau\).
The two extra facts are
\[
 \operatorname{pr}_1(\Gamma_\tau\cap(L_{\tau,\eta}\oplus H_\tau))
       =w^{-1}L_\tau,\qquad
 \operatorname{pr}_2(\Delta_\tau\cap(H_\tau\oplus O_{\tau,\eta}))
       =O_\tau.                                                    \tag{RC1}
\]

For the second graph calculation let \(V_i\) be the space of functions
vanishing on the span of the first \(i\) coordinate vectors of
\(\mathbb F^n\); hence \(V_n=0\). Let \(\Theta_\tau\) be the
original graph from the augmentation lattice to the dual Lie
lattice. Its first projection is \(M\cap V_0\), its second
projection is \(L_\tau^\vee\), and its second-coordinate-only
kernel is \(p L_\tau^\vee\). Define
\[
\begin{split}
 I_{\tau i}&=\{a:\ a\chi_\tau(x_i)\in M+V_{i+1}\},\\
 K_{\tau i}&=\operatorname{pr}_2(\Theta_\tau
                         \cap(V_i\oplus L_{\tau,\eta}^\vee)),\\
 B_{\tau i}&=\{(a,b):\
       (a\chi_\tau(x_i),b)\in\Theta_\tau+(V_{i+1}\oplus0)\}.
\end{split}
\]
Indices \(i\) for the coordinates run from \(0\) to \(n-1\);
\(K_{\tau i}\) is defined also at \(i=n\).
For \(A_\tau\) use the decreasing coordinate filtration given by
\(a_j=0\) for \(j<i\), and denote the image at coordinate \(i\) by
\(D_{\tau i}\). The other three pointwise facts required are
\[
 D_{\tau i}=I_{\tau i},\qquad
 \operatorname{pr}_2 B_{\tau i}=K_{\tau i},\qquad
 \ker(\operatorname{pr}_2 B_{\phi\tau,i})
            =\langle a^p:a\in I_{\tau i}\rangle_{\mathcal O_X}.
                                                                    \tag{RC2}
\]
In the last formula the kernel is identified with its first
coordinate. This is the \emph{span of powers}, not a map
\(a\mapsto a^p\) asserted to be additive in mixed characteristic.

\paragraph{Putting the pointwise facts on one auxiliary model.}
The following elementary flattening observation allows us to start
with the source's \texttt{TestFacts} instead of assuming (RC1)--(RC2)
on the entire original model.

\begin{lemma}[Simultaneous lattice diagram]
\label{lem:relative-chern-lattice-diagram}
Suppose that the original coherent lattice family has the indicated
base projections, kernels and inclusions, with the four bundles and
\(M\) locally free. Suppose the six facts
\(wT\subset A_\tau\), (RC1) and (RC2) hold after extension to
each of the algebraically closed, complete rank-one valued fields
in the source's injective test-point hypothesis, with every
lattice interpreted as its generated lattice there.
There is a common normal projective birational modification,
an isomorphism over \(X_\eta\), on which these facts hold
as sheaf identities and all bundles, filtered pieces and graph
projections needed below are locally free. The original four
bundles on it are their pullbacks.
\end{lemma}
\begin{proof}
For any coherent fractional module of constant generic rank,
take the closure of the generic identity quotient in its relative
Grassmannian of quotients of that rank, and normalize.
The tautological quotient is the pullback modulo torsion:
the kernel is generically zero and the quotient is locally free.
Thus any finite list of coherent image lattices can be made
locally free in succession. A surjection onto a locally free
image splits locally. Consequently its kernel (and the intersection
that this kernel represents) is locally free if its domain is
locally free, and these constructions now commute with every
extension to a valued test field. More explicitly, for a fixed
linear subspace use the image in the quotient by that subspace.
For the intersection of two lattices in one space use the image
of their direct sum under subtraction, i.e.\ their sum lattice.
Once that image is locally free, the kernel is the intersection
and the short exact sequence remains exact on pullback.
The image of a further coordinate projection is treated by
one more such modification. All subsequent pullbacks preserve
the locally split sequences already obtained.

For the present finite diagram apply this procedure first to
\(M+p\mathcal O^{\mathbb F^n}\), its quotient by the range of
the linear character map, and the coordinate flag in \(A_\tau\).
This constructs \(A_\tau,D_{\tau i},H_\tau\) with their actual
generic identifications. For \(I_{\tau i}\) use the restriction
image of \(M\) on the first \(i+1\) coordinates, followed by
intersection with the line spanned by \(\chi_\tau(x_i)\).
The latter line has coordinate \(\chi_\tau(1)=1\).
For \(K_{\tau i}\) and \(B_{\tau i}\) use the restriction
images of \(\Theta_\tau\) (retaining the second coordinate for
\(B_{\tau i}\)), the corresponding kernels and the intersections
with that character line times the dual Lie space.
For (RC1) take the sums with \(w^{-1}W_\tau\oplus H_\tau\)
and \(H_\tau\oplus w^{-1}O_\tau^{\rm HT}\), respectively, and
the indicated coordinate images. Also flatten the image of
the restricted \(\Gamma_\tau\) in its second coordinate.
These are finitely many coherent images and kernels.

The original graphs have the required locally free
presentations: for \(\Gamma_\tau\) use the first projection
with kernel \(T\); for \(\Delta_\tau\) use the first projection
with kernel \(O_\tau^{\rm HT}\); and for \(\Theta_\tau\) use the
first projection with kernel \(pL_\tau^\vee\). The augmentation
lattice \(M\cap V_0\) is a direct summand since \(1\in M\).
The second projection of \(\Delta_\tau\) is onto the locally
free \(w^{-1}O_\tau^{\rm HT}\), so its first-coordinate kernel
is also locally free. This is \(C_\tau^{\rm all}\), and hence
the second projection sequence of \(\Gamma_\tau\) is locally
split as well. Thus the base identities survive all pullbacks.
Over \(X_\eta\) the graphs are the coordinate direct products
and the function lattice is full; all the above images have
constant rank there. The modifications are therefore vertical.

On the final normal model take any height-one point over \(p\).
Its local ring is a DVR. Extend its fraction field, via completion
and completed algebraic closure, to an algebraically closed
complete rank-one valued field. The valuation ring is faithfully
flat over that DVR: it is torsion free and the local maximal
ideal has proper extension. All the indicated images, intersections
and duals commute with this extension by the split sequences.
The resulting test field contains the original function field
injectively and has positive valuation of \(p\); the original
\texttt{TestFacts} applies to the original chart and its
generated lattices. Faithful flatness yields the required equalities
or inclusions over the height-one local ring itself.
For the last equality in (RC2), a locally free fractional
line is principal on each trivialization; its span of \(p\)-th
powers is its ordinary product ideal of order \(p\), also
locally free. Thus both sides of each required comparison
are locally free lattices. On a normal scheme such lattices
are the intersections of their height-one localizations
inside their rational vector space. All facts follow globally,
including at points of codimension greater than one.
\end{proof}

The auxiliary modification in the lemma is used only for
intersection equalities. No uniform degree, height or number of
charts of it is needed. In particular it is not a replacement
for the prescribed tower or its sections.

\paragraph{Relative numerical classes.}
For two bundles \(E,F\) with the same generic identification,
write \([E:F]\) for the relative \(K\)-class supported on
\(Z\), with sign \(E-F\). If \(F\subset E\) it is represented
by the two-term resolution of \(E/F\). In general one can
use a common sublattice \(p^a E\subset E\cap F\) and subtract
the two quotient classes; additivity makes this independent
of \(a\). All exact sequences used below respect their generic
identifications. We use the localized Chern character of
perfect complexes (Fulton, \emph{Intersection Theory},
second edition, \S18.1, Theorem 18.1, Proposition 18.1 and
Corollary 18.1.2). Its additivity and projection formula apply
to the proper models over the DVR, including singular
special fibres. For an original complementary monomial \(Q\)
of degree \(d+1-k\) in line classes, write
\[
 c_k(E:F;Q)=k!\deg_Z
       \bigl(Q\,\operatorname{ch}^{Z}_k([E:F])\cap[X]\bigr).
\]
The degree is taken on the supporting fibre, \emph{before}
pushing forward to \(X\). Forgetting the support would make
the principal scalar divisors vanish incorrectly.

For a pair of lines \(E,F\) the expression is the vertical
divisor of their generic ratio intersected with
\(\sum_{j=0}^{k-1}c_1(E)^jc_1(F)^{k-1-j}Q\).
This follows by applying the localized character to the
two-term line inclusion (or subtracting two inclusions
through a common smaller line). Hence, by additivity along
matched flags, \(c_k(E:F;Q)\) is precisely the source's
mixed-difference power sum \texttt{ck}, including the sign
\(x-x'\). Two consequences used here are
\[
\begin{split}
 c_k(J^{\otimes p}:\mathcal O;Q)
   &=p^k c_k(J:\mathcal O;Q) \quad(J\text{ a fractional line}),\\
 c_k(w^{-1}E:E;Q)
   &= k\,\deg_Z\bigl(Q P_{k-1}(E)\cap[\operatorname{div}(w)]\bigr),
                                                               \tag{RC3}
\end{split}
\]
where \(P_j(E)=j!\operatorname{ch}_j(E)\).
For the second identity tensor the Cartier divisor resolution
by \(E\): its normal line \(\mathcal O(\operatorname{div}(w))\)
is trivial, so there are no higher self-intersection terms.
For the first identity use the line formula and
\(c_1(J^{\otimes p})=p c_1(J)\), with the ratio divisor
also multiplied by \(p\). In particular the normalized
power-sum factor \(k\) in the second line is necessary.

\begin{theorem}[Two relative Chern identities and the finite orbit formula]
\label{thm:relative-chern-two-identities}
Under the base and test hypotheses of
Lemma~\ref{lem:relative-chern-lattice-diagram}, set, for the original
complementary monomial \(Q\),
\[
 h_\tau=c_k(H_\tau:T;Q),\qquad
 \ell_\tau=k\deg_Z\bigl(QP_{k-1}(L_\tau)
                              \cap[\operatorname{div}(p)]\bigr).
\]
For every \(1\le k\le d+1\) the original two-level numerical
relation is
\[
\begin{split}
 c_k(L_\tau:W_\tau;Q)+c_k(O_\tau:O_\tau^{\rm HT};Q)
     &=h_\tau-\ell_\tau/(p-1),\\
 p^kh_\tau&=h_{\phi\tau}+\ell_{\phi\tau}.                 \tag{RC4}
\end{split}
\]
If the orbit of \(\tau\) under \(\phi\) has length \(m\), then
\[
 h_\tau=\frac{\sum_{a=1}^{m}p^{k(m-a)}
                                  \ell_{\phi^a\tau}}{p^{km}-1}.
                                                               \tag{RC5}
\]
These are the exact \texttt{Relations} and
\texttt{LabelChernComparison}, not estimates or asymptotic
formulas. They descend to the original model and hold with
any original homogeneous complementary polynomial by linearity.
\end{theorem}
\begin{proof}
Work first on the common model. Let
\[
 C_\tau=\operatorname{pr}_2\bigl(\Gamma_\tau
                  \cap(L_{\tau,\eta}\oplus H_\tau)\bigr).
\]
The restricted graph has first projection \(w^{-1}L_\tau\),
first-coordinate-only kernel \(W_\tau\), and
second-coordinate-only kernel \(T\). Taking the quotient by
\(W_\tau\oplus T\) in two ways gives
\[
 C_\tau/T\simeq w^{-1}L_\tau/W_\tau.
\]
The restricted second graph has first projection \(H_\tau\)
because \(H_\tau\subset w^{-1}T\); its second projection
is \(O_\tau\). Its two coordinate kernels are
\(C_\tau\) and \(O_\tau^{\rm HT}\). Taking the quotient
by their direct sum gives
\(H_\tau/C_\tau\simeq O_\tau/O_\tau^{\rm HT}\).
The three quotients are supported on \(Z\). Additivity in
relative \(K\)-theory therefore proves the exact identity
\[
 [H_\tau:T]=[L_\tau:W_\tau]+[O_\tau:O_\tau^{\rm HT}]
                         +[w^{-1}L_\tau:L_\tau].            \tag{RC6}
\]
Apply (RC3) and \((p-1)\operatorname{div}(w)=\operatorname{div}(p)\).
This gives the first equality of (RC4).

For the other identity, the first projection of \(B_{\tau i}\)
is \(I_{\tau i}\): a witness in \(M\) restricting to
\(a\chi_\tau(x_i)\) vanishes at zero, so has a preimage in
\(\Theta_\tau\). The second-coordinate-only kernel is
\(K_{\tau,i+1}\). By (RC2) the second projection is
\(K_{\tau i}\) and the first-coordinate-only kernel is
\(I_{\phi^{-1}\tau,i}^{\otimes p}\).
Taking the quotient of the graph by the two kernels gives
\[
 I_{\tau i}/I_{\phi^{-1}\tau,i}^{\otimes p}
                       \simeq K_{\tau i}/K_{\tau,i+1}.
\]
All these are quotients of locally free lattices, and the
isomorphism is induced by the original graph. Equivalently
one can compare its two locally split projection sequences.
Dualizing those sequences, with the generic coordinate
identifications unchanged, gives
\[
 [I_{\tau i}^{\vee}:(I_{\phi^{-1}\tau,i}^{\vee})^{\otimes p}]
                   =[K_{\tau i}^{\vee}:K_{\tau,i+1}^{\vee}].
                                                               \tag{RC7}
\]
This sign is important: both dual comparisons reverse the
respective inclusions. The coordinate filtration of \(A_\tau\)
has graded lines \(D_{\tau i}=I_{\tau i}\); the dual filtration
gives
\([H_\tau:T]=\sum_i[I_{\tau i}^\vee:\mathcal O]\).
Also \(K_{\tau0}=L_\tau^\vee\), \(K_{\tau n}=pL_\tau^\vee\),
by the two projection conditions of \(\Theta_\tau\).
Sum (RC7), apply the first identity of (RC3), and telescope
the right side. The result is
\[
 h_\tau-p^kh_{\phi^{-1}\tau}
   =c_k(L_\tau:p^{-1}L_\tau;Q)=-\ell_\tau,
\]
the second equality of (RC4) after relabelling.

Iterate this equality exactly \(m\) times; since \(p^{km}>1\),
division gives (RC5). No infinite tower of Chern identities
or convergence of model corrections is involved.
All auxiliary pullbacks of the original class comparisons
have unchanged numerical values by the localized projection
formula, since the birational map pushes the fundamental
cycle to the original one with degree one.
The same formula applies to \(\operatorname{div}(p)\)
and the old bundle roots. Hence (RC4)--(RC5) prove the
numerical identities on the original model.
\end{proof}

For clarity, the coefficient of \(\ell_{\tau'}\) in the
first line of (RC4) is zero outside the orbit and within
the orbit is exactly
\[
 \frac{p^{k(m-a)}}{p^{km}-1}
             -\frac{\mathbf1_{\tau'=\tau}}{p-1},
 \qquad \tau'=\phi^a\tau,\quad 1\le a\le m .
\]
This agrees with the source's \texttt{co}: group its geometric
series into complete periods. In particular the absolute row sum
is
\[
 \frac1{p-1}+\frac1{p^k-1}-\frac2{p^{km}-1}.
\]
Summing over a complete orbit gives the simpler identity
\[
 \sum_{\tau\ {\rm in\ orbit}}
  (c_k(L_\tau:W_\tau;Q)+c_k(O_\tau:O_\tau^{\rm HT};Q))
 =\left(\frac1{p^k-1}-\frac1{p-1}\right)
                      \sum_{\tau\ {\rm in\ orbit}}\ell_\tau .
\]
In particular the aggregate first correction is zero on
each complete orbit for \(k=1\). This cancellation does not
say that a single labelled first correction is zero.

With the original dual and sign identifications of the
corrected Hodge flags, the old Lie power sums in the labels
\(\tau,\bar\tau\) are exactly the two entries of \texttt{liePs},
and the left side of (RC4) is the original full flag power
sum difference for the corresponding label. Thus the formula
has precisely the coefficients of \texttt{FlagChernComparison}.
This final identification must use the actual labelled
bundle pairs at each correction prime, not merely bundles
with equal ranks. The previously established trivial
corrected-flag theorem handles the separate fibre-vanishing
condition. This calculation leaves the actual test-point
torsion-order facts, the global two-level model linkage,
the weighted tower mean and the section-rank losses to be
completed; it does not establish W10 or the final RH budget.

\subsection{Torsion-character flags by finite flat cotangent sequences}
\label{sec:geometry-torsion-character-flags}
The following local calculation separates the finite flat algebra from
the geometric inputs to the test-point theorem of
Section~\ref{sec:geometry-relative-chern-comparison}.
It applies to the original coordinate flag, not just a flag adapted to
a valuation-ring basis.

Let \(R\) be the valuation ring of an algebraically closed complete
rank-one mixed-characteristic field \(C\), with additive valuation
\(v(p)=1\). Put \(F=\mathbf F_{p^f}\), \(q=p^f\), and
\(N=F^n\), and let \(G\) be a finite locally free commutative
\(F\)-module scheme over \(R\), with its specified generic
identification \(G(C)=N\). Write \(M=\mathcal O(G)\subset R^N\)
for its evaluation lattice. We use the fundamental Teichm\"uller
characters \(\chi_\tau:F\to R\), extended by zero at zero, indexed
by the Frobenius orbit, so
\(\chi_{\phi\tau}=\chi_\tau^p\).
Write \(F_i\subset N\) for the first \(i\) coordinate subspace and
\(G_i\subset G\) for its schematic closure (\(0\le i\le n\)).
These closures are finite locally free: their coordinate algebras
are the finitely generated torsion-free images of \(M\) under
restriction to \(F_i\). Over a valuation ring such modules are free.
The group and field actions descend because the corresponding
generic identities hold in torsion-free algebras.
We use the classical finite locally free subgroup quotient theorem
(for commutative finite locally free groups), so \(G/G_i\) is
finite locally free and \(G\to G/G_i\) is a \(G_i\)-torsor.
In particular these assertions do not require a discrete valuation.

Let \(A_\tau\subset R^n\) consist of coefficient vectors \(a\)
such that \(\sum_j a_j\chi_\tau(x_j)\) is in \(M+pR^N\).
Let \(D_{\tau i}\) be its coordinate-\(i\) image on vectors with
coordinates \(<i\) zero. Let \(I_{\tau i}\subset R\) consist of
\(c\) such that \(c\chi_\tau(x_i)\) on \(F_{i+1}\) is the
restriction of a function of \(M\). These are the original
\texttt{af}, \texttt{dId} and \texttt{xId}, respectively.
All these coefficient lattices are finitely generated: the
conditions are images or kernels of finite matrices over a
valuation ring, for which finite Smith elimination applies.
All \(D_{\tau i}\) and \(I_{\tau i}\) are nonzero principal ideals.
For the first statement below, \(w\in R\) satisfies
\(w^{p-1}\in pR^\times\).

\paragraph{1. Annihilator and discriminant.}
One has \(wR^n\subset A_\tau\) for every \(\tau\).
In fact, for each \(\mathbf F_p\)-linear form
\(\mu:N\to\mathbf F_p\), the generic character
\(\zeta_p^{\mu(x)}\) extends to \(G\): it is a generic point of
the finite Cartier dual, and \(R\) is integrally closed.
Set
\[
 \ell(z)=\sum_{a=1}^{p-1}(-1)^{a+1}(z-1)^a/a,\qquad
 u=\zeta_p-1 .
\]
The formal logarithm identity through degree \(p-1\), whose
coefficients are \(p\)-integral, gives
\(\ell(\zeta_p^b)\equiv b\ell(\zeta_p)\pmod{pR}\)
for \(0\le b<p\): terms of degree at least \(p\) in \(u\)
have valuation \(>1\). Moreover \(v(\ell(\zeta_p))=1/(p-1)\).
Thus \(w\mu\) has a lift modulo \(p\) in \(M\), using integer
representatives of \(\mu\). In an \(\mathbf F_p\)-basis of \(F\),
each \(\chi_\tau(x_j)\) is a linear combination of such
coordinate functions modulo \(p\), since its lift is additive
modulo \(p\). This proves the inclusion. Evaluation at the
coordinate vectors proves \(A_\tau\subset R^n\).

For any of the groups \(P\) in this discussion, write
\(\lambda(P)=v(\det\mathcal O(P))\) in the full evaluation
lattice on its split generic points, and let
\(d(P)=\operatorname{length}_v\omega_{P/R}\).
Here the latter length of a finitely presented torsion module
is the sum of valuations of its nonzero Smith invariants.
There is an exact identity
\begin{equation}\label{eq:torsion-character-discriminant}
  2\lambda(P)=|P(C)|d(P).
\end{equation}
To see this without any tightness hypothesis, first recall why
cotangent lengths are additive. Finite locally free group schemes
are local complete intersections (SGA 3, Exp.~III, 4.15).
Their cotangent complexes at the identity are two-term complexes
of finite locally free modules in degrees \(-1,0\).
In the generically \'{e}tale case the degree-\(-1\)
differential is injective: its generic localization is an
isomorphism and its domain is torsion-free. In an exact finite
locally free group sequence \(H\to P\to Q\), transitivity for
the cotangent complex, and flat base change to the identity
fibre \(H\) of \(P\to Q\), give
\[
 0\longrightarrow\omega_Q\longrightarrow\omega_P
       \longrightarrow\omega_H\longrightarrow0 .
\]
All maps are the original differential maps. In particular
\(d(P)=d(H)+d(Q)\), also after applying any unit-group
character idempotent.

For a rank-one \(F\)-module scheme \(Q_i=G_{i+1}/G_i\),
the nonconstant augmentation character components all have rank
one. By Raynaud's classification (Cor.~1.5.1 of
\emph{Sch\'emas en groupes de type \((p,\ldots,p)\)},
Bull.~SMF 102 (1974), 241--280), if
\(I_{\tau i}=c_{\tau i}R\), the fundamental generators
\(X_\tau=c_{\tau i}\chi_\tau\) give the presentation
\[
 \mathcal O(Q_i)=
 R[X_\tau:\tau]/(X_\tau^p-\delta_{\tau i}X_{\phi\tau}:\tau),
 \qquad
 \delta_{\tau i}=c_{\tau i}^{p}/c_{\phi\tau,i}.
\]
The coefficient ideals here are indeed the specified
\(I_{\tau i}\): a function on \(G_{i+1}\) constant on the
generic \(G_i\)-cosets is invariant under \(G_i\) on the
whole product, by flatness. Faithfully flat descent identifies
it with a function on \(Q_i\). Conversely every function on
\(Q_i\) pulls back to \(G_{i+1}\) and lifts to \(M\).

The presentation implies, with the indicated coefficient
identification and not just up to isomorphism of lengths,
\begin{equation}\label{eq:torsion-character-cotangent}
 \omega_{Q_i,\tau}=I_{\tau i}/I_{\phi^{-1}\tau,i}^{\otimes p}.
\end{equation}
Here the tensor power is its product ideal in \(R\).
The monomials with exponents \(0,\ldots,p-1\) in the
\(X_\tau\) form a basis. Their unscaled evaluation matrix
is a unit matrix in determinant, because the nonconstant
monomials enumerate the \(q-1\) multiplicative characters
on \(F^\times\), extended by zero at zero. Consequently
\[
 \lambda(Q_i)=\tfrac12 q(p-1)\sum_\tau v(c_{\tau i}),
 \qquad d(Q_i)=(p-1)\sum_\tau v(c_{\tau i}).
\]
For a finite locally free \(H\)-torsor \(P\to Q\), the relative
discriminant ideal is the extension of the discriminant of
\(H/R\): this can be checked after the faithfully flat
base change trivializing the torsor. Transitivity of the
trace discriminant then makes
\(\lambda(P)/|P(C)|\) additive in \(H,Q\).
Indeed the generic trace pairing on the split coordinate
algebra is the dot product, whose Gram determinant in an
integral evaluation basis has valuation \(2\lambda(P)\).
Induction on the schematic coordinate flag, using the
rank-one formula, proves
\eqref{eq:torsion-character-discriminant} for \(G\) and every
coordinate quotient \(G/G_i\).

\paragraph{2. The two exact step-graph identities.}
We specify the still-required geometric interface to the
original graph \(\Theta_\tau\).
Assume its first projection is the augmentation ideal \(M^+\),
its second projection is \(L_\tau^\vee\), and its second-coordinate
kernel is \(pL_\tau^\vee\), as in \texttt{Base.th}.
Assume also the three original chart pins: the graph contains
\(((M^+)^2,0)\); it is equivariant under \(F^\times\) acting by
\(g(x)\mapsto g(ax)\) and by multiplication by \(\chi_\tau(a)\)
on the second coordinate; and on the \(\chi_\tau\)-eigenspace
\((g,0)\) in the graph implies \(g\in(M^+)^2\).
These are exactly the pins of
\texttt{chartD\_th\_sqIdeal\_mem\_and\_smul\_mem\_and\_eig\_mem\_sqIdeal}.
They imply the actual map of the graph modulo \(p\) is
\[
 M^+\longrightarrow\omega_{G,\tau}
               \xrightarrow{\ \sim\ } L_\tau^\vee/pL_\tau^\vee.
\]
Indeed the kernel over a fixed first coordinate is
\(pL_\tau^\vee\); apply the character projector (with the
unit denominator \(q-1\)) to obtain surjectivity and
injectivity on the indicated augmentation quotient.

Let \(K_{\tau i}\subset L_\tau^\vee\) be the second
projection of the graph restricted to functions vanishing on
\(F_i\). Since restriction \(M\to\mathcal O(G_i)\) is onto,
the image modulo \(p\) is precisely the kernel of
\(\omega_{G,\tau}\to\omega_{G_i,\tau}\). This follows directly
by lifting a sum of products in the augmentation square.
The preceding exact cotangent sequence gives
\[
 K_{\tau i}/K_{\tau,i+1}
       \simeq\omega_{Q_i,\tau}
       =I_{\tau i}/I_{\phi^{-1}\tau,i}^{\otimes p}.
\]
For the original step graph \(B_{\tau i}\), the map to this
quotient sends \(c\) to the differential class of
\(c\chi_\tau(x_i)\) restricted to \(G_{i+1}\).
Its first projection is \(I_{\tau i}\); for any lift of this
restricted function to \(M\) the class of its second coordinate
modulo \(K_{\tau,i+1}\) is independent of the lift.
Conversely all elements of \(K_{\tau,i+1}\) can be added to the
second coordinate using functions zero on \(F_{i+1}\).
The preceding quotient isomorphism therefore proves the
literal identities
\[
 \operatorname{pr}_2 B_{\tau i}=K_{\tau i},\qquad
 \ker(\operatorname{pr}_2:B_{\phi\tau,i}\to K_{\phi\tau,i})
     =I_{\tau i}^{\otimes p}.
\]
The last kernel is identified via its first coordinate.
This establishes two of the three identities of RC2 without
any product equality of cotangent lengths of the full order
or its intermediate flags. It does not replace the chart pins
by an arbitrary abstract isomorphism of cotangent modules.

\paragraph{3. The coordinate ideal equality from quotient generation.}
Call a function affine modulo \(p\) if, after subtraction of
its value at zero, it is additive on the split points modulo
\(pR\). Suppose the coordinate algebra of \(G/G_i\) is generated
over \(R\) by such functions, for every \(0\le i<n\). This is
precisely the quotient-generation premise \texttt{GenAF}:
invariant functions descend by the finite flat torsor, and
conversely pull back. Under this single premise we prove
\begin{equation}\label{eq:torsion-character-coordinate}
 D_{\tau i}=I_{\tau i}
\end{equation}
without a further full-order tightness hypothesis.

Let \(J_{\tau i}\subset I_{\tau i}\) consist of coefficients
of restrictions to \(Q_i\subset G/G_i\) of affine-modulo-\(p\)
functions in the augmentation \(\chi_\tau\)-eigenspace of
\(\mathcal O(G/G_i)\). Pullback to \(G\) shows
\(J_{\tau i}\subset D_{\tau i}\): such a function modulo
\(p\) is a linear combination of the \(\chi_\tau(x_j)\) with
\(j\ge i\). For clarity the fundamental characters give an
\(R/pR\)-basis of additive coordinate functions, since their
evaluation matrix on an \(\mathbf F_p\)-basis of \(F\)
is invertible modulo the maximal ideal. Its reduction is a
Moore matrix: a nonzero additive polynomial of degree at
most \(p^{f-1}\) cannot vanish at all \(p^f\) field elements.
Projection onto the chosen character is permitted since
\(q-1\) is a unit. The coefficient at \(i\) can be taken to
be the actual value on the \(i\)-th coordinate vector.

By generation and restriction to the closed subgroup \(Q_i\),
the differentials of these functions generate its cotangent
module. Subtract their values at zero and project onto the
\(\chi_\tau\)-component; projection preserves additivity
modulo \(p\). Formula \eqref{eq:torsion-character-cotangent}
gives
\[
 I_{\tau i}=J_{\tau i}+I_{\phi^{-1}\tau,i}^{\otimes p}.
\]
Whenever the second ideal is strictly smaller than the first,
this forces \(J_{\tau i}=I_{\tau i}\), since the first ideal is
principal over a local valuation ring. At an equality step
the result propagates along Frobenius: the \(p\)-th power of
an affine-modulo-\(p\) augmentation eigenfunction is again
one, at the successor character. If every step is an equality
then all \(I_{\tau i}=R\), by cycling the valuations. In that
case \(Q_i\) is constant by the presentation. At least one
\(J_{\tau i}\) must be the unit ideal: otherwise all the
generating functions would restrict to constants modulo the
maximal ideal on this constant subgroup. Indeed after removal
of constants their reductions are additive and hence sums
of their fundamental projections. This would contradict
generation of its full residue function algebra. Propagating
once around the orbit now proves \(J_{\tau i}=I_{\tau i}\)
in every case. We obtain \(I_{\tau i}\subset D_{\tau i}\).

For the reverse conclusion in degrees, let \(m=fn\),
\(h=|N|=p^m\), and let \(A\subset R^m\) be the full additive
coefficient lattice, expressed in any \(\mathbf F_p\)-coordinate
basis. It contains \(wR^m\), is contained in \(R^m\), and
decomposes into the \(A_\tau\) by the unit character projectors
and the invertible Moore change of basis above. Put
\[
 \alpha=v(\det A)=\sum_{\tau,i}v(D_{\tau i}).
\]
In a Smith basis \(\epsilon_j b_j\) for \(A\), the \(b_j\)
are the rows of a unit matrix, and
\(0\le v(\epsilon_j)\le 1/(p-1)\).
Choose lifts \(g_j\in M\) of these linear forms modulo \(p\).
For \(p>2\), the normalized functions \(g_j/\epsilon_j\)
are integral and reduce to the independent linear forms
\(b_j\). For \(p=2\) the same argument works unless
\(v(\epsilon_j)=1\); in this case take the exact lift by
linear combination of the functions \(2\mu(x)=1-(-1)^{\mu(x)}\),
which belong to \(M\) by Cartier duality.

The evaluation matrix of the monomials in the normalized
functions, with each exponent \(0,\ldots,p-1\), is invertible
over \(R\). In the residue field this follows by the invertible
linear change from the \(m\) original coordinates: the new
linear forms generate the full function algebra, and their
\(p\)-th powers are linear combinations of themselves, so the
\(p^m\) indicated monomials span and form a basis.
The unnormalized monomials belong to \(M\) and have
evaluation determinant valuation
\(\frac12 h(p-1)\alpha\).
It follows that
\[
 \lambda(G)\le\tfrac12 h(p-1)\alpha,\qquad
 \alpha\ge d(G)/(p-1)=\sum_{\tau,i}v(I_{\tau i})
\]
by \eqref{eq:torsion-character-discriminant} and cotangent
additivity. The inclusions already proved give the opposite
inequality term by term, and hence equality in every term.
Principal ideals in a valuation ring with the same valuation
are equal. This proves \eqref{eq:torsion-character-coordinate}.

The same determinant calculation gives a convenient sufficient
condition for the quotient-generation premise itself. For any
coordinate quotient \(P\), if its additive coefficient lattice
satisfies \((p-1)\alpha(P)=d(P)\), the above unnormalized
monomials have exactly the determinant valuation of
\(\mathcal O(P)\); their inclusion is therefore equality.
They generate the order by affine-modulo-\(p\) functions.
Thus just the original additive-coefficient degree equality
on these quotient orders suffices, without separate
discriminant identities. This last equality is not asserted
for arbitrary finite flat groups by the argument.

In summary, the annihilator inclusion is automatic for the
finite flat order, the two step-graph identities use only the
literal chart pins, and the remaining coordinate comparison
uses quotient generation (or its sufficient determinant
equality). The proof replaces the full \texttt{FlagConcl}
length-and-ideal comparison as an input to the relative
Chern calculation. At the original geometric test points one
must still realize the specified order, the chart pins and
quotient generation, and the exact Hodge image and kernel
comparisons RC1. In particular the existence of an induced
finite flat group alone does not prove all \texttt{TestFacts}.

\subsection{Flag nefness outside the correction primes}
\label{sec:geometry-uncorrected-flag-nefness}

\begin{lemma}[Degrees from the Hodge and conjugate filtrations]
\label{lem:geometry-two-filtration-nefness}
On a smooth projective connected curve \(C\) over an algebraically
closed field of characteristic \(p>2\), suppose the vector bundles
\(L,H\) have exact sequences
\[
 0\to L^\vee\to H\to L\to0,\qquad
 0\to F_C^*L\to H\to F_C^*L^\vee\to0.
\]
Here \(F_C\) is absolute Frobenius. Suppose these sequences respect
the split unramified CM coefficient action: \(L=\bigoplus_\sigma
L_\sigma\), duality pairs conjugate labels, and Frobenius permutes the
labels by \(\varphi(\sigma)=\sigma^p\). Put \(e=-\deg L\). Then
\[
 e\ge0,\qquad \deg W\le e/(p-2)\quad(W\subset L\text{ a subbundle}),
 \qquad \deg L_\sigma=\deg L_{\bar\sigma}.
 \tag{UN1}
\]
\end{lemma}
\begin{proof}
Let \(P\) be the maximum degree of saturated subbundles of \(L\),
including the zero subbundle. It exists: dualizing a generating map
for a sufficiently positive twist of \(L^\vee\) embeds \(L\) in a
sum of positive line bundles, and exterior powers give upper bounds
for the integer degrees in each rank. In particular \(P\ge0\).

Pull a subbundle of degree \(P\) back by Frobenius and inject it in
\(H\) by the second sequence. Its degree is \(pP\). Project to \(L\)
by the first sequence. The image has degree at most \(P\), since
saturation can only increase degree. The kernel in \(L^\vee\) has
degree at most \(e+P\): for a saturated subbundle \(K\subset L^\vee\),
\(\deg K=-\deg L+\deg (L^\vee/K)^\vee\le e+P\), since the last
dual quotient embeds as a subbundle of \(L\). This bound also holds
without saturation. Additivity of degree gives \(pP\le e+2P\).

Write \(d_\sigma=\deg L_\sigma\) and \(h_\sigma=\deg H_\sigma\).
The first sequence gives \(h_\sigma=d_\sigma-d_{\bar\sigma}\),
and the second gives \(h_{\varphi(\sigma)}=p h_\sigma\).
The permutation has finite orbits, so every \(h_\sigma=0\).
This proves (UN1).
\end{proof}

\begin{proposition}[The actual uncorrected finite-place flag bound]
\label{prop:geometry-uncorrected-flag-nefness}
For the coordinated original integral models, at each prime
\(p\notin S\), \(p>2\), \(p\nmid\operatorname{disc}(E)\), the root
model satisfies \(\mathrm{NefAt}(p,\mu,0)\) for every dominant
weight satisfying the original threshold
\[
 (p-2)(\mu_{\tau,j}-\mu_{\tau,j+1})
 \ge 4 f\|\mu\|_\infty\qquad(\text{every CM type label }\tau).
 \tag{UN2}
\]
Here \(f=[E:\mathbb Q]/2\), the weights are indexed from \(1\) to
\(N\), and \(j,N-j\) are the positive ranks of the paired tangent
eigenspaces. This supplies \(M2\) for the actual good tower,
for which the prime hypotheses follow from
\(p\notin\mathrm{InS}(B)\).
\end{proposition}
\begin{proof}
We first supply the exact sequences on the original test curves,
including curves contained in the boundary. The original unimodular
symplectic lattice with its \(O_E\)-action has a hyperspecial unitary
similitude model at this prime: \(O_E\otimes\mathbb Z_p\) is finite
etale, and the perfect pairing and involution give the reductive
similitude group after unramified splitting. The underlying principally
polarized PEL datum is unchanged (the primitive biextension's diagonal
factor \(2\) is a unit here). Forget the \(p\)-primary part of the
principal level, if present. The remaining level is neat by
\(\mathrm{GoodAt}(p)\). The coordinated level transitions and canonical
normalization construction identify the actual semiabelian family
with the pullback from this prime-to-\(p\) PEL toroidal model, also on
the global dominating model and flag modifications by their
cartesian comparisons.

Use Goldring--Koskivirta, \emph{Strata Hasse invariants, Hecke
algebras and Galois representations}, Invent. Math. 217 (2019),
887--984, Theorem 6.2.1 and Lemma 6.2.3, with their canonical
de Rham and Hodge bundles (Section 6.1). Refine the rational cones
to admit their log-integral Siegel compactification. For the
hyperspecial PEL embedding the theorem gives the Hodge and conjugate
filtrations with their coefficient tensors and Frobenius graded
isomorphisms on the whole toroidal model. Its canonical Hodge bundle
is the invariant cotangent bundle of the semiabelian family.
In de Rham homology the Hodge quotient is \(L=\operatorname{Lie}(G)\)
and the kernel is \(L^\vee\) with conjugate coefficient action by
the polarization. Dualizing the cohomological conjugate filtration
gives subbundle \(F_C^*L\) and quotient \(F_C^*L^\vee\).
The Frobenius maps raise residue embeddings to their \(p\)-th
powers. These are the two sequences in
Lemma~\ref{lem:geometry-two-filtration-nefness}.

For precision, the auxiliary fan refinement does not change the
arithmetic model or its degree. For an integral projective curve
of the original special fibre, base change to the algebraic closure
of the residue field and normalize a component. Pull back the proper
surjective refinement along its map to the prime-to-\(p\) model.
A closed point in the nonempty generic fibre of that pullback has
finite residue degree over the curve function field. Its closure,
normalized, is a smooth projective curve finite and surjective over
the chosen curve. The family, coefficient action and Hodge bundle
pull back along this map, and line-bundle degrees multiply by the
positive degree of the map. It suffices to prove the sign on this
test curve. In particular the special-fibre curve need not have an
abelian generic point. No uniform bound on the auxiliary curve
degree is needed for nefness.

Since \(p\notin S\), the patched tangent lattices agree with the
original tangent eigenlattices on the entire curve. The integral
idempotents of the etale coefficient algebra split the tangent
bundle and preserve the given subbundle flags. The root line
numbered \(i\le j\) is the \(i\)-th graded line of \(L_\tau\);
for \(i>j\) it is the dual of graded line \(N+1-i\) of
\(L_{\bar\tau}\), precisely the original
\(\mathrm{rootIdx},\mathrm{rootSgn}\) prescription.

Write \(q=N-j\); let \(a_{\tau,i},b_{\tau,i}\) be the degrees of
the prefixes of these two flags, starting at zero. By (UN1) all
are at most \(P=e/(p-2)\), with
\(a_{\tau,j}=b_{\tau,q}=d_\tau\) and \(\sum_\tau d_\tau=-e/2\).
Let \(S_\tau\) be the sum of root-line degrees weighted by
\(\mu_{\tau,i}\). Abel summation in the two flags, suppressing
\(\tau\) on the weights, gives
\[
\begin{split}
 S_\tau={}&(\mu_j-\mu_{j+1})d_\tau
 +\sum_{i=1}^{j-1}(\mu_i-\mu_{i+1})a_{\tau,i}\\
 &+\sum_{i=j+1}^{N-1}(\mu_i-\mu_{i+1})b_{\tau,N-i}.
\end{split}
\]
All displayed differences are nonnegative. Put
\(g=\min_\tau(\mu_{\tau,j}-\mu_{\tau,j+1})\) and
\(M=\|\mu\|_\infty\). Because \(d_\tau\le P\),
\[
 S_\tau\le g d_\tau+2MP,\qquad
 \sum_\tau S_\tau\le -ge/2+2fMP\le0
\]
by (UN2). The original numerical class has coefficients
\(-\mu_{\tau,i}\) at the root lines and zero at the extra line
when its last parameter is \(0\). Its curve degree is therefore
\(-\sum_\tau S_\tau\ge0\), exactly the required nefness.
\end{proof}

Together with the corrected-flag theorem at primes in \(S\), this
establishes both finite-place flag-positivity suppliers. It does not
establish the two-level arithmetic Chern-class comparison, which
requires the exact vertical correction sequences and their Frobenius
relation rather than only numerical degrees on curves.

\subsection{A height bound at every arithmetic point}
\label{sec:actual-point-height}

We now use the corrected fibres on the actual model to bound the
prescribed regularized height, including at boundary points. All
degrees and heights are normalized by the degree of the field of
definition. Auxiliary field extensions are allowed, but no bound for
their degrees is assumed or needed.

Retain the original host $E$, $f=[E^+:\mathbf Q]$, signature $(p,q)$,
selected rank $m\in\{p,q\}$, and cutoff $B$. Write
\[
 \Lambda_B=\log B+\log\operatorname{rd}(E)
       +\left|L'(1,\chi_{-d})/L(1,\chi_{-d})\right|,
 \qquad X_B=2\sum_{\ell\in S_B}\frac{\log\ell}{\ell-1}
                         +2\log\operatorname{rd}(E).
 \tag{PH1}
\]
Here $S_B=S_0\cup S_{\rm new}(B)$ is the original full set of
correction primes. The set $S_{\rm new}(B)$ is the prescribed index
set for the additional vertical factors. We use
its proved bound $X_B\le16\Lambda_B$ and assume
$B\ge\max\{128,4f(p+q)^2\}$, as well as the original initial-cutoff
conditions, in particular $|S_0|\le B$ and
$S_B\subset S_0\cup\{\ell\text{ prime}:\ell\le B\}$.
Choose the original admitted normal model with its locally free
corrected modules first, using
Proposition~\ref{prop:actual-hodge-free-model}; form its prescribed
determinant coordinates and vertical factors on that model, and
retain them throughout the argument. This order is part of the
original construction. Flattening a previously nonfree module is
not being asserted to preserve its old determinant: it can introduce
an exceptional vertical divisor. If these modules were already
locally free, their later pullbacks are exact by that proposition.

\paragraph{The point and its determinant.}
Fix a point, extend its field $L$ if necessary, and write $G_z$ for
the original semiabelian fibre and $A_z$ for its abelian quotient.
Let $\Lambda$ be the integral tangent lattice, let $I$ be the
cocharacter lattice of its split generic torus, embedded in the
tangent space by differentiation at one, and let $r$ be its
$\mathcal O_E$-rank. Thus $0\le r\le m$. The integral
arithmetic-point construction gives an exact sequence onto its image
\[
 0\longrightarrow I\otimes_{\mathbf Z}\mathcal O_L
   \longrightarrow\Lambda\longrightarrow V\longrightarrow0,
 \qquad V\subset\operatorname{Lie}(\mathcal A_z/\mathcal O_L),
 \tag{PH2}
\]
where $\mathcal A_z$ is the semistable N\'eron identity component.
This is the original sequence. Over each normal valuation ring,
the closure of the generic torus is a closed smooth torus, and its
generic splitting extends uniquely \cite[Propositions 3.3.1.5
and 3.3.1.7]{point-height-lan}. Its tangent map is injective on
the special fibre, so its integral tangent image is saturated.
The quotient map extends to the N\'eron model by its mapping
property. Its generic tangent kernel is the torus tangent space.
Saturation therefore identifies its integral kernel with
$I\otimes\mathcal O_L$, proving \textup{(PH2)}. Passing to an
eigenlattice takes the image actually obtained, without assuming
that the image is saturated. These identifications are preserved
by semistable base change \cite{conrad-semistable,actual-point-height-source}.

Let $\Lambda_\tau$ be a selected integral eigenlattice and
$V_\tau$ its image in the quotient. Choose $v_1,\ldots,v_r\in I$
independent over $E$ with finite index
\[
                  i_T=[I:\textstyle\sum_a\mathcal O_Ev_a].
\]
Choose $q_{\tau,1},\ldots,q_{\tau,m-r}\in\Lambda_\tau$
whose images form an $L$-basis of the quotient eigenspace. Then
$b_\tau=(e_\tau v,q_\tau)$ is an $L$-basis of the full
eigenspace, where $e_\tau$ is the coefficient projector.
For a full lattice $M$ and an $L$-basis $b$, put
\[
 J(M,b)=\log[M:C]-\log[\mathcal O_L b:C],
 \tag{PH3}
\]
where $C$ is any common sublattice of finite index. This signed
quantity is independent of $C$. With the prescribed quotient
metric $\pi H$, set
\[
 \begin{split}
 d_\tau(F)&=\frac{1}{[L:\mathbf Q]}
   \left(J(F_\tau,b_\tau)
     -\frac12\sum_{\iota:L\hookrightarrow\mathbf C}
                     \log\det\operatorname{Gram}_\iota(q_\tau)\right),\\
 \Phi_B(z)&=\sum_\tau d_\tau(F)-\frac12\log i_T
                                  +\frac{fr}{2}\log(4\pi).
 \end{split}
 \tag{PH4}
\]
Here $F_\tau$ is the fibre of the original corrected bundle.
These formulas are the original labelled determinant and point
quantity, including the signs of the toric index and $4\pi$ term.

\begin{proposition}[The actual determinant bound]
\label{prop:actual-point-determinant-bound}
For every original arithmetic point, including the boundary,
\[
 \Phi_B(z)\le fmX_B+T-fr\log\operatorname{rd}(E)
                                  +\frac{fr}{2}\log(4\pi)
                    \le52fm\Lambda_B,
 \qquad T=\sum_\tau\widehat{\deg}_{\rm n}V_\tau.
 \tag{PH5}
\]
The same bound covers $r=m$ and zero-dimensional abelian quotients.
\end{proposition}

\begin{proof}
We first compute the toric determinant, rather than discarding its
finite index. The trace-dual basis description of the coefficient
projectors identifies the saturated $\tau$-eigenlattice of
$I\otimes\mathcal O_L$ with the $\tau$-image of the determinant
ideal of $I$, multiplied by $\tau(\mathfrak D_E)^r$ on determinants.
To see the factor, first take $I=\mathcal O_E$. In the decomposition
$E\otimes L=\prod_\sigma L$, the integral elements supported on
the $\tau$ factor are precisely
$\tau(\mathfrak D_E)\mathcal O_L e_\tau$. This follows by writing
$e_\tau$ in a trace-dual integral basis; the inverse different is
the trace dual of $\mathcal O_E$. Localization proves the same
identity for every projective $I$, and taking determinants proves
the rank-$r$ assertion. Consequently
\[
 \frac{J((I\otimes\mathcal O_L)_\tau,e_\tau v)}{[L:\mathbf Q]}
       =-r\log\operatorname{rd}(E)+\frac{\log i_T}{2f}.
 \tag{PH6}
\]
The formula also holds for $r=0$, with $i_T=1$.

The exact sequence of lattices in \textup{(PH2)}, restricted to the
chosen eigenlattice and its actual image, gives an exact sequence
of determinant lines. Equivalently, choose local integral bases
for the kernel and image and lift the latter: the change-of-basis
matrix is block triangular. At infinity the quotient block has
the stated Gram determinant. Thus
\[
 d_\tau(\Lambda)
   =-r\log\operatorname{rd}(E)+\frac{\log i_T}{2f}
                              +\widehat{\deg}_{\rm n}V_\tau.
 \tag{PH7}
\]
Theorem~\ref{thm:actual-hodge-specialization} gives
$\mathfrak a_\tau F_\tau\subset\Lambda_\tau\subset F_\tau$
and $\log N\mathfrak a_\tau/[L:\mathbf Q]\le X_B$.
Locally a quotient of a rank-$m$ lattice killed by this ideal
has length at most $m$ times its valuation. Hence
\[
       0\le d_\tau(F)-d_\tau(\Lambda)\le mX_B.
\]
Sum \textup{(PH7)} over the $f$ labels. The term $\log i_T/2$
cancels the negative term in \textup{(PH4)}, proving the first
inequality of \textup{(PH5)}.

Corollary~\ref{cor:actual-selected-point-degrees} gives
$T\le36f(m-r)\Lambda_B$, with both integral indices retained.
Also $X_B\le16\Lambda_B$, $\log\operatorname{rd}(E)\ge0$,
and $\log(4\pi)\le\Lambda_B$. Therefore the right side of
\textup{(PH5)} is at most
$f\{16m+36(m-r)+r/2\}\Lambda_B\le52fm\Lambda_B$.
No slope estimate is applied when the quotient rank is zero.
\end{proof}

\paragraph{The boundary metric on the same determinant.}
We give the limiting calculation needed to compare \textup{(PH4)}
with the regularized metric. On an original normal-crossing chart,
choose the symplectic period coordinates of the semiabelian family.
Their period map has the form
\[
 M(z)(u+\tau(z)v),\qquad
 \tau(z)=T_0(z)+\sum_{a=1}^{h}\frac{\log z_a}{2\pi i}B_a,
 \tag{PH8}
\]
where $M(0)$ is invertible, $M,T_0$ are holomorphic at the centre,
and the $B_a$ are symmetric integral positive semidefinite matrices.
The polarization fixes the symplectic form. The extension of the
semiabelian family identifies the bounded periods with its original
central periods and the radical with its actual torus. The positive
definite form on the abelian quotient is the one from the same
polarization. To recall their construction, the central periods continue
holomorphically and span the central tangent space. Their rational
period subspace is coisotropic. Choose a Lagrangian in it fixed by
all monodromies and extend it to a symplectic basis. The resulting
monodromies have the upper triangular form with integral symmetric
blocks $B_a$. Positivity gives positive semidefiniteness of those
blocks; subtracting their logarithms gives the holomorphic $T_0$
by the removable-singularity argument for these weight-one periods.
Central spanning makes $M(0)$ invertible. Thus these are the period
coordinates of the original boundary construction
\cite{actual-point-height-source}.

Here is the exact cancellation that makes the limit valid for all
rates of approach. Put $K=\ker\sum_aB_a$ and let $\Pi$ be the real
orthogonal projector onto $K$. Positive semidefiniteness gives
$B_a\Pi=\Pi B_a=0$ for every $a$. In period coordinates the raw
Hermitian form is $Y(z)^{-1}$, where $Y=\operatorname{Im}\tau$.
Its quotient by the toric space $K^\perp$ is represented by
\[
 C(z)=\Pi\bigl(\Pi\operatorname{Im}T_0(z)\Pi+1-\Pi\bigr)^{-1}\Pi.
 \tag{PH9}
\]
Indeed $\Pi Y\Pi=\Pi\operatorname{Im}T_0\Pi$. For a vector $y$,
the vector $p=YC(z)y$ satisfies $p-y\in K^\perp$ and is
$Y^{-1}$-orthogonal to $K^\perp$; its pairing with $y'$ is
$y'^*C(z)y$. This proves the formula without an asymptotic error.
The restriction of $\operatorname{Im}T_0(0)$ to $K$ is positive
definite, since it is the Riemann form of the abelian quotient.
Thus \textup{(PH9)} and its transport by $M(z)$ extend
continuously to the centre, including $K=0$ and $K=\mathbf R^g$.

Let $s_\tau(z)$ be the actual selected tangent frames. Use the exact
moving toric period vectors, project them to the coefficient
eigenspaces, and complete them by continuous transverse vectors.
The coefficient projectors vary continuously, so this is an
eigenframe whose change-of-basis determinant from $s_\tau(z)$
is continuous and nonzero at the centre. The determinant of its
Gram matrix factors exactly as
\[
 \det\operatorname{Gram}(\hbox{toric vectors})\,
           \det\operatorname{Gram}_{C(z)}(\hbox{transverse vectors}).
 \tag{PH10}
\]
The coefficient eigenspaces are orthogonal. Thus the quotient by
the toric part in one eigenspace is the restriction of the full
toric quotient. This is the determinant formula for an orthogonal projection,
or the Schur-complement identity. If $n_z(I)$ denotes the original
product of toric Gram determinants divided by $i_T$, division by
$n_z(I)$ cancels the first factors in \textup{(PH10)} exactly.
The remaining factors converge by \textup{(PH9)} and the
nonzero frame determinants. No estimate of a small holomorphic
error times a divergent logarithm is being made.

At the centre write $b_\tau=C_\tau s_\tau$. A cocharacter
derivative has period $2\pi i$ times that derivative. Converting
the raw quotient form to the prescribed metric $\pi H$ therefore
gives the following limiting squared determinant norm:
\[
 \nu_\iota=i_T\prod_\tau
       \frac{(4\pi)^{-r}\det\operatorname{Gram}_\iota(q_\tau)}
                                {|\iota(\det C_\tau)|^2}>0.
 \tag{PH11}
\]
In detail, each selected label contributes $\pi^m$ from the
full tangent metric, $\pi^{-(m-r)}$ from the quotient Gram
matrix, and $|2\pi i|^{-2r}$ from the toric period vectors.
Their product is $(4\pi)^{-r}$. This verifies the normalization
in \textup{(PH11)} directly.

The original cusp correction $U$ satisfies
$U(z)\ge-\log n_z(I)$ for every saturated isotropic $I$.
For the zero lattice this says $U\ge0$. It is the original
first-minimum bound, with the same counting function.
The corrected interior logarithmic weight is the Hodge weight
plus $U/2$. Formulae \textup{(PH9)--(PH11)} show that its lower
limit at the point is at least the algebraic frame weight plus
$\tfrac12\log\nu_\iota$. Call this latter value $x_\iota$.
Changing projective coordinate charts only adds the prescribed
continuous algebraic frame term, so the statement holds in the
original homogeneous coordinates.

\begin{theorem}[The original all-point height floor]
\label{thm:actual-all-point-height}
Let $\overline{\mathcal L}_{B,a/b,t}$ be the prescribed smooth
regularization with positive rational perturbation $a/b$ by the
original ample line $\overline{\mathcal A}$. There is a finite
$K_{\mathcal A}\ge0$, independent of the point, such that for
every fixed $a/b>0$ and all sufficiently small $t>0$,
\[
 h_{\overline{\mathcal L}_{B,a/b,t}}(P)
        \ge-52fm\Lambda_B-(a/b)K_{\mathcal A}
                  \qquad(P\in X(\overline{\mathbf Q})).
 \tag{PH12}
\]
The threshold is chosen before $P$ and its field of definition.
All finite vertical factors and the original complex metrics
are retained. Thus \textup{(PH12)} supplies the all-point
hypothesis of Proposition~\ref{geo-reciprocal-height-limit}.
\end{theorem}

\begin{proof}
First verify the arithmetic metric identity. The determinant
tensor of the inverse corrected bundles is the original
coordinate line of degree $\mu_B$, with its prescribed vertical
correction factors. This can be checked on its local frames:
if $T_{\tau,v}$ changes a rational frame to an integral frame
of $F_\tau$ at a finite place, the transition of the tensor is
the product of $\det T_{\tau,v}^{-1}$. The original coordinate
and vertical transitions have exactly this product. Their ratio
on an arithmetic point is one element $g\in L^\times$, obtained
from the original algebraic frame factor multiplied by its
vertical factors. In particular this is a single algebraic
element at every embedding and every finite place.

More explicitly, choose nonzero homogeneous coordinates $w_{k,c_k}$
at the point and put
\[
 R_v=\prod_k\left(\frac{\max_j|w_{k,j}|_v}{|w_{k,c_k}|_v}
                         \right)^{\mu_{B,k}}.
\]
For the original finite factors write
$t_\ell=(v_{{\rm num},\ell}/v_{{\rm den},\ell})
\prod_k w_{k,c_k}^{-v_{\ell,k}}$ at the point. These are
nonzero: the normalized rational ratio is a unit at every
characteristic-zero point of the original chart. With its
algebraic Hodge frame factor $u_L$, which already includes the
initial $S_0$ correction, the original local
determinant presentation gives exactly
\[
 \left(\prod_\tau|\det T_{\tau,v}|_v\right)R_v=|g|_v,
 \qquad g=u_L\prod_{\ell\in S_{\rm new}(B)}t_\ell.
 \tag{PH13a}
\]
At infinity its vertical factor is
$|\iota(t_\ell)|\prod_k|\iota(w_{k,c_k})|^{v_{\ell,k}}$,
and $\mu_{B,k}=a_{\mu,k}+\sum_{\ell\in S_{\rm new}(B)}v_{\ell,k}$.
These identities verify the transition assertion for the exact
coordinate line, including every finite factor.

The signed index formula is
\[
 J(F_\tau,b_\tau)
   =\sum_{v\nmid\infty}\log|\det T_{\tau,v}|_v
       +\sum_{\iota:L\hookrightarrow\mathbf C}
                                      \log|\iota\det C_\tau|.
 \tag{PH13}
\]
Normalize the finite absolute values so the product formula
sums over all complex embeddings, counted separately.
The finite transition identity just described and
\textup{(PH11)} give
\[
                    h_B^0(P)+\frac1{[L:\mathbf Q]}
                         \sum_\iota x_\iota=-\Phi_B(P).
 \tag{PH14}
\]
Here $h_B^0$ consists of the finite homogeneous-coordinate height
of degree $\mu_B$ and the archimedean logarithms of the
additional vertical factors indexed by $S_{\rm new}(B)$. To check the signs, substitute
\textup{(PH13)} and
\[
 \tfrac12\log\nu_\iota
  =\tfrac12\log i_T-\tfrac{fr}{2}\log(4\pi)
    +\tfrac12\sum_\tau\log\det\operatorname{Gram}_\iota(q_\tau)
    -\sum_\tau\log|\iota\det C_\tau|.
\]
The terms in $C_\tau$ cancel; the remaining algebraic scalar
$g$ cancels by the product formula. What remains is precisely
the negative of \textup{(PH4)}. This proves
\textup{(PH14)} on the actual determinant line.

Let $\varphi$ be the original bounded upper semicontinuous
potential, relative to the fixed smooth reference metric. The
interior is dense. The preceding boundary lower-limit comparison
and upper semicontinuity give
\[
                 x_\iota\le\log h_{\rm ref}(w_\iota)
                                      +\varphi(P_\iota).
 \tag{PH15}
\]
The prescribed smooth regularizations satisfy
$\varphi\le\varphi_t$ at every point for all sufficiently small
$t>0$. For clarity, their threshold is global: a finite cover by
the original inner charts has a positive common margin;
submean inequalities for the circle-invariant smoothing kernel
and the prescribed positive correction give domination on each
inner chart. Their maximum construction gives domination on the
whole compact space. The same margin and smoothness conditions
are fixed before evaluating any point
\cite{actual-point-height-source}.

The exact degree and metric formula for the prescribed rational
line now reads
\[
 h_{\overline{\mathcal L}_{B,a/b,t}}(P)
   =h_B^0(P)+\frac1{[L:\mathbf Q]}\sum_\iota
        \{\log h_{\rm ref}(w_\iota)+\varphi_t(P_\iota)\}
                                      +(a/b)h_{\overline{\mathcal A}}(P).
 \tag{PH16}
\]
The ample height is bounded below: compare a positive tensor
power of $\mathcal A$ with its projective-coordinate height,
whose reference height is nonnegative; the logarithm of the
ratio of continuous positive metrics is bounded on the compact
complex space. This supplies $K_{\mathcal A}$, which may
depend on the fixed model and metric but not on $P$.
Combining \textup{(PH14)--(PH16)} with \textup{(PH5)} proves
\textup{(PH12)}. Properness realizes every algebraic point on the
model, after the harmless field extensions already allowed.
Normalized base change preserves all terms. If $m=0$, the
selected determinant is the unit line and the assertion follows
from the ample-height bound alone.
\end{proof}

The all-point floor is now an established input. The signed
intersection limit, the original weighted mean-height comparison,
and the finite rank losses must still be assembled on the same
family. No whole geometric proof-map node is closed by this
pointwise result alone.

\subsection{The actual determinant intersection}
\label{sec:geometry-height-assembly}

We retain the same arithmetic compactification and the same corrected
tangent bundles as in the point-height theorem.  Put $N=p+q$,
$f=[E:\mathbb Q]/2$, and $n_0=fpq$.  Throughout this section
$1\le q\le p$ and $3q+1\le N$.  The full flag dimension is
$d=fN(N-1)/2$; it is different from $n_0$.
The set of correction primes is the original
$S_B=S_0\cup S_{\mathrm{new}}(B)$, and
\[
 \Lambda_B=\log B+\log\operatorname{rd}(E)
              +\left|\frac{L'(1,\chi_{-d_E})}{L(1,\chi_{-d_E})}\right|.
\]
Here $-d_E$ denotes the host's original quadratic discriminant.
We take $B\ge\max(128,4fN^2)$ and $|S_0|\le B$.
First make the original cotangent modules and all their restricted
images locally free.  Then choose their determinant presentations,
coordinates and vertical factors.  No invariance of a nonfree
determinant under the preliminary modifications is used.

Write $\mathcal F_{B,\sigma}$ for these corrected tangent bundles.
The line underlying the base test is exactly
\[
       \mathcal L_B=
       \bigotimes_{\tau\in\Psi}
                    \det(\mathcal F_{B,\mathrm{sm}(\tau)})^{-1},
 \tag{HA1}
\]
where $\mathrm{sm}$ selects the original rank-$q$ half.  A constant
coefficient ideal in this identity is retained; its restriction to a
curve in a finite fibre has degree zero.  The source's coordinate
multidegree $\mu\mathrm{TwDeg}$ presents this line, including its
initial-prime and added-prime factors\cite{height-assembly-source}.

\begin{proposition}[Nefness on every finite fibre]
\label{prop:actual-finite-fibre-nefness}
The actual line $\mathcal L_B$ is nef on every finite fibre of the
chosen arithmetic model.
\end{proposition}

\begin{proof}
It suffices to prove nonnegative degree after pulling back to the
normalization of every integral projective curve in a finite fibre.
Finite extensions of its finite field do not change the sign.

Suppose first that its residue characteristic $\ell$ belongs to
$S_B$.  Put $n=\ell^{R_\ell}$, with the original $R_\ell\ge4$,
and retain the actual modules $\mathcal N_\ell$ and
$\mathcal Q_{\ell,\sigma}$ of
Section~\ref{sec:geometry-hodge-specialization}.  The chartwise
Hodge--Tate estimate proved there gives
\[
 \kappa_\ell\omega\subset\mathcal N_\ell,
 \qquad
 \kappa_\ell=\begin{cases}4,&\ell=2,\\
                         \zeta_\ell-1,&\ell>2.
             \end{cases}
\]
The algebraic integer $n/(\ell\kappa_\ell)$ is integral on the
original normal model.  For $\ell=2$ this is $2^{R_\ell-3}$;
for $\ell>2$ its valuation at every place above $\ell$ is
$R_\ell-1-1/(\ell-1)\ge0$, with the valuation normalized by
$v(\ell)=1$.  At other places $\kappa_\ell$ is a unit.
Consequently
\[
                        n\omega\subset\ell\mathcal N_\ell.
 \tag{HA2}
\]
Each original torsion differential
$\delta_t\in\mathcal N_\ell/n\omega$ therefore has a well-defined
global image in $\mathcal N_\ell/\ell\mathcal N_\ell$.
The labels here are global for a specific reason.  The original
level compactification has function field $\Omega$ and its
\texttt{FullTorsion} clause makes every geometric $n$-torsion
point $\Omega$-rational.  The generic-point clause identifies
$\Omega$ with the generic stalk; it is not merely a field of
constants containing roots of unity.  The construction with
$S_0\cup S_{\mathrm{new}}(B)$ retains this clause for every prime
in that union.  The local extensions of each original Weil
character, proved in the Hodge-specialization theorem, agree on
overlaps by uniqueness.  Their differentials consequently give
the asserted global sections\cite{height-assembly-source}.
These finitely many images generate it: locally $\mathcal N_\ell$
is generated by lifts of the $\delta_t$ and $n\omega$, and the
second summand disappears by (HA2).  Since
$\mathcal N_\ell\twoheadrightarrow\mathcal Q_{\ell,\sigma}$
is the original image map, its reduction is a surjection too.
Thus $\mathcal Q_{\ell,\sigma}/\ell\mathcal Q_{\ell,\sigma}$
is globally generated.  It is locally free by the chosen model
condition.  Its determinant has nonnegative degree on every such
curve.  The inverse determinant in (HA1) is precisely this determinant,
up to the prescribed constant different ideal.  Summing over the
selected labels proves the assertion at $\ell\in S_B$.

Now suppose $\ell\notin S_B$.  In particular
$\ell\nmid\operatorname{disc}(E)$ and the correction is absent.
The coefficient idempotents are integral after the existing splitting
base, so $\operatorname{Lie}(G)=\bigoplus_\sigma\mathcal T_\sigma$
and $\mathcal F_{B,\sigma}=\mathcal T_\sigma$ near this fibre.
The integral Hodge and conjugate filtrations used in
Theorem~\ref{thm:actual-flat-degree} give
\[
       \deg_C\det\mathcal T_\tau
           =\deg_C\det\mathcal T_{\bar\tau}
                         \quad(\tau\in\Psi).
 \tag{HA3}
\]
This is the finite-fibre character-determinant conclusion of that
proof, not an inference from the complex flat metric.  More explicitly,
the actual first-order extension has character determinant
\[
 \mathcal D_\tau=\det\mathcal T_\tau\otimes
                       (\det\mathcal T_{\bar\tau})^{-1}.
\]
Its Hodge and conjugate filtrations are those of
Theorem~\ref{geo-biextension-conjugate-filtration}.
Proposition~\ref{geo-frobenius-determinant-torsion} iterates their
Frobenius determinant identity around the finite coefficient orbit:
over a splitting field $\mathbb F_{\ell^a}$ it gives
$\mathcal D_\tau^{\otimes(\ell^a-1)}\simeq\mathcal O$.
It is this ratio, not either Hodge determinant separately, that has
degree zero\cite{height-assembly-determinant-record}.

The full Hodge determinant $\lambda_G=\det\omega_G$ is nef on
this fibre.  Here its extension is important.  The actual local
model is the normalization model with its original singleton
symplectic lattice, mapping to an auxiliary Siegel toroidal model;
the universal semiabelian family pulls back to the original family.
The Siegel Hodge line has a positive tensor power equal to the
pullback of an ample line on the minimal compactification
\cite[Theorem 7.2.4.1(2)--(3)]{height-assembly-lan}.
The normalization construction and its morphism to that auxiliary
model are \cite[(7.7)--(7.8), Proposition 7.11]{height-assembly-lan-ramified}.
The global dominating model and the subsequent vertical modifications
pull back the same family, so the equality of Hodge lines persists.
It holds along the boundary as well.  This proves the claimed
nefness of $\lambda_G$ without an assertion about every Hodge
subbundle in positive characteristic.

Finally (HA3) and the eigenspace decomposition give the exact identity
\[
             2\deg_C\mathcal L_B=\deg_C\lambda_G\ge0.
 \tag{HA4}
\]
This proves the proposition at every prime.  The general
characteristic-dependent nefness theorem for all flag weights is
not required for this particular determinant.
\end{proof}

\begin{proposition}[The original cusp potentials and their smoothing]
\label{prop:actual-cusp-smoothing}
Let $\varphi$ be the original logarithmic determinant metric relative
to a smooth reference metric, and let
$\psi=\varphi+U/2$ on the interior.  On the original smooth compact
complex model, $\varphi$ is smooth off the normal-crossing boundary
$D$, bounded above, and satisfies
\[
       \varphi\ge-C\log r-C,
       \qquad r=O(1)+\sum_{a}\log|z_a|^{-2}\ge1.
 \tag{HA5}
\]
The original $\psi$ extends to a bounded upper-semicontinuous
potential.  Both original metric potentials are plurisubharmonic
in holomorphic local frames.  There exist the original circle-invariant
chart smoothings and regularized maxima, with actual overlap errors
tending to zero, such that
\[
 \varphi_t\ge\varphi,\quad \psi_t\ge\psi,\quad
 \theta+dd^c\varphi_t\ge-\epsilon_t\omega_A,\quad
 \theta+dd^c\psi_t\ge-\epsilon'_t\omega_A,
 \qquad\epsilon_t,\epsilon'_t\longrightarrow0.
 \tag{HA6}
\]
Here $\theta$ is the reference curvature and $\omega_A$ is a fixed
positive reference form.  The first family converges locally in
$C^2$ off $D$, and the second pointwise and in $L^1$.  The bounded
family is invariant under complex conjugation, including its
original coordinate and vertical metric factors.
\end{proposition}

\begin{proof}
\emph{The actual local potentials.}
The period charts and the determinant identity of the point-height
theorem give, in a nonvanishing local frame,
\[
 \log h=\operatorname{Re}G+\tfrac12\log\det H(b,b),
 \qquad
 U=\max\bigl(0,\sup_I-\log\mathrm{nz}(I)\bigr).
 \tag{HA7}
\]
The determinant denotes the product over the original selected CM
labels; $G$ is holomorphic.  The supremum runs over the original
nonzero saturated isotropic coefficient sublattices.  Its index
normalization is the same as in the point-height theorem.

Here are direct reasons for the positivity and boundary bounds
needed below.  In a simply connected interior chart the polarized
period matrix is holomorphic, with imaginary part $Y>0$, and its
tangent metric is $Y^{-1}$ after a holomorphic change of frame.
For a holomorphic linearly independent frame $b$, the function
$\log\det(b^*Y^{-1}b)$ is plurisubharmonic.  This follows by
differentiating $Y^{-1}$, using $\partial\bar\partial Y=0$.
For an explicit check in one complex direction, make constant
changes of frame so $Y=1$ and $b=\binom{1}{0}$ at the point;
a holomorphic change of basis of $b$ makes
$\partial b=\binom{0}{C}$ there.  Write
$A=\partial Y=(A_{ij})$ in the same two blocks.  The Hessian
of this logarithm is
\[
       \|C-A_{21}\|_{\mathrm{HS}}^2
          +\|A_{11}\|_{\mathrm{HS}}^2
          +\|A_{12}\|_{\mathrm{HS}}^2\ge0.
\]
Since the direction was arbitrary this proves the assertion.

For a fixed rational isotropic subspace, complete its real period
basis to a symplectic basis.  Taking the quotient metric by its
period vectors is the Schur complement of their Gram block.  In
this basis it is the inverse of the imaginary part of the remaining
principal period block.  The same calculation therefore proves that
\[
          \log\det H(b,b)-\log\mathrm{nz}(I)
 \tag{HA8}
\]
is plurisubharmonic.  All index and trace factors in (HA8) are
constant in the period parameter.  This argument also applies
after restriction to the coefficient eigenspaces, which are
orthogonal and holomorphic; the isotropic coefficient subspace
and its symplectic quotient retain that action.

Only finitely many $I$ can have $\mathrm{nz}(I)\le1$ in a
fixed relatively compact interior chart.  Indeed the metrics
there are uniformly comparable.  The trace determinant formula
identifies $\mathrm{nz}(I)$, up to a fixed positive factor and
fixed positive power for each possible rank, with the norm of
the primitive exterior vector of its saturated integral sublattice.
In each exterior power of the fixed lattice,
a bounded set contains finitely many primitive integral vectors;
each such vector determines its rational subspace and saturation.
The formula for $U$ is consequently a finite maximum locally.
Equations (HA7)--(HA8) prove positivity and continuity off $D$.

Near the boundary use the exact original period form
\[
 \tau=T_0+\sum_a\frac{\log z_a}{2\pi i}B_a,
                    \qquad B_a\ge0,
 \tag{HA9}
\]
from the point-height proof.  On a sector with bounded arguments,
$\operatorname{Re}\tau$ is bounded and
$c\mathbf1\le\operatorname{Im}\tau\le Cr\mathbf1$ after
shrinking the polydisc.  To see the first inequality when
$\operatorname{Im}T_0$ is not positive on the whole space, split
off $\bigcap_a\ker B_a$: it is positive there, while the sum of
the other positive blocks tends to infinity on its complement.
A block completion of squares absorbs the bounded cross terms.
Thus, for real period coordinates $(u,v)$,
\[
 (u+\operatorname{Re}\tau v)^TY^{-1}
       (u+\operatorname{Re}\tau v)+v^TYv
                  \ge\frac c r(|u|^2+|v|^2).
 \tag{HA10}
\]
There is a matching upper bound $Cr(|u|^2+|v|^2)$.
The holomorphic changes of frame and their inverses stay bounded.
Consequently $|\varphi|\le C\log r+C$ and
$0\le U\le C\log r+C$: for the latter use (HA10) in each
exterior power and the positive minimum norm of a nonzero integral
Pl\"ucker vector.  These constants are allowed to depend on this
fixed datum; none enters the final intersection bound.

A plurisubharmonic function with upper growth $C\log r+C$ on a
punctured polydisc is bounded above on a smaller polydisc.  For
completeness, add $\eta\sum_a\log|z_a|$, $\eta>0$, and apply
the maximum principle successively in the punctured coordinates.
At an inner boundary the added negative linear logarithm dominates
$C\log r$; the outer distinguished boundary is a compact subset
of the interior.  Let the inner radii tend to zero and then
$\eta\downarrow0$.  This proves the assertion uniformly on a
smaller tangential compact set.  Apply it to both potentials in
(HA7).  The standard removable-singularity extension is their
upper-semicontinuous limit from the interior.

The corrected potential is also bounded below.  At a boundary
point take its actual saturated toric period sublattice $I$.
The exact toric Gram cancellation proved in the point-height
theorem makes (HA8) converge to a finite real value, along all
approaches to that point.  Since $U\ge-\log\mathrm{nz}(I)$,
it gives a lower bound for $\psi$.  This continuity assertion is
the analytic period and quotient-metric calculation from that
proof; it holds at every complex boundary point, independently
of whether the point is algebraic.  Compactness now makes $\psi$
bounded.  It gives (HA5) for $\varphi$ as well.

\emph{Uniform comparison of the actual local convolutions.}
Take finitely many holomorphic polydisc charts, with nested inner
and middle sets whose closures lie in the charts, and smooth
local reference potentials $q_j$.  On a common fixed neighbourhood
of every middle set, chart transitions and gauge differences have
bounded first derivatives.  Use the original nonnegative smooth
kernel $\rho$, of integral one and compact support in the unit
polydisc, invariant under complex rotations and conjugation.
Such a kernel is obtained by averaging a nonnegative bump over
the compact rotations and conjugation and normalizing its integral.

We give the comparison because boundedness alone does not assert
uniform approximation to an upper-semicontinuous function.  For
bounded plurisubharmonic $u$, let
$M_x(s)=\sup_{|z-x|_\infty\le s}u(z)$.  The maximum principle
on complex annuli shows that $M_x(e^s)$ is convex in $s$.
If $|u|\le B_0$ on the fixed polydisc of radius $\delta$, then,
for $L\ge2$ and sufficiently small $t$,
\[
 M_x(Lt)-M_x\bigl(t/(4L)\bigr)
 \le\frac{2B_0\log(4L^2)}{\log(\delta/(Lt))}
 \le\frac{4B_0\log(4L^2)}{\log(1/t)}.
 \tag{HA11}
\]
The first inequality is the adjacent-secant-slope inequality for
that convex function.

If $b_\rho=\sup\rho$ and
$c_\rho=b_\rho\pi^{n_0}(3/2)^{2n_0}$, the submean inequality
at a maximizer on the polydisc of radius $t/4$ gives
\[
  M_x(2t)-M_tu(x)
       \le c_\rho\bigl(M_x(2t)-M_x(t/4)\bigr).
 \tag{HA12}
\]
Indeed the polydisc of radius $3t/2$ about that maximizer contains
the kernel support and lies in the polydisc of radius $2t$ about
$x$.  The integral of the nonnegative function $M_x(2t)-u$ over
it is at most its volume times $M_x(2t)-M_x(t/4)$.  Multiplication
by the bound $b_\rho t^{-2n_0}$ for the kernel proves (HA12).
The reverse bound $M_tu(x)\le M_x(t)$ is immediate.

Apply (HA11)--(HA12) to $u_j=q_j+\psi$.  A change from chart
$j$ to $k$ sends a polydisc of radius $t$ into one of radius
$Lt$; its smooth gauge difference changes by at most $Lt$.
Subtracting the value of that gauge at the centre therefore gives
\[
 |M_tu_j-q_j-(M_tu_k-q_k)|
                 \le\frac C{\log(1/t)}+Ct
 \tag{HA13}
\]
on every middle-set overlap, with one $C$ for the finite atlas.
This is the required original overlap estimate.  It does not
assert that $M_tu_j-q_j$ converges uniformly to $\psi$.

For $\varphi$ use the original local truncations
$u_{j,K}=\max(q_j+\varphi,-K)$ with their removable
plurisubharmonic extensions.  Choose
$K(t)=A\log\log(1/t)$ with $A$ sufficiently large.
Every polydisc of radius $ct$ contains a point whose normal
coordinates have modulus at least $c't$.  Equation (HA5) thus gives
$\sup(q_j+\varphi)\ge-C_0\log\log(1/t)-C_1$ there.
Choosing $A>C_0$ absorbs the bounded gauge differences.  In the
transfer of maxima between charts, the added constant branch
$-K(t)$ is then below the transferred maximum.  The preceding
comparison applies without a constant gauge error from truncation.
Its oscillation bound is $O(K(t)+1)$, so the exact resulting error is
\[
 |M_tu_{j,K(t)}-q_j-(M_tu_{k,K(t)}-q_k)|
         \le C\frac{\log\log(1/t)}{\log(1/t)}+Ct.
 \tag{HA14}
\]
These are precisely the two original chart constructions and
their admissible overlap rates\cite{height-assembly-regularization-source}.

Choose each error at least $t$, and use the original regularized
maximum of the local expressions plus $4\epsilon(t)\eta_j$,
where $-1\le\eta_j\le0$, $\eta_j=0$ on the inner set, and
$\eta_j=-1$ near the edge of its middle set.  The overlap estimates
make the edge terms inactive.  The output is smooth.  The convex,
nondecreasing regularized maximum has derivative sum one, so its
curvature is bounded below by $-C\epsilon(t)\omega_A$.
An inner chart and the submean inequality prove domination of the
original potential.  At every point the bounded construction is
between its inner-chart convolution and that convolution plus
$2\epsilon(t)$; this proves its pointwise and $L^1$ convergence.
For the unbounded construction the truncation is inactive on each
compact subset off $D$ once $t$ is small.  There convolution
converges smoothly.  Differentiating the regularized maximum
twice adds errors $O(\epsilon(t)+t^2+t^4/\epsilon(t))$, which
tend to zero because $\epsilon(t)\ge t$.  This gives the asserted
local $C^2$ convergence and the uniform upper bound.

Finally include every chart with its conjugate chart, choose the
reference metric invariant under conjugation, and pair the cutoffs.
The original cusp norm and determinant metric are conjugation
invariant.  Change of variables in the kernel integral and
permutation invariance of the regularized maximum give exact
invariance of $\psi_t$.  Each original vertical factor has this
invariance too, as the norm of its rational arithmetic presentation.
The chosen smooth representative consequently agrees with the
original metric on the entire original complex model; no metric
averaging changes its values.
\end{proof}

\begin{theorem}[The original signed base intersection]
\label{thm:actual-signed-base-intersection}
For every original component $Z$, let $A_{B,Z}$ be the limiting
arithmetic top intersection of $\mathcal L_B$ with its original
determinant metric, and let $d_{B,Z}$ be its geometric top degree.
Use the original height and field normalization in both quantities.
Then these limits exist, $d_{B,Z}>0$, and
\[
 \frac{A_{B,Z}}{(n_0+1)d_{B,Z}}
                \ge-\frac{107}{2}fq\Lambda_B.
 \tag{HA15}
\]
For the actual full-flag test polynomial
\[
 \mu=\sum_{\tau}\sum_{i\ge p}X_{\tau i},\qquad
 V=\prod_\tau\prod_{\substack{a<b\\a,b<p\ \mathrm{or}\ a,b\ge p}}
                              (X_{\tau b}-X_{\tau a}),
 \qquad F_0=\mu^{n_0+1}V,
\]
write $J=\deg_{\mathrm{Fl}}(\mu^{n_0}V)>0$ and
$v=\sum_Zv_Z>0$, with the original component volumes $v_Z$.
The original normalized test intersection therefore satisfies
\[
 \beta_B:=\frac{\sum_Z I_{B,Z}(F_0)}{vfq(n_0+1)J}
                       \ge-\frac{107}{2}\Lambda_B.
 \tag{HA16}
\]
All initial-prime and added-prime corrections remain in these
intersections.
\end{theorem}

\begin{proof}
Fix one datum and component before taking any limit.  Write
$a_B=-52fq\Lambda_B$ for the proved all-point lower bound.
Let $A$ be the original auxiliary ample bundle, with positive
smooth metric and point-height lower bound $-K_A$.  For a fixed
positive rational number $z=a/b$, let $F_t(z)$ and $G_t(z)$ be
the original arithmetic intersections of the $(n_0+1)$st powers
of
\[
       (\mathcal L_B,\varphi_t)+zA,
                   \qquad (\mathcal L_B,\psi_t)+zA,
\]
respectively.  Equivalently use their genuine integral tensor
powers $\mathcal L_B^b\otimes A^a$ and divide the arithmetic
intersection by $b^{n_0+1}$, and the geometric degree by $b^{n_0}$.

Proposition~\ref{prop:actual-finite-fibre-nefness} gives relative
nefness at every finite place.  Proposition~\ref{prop:actual-cusp-smoothing}
makes the curvature of the second line positive when $t$ is
sufficiently small for this fixed $z$: its negative error is
absorbed by $z\omega_A$.  It is then ample on the generic fibre
and has its original conjugation-invariant smooth metric.
The point-height theorem gives, for every algebraic point and
the same sufficiently small $t$,
\[
                       h_{B,z,t}(P)\ge a_B-zK_A.
 \tag{HA17}
\]
The smoothing threshold is chosen before $P$.  Zhang's lower
intersection inequality, with precisely these finite and infinite
positivity hypotheses, gives
\[
         G_t(z)\ge(n_0+1)(a_B-zK_A)d(z),
 \qquad
 d(z)=\int_X(T+z\omega_A)^{n_0},\quad T=\theta+dd^c\varphi.
 \tag{HA18}
\]
The quoted inequality and its normalization are established in
the point-to-intersection argument\cite{height-assembly-height-record}.

The two metrics in this comparison have the same algebraic line,
finite metrics, rational factors and vertical factors.  The exact
metric-variation identity and the two energy convergence results
therefore give actual limits $F(z)$ and $G(z)$ and
\[
 G(z)-F(z)=(n_0+1)\{E_z(\psi)-E_z(\varphi)\}
       \le\frac{n_0+1}{2}\int_X U(T+z\omega_A)^{n_0}.
 \tag{HA19}
\]
The factor $1/2$ occurs only because $\psi-\varphi=U/2$;
the metric-variation formula has no additional half.
All the mixed integrals on the right are finite.  Indeed
$U=2(\psi-\varphi)$, $\psi$ is bounded, and the finite-energy
boundary theorem supplies every mixed integral weighted by
$|\varphi|$.  Polynomial interpolation in a positive auxiliary
curvature shows that $F(z)$ extends to $z=0$ with
$F(0)=A_{B,Z}$.  This is the original signed limit, not a new
positive part of that limit.

Taking $t\to0$ in (HA18), then subtracting the upper bound (HA19),
and finally taking $z\downarrow0$, gives
\[
 A_{B,Z}\ge(n_0+1)
       \left(a_B-\frac{\overline U_Z}{2}\right)d_{B,Z},
 \qquad
 \overline U_Z=\frac{\int_XUT^{n_0}}{\int_XT^{n_0}}.
 \tag{HA20}
\]
The geometric denominator is strictly positive: the original
interior curvature is the positive invariant volume density,
and the arithmetic component has positive finite volume.
The exact curvature-to-Haar dictionary, including the original
sheet multiplicities, gives
$\overline U_Z\le3fq\Lambda_B$ by the proved isotropic counting
and cusp-mean estimate\cite{height-assembly-cusp-record}.
This proves (HA15).  The order of inequalities is essential:
the bounded-metric intersection has a lower bound, and its
increase over the original metric has an upper bound.

We finish by checking that this is the actual flag test.
On the original flag model, $\mu$ is the first arithmetic Chern
class of the pullback of $\mathcal L_B$.  The determinant of a
metric flag is the product of the metrics on its successive
quotients, so this is an isometry as well as an algebraic line
identity.  It includes the original corrections at all primes.
The class $V$ has the fibre dimension $d-n_0$.
On a common projective model comparing the original flag and base
presentations, let $\delta>0$ be the original projection degree,
and let $c_{\mathrm{fib}}$ be the degree of $V$ on the generic
fibre.  The arithmetic projection formula for the smooth metrics,
and the geometric projection formula, give respectively
\[
 \delta I_{B,Z}(F_0)=c_{\mathrm{fib}}A_{B,Z},\qquad
 \delta v_ZJ=c_{\mathrm{fib}}d_{B,Z}.
 \tag{HA21}
\]
For the first equality take the just-proved metric limit after
the smooth projection formula.  Pullback of each already locally
free corrected bundle, and of its actual determinant presentation,
is exact on the common model.  Thus the vertical factors in the
two formulas are identical.  The second equality makes
$c_{\mathrm{fib}}>0$ and retains every finite projection and
component factor.  This is also the precise common-cover
calculation in the original test identification
\cite{height-assembly-source}.

Canceling $\delta$ and $c_{\mathrm{fib}}$, rather than setting
either to one, yields
\[
 I_{B,Z}(F_0)=
       \frac{A_{B,Z}}{(n_0+1)d_{B,Z}}\,(n_0+1)v_ZJ.
\]
Summing gives
$\beta_B=(fq)^{-1}\sum_Z(v_Z/v)
A_{B,Z}/((n_0+1)d_{B,Z})$, proving (HA16).
The initial corrections are already in the original determinant
presentation, while the new primes occur in its displayed
vertical product.  They have not been discarded because their
complex curvature vanishes.
\end{proof}

\subsection{Integral parallel twists on the actual cover}
\label{sec:actual-integral-parallel-twists}

\begin{proposition}\label{prop:actual-unramified-frames}
Use the original good cover, with $A=1$, coefficient field $K$,
and level $M$.  Put
\[
       B=\{\ell:\ell\text{ prime},\ \ell\mid
                                  2d\operatorname{disc}(E)M\}.
\]
For every original component and every finite place $v$ of $K$
above $\ell\notin B$, there is an original integral framed point
on that component over a finite unramified extension of $K_v$.
Its normalized valuation index over $K_v$ is one.  In particular,
the original $\mathrm{UnramPts}(B)$ assertion holds.
\end{proposition}

\begin{proof}
The original PEL order is the maximal order $\mathcal O_E$,
its alternating lattice is self-dual, and $\ell$ divides neither
its coefficient discriminant nor its level.  The datum is of type A,
so the extra invariant $I_{\rm bad}$ of Lan's definition is one.
These facts check every good-prime condition in
\cite[Definition 1.4.1.1]{theta-twist-lan}.  The maximal-order
condition in his Condition 1.4.3.10 is automatic here.
Choose a smooth projective admissible fan and the exact original
full level completed from the half-level.  The completion is the
one already proved in the actual level construction; since $\ell$
does not divide $M$, the prescribed cyclotomic fibre is an open and
closed finite-\'etale choice.  Theorem 6.4.1.1 of the same reference
gives the original smooth proper toroidal PEL model at $\ell$,
with relative normal-crossings boundary and the tautological family.
Base change to $R=\mathcal O_{K_v}$ preserves all these properties.

In a proper smooth family, the Stein algebra is finite \'etale and
its fibres are the connected-component algebra.  The selected
geometrically integral generic component therefore extends to an
open and closed proper smooth component over $R$; its generic
constant field is already $K_v$.  It has a nonempty special fibre.
A relative normal-crossings divisor contains no component of a
fibre, so this special fibre has a closed point outside the
boundary.  Its residue field is a finite extension of the finite
residue field of $R$.  Pass to the corresponding finite unramified
extension $R'/R$.  Smoothness and Hensel's lemma lift that point to
an $R'$-point of the same abelian interior.

The lifted object retains the original polarization, coefficient
action, both labelled level maps, and primitive rigidified line.
Since $\ell\nmid\operatorname{disc}(E)$ and $K$ contains every
specified image of $E$, the coefficient algebra over $R'$ splits
as the product of the original embedding factors.  Its tangent
eigenspaces are free with precisely the prescribed ranks; choose
bases.  The primitive principal line has rank-one space of sections
over $R'$, compatible with base change, so choose a generator and
its original rigidified generic comparison.  These choices give
exactly an original integral framed point.  They are evaluations
of the finite-type framed moduli charts, hence satisfy the original
test-frame condition.  The residue field remains finite: the
source's $\mathrm{IntFr.IsAlg}$ means precisely this condition,
not algebraicity of the complete fraction field over $K$
\cite{theta-twist-source}.

The extension $R'/R$ is unramified, so restriction of its integer
valuation is the integer valuation of $K_v$.  Every evaluated rule
has value zero or an integral normalized valuation.  This is the
literal index-one criterion defining $\mathrm{UnramPts}(B)$.
\end{proof}

\begin{theorem}\label{thm:actual-short-parallel-twist}
Use the actual flat determinant sections of
Theorem~\ref{thm:actual-flat-degree}.  On each original component,
for every integer parallel shift $w=(w_\tau)_\tau$, there is a
nonzero section $t_w$ of the original parallel weight
$\operatorname{parWt}(w)$, with theta tag zero, such that
\[
 v(t_w)\ge0\quad\hbox{at every original integral test frame},
\]
and, at every complex embedding and every frame on the component,
\begin{equation}\label{eq:actual-parallel-twist-bound}
 \log\|t_w\|\le
 \frac{|p-q|}{2}|P_E|\sum_\tau|w_\tau|
       +\frac12\log\operatorname{rd}(K)
       +\sum_{\ell\in B}\log\ell+1.
\end{equation}
Both $t_w$ and its generic inverse have constant complex norm and
preserve the original cusp spaces under multiplication, with the
corresponding weight shifts.  The inverse is not asserted integral.

Under $|P_E|\le3L_b$, $\log\operatorname{rd}(E)\le3L_b$,
$\log M\le30L_b$, $\log d\le L_b$, $f=[E:\mathbf Q]/2\ge1$,
and $L_b\ge1$, the right side is at most
\[
       \tfrac32 rL_b\sum_\tau|w_\tau|
                  +\tfrac12\log\operatorname{rd}(K)+40fL_b+1.
\]
\end{theorem}

\begin{proof}
Let $s_\tau$ be the actual uncorrected determinant generator,
with the original tangent-frame norm and finite valuations.  Form
\[
                       t=\prod_\tau s_\tau^{-w_\tau}.
\]
The identity $\operatorname{parWt}(w)=
\sum_\tau(-w_\tau)\operatorname{flatWt}(\tau)$ shows that this
is the required original weight, with theta tag zero.
Let $r_\sigma>0$ be its constant complex norm and $\rho_v$ its
point-independent normalized finite valuation.  The latter have
finite support.  By Proposition~\ref{prop:actual-unramified-frames},
$\rho_v$ is an integer whenever $v$ lies above a prime outside $B$:
evaluate the same constant valuation at the index-one point.

The uncorrected degree of $s_\tau$ is $(p-q)P_E/2$; hence
\[
 D(t)=\sum_v\rho_v\frac{\log N v}{[K:\mathbf Q]}
         -\frac{1}{[K:\mathbf Q]}\sum_\sigma\log r_\sigma
       =-\frac{p-q}{2}P_E\sum_\tau w_\tau.
\]
We use the uncorrected frame here deliberately.  The correction
$q(\log d-\log\operatorname{rd}(E))$ in
$\mathrm{FlatSec}$ changes its degree bookkeeping, not the
actual evaluated section or its integral tangent lattice.

For completeness, put
$I=\prod_v\mathfrak p_v^{\lceil-\rho_v\rceil}$, a nonzero
fractional ideal.  Its logarithmic norm satisfies
\[
 \frac{\log N I}{[K:\mathbf Q]}
 \le-\sum_v\rho_v\frac{\log N v}{[K:\mathbf Q]}
                     +\sum_{\ell\in B}\log\ell.
\]
Indeed rounding costs zero outside $B$, and for each remaining
rational prime
$\sum_{v\mid\ell}\log N v\le[K:\mathbf Q]\log\ell$.
Minkowski's theorem, applied to the embedding box with radii
$R/r_\sigma$ (discs at complex places), gives $b\in I\setminus\{0\}$
with
\[
 \log(r_\sigma|\sigma b|)
 \le -D(t)+\tfrac12\log\operatorname{rd}(K)
                 +\sum_{\ell\in B}\log\ell+1
 \quad\hbox{for all }\sigma.
\]
The additive constant $1$ makes the volume strictly larger than
$2^{[K:\mathbf Q]}$ times the lattice covolume; the conjugate
radii agree because the original metric respects complex
conjugation.  This is the exact scalar construction, with the
field-degree normalization retained.

Set $t_w=bt$.  The inequalities
$v(b)\ge\lceil-\rho_v\rceil$ prove integrality at every original
test frame, and $-D(t)\le |p-q||P_E|\sum|w_\tau|/2$ proves
\eqref{eq:actual-parallel-twist-bound}.  The determinant generator
and its inverse extend as units on the generic compactification;
equivalently their original normalized norms are constant.
Multiplication by them therefore preserves holomorphy, the row
weight conditions, and cusp vanishing, in both directions.  This
argument does not require the separate theta key section.

Finally
\[
 \sum_{\ell\in B}\log\ell
 \le\log2+\log d+2f\log\operatorname{rd}(E)+\log M
 \le(6f+32)L_b\le40fL_b.
\]
\end{proof}

This supplies the actual unramified-point and integral parallel-shift
inputs.  Section~\ref{sec:actual-key-short-span} constructs the separate
key-weight shift with its finite-place normalization and applies both
shifts to the exact short-family consumer.

\subsection{A peak section on the actual thick quotient}
\label{shared-bergman-actual-thick}

The following argument concerns the original good cover and its
Hermitian line, at one coefficient embedding.  It uses a complete
quotient metric, so no estimate for a separately periodized peak or for
an arithmetic fundamental volume is needed.  Write
$f=[E:\mathbb Q]/2$, $A=fr$, and $D=f\binom r2$.
At an orthogonal frame the sign matrix is
$J=\operatorname{diag}(1_p,-1_q)$, where $p,q\geq1$ and $p+q=r$.
For an integral weight $\lambda$ put
$S=\sum_{\tau,i}|\lambda_{\tau i}|$.
The original norm of a section is denoted by $\|P\|_\lambda$.

\begin{theorem}[The actual thick peak]
\label{bergman-actual-thick-peak}
There is an absolute constant $C\geq1$ with the following property.
On any component of the original good cover, suppose
\[
 \begin{gathered}
 \lambda_{\tau i}-\lambda_{\tau,i+1}\geq CA^8
       \quad(i+1\ne p),\\
 \lambda_{\tau,r-1}-\lambda_{\tau0}\geq CA^8
       \quad\text{for every }\tau.
 \end{gathered}
\]
Let $x$ be an orthogonal frame with
\begin{equation}
 \label{bergman-actual-thick-test}
  M\,\operatorname{firstMinimumSq}(x)\geq100A^4,
\end{equation}
where $M$ is the original level and the auxiliary level factor is one.
Then the original cusp space contains a section $P$ with $P(x)\ne0$ and
\begin{equation}
 \label{bergman-actual-thick-peak-bound}
 e^{-C(S+A^8)}\int_{\mathcal F}\|P\|_\lambda^2\,d\mu_F
 \leq \|P(x)\|_\lambda^2\,\mu_F(\mathcal F)
\end{equation}
for every original fundamental set $\mathcal F$.
The same conclusion holds whenever the original chart ball of radius
$1/(100r^2)$ injects in the quotient, without assumption
\eqref{bergman-actual-thick-test}.
\end{theorem}

All arithmetic hypotheses here are those of the constructed good cover:
the original self-dual lattice, full level markings, fine moduli
uniformization and its exact automorphic dictionary.  The assertion is
not an extension to an arbitrary analytic lattice with an unspecified
multiplier.  In particular it applies to the untagged weights in the
original thick Bergman consumer.\cite{bergman-thick-source}

\begin{proof}
\emph{A tensor test for a short return.}
For Euclidean lattices $L_1,L_2$ and a nonzero integral tensor
$T\in L_1\otimes L_2$ of rank $k$,
\begin{equation}
 \label{bergman-tensor-minimum}
 \|T\|_{\mathrm{HS}}\geq
      \frac{\lambda_1(L_1)\lambda_1(L_2)}{\sqrt{k}}.
\end{equation}
Indeed, choose a basis $u_1,\ldots,u_k$ of the saturated first support
in $L_1$ and write $T=\sum_i u_i\otimes v_i$.  The $v_i$ are independent
elements of $L_2$.  If $U,V$ are their Gram matrices, the arithmetic--
geometric mean inequality for the positive eigenvalues of
$U^{1/2}VU^{1/2}$ gives
\[
 \|T\|_{\mathrm{HS}}^2
 \geq k(\det U\det V)^{1/k}.
\]
The cube form of Minkowski's theorem gives
$\operatorname{covol}(L)^{1/k}\geq\lambda_1(L)/\sqrt{k}$
for every rank-$k$ Euclidean lattice.  Apply it to the two support
lattices, whose first minima are at least those of $L_1,L_2$.
This proves \eqref{bergman-tensor-minimum} without a restriction on rank.

For the original lattice $\Lambda$, use the actual majorant
\[
 Q_x(v)=\sum_\tau2|\tau\xi_0|\,
                 \operatorname{maj}_{x,\tau}(v_\tau).
\]
Its real dimension is $N=2fr$.  If $L$ is its Euclidean realization,
the original alternating trace self-duality identifies $L^*$
isometrically with $L$.  More explicitly, the two trace-dual
realizations $\phi,\psi$ satisfy
\[
 \langle\phi u,\psi w\rangle=\operatorname{Tr}_{E/\mathbb Q}h(u,w),
 \qquad \|\psi(-\xi_0w)\|^2=Q_x(w),
 \qquad L^*=\psi(-\xi_0\Lambda).
\]
Thus $\lambda_1(L)^2=\lambda_1(L^*)^2$
is exactly $\operatorname{firstMinimumSq}(x)$; there is no division by
$f$ and no replacement by a rank-one covolume.

Suppose $\gamma$ preserves $h,\Lambda$, is congruent to one modulo $M$,
and takes two points of the stated small chart ball to the same flag.
The two chart majorants differ from $Q_x$ by factors between $1/2$ and
$2$, so $\|\gamma\|_{\mathrm{op}}\leq2$ in $Q_x$.
For completeness, write a chart frame as $x(I+N(z))$, with $N$ strictly
triangular.  Each coordinate is bounded by $\|z\|$, and
$\|N\|_{\mathrm{op}}\leq r\|z\|\leq1/(100r)$.
The leading Gram matrices are perturbations of the corresponding
blocks of $J$.  Indefinite Gram--Schmidt, or its successive Schur
complements, gives a unitary representative with norm and inverse norm
at most $\sqrt2$.  Its ambiguity is the compact diagonal torus.
This proves the asserted majorant comparison and return bound.

If $\gamma\ne1$, then
$(\gamma-1)/M\in\operatorname{End}_{\mathbb Z}\Lambda=L\otimes L^*$
is nonzero.  Equation \eqref{bergman-tensor-minimum} yields
\[
 M\operatorname{firstMinimumSq}(x)
 \leq \sqrt N\,\|\gamma-1\|_{\mathrm{HS}}
 \leq3N=6A.
\]
This contradicts \eqref{bergman-actual-thick-test}.

One must still pass from this principal subgroup to the original
uniformizing group.  The actual moduli comparison supplies this step.
An identification of two flags by an original return gives an
$\mathcal O_E$-linear, polarization-preserving isomorphism of their
marked complex tori.  On their homology lattice it is a unitary
$\gamma'$ preserving $\Lambda$.  Preservation of the two original
level markings gives congruence on the two conjugate level ideals.
Their coprimality and trace self-duality give
$(\gamma'-1)\Lambda\subset M\Lambda$.  In the general mixed level the
modulus is $\operatorname{lcm}(M,A_{\rm lev})$; here
$A_{\rm lev}=1$, so it is precisely $M$.
The two lifted flags differ by the original upper triangular frame
change.  The preceding argument makes $\gamma'=1$; faithfulness of
the fine level uniformization then makes the original return trivial.
This uses the actual orbit comparison, not an assertion that its
original group is itself principal.\cite{bergman-thick-source}

\emph{A complete metric with a uniform local chart.}
Let $\mathscr D=\prod_\tau U(p,q)/T$ be the full flag domain and
$Y=\Gamma\backslash\mathscr D$ its actual component.
For one embedding put
\[
 \lambda^0_i=\begin{cases}-i&i<p,\\ r-i&i\geq p.
                        \end{cases}
\]
The original squared line metric in the chart $I+N(z)$ is
$e^{-\varphi_\lambda}$, where
\begin{equation}
 \label{bergman-exact-metric}
 \varphi_\lambda(z)=\sum_{\tau,i}\lambda_{\tau i}
       \bigl(\log|\Delta_{\tau,i+1}(z)|-
             \log|\Delta_{\tau,i}(z)|\bigr).
\end{equation}
Here $\Delta_{\tau,j}$ is the leading $j$ by $j$ Gram determinant,
and $\Delta_{\tau,0}=1$.
Differentiation at zero gives the coefficient
\[
 \epsilon_i\epsilon_j(\lambda_{\tau i}-\lambda_{\tau j})
       \quad(i<j),\qquad
 \epsilon_i=\begin{cases}1&i<p,\\-1&i\geq p.\end{cases}
\]
For $\lambda^0$ these coefficients are between $1$ and $r$.
Write $\omega_0=i\partial\bar\partial\varphi_{\lambda^0}$.
It is a positive invariant K\"ahler form.  Its homogeneous Riemannian
metric is complete, and so is the quotient metric on $Y$.
If $m$ is the least of the gaps in the theorem, telescoping the gaps,
including the gap from the last to the first weight, gives globally
\begin{equation}
 \label{bergman-curvature-positive}
 i\Theta(L_\lambda)\geq m\omega_0.
\end{equation}
The invariant canonical weight is the sum of the flag roots; its
successive coefficients have absolute value at most $2r$.
The same root calculation therefore gives
$\operatorname{Ric}(\omega_0)\geq-Cr\omega_0$.

There is an absolute $c>0$ such that, on the coordinate ball of radius
\[
 \eta=cA^{-3},
\]
the metric $\omega_0$ is between one half and $2r$ times the Euclidean
metric, and $|\varphi_\lambda|\leq CS$.
Here and below $c$ can be decreased by an absolute factor.
These are uniform matrix estimates, not an appeal to compactness in
varying dimension.  On this ball each leading Gram matrix has inverse
norm at most $2$.  Differentiate
$d\log\det G=\operatorname{tr}(G^{-1}dG)$ three times.
For unit Hilbert--Schmidt directions every resulting trace is bounded
by $Cr$; the sum of the absolute coefficients of $\lambda^0$ is at
most $r^2$.  Thus the third derivative is bounded by $Cr^3$ on each
embedding block.  The blocks are orthogonal, so this is also an
operator bound for their product.  It gives the stated Hessian
comparison for $\eta=cA^{-3}$.  The same determinant formula bounds
\eqref{bergman-exact-metric} by $CS$.
This smaller ball lies in the injective original chart.

\emph{Solve once on the quotient.}
Choose a holomorphic line frame $e$ with $\|e(x)\|=1$ and a smooth
cutoff $\chi$, equal to one for $|z|\leq\eta/8$, supported in
$|z|<\eta/4$, with $|\bar\partial\chi|_{\omega_0}\leq C/\eta$.
Let $\chi_0$ be one on $|z|\leq\eta/2$, zero for $|z|\geq\eta$,
and define on the quotient
\[
 \psi_\varepsilon=(D+1)\chi_0(z)
       \log\frac{|z|^2+\varepsilon^2}{\eta^2+\varepsilon^2},
       \qquad0<\varepsilon<\eta/16.
\]
Extend it by zero outside the chart.  Then $\psi_\varepsilon\leq0$ and
\[
 i\partial\bar\partial\psi_\varepsilon
       \geq-C(D+1)\eta^{-2}\omega_0\geq-CA^8\omega_0.
\]
The logarithm is plurisubharmonic where $\chi_0=1$; on the cutoff
annulus it and its first two scaled derivatives have absolute bounds.
Since $D\leq A^2$, the power $A^8$ includes every dimension factor.

Regard $\bar\partial(\chi e)$ as a $(D,1)$ form with values in
$L_\lambda\otimes K_Y^{-1}$.  Its canonical identification with an
$L_\lambda$-valued $(0,1)$ form is isometric for the induced metrics.
By \eqref{bergman-curvature-positive}, the Ricci estimate and a large
enough absolute $C$ in the gap hypothesis, this bundle with weight
$e^{-\psi_\varepsilon}$ has curvature at least $\omega_0$.
The complete K\"ahler $L^2$ theorem\cite{bergman-demailly}
therefore solves
$\bar\partial u_\varepsilon=\bar\partial(\chi e)$ with
\begin{equation}
 \label{bergman-solved-norm}
 \int_Y\|u_\varepsilon\|_\lambda^2e^{-\psi_\varepsilon}\,dV_0
 \leq C\eta^{-2}e^{C(S+D)}\operatorname{vol}_0(B_\eta).
\end{equation}
The right side follows on the support of $\bar\partial\chi$, where
$e^{-\psi_\varepsilon}\leq64^{D+1}$.
No pseudoconvexity assumption on the quotient is used.

Take a weak $L^2$ limit.  On every closed annulus about $x$ the weights
converge uniformly, and weak lower semicontinuity followed by monotone
exhaustion gives \eqref{bergman-solved-norm} with
$\psi=(D+1)\chi_0\log(|z|^2/\eta^2)$.
The limit $u$ is holomorphic near $x$.  Its weighted square integrability
forces $u(x)=0$.  Consequently $P=\chi e-u$ is global and holomorphic,
$P(x)=e(x)$, and
\begin{equation}
 \label{bergman-quotient-l2}
 \int_Y\|P\|_\lambda^2\,dV_0
       \leq e^{C(S+A^8)}\operatorname{vol}_0(B_\eta).
\end{equation}

\emph{Membership in the original cusp space.}
This last point requires the actual arithmetic quotient and the deep
weight range.  Restriction to a compact flag fibre gives the fixed
Borel--Weil representation of $U(p)\times U(q)$.
The resulting automorphic function is $K$-finite, and holomorphicity
gives its holomorphic discrete-series infinitesimal character.
One convention for its lowest compact highest weight is
$\mu=-w_{0,K}\lambda$.
The associated finite-dimensional highest weight is
\[
 \nu=\mu-(\underbrace{q,\ldots,q}_{p},
                   \underbrace{-p,\ldots,-p}_{q}).
\]
The adjacent gaps of $\nu$ inside the two blocks are at least $m$;
its central gap is at least $m-r$.  Thus $\nu$ is integral and strictly
dominant at every real place.  A simultaneous dual convention has the
same regularity property.

The $L^2$ holomorphic section defines a square-integrable automorphic
form in the corresponding coherent degree-zero space.
To check the growth qualification, a local holomorphic mean-value
estimate on a fixed translated chart bounds its value by its local
$L^2$ mass.  The overlap multiplicity is bounded by the number of
integral matrices in a ball of radius polynomial in
$\|g\|+\|g^{-1}\|$.  This is a polynomial count on the fixed
arithmetic lattice.  It gives moderate growth; the highest-weight
relations give finiteness under the centre of the enveloping algebra.
These are qualitative assertions at each fixed setting and introduce
no constant into \eqref{bergman-quotient-l2}.
The actual dictionary supplies the algebraic automorphic bundle.
Its multiplier on the present tag-zero weight is trivial, as follows
directly from the constructed good-cover law.  For a nonnegative
antitone exponent array $a$, that law gives
$\chi(\operatorname{laU}(a,0),\gamma)=1$.
At tag zero, $\operatorname{laU}$ is the signed permutation taking
the first block in reverse order with a minus sign and the second
block in its original order with a plus sign, with the prescribed
conjugate embedding relabeling.  Given any integral flag weight,
pull it back to an integer array $c$.  Choose an integer $B$ larger
than every $|c_j|$ and $|c_j-c_{j+1}|$, including the last jump to
zero.  The arrays $b_j=B(r-j)$ for $j<r$ and zero afterward, and
$a=b+c$, are nonnegative and antitone.  Since
$\lambda=\operatorname{laU}(a,0)-\operatorname{laU}(b,0)$,
the original weight-addition law gives $\chi(\lambda,\gamma)=1$.
Thus no inference from finite character image to a congruence kernel,
and no new level cover, is needed.

For unitary groups, square-integrable coherent cohomology at the
infinitesimal character of an integral strictly dominant weight equals
cuspidal coherent cohomology.\cite{bergman-harris}
The cited formulation uses the unitary similitude Shimura variety.
Use its connected component with a compact finite level intersecting
the unitary group in the given level and with similitude in
$\widehat{\mathbb Z}^{\times}$.  A positive rational similitude in
that finite level is one.  Consequently its component arithmetic
group is exactly the original unitary arithmetic group.
An integral unitary weight extends to the similitude group over
$\mathbb C$; its complexification is a product of the original
general linear factors and a multiplicative group.  This identifies
the same coherent bundle on the same component, with no change of
the section or its metric.  The preceding calculation checks the
regularity hypothesis of the cited result, so $P$ is cuspidal.
Rapid decay of a cuspidal automorphic form of moderate
growth, applied to this finite compact type, makes its original
invariant point norm bounded.\cite{raw-theta-rapid-decay}
The constructed good-cover dictionary therefore puts it in the
literal space $X.\mathrm{SS}(\lambda)$.
In particular, square integrability alone has not been used to infer
cuspidality.

Finally $dV_0$ and the original flag Haar measure differ by a positive
constant, since both are invariant.  The injective ball has volume at
most that of the quotient.  These two facts turn
\eqref{bergman-quotient-l2} into
\eqref{bergman-actual-thick-peak-bound}, with every original Haar
normalization intact.
\end{proof}

\paragraph{The exact Bergman consequence.}
On the actual component extend $P$ by zero to the other components.
The original dictionary identifies its invariant point norm and its
component mean square.  A component's probability mass in the original
global Petersson norm is at most one.  The variational characterization
of the squared norm of evaluation therefore gives
\[
 \operatorname{bergOn}(y)\geq e^{-C(S+A^8)}
\]
at every point satisfying \eqref{bergman-actual-thick-test}; orthogonal
frame transfer preserves both sides.  This is the exact comparison
used by the original component Bergman consumer, rather than a change
to an independently normalized $L^2$ space.
For its untagged weight $u_1$, the model permutation preserves the
absolute weight sum, so $S=\ell_1(u_1)$.
If the original numerical data give
$\operatorname{lenT}\ell_1(u_1)\leq C_2\operatorname{unitE}$ and
$\operatorname{lenT}A^8\leq\operatorname{unitE}$, then the same
$B=C(\ell_1(u_1)+A^8)$ satisfies
$\operatorname{lenT}B\leq C(C_2+1)\operatorname{unitE}$.
These numerical inequalities must be checked in their original host
range when invoking that asymptotic wrapper.
The theorem supplies the thick analytic estimate and its actual
arithmetic injectivity.  The thin estimate, global upper section-frame
bound, point-height estimate and final section remainder remain
separate obligations.

\subsection{Theta jets and the lower bound in the thin region}
\label{shared-actual-thin-bergman}

The issue in the thin region is quantitative nonvanishing of a
prescribed determinant differential operator.  A short spanning set
does not give this.  We construct a section of the original period
torus, and then use the finite theta basis to obtain an actual cusp
section~\cite{bergman-thin-source}.  Maximizing a determinant over a bounded set of lattice
frequencies avoids any assumption on the condition number of their
projections to one complex embedding.

\paragraph{The period torus and its parameters.}
Let $E$ be a CM field of degree $2f$, put $R=\operatorname{rd}(E)$,
and let
\[
 V=\bigoplus_{\tau:E\hookrightarrow\mathbb C}\mathbb C^{s_\tau},
 \qquad s_\tau+s_{\bar\tau}=r,\qquad g=\dim_{\mathbb C}V=fr.
\]
Use the standard Hermitian form $H$, linear in its first argument.
The original period lattice $\Gamma\subset V$ is stable under
$\mathcal O_E$, and $\omega=\operatorname{Im}H$ is integral and
unimodular on $\Gamma$.  Its semicharacter
$\varepsilon:\Gamma\to\{1,-1\}$ satisfies
\[
 \varepsilon(\gamma+\eta)
   =\varepsilon(\gamma)\varepsilon(\eta)e^{\pi i\omega(\gamma,\eta)}.
\]
We retain its original isotropic normalization:
$\varepsilon(\gamma)=1$ whenever
$\omega(a\gamma,\gamma)=0$ for every $a\in\mathcal O_E$.
Write $\lambda=\min_{0\ne\gamma\in\Gamma}\|\gamma\|$.

For an integer $q\ge1$, let $\mathcal H_q$ consist of the entire
functions with transformation law
\[
 F(z+\gamma)=\varepsilon(\gamma)^q
  e^{\pi q H(z,\gamma)+\pi q\|\gamma\|^2/2}F(z).
\]
Its squared norm is the integral of
$|F(z)|^2e^{-\pi q\|z\|^2}$ on $V/\Gamma$.
This is also the probability norm: unimodularity of $\omega$
gives Euclidean covolume one.  On $V^n$ use the exterior tensor
power of this line and the product probability measure.

Let $P$ be a product of powers of determinants.  Each determinant
uses $i\le s_\tau$ rows at one embedding $\tau$, and $i$
distinct columns among the $n$ copies of $V$.  Its exponent is an
integer $c\ge1$.  Different determinants have disjoint column
sets.  Write $D=\sum ic\ge1$ for its degree.  The rows can be any
specified orthonormal coordinate rows.  The differential operator
$P(\partial)$ uses ordinary derivatives, with no factorial
inserted into its definition.

Choose $m\ge1$ with $s_\tau\le m$ at at least one embedding, and
put
\begin{gather*}
 d_*=2fm,\qquad \Omega=2f^{3/2}R,\qquad
 \theta^2=\frac{fR_0^2R^{21}}q,\qquad
 b_0=\frac q{R_0^2R^{28}},\qquad t_0=qb_0,\\
 C=8100m^4d_*^4R^7,\qquad
 \phi=\begin{cases}2f/(R^2\lambda^2),&\lambda<\theta,\\
                    0,&\lambda\ge\theta.
       \end{cases}
\end{gather*}
The following follow from the original numerical hypotheses and rank
conditions:
\begin{gather}
 r\ge2,\quad m^4\le f,\quad f^2\le R,\quad R\ge16,
 \quad q\ge6,\quad R_0^2\ge90(g+D),\label{eq:thin-jet-parameters-a}\\
 \Omega\theta^2<1,\qquad b_0\ge1,\qquad
 t_0\ge4D\bigl(7+\log(qD)\bigr)+d_*+1.
 \label{eq:thin-jet-parameters-b}
\end{gather}
Indeed the original ranks have $r=k+m-1$, $m\ge2$ and
$k\ge m+3$, so $r\ge6$.  The original inequalities include
$d_*^2\Omega\theta^2<1$ and $90m^2d_*^2R\le q/2$.
Since $d_*=2fm\ge2$, they imply $\Omega\theta^2<1$ and
$q\ge180m^2d_*^2R\ge11520$.  The other displayed inequalities
are retained verbatim.  In particular this argument does not infer
$r\ge2$ merely from $fr\ge2$.

\begin{theorem}[The required determinant jet]
\label{thm:actual-thin-determinant-jet}
Under these hypotheses, there is a section $F$ of the original
line on $(V/\Gamma)^n$ such that
\begin{equation}\label{eq:thin-jet-lower}
 \frac{|P(\partial)F(0)|^2}{\|F\|_2^2}
   \ge 2^{-3n}16^{-D}
                  \exp\left(-\frac{2\pi nC\phi}{q}\right).
\end{equation}
In particular this implies the prescribed lower bound
\[
 (4K_6)^{-n}2^{-(r+2)D}
                  \exp\left(-\frac{2\pi nC\phi}{q}\right),
 \qquad K_6=6(2d_*)^{d_*/2}.
\]
The theorem concerns the original principally polarized period
lattice, including its semicharacter and CM action.
\end{theorem}

\begin{proof}
\emph{An isotropic sublattice and its dual frequencies.}
There are $2f$ rationally independent elements of
$\mathcal O_E$ whose absolute values at all embeddings are at
most $\Omega$.  Indeed the trace-dual minimum is at least
$\sqrt{2f}/R$; the transference argument proved in
Section~\ref{sec:actual-short-theta-generators} gives Euclidean
length at most $(2f)^{3/2}R$.  A single CM embedding accounts for
two equal summands in its squared Euclidean norm.

Start with $U=0$ and $I=\Gamma\cap U$.  We build an
$E$-stable, real isotropic subspace $U$ with the following property:
if $b$ is the vector in $U$ representing the functional
$v\mapsto\omega(\gamma,v)$ on $U$, for $\gamma\in\Gamma$,
then
\begin{equation}\label{eq:thin-dual-separation}
             b\ne0\quad\Longrightarrow\quad
                         \|b\|\ge\frac1{\Omega\theta}.
\end{equation}
Suppose $x\in\Gamma\setminus U$ has distance less than
$\theta$ from $U$.  The property forces $\omega(x,U)=0$,
since $\Omega\theta^2<1$.  Write $x=u+w$, with $u\in U$;
then $w$ is orthogonal to $U+iU$ and $\|w\|<\theta$.
For each of the short rational basis elements $a$, the integer
$\omega(ax,x)=\omega(aw,w)$ has absolute value less than one.
It is zero.  Rational linearity gives the same conclusion for
every $a\in E$.  Thus adjoining the $E$-span of $x$ preserves
isotropy.

Property~\eqref{eq:thin-dual-separation} is preserved as well.
If a pairing is already nonzero on the old $U$, use its old
orthogonal projection.  Otherwise some pairing with $ax$, for
one of the short rational basis elements $a$, is a nonzero
integer.  The old component does not contribute, so its absolute
value is at most $\|b\|\|aw\|\le\|b\|\Omega\theta$.
This proves the asserted separation.

The process stops after at most $m$ steps.  More explicitly, the
complexification of an isotropic $E$-subspace of rank $\rho$
injects into $V$: $U\cap iU=0$.  At every embedding its
complex dimension is $\rho$, whence $\rho\le s_\tau$ for
every $\tau$.  At termination put $d=\dim_{\mathbb R}U=2f\rho
\le d_*$.  The lattice $I=\Gamma\cap U$ is saturated, and
\begin{equation}\label{eq:thin-quotient-separation}
 x\in\Gamma\setminus I\quad\Longrightarrow\quad
                         \operatorname{dist}(x,U)\ge\theta.
\end{equation}
The projected vectors used at each step, multiplied by the $2f$
short coefficient elements, also give a full sublattice of $I$
with successive Gram--Schmidt lengths at most $\Omega\theta$.
Consequently
\begin{equation}\label{eq:thin-isotropic-covolume}
                         \operatorname{covol}(I)
                                      \le(\Omega\theta)^d.
\end{equation}
The semicharacter is one on $I$ by its retained normalization.
The space $U$ is zero exactly when $\lambda\ge\theta$.

Write $W=U+iU$ and $A=W^\perp$.  Each period is written
\[
                   \gamma=t+ib+a,
             \qquad t,b\in U,\quad a\in A.
\]
Its frequencies $b$ constitute exactly the Euclidean dual lattice
$L=I^*$ in $U$.  A homomorphism $I\to\mathbb Z$ extends to
$\Gamma$ because $I$ is saturated, and unimodularity of
$\omega$ then supplies its representing period.  Thus this is
an equality of lattices, not an inclusion of unspecified index.
The periods with $b=0$, modulo $I$, project to a lattice in $A$
with minimum at least $\theta$.  Every other frequency fibre is
a translate of that lattice.  By
\eqref{eq:thin-dual-separation},
\begin{equation}\label{eq:thin-frequency-minimum}
 \lambda_L\ge(\Omega\theta)^{-1}\ge\sqrt{b_0}.
\end{equation}
The second inequality is equivalent to $f\Omega^2\le R^7$,
which follows from $f^2\le R$ and $R\ge16$.

\emph{A bounded set controlling every relevant determinant.}
Let $\mathcal V$ be the closed Voronoi cell of $L$ and put
$S=L\cap3\mathcal V$.  Then
\begin{equation}\label{eq:thin-voronoi-convex}
                            \mathcal V\subset\operatorname{conv}S.
\end{equation}
To see this, for $v\in\mathcal V$ choose $\ell\in L$ with
$2v-\ell\in\mathcal V$.  Then $\ell\in3\mathcal V$ and
$v=(\ell+(2v-\ell))/2$.  Iterate this identity.  Its remainders
stay in the compact cell and their coefficients tend to zero;
the finite set $S$ includes zero, so the limit belongs to its
closed convex hull.

For $v\in U$ its Voronoi gauge is
\[
 \mu(v)=\sup_{0\ne\ell\in L}
                  \frac{2|\langle v,\ell\rangle|}{\|\ell\|^2}
                      \le\frac{2\|v\|}{\lambda_L}.
\]
Thus $v$ belongs to $\mu(v)\operatorname{conv}S$.
Moreover the covering radius of $L$ is at most
$d\lambda_d(L)/2\le d^3/(2\lambda)$, by transference and
$\lambda_1(I)\ge\lambda$.  Every $u\in S$ consequently has
\begin{equation}\label{eq:thin-frequency-radius}
 \|u\|^2\le\frac{9d^6}{4\lambda^2}\le C\phi.
\end{equation}
For the last inequality substitute $d\le2fm$ and
$\phi=2f/(R^2\lambda^2)$; it is implied by
$f\le1800m^2R^5$.  When $U=0$, set $S=\{0\}$ and only
use the abelian directions below.

Fix a block of size $i$ and exponent $c$ at $\tau$.  Consider
the projections to $V_\tau$ of $S$, together with the ball of
radius $\sqrt{qc}$ in $A\cap V_\tau$.  Maximize the absolute
value of the specified $i$-row determinant on $i$ choices from
this compact set, and call its maximum $\Delta$.  The maximum
controls determinants with columns in the complex balanced convex
hull as well, by multilinearity and the triangle inequality.
Give this span the norm whose unit ball is that hull.  For symmetric
tensors we use its projective tensor norm: the infimum of
$\sum_j\prod_k\|v_{jk}\|$ over representations as sums of
symmetrized pure tensors.  Symmetrization is the average over
permutations.  This specifies the tensor estimates used below.

The real form $U$ pairs $W\cap V_\tau$ with
$W\cap V_{\bar\tau}$ by conjugation.  A unit vector in the
first summand is the projection of a vector in $U$ of norm
$\sqrt2$.  Since $\mathcal V$ contains the radius
$\lambda_L/2$ ball, \eqref{eq:thin-voronoi-convex} and
$\lambda_L\ge1$ show that this balanced convex hull contains
the radius $1/4$ ball in all of $V_\tau$.  Hence
\begin{equation}\label{eq:thin-leading-determinant}
                              \Delta\ge4^{-i}.
\end{equation}
No lower singular-value bound for a chosen lattice basis is used.

\emph{Actual Appell--Humbert sections.}
The real form $U\subset W$ defines the complex bilinear form
$B$ whose restriction to $U$ is its Euclidean form.  Conjugation
in this real form interchanges $\tau$ and $\bar\tau$.
Consequently $B(z,z)=0$ on every one-embedding subspace
$V_\tau$.  For $u\in L$ set
\[
 e_u(z)=\exp\left(\frac{\pi q}2B(z_W,z_W)
                                  +2\pi iB(z_W,u)\right).
\]
It has the prescribed transformation law on $I$.
For a toric column selected above use
$h=e_u/(2\pi i)^c$.  For an abelian column
$\sqrt{qc}\,\xi$, where $\xi\in A\cap V_\tau$ is a unit
vector, use
\[
                       h(z)=e_0(z)
                     \frac{(\sqrt{qc}\,H(z,\xi))^c}{c!}.
\]
For an unused column use $h=e_0$.
Periodize these functions by the original lattice:
\begin{equation}\label{eq:thin-periodization}
 \Theta_h(z)=\sum_{\gamma\in\Gamma/I}
   \varepsilon(\gamma)^q
   e^{\pi qH(z,\gamma)-\pi q\|\gamma\|^2/2}h(z-\gamma).
\end{equation}
The summand is unchanged by changing the representative modulo
$I$.  The semicharacter identity makes the translation operators
a representation of $\Gamma$, so the sum belongs to
$\mathcal H_q$.  To justify the sum, choose the real $U$
coordinate of each representative in a bounded fundamental
parallelepiped.  In the remaining coordinates its modulus and
every derivative on a compact set are bounded by a polynomial
times a decaying Gaussian.  The frequency lattice and the
separated affine lattices in $A$ give locally uniform absolute
convergence.  This also justifies differentiation term by term.

For $\gamma=t+ib+a$, the scalar in the toric summand at zero
has modulus
\begin{equation}\label{eq:thin-summand-modulus}
  \exp\left(-\pi q\|b\|^2-\frac{\pi q}2\|a\|^2
                                          +2\pi\langle b,u\rangle\right).
\end{equation}
Because $u\in3\mathcal V$, this is at most
$\exp(-\pi(q-3)\|b\|^2-\pi q\|a\|^2/2)$.
Along $V_\tau$ its linear exponential is exactly
\[
            2\pi iH(z,u-qb)+\pi qH(z,a).
\]
The toric $c$-jet, after its normalization, is therefore the
$c$-th power of the functional with column
$u-qb-(q/(2i))a$, with the conjugation in the second argument
of $H$ understood.  The identity summand gives the selected column
$u$.

For the abelian column put $v=\sqrt{qc}\,\xi$.  Its $c$-jet
is the sum, for $0\le l\le c$, of the symmetric tensors
\begin{equation}\label{eq:thin-abelian-jet}
       \frac{\binom cl}{l!}(-H(a,v))^l
                  H(\,\cdot\,,v)^{c-l}\,\mathcal L_\gamma^l,
 \qquad \mathcal L_\gamma(z)=-2\pi iqH(z,b)+\pi qH(z,a),
\end{equation}
times its scalar Gaussian.  This identity follows by expanding
the polynomial and exponential before taking the $c$-th
ordinary derivative.  At $\gamma=0$ only the chosen column
$v$ remains.  The factorial in this formula is essential.

\emph{Uniform bounds for the error terms.}
Here are the estimates including their constants.  Put
\[
              A_0=q\theta^2\ge90(g+D),\qquad B_0=qb_0=t_0.
\]
For a nonzero frequency put $x=\|b\|/\lambda_L\ge1$, and
put $y=\sqrt q\|a\|$.  Measure a column by the balanced convex
hull used to define $\Delta$.  A toric jet has total tensor
coefficient norm at most
$(1+2qx+y/(2\sqrt c))^c$.
Formula~\eqref{eq:thin-abelian-jet} bounds the abelian coefficient
norm by
\[
 \sum_{l=0}^c\frac{\binom cl}{l!}
                  (4\pi q\sqrt c\,xy+\pi y^2)^l.
\]
Both are at most
\begin{equation}\label{eq:thin-nonzero-frequency-polynomial}
                    (20q\sqrt c\,x)^c(1+y)^{2c}
                         \quad(b\ne0).
\end{equation}
Indeed $1+4\pi q\sqrt c\,xy+\pi y^2$ is at most
$4(1+4q\sqrt c\,x)(1+y)^2$, and $l!\ge1$.
For $b=0$ the sharper bound is
\begin{equation}\label{eq:thin-zero-frequency-polynomial}
                         \exp(2\sqrt{\pi c}\,y).
\end{equation}
For the abelian jet use $\binom cl\le c^l/l!$ and
$\sum_{l\ge0}t^l/(l!)^2\le e^{2\sqrt t}$; the latter follows
by comparing with the even terms of the exponential using
$\binom{2l}l\le4^l$.  For the toric jet use
$\log(1+y/(2\sqrt c))\le y/(2\sqrt c)$.

We record the elementary packing bounds used here.  In a real
lattice of dimension $a$ and minimum $\delta$, at most
$(1+2R/\delta)^a$ points lie in a ball of radius $R$; this
also holds for any translate.  Split into shells of width
$\delta/2$.  Since $A_0\ge90(g+D)$, for every affine lattice
in $A$ with minimum at least $\sqrt{A_0}$ one obtains
\[
                         \sum_z e^{-\pi\|z\|^2/4}\le2.
\]
There is at most one point in the open radius $\sqrt{A_0}/2$
ball.  The remaining shells are bounded by
$\sum_{j\ge1}(j+2)^{2g}e^{-\pi A_0j^2/16}<1$.
For the lattice itself, excluding zero, the same shell calculation
starts at $j=2$ and gives
\[
                  \sum_{z\ne0}e^{-\pi\|z\|^2/4}
                         \le4\,4^{2g}e^{-\pi A_0/4}.
\]
For example $j+2\le4e^{(j^2-4)/8}$ for $j\ge2$,
and $j+2\le3e^{(j^2-1)/6}$ for $j\ge1$ prove these displayed
bounds by a geometric series.

When $b=0$ and $a\ne0$, $y\ge\sqrt{A_0}$, and
$2\sqrt{\pi c}\,y\le\pi y^2/4$.
Thus \eqref{eq:thin-zero-frequency-polynomial} and the Gaussian
give total error at most
\[
                 4\,4^{2g}e^{-\pi A_0/4}\le e^{-30D}.
\]
For $b\ne0$, use
\[
 (1+y)^{2c}e^{-\pi y^2/2}
                    \le(4c)^c e^{-\pi y^2/4}
\]
and the uniform affine-lattice bound.  Summing the frequency
shells, whose dimension is $d\le d_*$, gives error at most
\begin{align*}
 2(80qD^{3/2})^D
   \sum_{b\ne0}e^{-\pi q\|b\|^2/2}
                            (\|b\|/\lambda_L)^D
 &\le4(120qD^{3/2})^D4^d e^{-\pi B_0/2}\\
 &\le e^{-30D}.
\end{align*}
For the first inequality use the shells $j\ge2$ and
\[
 (j+2)^d((j+1)/2)^D
 \le4^d(3/2)^D e^{(d+D)(j^2-4)/8}.
\]
The resulting geometric sum is less than two since
$B_0\ge d+D+1$.  For the second inequality substitute
\eqref{eq:thin-jet-parameters-b}; its logarithmic margin is at
least $30D$, using $\log120<5$, $\log4<1.4$ and $\pi>3$.

Accordingly, every selected column has total nonidentity error,
in the indicated tensor norm, at most
\begin{equation}\label{eq:thin-column-error}
                            \delta=2e^{-30D}\le(8D)^{-1}.
\end{equation}
The same upper bound applies to the unused zero-degree column
by omitting its polynomial factors.  Polarizing the product of
$c$ determinants expresses its contraction with symmetric
column tensors as an average of products of $c$ determinants.
Each factor is bounded by $\Delta$.  Thus the relative error in
an $i$-column block is at most $(1+\delta)^i-1<1/2$, since
$i\le D$.  Its jet is at least $\Delta^c/2$ in absolute value.
Each unused column has value at least $1/2$.

\emph{The original $L^2$ norm.}
The periodization also has a uniform norm bound; a pointwise
estimate alone would not suffice.  Absolute convergence of the
unfolding can be checked before using orthogonality.  Write
$N(h)=\int_{V/I}|h|^2e^{-\pi q\|z\|^2}$ and let $T_\gamma$
denote the translation in \eqref{eq:thin-periodization}.
For a toric column, completing the square gives exactly
\[
 \frac{\int_{V/I}|T_\gamma h\,\overline h|
                            e^{-\pi q\|z\|^2}}{N(h)}
       =e^{-\pi q\|b\|^2/2-\pi q\|a\|^2/4}.
\]
For an abelian column the same ratio is at most this Gaussian
times $\sum_{l=0}^c(\pi q|H(a,\xi)|^2/4)^l/l!$.
Indeed translate the $A$ coordinate by $a/2$, bound the product
of its two polynomial factors by
$(|H(z,\xi)|^2+|H(a,\xi)|^2/4)^c$, and integrate its monomials.
These Gaussian-polynomial sums converge by packing, uniformly
on the finite-volume real $U$ domain.  Thus the absolute double
sum can be unfolded to $V/I$.

Now integrate first over $U/I$.  The characters are
$u-qb$, so orthogonality kills every term with $b\ne0$.
The remaining terms are ordinary Gaussian overlaps in $A$.
For a toric column their normalized absolute value is
$e^{-\pi q\|a\|^2/2}$.  For an abelian column it is
\[
 e^{-\pi q\|a\|^2/2}
       \left|\sum_{l=0}^c\binom cl
              \frac{(-\pi q|H(a,\xi)|^2)^l}{l!}\right|
 \le e^{-\pi y^2/2+2\sqrt{\pi c}\,y}.
\]
This formula follows by integrating monomials in the one
coordinate $\xi$; its Gaussian moment is
$c!/(\pi q)^c$.  The preceding zero-frequency tail estimate
shows that the sum of the absolute nonidentity overlaps is
less than one.

The squared norm of the unperiodized toric function on $V/I$ is
\[
 \operatorname{covol}(I)(2q)^{-d/2}q^{-(g-d)}
                         e^{2\pi\|u\|^2/q}.
\]
Its normalization $(2\pi i)^{-c}$ only decreases this.
For the abelian polynomial the additional factor is
$c^c/(\pi^c c!)\le(e/\pi)^c\le1$.
By \eqref{eq:thin-isotropic-covolume} and $f\Omega^2\le R^7$,
\[
 \frac{\Omega\theta}{\sqrt q}\le t_0^{-1/2}\le1.
\]
The prefactor of the Gaussian norm is therefore at most one.
We have proved
\begin{equation}\label{eq:thin-column-l2}
                  \|\Theta_h\|_2^2
                           \le2e^{2\pi\|u\|^2/q},
\end{equation}
where $u=0$ for an abelian or unused column.

Take the tensor product of all the selected columns.  The
determinant blocks and unused columns give
$|P(\partial)F(0)|\ge2^{-n}4^{-D}$, because their number is at
most $n$.  Equation~\eqref{eq:thin-column-l2} gives
$\|F\|_2^2\le2^n e^{2\pi nC\phi/q}$ by
\eqref{eq:thin-frequency-radius}.  This proves
\eqref{eq:thin-jet-lower}.  Finally $K_6\ge2$ and $r\ge2$
give the asserted comparison with the prescribed bound.
\end{proof}

\paragraph{Return to the actual cusp sections.}
The required polynomial really has disjoint column sets.  For
the original antitone exponents, give the minor of size $i$ at
$\tau$ its own $i$ columns and raise it to
$a_{\tau,i-1}-a_{\tau,i}$, omitting zero exponents.  This uses at
most
$\sum_\tau s_\tau(s_\tau+1)/2\le fr(r+1)/2\le n$ columns.
Its row degrees telescope to the prescribed exponents.  Different
embeddings, including conjugate embeddings, use disjoint columns,
so every harmonic contraction is zero.  The leading minors give
the required highest-weight relations, and their product is
nonzero with integral coefficients.  Its degree is $D$, and its
coefficient sum is at most $r^D$.  This is the actual normalized
block polynomial; the construction does not use overlapping
columns of nested minors.

The principally polarized theta section has original probability
norm one.  Pulling it back by the prescribed half-level generator
and translating through the finite half-level group gives an
orthonormal basis of $\mathcal H_M$; on $n$ copies its cardinality
is $M^{gn}$.  Indeed $\alpha\bar\alpha=M$, and the original
semicharacter compatibility makes pullback by $\alpha$ a section
of this very line.  Pullback preserves its probability integral
by Haar invariance under the finite isogeny; the translations are
unitary, so every vector has squared norm exactly one.
The half-level ideal has norm $M^f$, and the original projective
rank-$r$ lattice has quotient cardinality $M^{fr}$.
To check completeness, the translations are distinct
characters for the complementary finite Heisenberg subgroup, so
they are orthogonal; their number is $M^g=\dim\mathcal H_M$.
The dimension identity follows directly by writing a theta section
in Fourier series along a symplectic integral basis: its
coefficients are determined freely by the $M^g$ residue classes.
The original normalized theta dictionary identifies these basis
vectors with its half-level rules.  Before taking absolute values
the evaluation factor is
$M^{-D}\kappa(\sqrt{2\xi_R})^D$, with
$|\kappa|^2=M^D$.  Thus its squared absolute value is exactly
$(2\xi_R/M)^D$.  The Lie-norm factor is
$(2\pi\xi_R)^{-D/2}$.  Hence Parseval and
Theorem~\ref{thm:actual-thin-determinant-jet} give, for the original
rules $s_\zeta$,
\begin{equation}\label{eq:thin-actual-theta-sum}
 M^D\sum_{\zeta}
       \bigl((2\pi\xi_R)^{-D/2}\|s_\zeta(x)\|\bigr)^2
             \ge\pi^{-D}\,2^{-3n}16^{-D}
                           e^{-2\pi nC\phi(x)/M}.
\end{equation}
The integrality, harmonicity, cuspidality and original evaluation
normalization of these rules were established in the actual theta
construction.  The auxiliary section just constructed selects a
nonzero value among them; it does not replace their arithmetic
lattice or their metric.

\begin{theorem}[The original thin-region Bergman bound]
\label{thm:actual-thin-bergman}
Retain the original sized good cover, anchored divided weight
$u_1$, flat-weight comparison and their numerical hypotheses.
At every coefficient embedding and every point of an original
component $Z$ satisfying
$M_{\mathrm{deep}}\lambda(x)^2<b$, one has
\begin{equation}\label{eq:actual-thin-bergman}
                  B_Z(x)\ge
                  \exp\left(-\frac{2f}{625R^2\lambda(x)^2}\right).
\end{equation}
Here $\lambda(x)^2$ is exactly the original first minimum
squared, and $B_Z$ uses the original component probability norm.
\end{theorem}

\begin{proof}
We spell out the numerical and metric interfaces.  On the original
minimal level $M$ the numerical hypotheses give
\[
 b\ge1000,\quad gn\le b,\quad r\le n,\quad R^7\le b,
 \quad b^{16}\le M\le b^{70},\quad 1\le D\le b^2,
 \quad M_{\mathrm{deep}}\ge bM.
\]
They give \eqref{eq:thin-jet-parameters-a}--\eqref{eq:thin-jet-parameters-b}
for this very period torus.  The deep-to-minimal-level dictionary
preserves the first minimum and the point metric.  Thus a point
under consideration has $M\lambda(x)^2<1$ and $\lambda(x)<\theta$, and
\[
                     \phi(x)=\frac{2f}{R^2\lambda(x)^2}
                                      \ge b^{14}.
\]
The normalized integral block polynomial has coefficient sum at
most $e^{Db}$.  The already proved actual raw-theta upper bound,
with $L=3\log b$ and size bound $b^2$, gives
\[
 (2\pi\xi_R)^{-D/2}\sup_z\|s_\zeta(z)\|
   \le\exp\bigl(Db+9000(3fr\log b)^2b^2\bigr).
\]
There is no determinant-index factor for this original principal
theta family.  Averaging \eqref{eq:thin-actual-theta-sum} over
its $M^{gn}$ indices gives at least one nonzero value with
\[
 -\log\bigl((2\pi\xi_R)^{-D/2}\|s_\zeta(x)\|\bigr)
 \le3b^3D+\frac{\phi(x)}{2500}+(D+gn)\log M.
\]
Indeed the nonconstant coefficient is
$\pi nC/M\le1/2500$: use $n,m\le b$, $d_*\le2b$,
$R^7\le b$ and $M\ge b^{16}$.
The remaining logarithms from
\eqref{eq:thin-actual-theta-sum} are less than $3b^3D$.

The total terms not containing $\phi$ are at most $10^5b^6$,
using $fr\le b$, $\log b\le b$ and $\log M\le70b$.
Since $b\ge1000$ and $\phi\ge b^{14}$, they are at most
$\phi/2500$.  The same original cusp section therefore satisfies
\[
                    \|s_\zeta(x)\|
                       \ge e^{-\phi(x)/1250}
                                      \sup_z\|s_\zeta(z)\|.
\]
Pull it to the prescribed cover.  The established flat-weight
comparison from the reference theta weight to $u_1$ gives a
section $g$ supported on the same $Z$ and a constant $K>0$ with
$\|g(z)\|=K\|s_\zeta(z)\|$ for every $z$ on $Z$.
It is the product of the already constructed key and parallel
twisting sections.  Thus it preserves the ratio on this very
component.  Take this section supported on $Z$.
Its original probability $L^2$ norm is at most its supremum norm.
Squaring the ratio proves \eqref{eq:actual-thin-bergman}.
\end{proof}

The theorem supplies the lower thin-region norm used in the
geometric section argument.  The construction of the original
numerical family remains a separate obligation wherever it is
still an input.  Point heights and the final arithmetic section
remainder are not supplied by this analytic theorem.

\subsection{A uniform upper bound for the original cusp norms}
\label{shared-bergman-actual-upper}

We retain the integral model, its complex uniformization and its
Hermitian metric.  The probability measure on a component is the one
induced by the original adelic measure~\cite{bergman-upper-source}.
In particular, normalizing a
component does not permit us to discard its arithmetic volume.

Let $E$ be the original CM field, let $F=E^+$, and put
\[
 f=[F:\mathbb Q],\qquad \Delta=|\operatorname{disc}E|,
 \qquad \operatorname{rd}(E)=\Delta^{1/(2f)}.
\]
The original hypotheses include that $E/F$ is unramified at all finite
primes.  The rank-$r$, $r\ge2$, Hermitian lattice used in the principal-level
description is self-dual over $\mathcal O_E$; it need not be free.
Write $M\ge3$ for the level and $d_H\ge1$ for the original positive
quadratic parameter.  At an orthogonal frame $x$, write
\[
 q_x(v)=\sum_{\tau\in\Psi}\operatorname{maj}_{x,\tau}(v)
\]
for the sum of the positive Hermitian majorants, where $\Psi$ contains
one embedding from each conjugate pair.  If $\lambda_{\tau,i}$ are
the model weights, assume
\[
 \sum_{\tau\in\Psi}
 \bigl(\max_i\lambda_{\tau,i}-\min_i\lambda_{\tau,i}\bigr)
 \le w,\qquad w\ge1.
\]
The weights in this statement are the actual model weights after the
prescribed signed permutation.  Let $\mathcal H_Z$ be the original
finite-dimensional cusp section space on a component $Z$, with its
component probability norm $\|\cdot\|_{2,Z}$.  Its Bergman function is
\[
 B_Z(x)=\sup_{0\ne s\in\mathcal H_Z}
              \frac{\|s(x)\|^2}{\|s\|_{2,Z}^2},
\]
with value zero when $\mathcal H_Z=0$.

Set
\begin{align}
 D&=fr(r-1)/2,&
 \rho&=\frac1{200r^3(w+f)},\label{eq:actual-upper-radius}\\
 \epsilon&=\min\left\{1,
 \frac{\pi}{2M\sqrt{d_H}(fw+2)\log2}\right\},&
 A_\epsilon&=
 \left(1+\frac{2(2fr)^2\Delta}{\epsilon}\right)^{2fr^2},
 \label{eq:actual-upper-multiplicity}\\
 V&=2h_E M^{2f(r^2+1)}
       \operatorname{rd}(E)^{f(r^2-1)/2},&
 C_{\mathrm{sup}}&=
       \frac{8D!}{(\pi\rho^2)^D}A_\epsilon V,
 \label{eq:actual-upper-constant}
\end{align}
where $h_E$ denotes the ideal class number of $E$.

\begin{theorem}[The original upper Bergman bound]
\label{thm:actual-upper-bergman}
For the original good principal-level cover, every coefficient
embedding, every component $Z$ and every $s\in\mathcal H_Z$ satisfy
\[
 \sup_{x\in Z}\|s(x)\|^2
       \le C_{\mathrm{sup}}\|s\|_{2,Z}^2,
 \qquad B_Z(x)\le C_{\mathrm{sup}},
 \qquad \dim\mathcal H_Z\le C_{\mathrm{sup}}.
\]
The statement uses the classical Tamagawa-number theorem for special
unitary groups.  All comparisons from its invariant differential
measure to the displayed constant are given below.  No bound is
asserted here for a ramified extension $E/F$ or a different integral
lattice.

If $L\ge1$ and
\begin{gather}
 \log(fr)\le L,\quad \log M\le30L,
 \quad \log\Delta\le4fL,\quad
 \log d_H\le4L,\quad \log w\le2L,
 \label{eq:actual-upper-size}
\end{gather}
then
\begin{equation}\label{eq:actual-upper-logbound}
            1\le C_{\mathrm{sup}},\qquad
            \log C_{\mathrm{sup}}\le250f^2r^2L.
\end{equation}
In the original family, let $L_b\ge1$.  Suppose
\[
 \log M\le30L_b,\qquad
 \log(2f)\le3L_b,\qquad \log\operatorname{rd}(E)\le3L_b,\qquad
 \log(fr)\le L_b,
\]
and put $S=\sum_{\tau,i}|u_{\tau,i}|$.  If $S=0$ or
$\log S\le20L_b$, then one may take
\begin{equation}\label{eq:actual-upper-family}
                   B_Z(x)\le \exp\bigl(4000(frL_b)^3\bigr).
\end{equation}
Here $u$ is the original integral weight and the original host
discriminant relation $d_H^f\le\Delta$ is retained.
\end{theorem}

\begin{proof}
\emph{A half-maximum point has a lower lattice minimum.}
The zero section requires no argument.  For a nonzero section choose
a frame $x$ with point norm greater than half its supremum, then
orthogonalize the frame.  The original norm is unchanged.  The
integral Hermitian lattice $L'$ and principal-level transvections
constructed in the short-theta proof give
\begin{equation}\label{eq:actual-upper-minimum}
                    q_x(v)\ge\epsilon
                  \qquad(0\ne v\in L').
\end{equation}
We recall why this estimate has precisely the displayed level
dependence.  If all embedding majorants of $v$ are less than
$\epsilon\le1$, its integral Hermitian norm has all conjugates of
absolute value less than one and is therefore zero.  Its rank-one
unipotent orbit, with complex parameter $\zeta$, remains in the
domain for $\Im\zeta>-T$, where
$T=(k_0\epsilon)^{-1}$.  The period $M_*$ of the original
principal-level action obeys
\[
                       M_*k_0=2M\sqrt{d_H}.
\]
The actual arithmetic character is trivial on this action, so the
evaluated section is periodic with period $M_*$.  Its constant
Fourier coefficient is zero by cuspidality.  On the line
$\Im\zeta=-T/2$, its norm differs from the norm at $x$ by at most
$2^{w/2}$ in the frame factor.  Applying the disk Schwarz lemma
after $q=\exp(2\pi i\zeta/M_*)$ gives
\[
              1\le2\,2^{w/2}\exp(-\pi T/M_*).
\]
But the definition of $\epsilon$ makes the right side at most
$2^{1+w/2-(fw+2)}<1$.  This contradiction proves
\eqref{eq:actual-upper-minimum}.  The arithmetic character assertion
is the one for the constructed bundle: its character law on all
nonnegative antitone exponent arrays extends, by differences, to
all integral model weights.  Finite image alone would not suffice.

\emph{The actual chart has bounded multiplicity.}
Put $n=2fr$.  The real quadratic form $q_x$ has dimension $n$.
For its Euclidean dual lattice $(L')^*$, the trace-pairing identity
and Hermitian self-duality imply
\[
             \frac{\Delta}{2}J_x(L')^*\subseteq L',
\]
where $J_x$ is the orthogonal involution of the Hermitian majorant.
Indeed, trace duality gives the inverse different, the integer
$\Delta$ clears that different, and the sum over a CM type is
half the full trace.  Thus the dual minimum is at least
$2\sqrt\epsilon/\Delta$.

The elementary transference argument already proved for the short
theta span gives
$\lambda_n(L')\lambda_1((L')^*)\le n^2$.
Consequently $L'$ contains $r$ vectors $v_1,\ldots,v_r$ independent
over $E$, with
\begin{equation}\label{eq:actual-upper-short-basis}
                  q_x(v_i)\le\frac{n^4\Delta^2}{4\epsilon}.
\end{equation}
The argument first gives $n$ vectors independent over $\mathbb Q$;
one then selects an $E$-basis.  It does not select a free
$\mathcal O_E$-basis or change the integral lattice.

Use the original lower-triangular flag chart of radius $\rho$ at
$x$.  Comparing the two positive/negative orthogonal projections
in its Gram matrices gives
\[
 \tfrac12q_x(v)\le q_z(v)\le2q_x(v)
       \qquad(\|z\|\le\rho).
\]
For clarity, write the chart matrix as $I+N$ in an orthogonal frame.
Its entries satisfy $|N_{ab}|\le\rho$ and
$\|N\|_{\mathrm{op}}\le r\rho$; the Gram perturbation has norm
at most $2r\rho+r^2\rho^2$.  The positive block and its inverse
remain within the same bounds by the Neumann series.  In the
formula for the positive projection this gives an operator norm
error at most $8r^2\rho<1/2$.  This proves the displayed
comparison, for the indefinite signatures as well as for the
definite endpoints.

If an arithmetic element $\gamma$ returns this chart to itself,
the intermediate upper-triangular change of flag leaves the
majorant invariant.  Hence
$q_x(\gamma v)\le4q_x(v)$.  The map
\[
               \gamma\longmapsto(\gamma v_1,\ldots,\gamma v_r)
\]
is injective.  Balls of radius $\sqrt\epsilon/2$ around lattice
points are disjoint.  Comparing their volumes shows that a ball
of radius $R$ contains at most
$(1+2R/\sqrt\epsilon)^n$ lattice points.  Taking
$R=n^2\Delta/\sqrt\epsilon$ in
\eqref{eq:actual-upper-short-basis}, and then taking the $r$-th
power, bounds the number of returning elements by $A_\epsilon$.

\emph{A local mean inequality retains the metric and measure.}
In the same chart let $\delta_{\tau,j}$ be the leading Hermitian
Gram determinants, with $\delta_{\tau,0}=1$.
The triangular chart has determinant one, so
$|\delta_{\tau,r}|=1$.  The same matrix estimates give
\[
        \bigl|\log|\delta_{\tau,j}|\bigr|\le9r^2\rho.
\]
In the exact metric formula, subtracting the midpoint of the
largest and smallest weights costs nothing because the two end
determinants have absolute value one.  Summation by parts therefore
gives a squared frame factor at least $\exp(-1/20)$.
The invariant flag density is the product of
$|\delta_{\tau,j}|^{-2}$, and is at least $\exp(-1/10)$.
Both assertions follow directly from
$200r^3\rho(w+f)=1$.  Their product exceeds $1/2$.

The mean-value inequality for the squared absolute value of the
holomorphic coefficient, on a ball in $\mathbb C^D$, now gives
\[
 \|s(x)\|^2\le
 \frac{2D!}{(\pi\rho^2)^D}
           \int_{\|z\|<\rho}\|s(z)\|^2\,d\mu_{\mathrm{flag}}(z).
\]
Unfolding this nonnegative integral to a fundamental set multiplies
it by at most $A_\epsilon$.  The half-maximum choice costs a factor
four.  Thus, on this uniformizing component,
\begin{equation}\label{eq:actual-upper-raw}
 \sup\|s\|^2\le
 \frac{8D!}{(\pi\rho^2)^D}A_\epsilon
       \int_{\mathcal F}\|s\|^2\,d\mu_{\mathrm{flag}}.
\end{equation}
Tonelli's theorem applies throughout.  The argument needs neither
an injective chart nor an unspecified compactness constant.

\emph{The arithmetic volume in the probability normalization.}
The required bound for the sum of the original fundamental-set
volumes is
\begin{equation}\label{eq:actual-upper-mass}
             \sum_k\mu_{\mathrm{flag}}(\mathcal F_k)\le V.
\end{equation}
Here is the normalization calculation.  At a finite place of $F$
with residue size $q$, the self-dual lattice in the unramified
quadratic algebra defines a smooth special unitary group scheme.
Its volume for an integral invariant differential is
\[
 a_v=\prod_{j=2}^r(1-\chi_v^j q^{-j}),
 \qquad
 \chi_v=\begin{cases}1&\text{at a split place},\\-1&\text{at an inert place}.
                    \end{cases}
\]
The formula follows from the orders of the split and unitary
finite groups and smooth lifting.  It holds at primes over two
and at primes dividing $r$; no division by these primes is used.
It also gives
\[
                 \prod_v a_v^{-1}
                   \le\prod_{j=2}^r\zeta_F(j)
                   \le\left(\prod_{j=2}^r\zeta(j)\right)^f.
\]

The only global volume theorem used is that the Tamagawa number
of the simply connected group $\operatorname{SU}(h)$ is one.
For example, this is the outer type-$A$ case recorded in
Prasad~\cite[\S3.3]{bergman-upper-prasad}; his \S3.7 gives the corresponding local-factor
formula.  Summing over all finite double cosets gives the same
identity when all real factors are compact, so no isotropy or
strong-approximation assumption is imposed at this step.

To compare its differential measure with the original Riemannian
measure, choose any integral $F$-basis of the trace-zero
anti-Hermitian Lie algebra, spanning a sublattice of index $a$.
The trace form on the full endomorphism lattice is unimodular
locally because $E/F$ is unramified and the Hermitian lattice is
self-dual.  Passing to trace zero contributes the determinant
$r$: in diagonal coordinates its Gram matrix is
$I_{r-1}+\mathbf1\mathbf1^t$, of determinant $r$.
The off-diagonal pairs have unit determinant.  This computation
may be made after unramified faithfully flat extension and hence
gives the same determinant ideal before extension.  It proves
that the real volume conversion is $r^{f/2}a$.
The same $a$ occurs inversely in the finite differential volumes,
so it cancels.  The sum of the Riemannian covolumes at level one is
therefore
\[
        r^{f/2}|\operatorname{disc}F|^{(r^2-1)/2}
                             \prod_v a_v^{-1}.
\]
This computation allows a projective integral Lie lattice and
does not assume it has a global integral basis.

The compact diagonal torus has Riemannian volume
$(2\pi)^{r-1}\sqrt r$ at each real place; the determinant-one
condition gives the Gram matrix
$I_{r-1}+\mathbf1\mathbf1^t$.  Each complex off-diagonal flag
coordinate contributes the factor two.  Consequently the exact
conversion from a flag fundamental-set volume to its special
unitary covolume, for a torsion-free arithmetic group, is
\[
       \left(2^{r(r-1)/2}(2\pi)^{r-1}\sqrt r\right)^f.
\]
Since $\zeta(j)<2$ for $j\ge2$, the ratio of the sum of
level-one group covolumes to this factor is at most
$|\operatorname{disc}F|^{(r^2-1)/2}$.  At level one this is a
volume with stabilizer weights.  We use the ordinary flag
fundamental-set volume only after passing to the principal level.

It remains to retain the determinant classes and level.
An element of the principal congruence group has determinant
$a\in\mathcal O_E^\times$ with $a\bar a=1$ and
$a-1\in M\mathcal O_E$.  If $a\ne1$, the norm of $a-1$ is
both at most $2^{2f}$ and at least $M^{2f}$, a contradiction.
Thus these rational stabilizers really lie in the special unitary
group; this is where $M\ge3$ is used.
The principal group is also torsion-free.  Represent it on the
underlying integral lattice.  If an odd prime $p$ divides $M$,
the group $1+p\operatorname{Mat}_n(\mathbb Z_p)$ is torsion-free;
if $M$ has no odd prime factor then $4\mid M$, and one uses
$1+4\operatorname{Mat}_n(\mathbb Z_2)$.
For completeness, if $A-1=p^aB$ has least entry valuation $a$,
the binomial expansion shows that $A^p-1$ has least entry
valuation $a+1$ in these ranges.  Reduction shows that any
finite order would be a power of $p$, which is impossible.
The action on the flag domain is therefore free, since a
discrete subgroup meets its compact point stabilizer in a
finite group.  This justifies the preceding measure conversion
at the actual level, including the central stabilizers.
Hilbert 90 writes a finite norm-one idele as $b/\bar b$.
Choose the ideal class of $b$ among $h_E$ representatives; the
remaining integral units, modulo $M$, have at most $M^{2f}$
residue classes.  Their quotient has norm one and is congruent
to one modulo $M$.  Locally it is the determinant of an integral
unitary transformation congruent to one: diagonalize the
unramified unimodular Hermitian lattice and act on one line.
Descent gives the assertion for its original lattice.
Thus the number of determinant classes is at most
$h_E M^{2f}$.  In each determinant class the principal level
index is at most $M^{2fr^2}$, by its injection into endomorphisms
of the rank-$r$ module modulo $M$.  These two bounds, with the
harmless factor two retained in $V$, prove
\eqref{eq:actual-upper-mass}.  Finally,
$\Delta=|\operatorname{disc}F|^2$ follows from finite
unramifiedness, giving the exponent in
\eqref{eq:actual-upper-constant}.

The constructed component dictionary identifies its probability
norm with the sum of these integrals divided by the mass of the
corresponding components.  That mass is at most the full sum in
\eqref{eq:actual-upper-mass}.  Applying
\eqref{eq:actual-upper-raw} and then this exact normalization proves
the asserted pointwise norm inequality.  Taking its supremum over
unit sections proves the Bergman bound.  Integrating an orthonormal
basis gives $\int_Z B_Z\,d\mu_Z=\dim\mathcal H_Z$, and proves the
dimension bound.

\emph{The numerical estimate.}
Minkowski's ideal-class theorem gives an integral ideal of norm
at most $\sqrt\Delta$ in each class.  Indeed its CM constant is
at most $(2/\pi)^f<1$, by bounding the first $f$ factors in
$(2f)!/(2f)^{2f}$ by $1/2$.  Counting ideals with the Dirichlet
series at two gives
\[
                 h_E\le\Delta\zeta_E(2)
                        \le4^f\Delta.
\]
Under \eqref{eq:actual-upper-size} this implies
$\log V\le80fr^2L$.
Also
\[
       \epsilon^{-1}\le3fwM\sqrt{d_H},\qquad
       \log\epsilon^{-1}\le2+35L.
\]
Using $1+2n^2\Delta/\epsilon\le3n^2\Delta/\epsilon$
and $n=2fr$, we obtain
\[
          \log A_\epsilon
                  \le86fr^2L+8f^2r^2L.
\]
Since $D!\le D^D$, $D\le fr^2/2$, $\log D\le2L$,
$\log(w+f)\le\log2+2L$ and $\log2<3/4$, the chart factor obeys
\[
        \log\frac{8D!}{(\pi\rho^2)^D}\le14fr^2L.
\]
Adding gives at most $188f^2r^2L$, and hence the stated bound
$250f^2r^2L$.  Each factor in $C_{\mathrm{sup}}$ is at least
one, so its lower bound also holds.

For the family statement put $S=\sum_{\tau,i}|u_{\tau,i}|$ and
$w=4fS+1$.  The exact model-weight formula gives total spread at
most $4fS$.  If $S=0$, then $w=1$.  Otherwise $S\ge1$ because
the weights are integral; the displayed hypotheses imply
$\log w\le25L_b$.  The root-discriminant bound gives
$\log\Delta\le6fL_b$, and $d_H^f\le\Delta$ gives
$\log d_H\le6L_b$.  Apply
\eqref{eq:actual-upper-logbound} with $L=25L_b/2$.
It gives
\[
      \log C_{\mathrm{sup}}\le3125f^2r^2L_b
                                 \le4000(frL_b)^3.
\]
All constants are independent of the component, embedding and
growing arithmetic field.
\end{proof}

\begin{corollary}[The multiplication norm for the original sections]
\label{cor:actual-upper-multiplication}
If $s$ has the first weight and $t$ the second weight in the
original section multiplication, then
\[
       \|st\|_{2,Z}\le\sqrt{C_{\mathrm{sup}}}
                              \|s\|_{2,Z}\|t\|_{2,Z}.
\]
This holds for the original integral multiplication map whenever
$s,t$ are integral sections, since their pointwise metrics multiply.
\end{corollary}
\begin{proof}
Integrate $\|s(x)t(x)\|^2\le
(\sup\|s\|^2)\|t(x)\|^2$ against the same component probability
measure and apply Theorem~\ref{thm:actual-upper-bergman}.
\end{proof}

The theorem supplies the global upper Bergman, supremum and rank
estimates with their original numerical family bound, and the
corollary supplies the corresponding multiplication norm.  The
thin-region lower bound, point-height comparison and final arithmetic
section remainder remain separate.  Generation of the entire
integral cusp lattice from the short coefficient-field spanning
family is treated in Section~\ref{shared-actual-integral-cusp-frame}.

\subsection{Generators of the original integral cusp lattice}
\label{shared-actual-integral-cusp-frame}

The short theta sections span over the coefficient field.  The
arithmetic section argument also uses a finite family generating the
whole integral lattice~\cite{bergman-integral-frame-source}.
Its index in the lattice spanned by the
known sections can be arbitrarily large.  The following construction
retains that index without putting it into the norm bound.

\begin{lemma}[Integral generators from a short spanning family]
\label{lem:actual-integral-generator-conversion}
Let $K$ be a number field of degree $a$, let $V$ have dimension
$q\ge1$ over $K$, and let $\mathcal M\subset V$ be a full,
finitely generated torsion-free $\mathcal O_K$-module.  Give every
$V\otimes_{K,\sigma}\mathbb C$ a Hermitian norm.  Suppose
$\mathcal M$ contains a $K$-spanning family with norm at most
$R\ge1$ at every embedding.  Then $\mathcal M$ has at most $2aq$
generators over $\mathbb Z$, each with every embedding norm at most
\begin{equation}\label{eq:actual-integral-generator-bound}
                      a^3q\operatorname{rd}(K)R.
\end{equation}
No freeness over $\mathcal O_K$ or bound for a lattice index is
assumed.
\end{lemma}

\begin{proof}
Use the canonical Euclidean norm
$\|x\|^2=\sum_{\sigma:K\hookrightarrow\mathbb C}|\sigma x|^2$,
where both members of a conjugate pair are included.  The Euclidean
dual of $\mathcal O_K$ is the conjugate of the inverse different.
For $0\ne b\in\mathfrak D_K^{-1}$, its ideal norm gives
$|N_{K/\mathbb Q}b|\ge|\operatorname{disc}K|^{-1}$.
Arithmetic--geometric mean therefore gives dual minimum at least
$\sqrt a/\operatorname{rd}(K)$.
The elementary transference bound
$\lambda_a(L)\lambda_1(L^*)\le a^2$ gives integral elements
$b_1,\ldots,b_a$, independent over $\mathbb Q$, such that
\[
           |\sigma b_j|\le a^{3/2}\operatorname{rd}(K)
                         \le a^2\operatorname{rd}(K)
                      \quad\text{for every }\sigma,j.
\]
Select a $K$-basis $s_1,\ldots,s_q$ from the given short family.
The $aq$ elements $b_js_i$ form a $\mathbb Q$-basis and belong to
$\mathcal M$.  Their integral span $\mathcal N$ has finite index
in $\mathcal M$.

The finite abelian group $\mathcal M/\mathcal N$ has at most
$aq$ generators: the images of any integral basis of the rank-$aq$
abelian group $\mathcal M$ suffice.  Reduce a representative of
each chosen generator in the parallelepiped of the basis $b_js_i$,
so that all its real coordinates lie in $[-1/2,1/2]$.
The reduced representative remains in $\mathcal M$ and has norm
at most
\[
          \frac{aq}{2}\,a^2\operatorname{rd}(K)R
                          \le a^3q\operatorname{rd}(K)R
\]
at every embedding.  The basis vectors of $\mathcal N$ satisfy
the same bound.  These two finite families generate
$\mathcal M$ over $\mathbb Z$ and have at most $2aq$ members.
The reduction is purely in the underlying real vector space;
it subtracts integral multiples of elements of $\mathcal N$
and so introduces no denominators into $\mathcal M$.
\end{proof}

\begin{theorem}[The prescribed integral cusp frame]
\label{thm:actual-integral-cusp-frame}
Use the original good cover and coefficient field $K$, with all
the hypotheses of Theorem~\ref{thm:actual-upper-bergman} for the
first divided weight $u_1$.  Put
\[
 X=frL_b,\qquad Y=\kappa/T+1,
 \qquad T\ge1,\quad\kappa\ge0.
\]
Suppose the original numerical conditions give
\begin{gather}
 X\ge8,\qquad P\ge X^{12},\qquad
 \log[K:\mathbb Q]\le391X^2,\qquad
 \log\operatorname{rd}(K)\le391X^2,\label{eq:actual-cusp-frame-size}\\
 \sum_{\tau,i}|u_{1,\tau,i}|\le X^2Y,
 \label{eq:actual-cusp-frame-weight}
\end{gather}
and give the antitonicity and central gap $8fr+10$ required by
Theorem~\ref{thm:actual-shifted-short-span}.
Then the original integral cusp lattice on each component admits
a finite family of integral generators $w_1,\ldots,w_n$ and
numbers $\Lambda,\beta\ge1$ with
\begin{gather}
 \log n\le P,\qquad \log\beta\le P,\qquad
 \log\Lambda\le P(\kappa/T+1),
 \label{eq:actual-cusp-frame-budget}\\
 \|w_i\|_{2,Z,\sigma}\le\Lambda\quad\text{for all }i,\sigma,
 \qquad B_{Z,\sigma}(x)\le\beta.
\end{gather}
Here the norms are the original integral cusp-lattice Hermitian
norms, as identified with the component Petersson norms by the
constructed dictionary.  The coordinate matrix of these generators,
together with that Hermitian Gram matrix, is a frame for the
original lattice.  Its integral span, even over $\mathbb Z$, is
the whole lattice.

If a component's section space is zero, its integral lattice is
zero and the assertion is interpreted using one zero generator;
then $n=\Lambda=1$ is allowed.  The positive-dimensional spaces
in the original section comparison are covered without this
convention.
\end{theorem}

\begin{proof}
For a nonzero section space write $q$ for its dimension over $K$
and $a=[K:\mathbb Q]$.  The integral lattice is full and finite:
the integral model is projective, the cusp condition is the
prescribed linear restriction, and its ambient section module is
torsion-free.  These are the same finite modules used in the
short-section construction; no saturation at a finite prime is
changed here.

Theorem~\ref{thm:actual-shifted-short-span} supplies a finite
integral $K$-spanning family with all its original Hermitian norms
at most
\[
 R=\exp\left(400000\left[X^3
                   \sum_{\tau,i}|u_{1,\tau,i}|+X^5\right]\right).
\]
Apply Lemma~\ref{lem:actual-integral-generator-conversion} to
this very lattice and these norms.  Theorem
\ref{thm:actual-upper-bergman} gives
\[
           q\le\beta,\qquad \beta=\exp(4000X^3).
\]
We may therefore choose $n\le2aq$ and
$\Lambda=a^3q\operatorname{rd}(K)R$.
Writing the generators in any $K$-basis gives their required
coordinate matrix; the Gram matrix is that of the original
Hermitian norm in this basis.  The isometry and generation
properties follow from these definitions.  This construction
does not require a free integral basis.

Since $Y\ge1$,
\begin{align*}
 \log n&\le\log2+391X^2+4000X^3\le5000X^3,\\
 \log\Lambda
 &\le1564X^2+4000X^3
                      +400000X^5(Y+1)\\
 &\le806000X^5Y.
\end{align*}
For $X\ge8$ one has $5000X^3\le X^{12}$,
$4000X^3\le X^{12}$ and $806000X^5\le X^{12}$.
Together with $P\ge X^{12}$, these are precisely
\eqref{eq:actual-cusp-frame-budget}.
\end{proof}

In the original section construction, $P$ is its prescribed size
parameter and $T$ its prescribed integer length.  The existing
weight and size calculations give
\eqref{eq:actual-cusp-frame-size}--\eqref{eq:actual-cusp-frame-weight}
for the actual anchored division of the weight.  Thus the theorem
supplies the upper cusp frame required there, without a new
coefficient field, a free-lattice hypothesis or a hidden index
factor.  The second divided weight already has its finite
Hermitian lattice by the constructed positivity and section
theorems, and requires no additional quantitative frame bound.
The thin-region lower norm, point heights and final section
remainder remain open.

\subsection{A determinant bound for the section lattice}
\label{sec:arithmetic-determinant-degree}

The arithmetic degree can be estimated before constructing integral
lifts of individual jets.  We develop the determinant argument
of~\cite{arithmetic-determinant-source}.  We use normalized degrees: all finite-place
sums and sums over the complex embeddings of a number field $K$ are
divided by $[K:\mathbf Q]$.  Finite absolute values are normalized by
$|\pi|_v=(\#k_v)^{-1}$, and every complex embedding occurs in the
archimedean sum.

Let $1\leq N\leq n$, and let $w_1,\ldots,w_n$ span $K^N$.
At every embedding $\sigma$ fix a
Hermitian norm $|\cdot|_\sigma$ on $\mathbf C^N$, with
$|w_i|_\sigma\leq\Lambda$, where $\Lambda\geq1$.  For a nonzero
multihomogeneous polynomial $P$ in blocks of $N$ variables, write
$H(P)$ for the sum of its finite Gauss heights after substituting
$u=\sum_i z_iw_i$ in each block, and its mean logarithmic absolute
value on the product of the Hermitian unit spheres at infinity.
Thus no maximum norm replaces the spherical mean.  Put
\[
 c=\log(n\Lambda)+\frac{\log n+1}{2}.
\]

\begin{lemma}\label{lem:arithmetic-polynomial-determinant}
For nonzero multihomogeneous polynomials,
$H(PQ)=H(P)+H(Q)$, $H(a)=0$ for $a\in K^\times$, and
$H(P)\geq-c\deg P$.  Let $M$ be a full lattice in a Hermitian
$K$-space $V$ of dimension $d$.  Suppose $p_1(u),\ldots,p_d(u)$ are
polynomial vectors, their coefficients after the above substitutions
belong to $M$, and their determinant $P$ is nonzero and
multihomogeneous.  If $\|p_j(u)\|_\sigma\leq B_j$ on every product
of unit spheres, then
\begin{equation}\label{eq:arithmetic-polynomial-determinant}
             \widehat{\deg}M\geq H(P)-\sum_{j=1}^d\log B_j.
\end{equation}
\end{lemma}

\begin{proof}
Multiplicativity follows from the Gauss lemma and addition of
integrable logarithms; scalar invariance is the product formula.
For a homogeneous polynomial of degree $a$ in $N$ complex variables,
iterated Jensen inequalities on coordinate circles give
\[
 \log|P(u)|\leq\int_{|z|=1}\log|P(z)|\,d\mu(z)
                    +\frac a2\sum_{j=1}^{N-1}\frac1j
                 \qquad(|u|=1).
\]
Indeed rotate $u$ to the first coordinate, average the phases of the
other coordinates, and use
$\int\log|z_1|\,d\mu=-\frac12\sum_{j=1}^{N-1}j^{-1}$.
Apply this in each block.  On the coefficient polytorus the substituted
vectors have norm at most $n\Lambda$.  The Cauchy coefficient inequality
and the product formula, applied to any nonzero coefficient, now give
the asserted lower bound, since $N\leq n$ and
$\sum_{j=1}^{N-1}j^{-1}\leq\log n+1$.

For the determinant assertion, choose a $K$-basis $b$ of $V$, and let
$I$ be the fractional ideal representing $\bigwedge^dM$ in that basis.
Every coefficient of $P$ belongs to $I$, by multilinearity of the
determinant.  Its finite Gauss contribution is therefore at most
$-[K:\mathbf Q]^{-1}\log\operatorname{N}I$.
At infinity, Hadamard's inequality gives
\[
 \log|P(u)|\leq\sum_j\log B_j
                   -\tfrac12\log\det\operatorname{Gram}_\sigma(b).
\]
Average and add the finite contributions.  This is exactly the
definition of $\widehat{\deg}M$.  A basis of $M$ over the ring of
integers is neither chosen nor needed.
\end{proof}

We specify the geometric input to the next theorem.  Let $D\geq1$,
$T\geq D+1$, and let $V_j$ be finite dimensional $K$-spaces in a
common space with commuting multiplication operators $m(u)$, linear
in $u\in K^N$, taking $V_j$ to $V_{j+1}$.  Let $M_j\subset V_j$
be full lattices, stable under each $m(w_i)$.  Assume at infinity
\[
 \|m(u)x\|_{j+1,\sigma}\leq\sqrt\beta\,|u|_\sigma\|x\|_{j,\sigma},
                  \qquad\beta\geq1,
\]
and assume that $M_0,M_D,M_T$ contain $K$-spanning families of norms
at most $\Lambda_0,\Lambda_D,\Lambda_T$, respectively, all at least one.

Let $f(s_1,\ldots,s_D,l)$ be irreducible, of degree $\delta\geq1$
in every block.  On a nonempty principal open set of $D$-tuples $s$
over $K$, assume the following finite-dimensional description, after
a field extension if necessary.  There are nonzero linear forms
$L_1,\ldots,L_\delta$, a surjection
$e:V_0\otimes_K E\to E^\delta$, and a decreasing filtration $F_a$ in $V_T\otimes_K E$,
$0\leq a\leq T-D$.  The middle two identities below hold for
$0\leq a<T-D$; the last describes the deepest space:
\begin{equation}\label{eq:arithmetic-generic-filtration}
 \begin{split}
 f(s,l)&=c_s\prod_{p=1}^{\delta}L_p(l),\qquad c_s\ne0,\\
 F_a/F_{a+1}&\simeq\bigoplus_{|\alpha|=a}E^\delta,\\
 m(s)^\alpha m(l)^{T-a}x&\longmapsto
           \bigl(\mathbf1_{\alpha'=\alpha}L_p(l)^{T-a}e_p(x)\bigr)_{\alpha',p},\\
 F_{T-D}&=\sum_{|\alpha|=T-D}m(s)^\alpha(V_D\otimes E).
 \end{split}
\end{equation}
Here the product vectors in the third line belong to $F_a$; the
isomorphism in the second line has kernel $F_{a+1}$.  These are
geometric vector-space conditions.  They do not assert that arbitrary
integral jets lift with controlled norm.

\begin{theorem}\label{thm:arithmetic-determinant-degree}
Under these hypotheses put
\[
 R=\binom{T-1}{D},\quad q=R\delta,\quad d=\dim_KV_T,\quad
 h_f=\frac{H(f)}{(D+1)\delta},\quad\theta=\frac qd.
\]
Then $d\geq q>0$ and
\begin{equation}\label{eq:arithmetic-determinant-slope}
 \begin{split}
 \frac{\widehat{\deg}M_T}{d}\geq{}&
 \theta\left(Th_f-D(D+1)c-\log\Lambda_0-\frac T2\log\beta\right)\\
 &-(1-\theta)\left(Tc+\frac T2\log\beta
                          +\log\Lambda_D+\log\Lambda_T\right).
 \end{split}
\end{equation}
If the geometric height $h$ satisfies
$h_f\geq h-(\log n+1)/2$, the same bound holds with $Th_f$ replaced
by $Th-T(\log n+1)/2$.
\end{theorem}

\begin{proof}
Choose a good $s$ for which the dimension $y$ of the deepest space is
maximal.  From the original $K$-vectors in the short family in $M_0$
select $\delta$ vectors whose residues are independent over $E$.
Choose $l\in K^N$ avoiding the finitely many hyperplanes $L_p(l)=0$;
this is possible because $K$ is infinite and each $L_p$ is nonzero.
For every $a<T-D$ and $|\alpha|=a$, their products
$m(s)^\alpha m(l)^{T-a}x$ give a basis of that graded piece.
Choose $y$ independent deepest vectors from the products of the original
short family in $M_D$, and complete with vectors from the original
short family in $M_T$.  All chosen vectors are over $K$; independence
may be checked after extension to $E$.  There are exactly $q$ vectors
of the first kind.

Replace $s,l$ by variables in these fixed columns and take their
determinant $P$.  It is nonzero, since the chosen specialization is a
basis.  All substituted coefficients are in $M_T$ by multiplication
stability.  The three types of column have bounds
\[
       \beta^{T/2}\Lambda_0,\qquad
       \beta^{(T-D)/2}\Lambda_D,\qquad\Lambda_T,
\]
and occur $q,y,d-q-y$ times.

Write
\[
 b=\sum_{a=0}^{T-D-1}\binom{a+D-1}{D-1}(T-a),\qquad
                 (D+1)b=R\{T+D(D+1)\}.
\]
At any good specialization the selected deepest columns either become
dependent, in which case $P(s,l)=0$, or form a basis of the deepest
space by maximality of $y$.  In the latter case the determinant on the
graded pieces is
$\gamma_s(\prod_pL_p(l))^b$, with $\gamma_s$ independent of $l$.
Thus $P(s,l)$ is a scalar multiple of $f(s,l)^b$ in both cases.
To obtain divisibility over $K$, choose $l_0$ such that
$f(s,l_0)$ is not the zero polynomial.  Zariski density gives
\[
       P(s,l)f(s,l_0)^b=P(s,l_0)f(s,l)^b.
\]
The irreducible polynomial $f$ cannot divide the nonzero polynomial
$f(s,l_0)$, which has degree zero in $l$.  Hence $P=f^bG$ for a
nonzero multihomogeneous polynomial $G$.  Its total degree is
\[
           e=y(T-D)-qD(D+1)\geq0.
\]
Lemma~\ref{lem:arithmetic-polynomial-determinant} proves the sharper
intermediate bound
\begin{equation}\label{eq:arithmetic-determinant-exact}
 \begin{split}
 \widehat{\deg}M_T\geq{}&bH(f)-ce-q\log(\beta^{T/2}\Lambda_0)\\
 &-y\log(\beta^{(T-D)/2}\Lambda_D)
   -(d-q-y)\log\Lambda_T.
 \end{split}
\end{equation}
Finally $h_f\geq-c$, $e\leq T(d-q)$, and $y\leq d-q$.
Together with the displayed formula for $b$, these inequalities give
\eqref{eq:arithmetic-determinant-slope}.  The geometric-height version
follows by substitution.  This estimates the degree of the full
lattice $M_T$, including its finite-place contribution.
\end{proof}

The required geometric filtration, Chow form and height identity, as
well as the full lattices and their short spanning families, remain
inputs to an application.  Once these data are supplied, the argument
uses only the short generators themselves; separate arithmetic jet
lifts and their extension estimates are unnecessary.

\subsection{Constructing the geometric filtration}
\label{sec:geometric-section-filtration}

The vector-space filtration in the determinant argument follows from
ordinary vanishing and the Koszul complex.  It does not require an
arithmetic extension estimate.

Let $Y$ be a smooth geometrically integral projective variety of
dimension $D\geq1$ over a number field $K$, and fix an integer
$T\geq D+1$.  Let $L$ be generated by
a fixed space of sections, let $A$ and $B$ be nef and big line bundles,
and fix
\[
 e\in H^0(Y,A\otimes B^{-1}\otimes L^{-D}),\qquad
 U\subset\{e\ne0\}
\]
with $U$ a nonempty open set.  Write $\mathcal G_j=\omega_Y\otimes A\otimes L^j$
and $G_j=H^0(Y,\mathcal G_j)$.  For $0\leq j\leq T$, suppose that the original spaces
$V_j\subset H^0(U,\mathcal G_j)$ are finite-dimensional and contain
$G_j$, that $V_0=G_0$, and
that multiplication by the fixed sections of $L$ maps $V_j$ into
$V_{j+1}$ for $j<T$.  Global sections are identified with their restrictions to
$U$, an injective map because $Y$ is integral.

\begin{proposition}\label{prop:geometric-section-filtration}
Let $s_1,\ldots,s_D$ in the fixed linear system
have a nonempty finite \'etale common zero scheme $Z$ contained in $U$.
Let $I$ be its ideal, set $a_0=T-D$, and put
\[
 P_s=\sum_{|\alpha|=a_0}s^\alpha V_D,\qquad
 W_s=G_T+P_s\subset V_T,\qquad
 F_a=W_s\cap H^0(U,I^a\mathcal G_T).
\]
After a finite extension splitting $Z$, with $\delta=\deg Z$, these
spaces and their natural symbol maps satisfy the filtration and product
identities in~\eqref{eq:arithmetic-generic-filtration}.  In particular,
\[
             F_{a_0}=P_s,
 \qquad F_a/F_{a+1}\simeq
           \bigoplus_{|\alpha|=a}E^\delta\quad(0\leq a<a_0).
\]
The evaluation map from $V_0$ onto the $\delta$ residue coordinates is
surjective.
\end{proposition}

\begin{proof}
Since $L$ is generated, it is nef.  For $C=A$ or $B$ and every $j\geq0$,
$C\otimes L^j$ is nef and big.  Kawamata--Viehweg vanishing therefore
gives
\[
       H^i(Y,\omega_Y\otimes C\otimes L^j)=0
                      \qquad(i>0,\ j\geq0).
\]
One may extend to an algebraic closure and descend this statement by
flat cohomology base change~\cite[Tag 02KH]{stacks-filtration-regular}.
We use precisely the line-bundle case of
\cite[Theorem 2.28]{viehweg-filtration-vanishing}.

At a point of $Z$, the $s_i$ are a regular system of parameters after
trivializing $L$.  Off $Z$ one is a unit.  Their Koszul complex is
therefore exact, and the associated graded algebra of $I$ is the
symmetric algebra on $I/I^2$; its degree-$a$ basis consists of the
monomials $s^\alpha$ with $|\alpha|=a$
\cite[Lemmas 31.21.1 and 31.21.4]{stacks-filtration-regular}.

First twist the Koszul complex by $\omega_Y\otimes B\otimes L^D$.
All its terms have the form $\omega_Y\otimes B\otimes L^{D-i}$,
where $0\leq i\leq D$, so the vanishing just proved makes restriction
onto $Z$ surjective.  Multiplication by $e$ maps these global sections
into $G_0=V_0$.  On $Z$, this multiplication is invertible.
Adjunction gives $(\omega_Y\otimes L^D)|_Z\simeq\omega_Z$, and
$\omega_Z\simeq\mathcal O_Z$ because $Z$ is \'etale.  Thus the actual
residue map $V_0\to H^0(Z,B|_Z)$ is onto.  No negative twist of $A$
has been used.

Next twist the Koszul complex by $\mathcal G_n$, with $n\geq D$.
Its terms have nonnegative exponents $n-i$, so it gives first-order
division:
\[
 H^0(Y,I\mathcal G_n)=\sum_{i=1}^D s_iG_{n-1}.
\]
Inductively, for every $a\geq0$ and $j\geq D$,
\begin{equation}\label{eq:geometric-global-division}
 H^0(Y,I^a\mathcal G_{j+a})
                  =\sum_{|\alpha|=a}s^\alpha G_j.
\end{equation}
Indeed write a section vanishing to order $a+1$ as
$\sum_{|\alpha|=a}s^\alpha g_\alpha$ with $g_\alpha\in G_{j+1}$,
using the induction hypothesis.  Its degree-$a$ symbol is zero.
Independence of the monomials in $I^a/I^{a+1}$ forces each
$g_\alpha|_Z=0$.  First-order division at $n=j+1$ now divides every
coefficient by one further $s_i$, with coefficients in $G_j$.
Collecting equal monomials proves the next step.  This proves the
required higher division without assuming vanishing for powers of $I$.

All elements of $P_s$ vanish to order $a_0$.  If $g+p\in W_s$ has
that order, with $g\in G_T$ and $p\in P_s$, then so does $g$.
Equation~\eqref{eq:geometric-global-division}, with $j=D$, puts $g$
in $\sum s^\alpha G_D\subset P_s$.  This proves $F_{a_0}=P_s$.

Finally split $Z$ and choose fibre coordinates for $B$ and $L$.
Write $L_p(l)$ for evaluation of a section $l$ at its $p$th point.
The local regular-sequence calculation identifies the degree-$a$
symbol of $s^\alpha l^{T-a}x$, for $x\in V_0$, with
\[
       (\mathbf1_{\alpha'=\alpha}
                       L_p(l)^{T-a}\operatorname{res}_p(x))_{\alpha',p}.
\]
Choose one $l$ over $K$ nonzero at every point of $Z$; this is possible
over the infinite field $K$ by avoiding finitely many hyperplanes.
The residue map is
onto, so these products supply every symbol.  They lie in $G_T$ and
hence in $W_s$.  The kernel of the symbol map is exactly $F_{a+1}$
by the definition of the $I$-adic filtration.  This supplies all the
claimed isomorphisms and identities.
\end{proof}

If the fixed linear system is birational onto its image $X$, good
sections as above exist on a nonempty open set of tuples.  Here
$\delta=\deg X$.  The projection construction used in
Proposition~\ref{prop:finite-section-jets} gives a morphism
$f:Y\to\mathbf P^D$ of generic degree $\delta$.  Let $U^\circ$ be
the intersection of $U$ with the birational isomorphism locus.
The image $f(Y\setminus U^\circ)$ is closed and has dimension at most
$D-1$.  Choose $q$ outside it and in the finite \'etale locus of $f$.
Its fibre consists of $\delta$ distinct points in $U^\circ$, and
its equations are $D$ sections from the same fixed linear system.

The universal common zero scheme over the affine tuple space is
proper.  At this tuple its relative Jacobian is invertible.  Remove
the closed images of its non-\'etale locus and of its intersection
with $Y\setminus U^\circ$.  Over the resulting nonempty open set,
the universal zero scheme is proper and \'etale, hence finite
\'etale, and has constant degree $\delta$.  Restrict to a nonempty
principal open.  All constructions are over $K$; since $K$ is
infinite, this open has $K$-rational tuples.

For the original Chow form $f$, its specialization at any such tuple
is nonzero and vanishes precisely on the hyperplanes through these
$\delta$ distinct image points.  Its degree is $\delta$, so
$f(s,l)=c_s\prod_pL_p(l)$ with $c_s\ne0$.  Thus the filtration and
Chow specialization required by the determinant theorem are supplied
by the original geometric hypotheses.  The arithmetic model, its
metrics, multiplication bounds, short integral families and geometric
height identity remain the data needed for the arithmetic application.

\subsection{The arithmetic degree of the original component lattice}
\label{sec:actual-component-section-degree}

We now apply the determinant theorem to the actual cusp lattices.
The geometric realization, its canonical line, the finite determinant
ideals and the Petersson norms are all retained.  In particular the
argument does not require integral lifts of arbitrary jets or the
selection of algebraic points of controlled height.

\paragraph{The original spaces.}
Fix one of the constructed components over its original totally complex
coefficient field $K$.  Let $X$ be its smooth projective toroidal full
flag compactification, $B_X$ its reduced boundary, and $U=X\setminus B_X$.
Write $D=\dim X\ge1$, and use the original automorphic lines
$\mathcal L(a)$ and canonical weight $\kappa_{\mathrm{can}}$.  Their
prescribed multiplicative identifications and
\[
                  \omega_X(B_X)=\mathcal L(\kappa_{\mathrm{can}})
\]
are those proved in Section~\ref{actual-geometric-positivity}.
Take the two original divided weights $u_1,u_2$ and the original integer
$T\ge D+1$.  Retain their original strict gap bounds
$\operatorname{gap}(u_i)\ge(D+3)^2+2r$, $i=1,2$, where
$r=p+q$ is the original rank.  Put
\[
 a_j=\kappa_{\mathrm{can}}+u_2+(D+j)u_1,\qquad 0\le j\le T.
\]
For any weight $a$ used here, let $C(a)$ be its original $K$-space of
cuspidal sections supported on this component, with its original
integral lattice $M(a)$ and component Petersson norms
$\|\cdot\|_{\sigma,a}$, for all embeddings $\sigma:K\hookrightarrow
\mathbb C$.  These norms use the original component probability
measure.

Choose the full integral frame of $C(u_1)$ constructed in
Section~\ref{shared-actual-integral-cusp-frame}.  Denote its generators
by $w_1,\ldots,w_n$, its dimension by $N$, its generator bound by
$\Lambda\ge1$, and its upper Bergman bound by $\beta\ge1$.
Thus the $\mathcal O_K$-span of the $w_i$ is exactly $M(u_1)$,
$1\le N\le n$, and
\begin{equation}\label{eq:actual-section-frame}
 \|w_i\|_{\sigma,u_1}\le\Lambda,\qquad
 \|s(x)\|^2\le\beta\|s\|_{\sigma,u_1}^2.
\end{equation}
Let $\Lambda_0,\Lambda_D,\Lambda_T\ge1$ be the supplied short
spanning bounds for $M(a_0),M(a_D),M(a_T)$.

Write $\delta$ and $h$ for the original image degree and normalized
image height of the system $C(u_1)$ on this component.  These use
the full lattice $M(u_1)$ and its stated Hermitian norms.  The precise
height formula is recalled and verified below.  Put
\begin{gather*}
 d=\dim_K C(a_T),\qquad
 q=\binom{T-1}{D}\delta,\qquad \vartheta=\frac qd,\\
 c=\log(n\Lambda)+\frac{\log n+1}{2},\qquad
 a_n=\frac{\log n+1}{2}.
\end{gather*}

\begin{theorem}[The original component section estimate]
\label{thm:actual-component-section-degree}
For the constructed geometric data and the original weight-size
conditions, one has $d\ge q>0$ and
\begin{equation}\label{eq:actual-component-section-strong}
\begin{split}
\frac{\widehat{\deg} M(a_T)}d\ge{}&
 \vartheta\left(Th-Ta_n-D(D+1)c-\log\Lambda_0
                                      -\frac T2\log\beta\right)\\
&-(1-\vartheta)\left(Tc+\frac T2\log\beta
                             +\log\Lambda_D+\log\Lambda_T\right).
\end{split}
\end{equation}
Arithmetic degree is divided by $[K:\mathbb Q]$ and is the original
component degree.

In particular, set
\[
 \rho=\log\Lambda+
 26\log\!\left(nD\delta(2+|h|+\log\Lambda)\right),
 \qquad G=\frac{(D+1)(\log n+1)}2+1,\qquad z=1-\vartheta .
\]
Then $0\le z\le1$ and the literal numerical estimate required by
the component construction holds:
\begin{equation}\label{eq:actual-component-section-original}
\begin{split}
\frac{\widehat{\deg}M(a_T)}d\ge{}&
(1-z)\bigl(Th-3(D+1)|h|-\log\Lambda_0
 -26D^2\log(e^\rho T)\\[-2pt]
&\hspace{36mm}-(T+D+1)\log\beta-(T+D^2)(G+2)\bigr)\\
&-z\bigl(T\log(2e^\rho\sqrt\beta)
                              +\log\Lambda_D+\log\Lambda_T\bigr).
\end{split}
\end{equation}
At the original family size $P\ge0$ and weight size $\kappa\ge0$, the already
constructed short families give
\[
 \Lambda_0=\exp\!\left(DP(\kappa/T+1)\right),\qquad
 \Lambda_D=\Lambda_T=\exp(P\kappa).
\]
Thus the original remainder coefficient is exactly
$J=DP(\kappa/T+1)$.  The theorem supplies the component arithmetic
section inequality; the separate all-point height comparison and
any still unconstructed numerical family are not assumptions
silently discharged by this application.
\end{theorem}

\begin{proof}
\emph{The cusp space and the canonical extension.}
For each weight under consideration there is an equality over $K$
\begin{equation}\label{eq:actual-section-cusp-extension}
             C(a)=H^0(X,\mathcal L(a)(-B_X)).
\end{equation}
Here the line on the right is the subcanonical extension.  The
analytic--algebraic toroidal comparison identifies its automorphic
bundle and Fourier--Jacobi coordinates with those of the original
moduli problem; the subcanonical comparison identifies its sections
with holomorphic cuspidal forms
\cite{positive-lan-analytic,positive-lan-stroh}.
The passage from the automorphic bundle to its full flag line is the
usual highest-weight realization, with the same ordered characters.
The original weights satisfy the needed dominance conditions.

We spell out why no growth condition changes this space.  In a smooth
toroidal boundary chart, a canonical Hodge frame and its inverse have
at most a fixed power of the logarithmic cusp coordinates as growth.
The same holds for every fixed tensor and flag character.  Rapid decay
of a holomorphic cuspidal form therefore makes its coefficient tend
to zero at a generic point of each boundary divisor
\cite{positive-cusp-decay}.  Its algebraic meromorphic coefficient
has positive order there.  Smoothness, hence normality, extends it
across the remaining codimension at least two subset.  It is a
section of the negative reduced boundary.  Conversely such a section
has vanishing constant Fourier--Jacobi coefficients, so its
automorphic reading is cuspidal; the same rapid-decay theorem gives
the original bounded Petersson reading.  The forms have the
prescribed holomorphic highest weight and hence finite infinitesimal
character, as in the earlier cusp verification.  Finally the
identifications are algebraic over the given $K$, and cohomology
base change descends their equality from an embedding into
$\mathbb C$.  This is an equality of rational spaces, and does not
replace their integral lattices.

\emph{The exact geometric realization.}
Principalize the base ideals of $C(u_1)$ and $C(u_2)$
simultaneously over $K$, using their prescribed birational
modification $\pi:Y\to X$~\cite{WlodarczykResolution}.
One can principalize their product; the two factors of an invertible
product ideal are invertible fractional ideals, so both original
ideals become invertible.  Both systems are generated on $U$, so
the modification is an isomorphism there.  The constructed $X$
is smooth projective, and the principalization has smooth projective
source $Y$.  Let $F_1,F_2$ be the actual effective fixed divisors of
the pulled-back systems.  They are supported outside $U$.  Set
\begin{equation}\label{eq:actual-section-lines}
 L=\pi^*\mathcal L(u_1)(-F_1),\qquad
 L_2=\pi^*\mathcal L(u_2)(-F_2),\qquad
 A=\pi^*\mathcal L(u_2+Du_1).
\end{equation}
The two divided systems generate $L,L_2$ everywhere and are
birational.  Thus $L,L_2$ are nef and big.  The original weight gaps
make $\mathcal L(u_2+Du_1)$ nef and big by
Theorem~\ref{actual-positive-cusp-theorem}; pullback under $\pi$
preserves both properties.  This proves the required assertion for
$A$.  The canonical section with divisor $F_2+DF_1$ is a section
\[
            e\in H^0(Y,A\otimes L_2^{-1}\otimes L^{-D})
\]
which is nonzero throughout $U$.

Put $\mathcal G_j=\omega_Y\otimes A\otimes L^j$.
Use the prescribed Kodaira--Spencer canonical identification on $U$
and the actual sections of the $F_i$ to identify $C(a_j)$ with
a space $V_j\subset H^0(U,\mathcal G_j)$.
For clarity the needed inclusions follow without a vanishing
theorem for higher direct images.  One has
\begin{equation}\label{eq:actual-section-canonical-pushforward}
                         \pi_*\omega_Y=\omega_X.
\end{equation}
Indeed a proper birational morphism to the smooth $X$ is an
isomorphism above the generic point of every divisor.  A regular
top form on $Y$ is therefore a regular top form on $X$ outside
codimension at least two, and extends uniquely because $\omega_X$
is locally free and $X$ is normal.  Conversely the pullback of a
regular top form is regular: locally its coefficients are multiplied
by the Jacobian determinant of a morphism between smooth varieties.
These two operations are inverse on a dense open.

By the projection formula, \eqref{eq:actual-section-canonical-pushforward}
and \eqref{eq:actual-section-cusp-extension},
\[
 H^0(Y,\omega_Y A)
   =H^0(X,\omega_X\mathcal L(u_2+Du_1))=C(a_0).
\]
Thus $V_0=H^0(Y,\mathcal G_0)$.  Multiplication by the section
of $jF_1$ embeds
$\mathcal G_j$ into
$\omega_Y\otimes\pi^*\mathcal L(u_2+(D+j)u_1)$.
Its global sections descend, by the same argument, to $C(a_j)$.
Hence
\[
                    H^0(Y,\mathcal G_j)\subset V_j
                         \quad(0\le j\le T).
\]
Multiplication by the divided section of $L$ corresponds exactly
to multiplication by the original cusp section of weight $u_1$:
the additional factor of $F_1$ is cancelled by the next division.
These are precisely the spaces and products in
Proposition~\ref{prop:geometric-section-filtration}.  Its finite
section and Koszul argument supplies all of the required geometric
filtration, including the surjection at degree zero and the equality
at the deepest step.  No arithmetic estimate on $e$ or the $F_i$
enters that argument.

\emph{The original arithmetic data.}
Transport $M(a_j)$ and its Hermitian norms through the preceding
$K$-linear identification.  This is exactly how the original
component section system assigns its lattice and metric
\cite{actual-section-source}.  Thus lattice membership and every
archimedean norm are equal before and after transport.
In particular no lattice of all integral global sections on $Y$
is substituted for $M(a_j)$.

Multiplication of two original integral cusp sections remains
integral at every original integral test object, and preserves the
component support.  The exact weight identity
$u_1+a_j=a_{j+1}$ proves the required lattice stability.
Pointwise metric multiplicativity and the original probability
measure, together with \eqref{eq:actual-section-frame}, give
\[
 \|s\,v\|_{\sigma,a_{j+1}}^2
   \le \|s\|_{\sup,\sigma,u_1}^2\|v\|_{\sigma,a_j}^2
   \le\beta\|s\|_{\sigma,u_1}^2\|v\|_{\sigma,a_j}^2.
\]
The short families at $0,D,T$ transport without a change in bound.

There is also no factor of two hidden in the degree.  If
$a=[K:\mathbb Q]$ and the covolume uses one complex embedding
per conjugate pair with its underlying real Euclidean measure,
then for a rank-$d$ full Hermitian lattice $M$
\[
 \widehat{\deg}M
   =-\frac1a\log\operatorname{covol}(M)
      -\frac d2\log2+\frac d2\log\operatorname{rd}(K).
\]
To verify this, express $\bigwedge^dM$ by a fractional ideal
$I$ in any $K$-basis.  Its real covolume is
$2^{-ad/2}|\Delta_K|^{d/2}\operatorname NI$
times the product of the Hermitian Gram determinants over the
complex places.  Substitution gives
$-a^{-1}(\log\operatorname NI+
\tfrac12\sum_\sigma\log\det\operatorname{Gram}_\sigma)$,
the normalized arithmetic degree used in the determinant theorem.
The original squared-norm slope is twice this degree divided by
rank.  Hence the degree just transported is the stated original
component degree.

\emph{The exact image height.}
Use any $K$-basis of $C(u_1)$, and let $f$ be the irreducible Chow
form of its image, of degree $\delta$ in each of $D+1$ blocks.
Its finite height is the Gauss height after substituting the full
integral generators $w_i$.  Its archimedean height is the mean of
$\log|f|$ over the product of the intrinsic Hermitian unit spheres.
Denote their normalized sum by $H(f)$, exactly as in
Section~\ref{sec:arithmetic-determinant-degree}.  The original image
height satisfies
\begin{equation}\label{eq:actual-section-height}
       h=\frac{H(f)}{(D+1)\delta}
                   +\frac12\sum_{i=1}^{N-1}\frac1i .
\end{equation}
This is its defining Chow-height normalization, not an estimate on
selected algebraic points.

To check independence of the actual frame, two finite generating
families of the same full lattice express one another by integral
matrices.  Integral linear substitution cannot increase a finite
Gauss norm, so the two opposite substitutions prove equality at
every finite place.  At infinity the change between the two coordinate
realizations, each with its transported Hermitian Gram matrix, is
an isometry of the same Hermitian space.  It preserves the spherical
probability measure; the basis-change matrix need not be unitary
for an unrelated standard form.  Chow forms in two bases correspond by that
linear substitution, up to a nonzero scalar; the product formula
cancels the scalar.  This proves \eqref{eq:actual-section-height}
for the original frame and the original $h$, without a projective
maximum norm or an extra period factor.  Birationality identifies
the image degree with the degree of the good finite section in
the geometric filtration.  Since $N\le n$,
\[
        \frac{H(f)}{(D+1)\delta}\ge h-a_n.
\]
All hypotheses of Theorem~\ref{thm:arithmetic-determinant-degree}
are now supplied for the original component spaces and lattices.
It gives \eqref{eq:actual-component-section-strong} and $d\ge q>0$.

\emph{The literal numerical remainder.}
The argument of the logarithm in $\rho-\log\Lambda$ is at least
$2n$.  As $n,\Lambda\ge1$, this implies $0\le c\le\rho$.
Also $D(D+1)\le2D^2$ and $T\ge2$, so
\[
 D(D+1)c\le26D^2(\rho+\log T),\qquad
 \frac T2\log\beta\le(T+D+1)\log\beta,\qquad
 Ta_n\le(T+D^2)(G+2).
\]
The added $3(D+1)|h|$ is nonnegative.  The complementary term
is bounded because
$c+\tfrac12\log\beta\le
\log(2e^\rho\sqrt\beta)$.
Multiplication by the nonnegative factors $\vartheta$ and
$1-\vartheta$ now proves \eqref{eq:actual-component-section-original}.
Finally Corollary~\ref{cor:actual-hs-short-families} supplies the displayed
$\Lambda_0,\Lambda_D,\Lambda_T$ in the original family, while
the full frame theorem supplies $n,\Lambda,\beta$ with their
original bounds.  This proves the stated value of $J$.
\end{proof}

The stronger estimate also combines linearly over components.
If $d_Z,q_Z,h_Z$ are the quantities above, put
$d_{\mathrm{tot}}=\sum_Zd_Z$,
$\Theta=\sum_Zq_Z/d_{\mathrm{tot}}$ and
$\bar h=(\sum_Zq_Zh_Z)/(\sum_Zq_Z)$.
With common upper bounds for the logarithmic costs, summation gives
\[
\begin{split}
\frac{\sum_Z\widehat{\deg}M_Z(a_T)}{d_{\mathrm{tot}}}
\ge{}&\Theta T\bar h-\frac T2\log\beta\\
&-\Theta\bigl(Ta_n+D(D+1)c+\log\Lambda_0\bigr)\\
&-(1-\Theta)\bigl(Tc+\log\Lambda_D+\log\Lambda_T\bigr).
\end{split}
\]
All dependence on heights is through the original signed weighted
mean.  Bounds for that mean, the rank ratios and the remaining
original numerical family are still required at the final
geometric deduction.


\subsection{A finite section bound from one fibre}
\label{sec:finite-section-jets}

The upper section-rank estimate has a short geometric proof. It does
not require a uniform threshold in an asymptotic Hilbert--Samuel theorem.

\begin{proposition}\label{prop:finite-section-jets}
Let $X$ be an integral projective variety of dimension $D\ge1$ over
an algebraically closed field of characteristic zero. Let $A$ be nef,
and suppose a finite linear system in $H^0(X,A)$ gives a birational
map onto a projective image of degree $\delta>0$. Let $B$ be a line
bundle and $n\ge0$ an integer such that $nA-B$ is nef. For an integer $a\ge0$,
\[
 nA^D<(a+1)\delta
 \quad\Longrightarrow\quad
 h^0(X,B)\le\delta\binom{a+D}{D}.
\]
\end{proposition}

\begin{proof}
Blow up the base ideal, obtaining $\pi:Y\to X$ and an effective
Cartier divisor $E$ such that $L=\pi^*A-E$ is generated by the chosen
sections. Then $Y$ is integral and projective, $L$ is nef, and
$L^D=\delta$. A general linear projection of the image gives
$f:Y\to\mathbf P^D$ with $L=f^*\mathcal O(1)$, generically finite
of degree $\delta$. Choose $q$ in its finite \emph{\'etale} locus.
The fibre $Q=f^{-1}(q)$ consists of $\delta$ regular points. We claim
that the pullback of every section of $B$ is determined by its jets
through order $a$ at $Q$.

Suppose a nonzero section $s$ has all these jets zero. Its pullback
has an effective zero divisor $Z_s$. The nonfinite-fibre image of
$f$ has codimension at least two: positive-dimensional fibres over
a divisor would force a proper closed subset of the integral
$D$-fold $Y$ to have dimension at least $D$. A general line $\ell$
through $q$ therefore avoids that image and is not contained in
$f(Z_s)$. For $D=1$ take $\ell=\mathbf P^1$. The curve cycle
$C=f^{-1}(\ell)$ has class $L^{D-1}$. Every curve component dominates
$\ell$, hence none lies in $Z_s$; near $Q$, the curve is smooth
with multiplicity one. Each point of $Q$ contributes at least
$a+1$ zeros to $s|_C$. Consequently,
\[
 (a+1)\delta\le(\pi^*B)L^{D-1}
       \le n(\pi^*A)L^{D-1}\le n(\pi^*A)^D=nA^D,
\]
a contradiction. The middle inequalities use nefness. For the
last one, expand $(\pi^*A)^D-(\pi^*A)L^{D-1}$ as
$E(\pi^*A)\sum_{j=0}^{D-2}(\pi^*A)^{D-2-j}L^j$ if $D\ge2$;
each summand is nonnegative. For $D=1$ it is an equality.

Thus evaluation embeds $H^0(X,B)$ in the direct sum of the local
quotients $(\pi^*B)_z\otimes\mathcal O_{Y,z}/\mathfrak m_z^{a+1}$,
for $z\in Q$. Their dimensions are $\binom{a+D}{D}$, by counting
monomials in regular local parameters. This proves the bound.
The blowup, codimension-one finiteness and intersection facts used
here are recorded in the Stacks Project, tags
\href{https://stacks.math.columbia.edu/tag/01OF}{01OF},
\href{https://stacks.math.columbia.edu/tag/0AB5}{0AB5} and
\href{https://stacks.math.columbia.edu/tag/0BEL}{0BEL}.
\end{proof}

Suppose now that $t>D$ is an integer, $0\le\rho<1$ and
$\delta\ge(1-\rho)A^D$. Set
\[
 n=t+D+3,\qquad
 a=\left\lfloor\frac{t+D+3}{1-\rho}\right\rfloor,\qquad
 w=\delta\binom{t-1}{D},\qquad
 \eta=1-\frac{\binom{t-1}{D}}{\binom{a+D}{D}}.
\]
The proposition gives $(1-\eta)r\le w$ for any section subspace
of dimension $r$. The ratio defining $1-\eta$ is
$\prod_{j=1}^D(t-j)/(a+j)$. Its $j$th factor is at least $1-e_j$,
where
\[
 1-e_j=\frac{(1-\rho)(t-j)}{t+D+3+j(1-\rho)},\qquad
 0\le e_j<1,\qquad
 e_j\le\rho+\frac{D+3+2j}{t}.
\]
Using $1-\prod_j(1-e_j)\le\sum_j e_j$ gives
\[
 \boxed{\eta\le D\rho+\frac{2D^2+4D}{t}
               \le D\rho+\frac{4D^2}{t}\quad(D\ge2).}
\]
If a separate lower estimate gives $w\le r$ and the first majorant
$\Theta$ is less than one, the actual excess rank satisfies
\[
             0\le\frac{r-w}{w}\le\frac{\eta}{1-\eta}
                                  \le\frac{\Theta}{1-\Theta}.
\]
This is a relative rank comparison, not a bound for an arithmetic
degree. Its application retains the original projective model,
section realization, nefness, birationality and degree comparison.
The signed cusp mean and the other arithmetic section remainders
remain separate requirements.


\subsection{Keep the component degrees when averaging}

\begin{lemma}[Cancellation of the degree weights]
\label{rh-degree-weights}
Let $0\le\rho<1$, $\delta_*>0$, and
$\delta_*\le\delta_Z\le\delta_* /(1-\rho)$.
Put $w_Z=B_0\delta_Z$ with $B_0>0$, $W_0=\sum_Zw_Z$, and
$N=|\mathcal Z|$. Suppose $c,\gamma\ge0$ and
\[
 \delta_*(s_Z+c)\le\delta_Z(h_Z+c+\gamma).
\]
With $q_0=NB_0\delta_*/W_0$, $\bar s=N^{-1}\sum_Zs_Z$ and
$M=W_0^{-1}\sum_Zw_Zh_Z$, one has
\begin{equation}\label{rh-degree-weights-result}
 1-\rho\le q_0\le1,\qquad
 M\ge q_0\bar s-(1-q_0)c-\gamma.
\end{equation}
Thus $\bar s\ge-X$ gives $M\ge-H$, where
$H=q_0X+(1-q_0)c+\gamma\le X+\rho c+\gamma$.
\end{lemma}

\begin{proof}
The degree inequalities give the bounds on $q_0$. Multiply the
component comparison by $B_0$, sum it, and divide by $W_0$.
\end{proof}

The degree at the smaller weight is the common $\delta_*$. The
image degrees $\delta_Z$ may vary. Summing before discarding them
therefore retains the geometric averaging measure exactly.

\subsection{Short sections replace a negative-height moment}

\begin{proposition}[A degree bound from one mean height]
\label{rh-short-sections}
Let $E_Z$ be nonzero Hermitian lattices, with ranks $r_Z$ and
normalized degrees $d_Z$. Suppose integral vectors of norm at most
$e^L$ at every embedding span each lattice over the field, where
$L\ge0$. Let $0<w_Z\le r_Z$ and suppose
\[
 d_Z\ge w_Z\{t h_Z-a|h_Z|-A\log(2+|h_Z|)-R_0\},
 \quad a,R_0\ge0,\quad A>0,\quad t\ge a+1.
\]
Put $W_0=\sum_Zw_Z$, $R=\sum_Zr_Z$, $\vartheta=W_0/R$,
$M=W_0^{-1}\sum_Zw_Zh_Z$, and
$\lambda=t-a-1$, $b_0=t+a$, $K_0=\max(2,A)$. Then
\begin{equation}\label{rh-one-mean-exact}
 \frac{\sum_Zd_Z}{R}
 \ge\lambda\vartheta M-\frac{2a+1}{b_0}L-\vartheta R_0
       -\vartheta A\log\!\left(K_0+\frac{L}{b_0\vartheta}\right).
\end{equation}
In particular, $M\ge-H$ with $H\ge0$ implies
\begin{equation}\label{rh-one-mean-bound}
 \frac{\sum_Zd_Z}{R}
 \ge-\lambda\vartheta H-\vartheta\{R_0+A\log\max(A,1)\}
       -2\log2-\frac{2a+2}{t+a}L.
\end{equation}
The term $2\log2$ is unnecessary when $A\ge2$.
\end{proposition}

\begin{proof}
Choose a field basis among the short integral vectors. The determinant
norm inequality, including the nonnegative integral-index contribution,
gives $d_Z\ge-r_ZL$.
For one component put $P=r_ZL/w_Z$. If $h_Z=-y\le0$ and
$y\le P/b_0$, its assumed degree bound yields
\[
 d_Z/w_Z\ge-\lambda y-\frac{2a+1}{b_0}P
                           -R_0-A\log(K_0+P/b_0).
\]
If $y\ge P/b_0$, the same conclusion follows from $d_Z/w_Z\ge-P$,
since $b_0-\lambda=2a+1$. For $h_Z\ge0$, use
$A\log(2+h_Z)\le h_Z+A\log K_0$. Thus in every case
\[
 d_Z/w_Z\ge\lambda h_Z-\frac{2a+1}{b_0}P
                           -R_0-A\log(K_0+P/b_0).
\]
The weighted mean of $P$ is $L/\vartheta$; concavity of the logarithm
proves \eqref{rh-one-mean-exact}. Finally,
\[
 \vartheta A\log\!\left(K_0+\frac{L}{b_0\vartheta}\right)
 \le\vartheta A\log K_0+L/b_0,
 \quad A\log K_0\le A\log\max(A,1)+2\log2.
\]
This proves \eqref{rh-one-mean-bound}, also when $L=0$ or $\lambda=0$.
\end{proof}

Only the weighted mean $M$ is needed. The short sections already give
the independent lower degree bound used on components of large negative
height. Without that bound a mean alone is insufficient: equal weights
at heights $H$ and $-H$ keep the mean zero while the logarithmic and
absolute-value losses grow without bound. Short sections must belong
to these same component lattices, with these metrics.

\subsection{Descent preserves the actual rank ratio}

\begin{lemma}[Determinant descent]\label{rh-descent}
Let $E,F$ be nonzero Hermitian lattices and let $n\ge1$ be an integer.
Suppose an integral injection $E\to F^{\oplus n}$ has operator norm
at most $e^\beta$ at every embedding. Suppose $F$ is spanned by
integral vectors of norm at most $e^\Lambda$ at every embedding.
With $a=\rank E$, $b=\rank F$ and $\tau=a/(nb)$,
\[
                  \mu(F)\ge\tau\{\mu(E)-\beta\}
                                      -(1-\tau)\Lambda.
\]
\end{lemma}

\begin{proof}
Saturate the image. Its degree is at least
$\widehat{\deg}(E)-a\beta$ by the determinant norm inequality.
The quotient has rank $nb-a$ and is spanned by the images of the
short vectors, so its degree is at least $-(nb-a)\Lambda$.
Add the degrees and divide by $nb$.
\end{proof}

A signed section-degree error is multiplied by its actual ratio
$\tau$ in this formula. Replacing it prematurely by an absolute
value can increase the coefficient unnecessarily. Flat twists likewise
enter once, through their actual rank-weighted degree costs. The
contribution of this section to the final arithmetic budget still
requires a compatible mean-volume theorem, quantitative section-degree
bound, short spanning sections, and the stated integral descent maps.
